\documentclass[11pt,a4paper]{article}
\usepackage{amsmath,amssymb,amsthm}
\usepackage[margin=23mm]{geometry}
\usepackage{color}
\usepackage[colorlinks=true,linkcolor=blue,urlcolor=blue,filecolor=blue,bookmarksopen=true]{hyperref}
% Scale existing Computer Modern bold math fonts where a design-size TFM
% is absent from the existing, deliberately small offline bundle.
\DeclareFontShape{OT1}{cmr}{bx}{n}{<-> cmbx8}{}
\DeclareFontShape{OT1}{cmr}{m}{n}{<-> cmr10}{}
\DeclareFontShape{OML}{cmm}{m}{it}{<-> cmmi10}{}
\DeclareFontShape{OML}{cmm}{b}{it}{<-> cmmi8}{}
\DeclareFontShape{OMS}{cmsy}{m}{n}{<-> cmsy10}{}
\DeclareFontShape{OMS}{cmsy}{b}{n}{<-> cmsy10}{}
\DeclareFontShape{OMX}{cmex}{m}{n}{<-> cmex10}{}
\input{umsa.fd}
\input{umsb.fd}
\DeclareFontShape{U}{msa}{m}{n}{<-> msam10}{}
\DeclareFontShape{U}{msb}{m}{n}{<-> msbm10}{}
% The two AMSa glyphs in these sources are only proof-end squares.
\renewcommand{\square}{\mathord{\vcenter{\hbox{\fboxsep=0pt\fboxrule=.4pt\fbox{\rule{0pt}{.7ex}\rule{.7ex}{0pt}}}}}}
\renewcommand{\blacksquare}{\mathord{\vcenter{\hbox{\rule{.8ex}{.8ex}}}}}
% This cache contains Roman faces but no LM Mono. Use the cached Unicode
% face even for path styling; a full TeX installation is not required.
\makeatletter
\input{tulmtt.fd}
\DeclareFontShape{TU}{lmtt}{m}{n}{<-> \UnicodeFontFile{lmroman10-regular}{}}{}
\DeclareFontShape{TU}{lmtt}{m}{it}{<-> \UnicodeFontFile{lmroman10-italic}{}}{}
\DeclareFontShape{TU}{lmtt}{m}{sl}{<-> \UnicodeFontFile{lmroman10-italic}{}}{}
\DeclareFontShape{TU}{lmtt}{b}{n}{<-> \UnicodeFontFile{lmroman10-bold}{}}{}
\g@addto@macro\UrlBreaks{\do\0\do\1\do\2\do\3\do\4\do\5\do\6\do\7\do\8\do\9\do\a\do\b\do\c\do\d\do\e\do\f}
\makeatother
\setlength{\parindent}{0pt}
\setlength{\parskip}{5pt plus 2pt minus 1pt}
\setlength{\emergencystretch}{3em}
\setcounter{secnumdepth}{5}
\allowdisplaybreaks[2]
\urlstyle{same}
\providecommand{\tightlist}{\setlength{\itemsep}{0pt}\setlength{\parskip}{0pt}}
\providecommand{\PaperSixMissingEvidence}[2]{#2 (unbundled archive path: \path{#1})}
\providecommand{\SubtitleForOfflineBuild}[1]{\par\medskip{\large #1\par}}
\newlength{\OfflineResearchTableWidth}
\newcommand{\OfflineResearchHeaderRule}{\noindent\rule{\linewidth}{0.45pt}\par}
\newcommand{\OfflineResearchRow}{\par}
\newenvironment{OfflineResearchTable}{%
  \par\medskip\begingroup\small\setlength{\parindent}{0pt}%
  \setlength{\parskip}{3pt plus 1pt minus .5pt}%
  \setlength{\OfflineResearchTableWidth}{\linewidth}%
}{%
  \par\endgroup\medskip%
}

% Only simple Latin multi-indices n, j, a use \boldsymbol in the frozen draft.
\font\OfflineBoldText="[KaTeX_Math-BoldItalic.ttf]" at 10.95pt
\font\OfflineBoldScript="[KaTeX_Math-BoldItalic.ttf]" at 7pt
\font\OfflineBoldScriptScript="[KaTeX_Math-BoldItalic.ttf]" at 5pt
\renewcommand{\boldsymbol}[1]{\mathord{\mathchoice
  {\hbox{\OfflineBoldText #1}}{\hbox{\OfflineBoldText #1}}
  {\hbox{\OfflineBoldScript #1}}{\hbox{\OfflineBoldScriptScript #1}}}}
\hypersetup{pdftitle={Native Information and Nonlinear Reconstruction in Interpolation Counting},pdfauthor={Chuhang Ding},pdfsubject={Evaluation capacity and native interpolation},pdfcreator={Offline XeTeX proof replay}}
\title{Native Information and Nonlinear Reconstruction in Interpolation Counting\SubtitleForOfflineBuild{Exact Decoder Capacities, Native Interfaces, and Escape Mechanisms}}
\author{Chuhang Ding (Peking University)}
\date{26 November 2026}
\newcommand{\PaperSixCatalogue}[1]{\csname hyper@linkurl\endcsname{#1}{EVIDENCE_CATALOG.md}}
\newcommand{\PaperSixEvidence}[2]{\PaperSixCatalogue{#2}}
\begin{document}
\maketitle
\begin{abstract}
We distinguish the information determined by finite interpolation data from
the cost of acquiring those data and from the cost of producing the final
count. Newton--Smith coordinates describe exact reconstruction capacities.
We compare these algebraic interfaces with genuine bounded-occurrence
counting instances and actual root and child registers. Low-precision
whole-menu limits coexist with constructive higher-precision recovery.
The bounds apply to the interfaces specified in each theorem rather than
to all optimized implementations of interpolation counting.
\end{abstract}
\tableofcontents
\section*{How the two parts fit together}
Part I asks what data determine, for a fixed observation interface and a
specified receiver precision. Part II asks which of those data are
obtainable from the original source registers. These are distinct questions:
a bound on the number of possible transcripts is not an acquisition bound,
and a polynomial number of queries is not a polynomial bit-time bound.
The separately constructed physical witness families in the two parts
retain their own hypotheses; no theorem about the five-copy family is
invoked on the four-copy family without a construction.
See the \hyperref[paper6-current-coverage]{current theorem coverage and
release evidence}, distinct from the original versioned appendices.
\subsection*{Reading the source comparisons}
\textbf{Same addresses, different precision.} In the four-clone family in J9, the public probes and modes are fixed on one
physical layout before the pin signs vary. The prescribed 10M-bit observations
recover all signs from N(10M+1) requested products. Projecting the same
original registers to M/12 bits makes every value in the common low-precision
menu zero. That comparison concerns available observations. It does not price
their acquisition or authorize separate runs at the two precisions.

\textbf{Different locality rules.} C3 uses a different, five-copy family. The MAIN calculation separates
individual switch factors across a cut even when a MAIN group meets both
sides. The Appendix D locality certificate concerns one child return as a
function of its whole group, and requires the specified cut to be a union of D
groups. For the fixed lists, the two precision proofs permit replacing
acquired complete slots by the indicated child lifts while preserving the
terminal residue. They do not establish the bit time of an arbitrary adaptive
algorithm.

Start with the \hyperref[paper6-featured-index]{featured theorem index} for precise evidence and limits.
More recent checks appear in the \hyperref[paper6-late-coverage]{later source proof map}.
The \hyperref[r47-v9-review-consequences]{J2 supplement} gives the general
seed-law proof. The \hyperref[sec:retained-physical-execution]{retained-execution
appendix} specifies the J6 source inputs and the checked count, clock
and named-buffer results.

\paragraph*{Historical evidence.}
Links to historical records open the \PaperSixCatalogue{public
evidence catalogue}. It lists original paths, SHA-256 hashes and
public copies where available. The TeX source retains the paths.


\part{Exact capacities and limits of native observations}
\section{Research question and principal
results}\label{v8-research-question-and-principal-results}

The obstruction to a nonlinear extension of a rank budget is already
visible in the equality function: a product of local equality factors
has exponential value-matrix rank and a linear-size arithmetic circuit.
Consequently a proof that an interpolation program encounters a
high-rank object is not a running-time lower bound. The present work
instead fixes the messages, operations, coefficient ring and precision
at an actual interface, then determines the functions that interface can
produce.

The organizing object is an evaluation module. Integral polynomial
operations generate a module of functions; a denominator budget
saturates it, and finite output precision takes a quotient. Its Smith
profile gives both exact decoder capacity and a target-specific minimum
denominator loss. The latter can be large for a rank-one target. We
compute the profile explicitly for scalar messages, give an exact finite
matrix for indivisible jet messages over every finite residue field, and
prove sharp witnesses. The proofs allow repeated messages and arbitrary
polynomial DAGs without expanding their computational paths.

There are two complementary barriers beyond this polynomial receiver.
For linear observations on a full input module, arbitrary decoding is
governed exactly by a kernel condition; finite Frobenius duality
supplies a row criterion, including unequal precision rectangles. For a
genuine bounded-occurrence CNF, positive origin contraction needs
exponentially many actual scalar gates on expander count slices,
including overlapping products and cache reuse. Its proof uses positive
support contexts and classical monotone-circuit arguments, not a
rank-as-time inference.

We then examine the supplied OpenAI manuscript, \emph{The Exponential
Time Hypothesis is False}, in its revision dated 25 October 2026. We
call its main-text rational route MAIN and keep its Appendix D label for
the other route. The archived source is pinned by content hashes; the
online reference identifies the supplied version.
\href{https://github.com/openai/math/tree/main/preprints/The-Exponential-Time-Hypothesis-is-False-October-25-2026}{Source
manuscript}

For these interfaces, a complete low-precision D-coordinate menu has too
little row information even for an arbitrary decoder. A sparse native
MAIN/D cross-product has a work barrier under an explicit acquisition
contract, with the published MAIN exponent for one D coordinate.
Ordinary canonical slot lifts obey the requisite loss inequality. These
are source-specific extensions of the v7 frontier.

The operation boundary is equally explicit. Two legal raw clause entries
give a Boolean variable and its complement; nonnegative two-rail gates
simulate every Boolean circuit. Contraction to the positive origin
tensor removes that complementary entry. Higher-digit tests escape
low-loss polynomial capacities, and freely designed field encodings
attain the one-way width bound on the older, easily computed graph
family. None of these facts gives a polynomial-time solver for general
SAT. They prevent incorrect generalization of the valid interface
barriers.

The low-precision result also survives pooling every completed root
coordinate and every probe, with arbitrary joint decoding. On one fixed
cubic expander topology, the same physical CNF attains exponentially
many exact counts, while the entire pooled menu has a sublinear bit
encoding. An exact decoder must obtain a linear number of additional
bits. This argument compares output values and joint transcripts, so it
does not need a common fixed matrix cut across the coordinates. It still
requires the stated low-precision bound; the source's full-precision
recovery formula is outside that restriction.

The complete source, formal modules, exact finite checks, independent
reviews, coverage matrix and reproduction instructions accompany this
paper. Mathematical prerequisites with established provenance are
identified individually. Research hypotheses, counterexamples and model
limits are retained. The paper establishes no priority claim for its
classical ingredients and no complete obstruction for unrestricted
signed native hybrid computation.

\section{The observable-module framework}\label{v8-section-framework}

This framework computes what polynomial arithmetic can recover from a
fixed message interface at a given precision and denominator budget. On
a common family of source CNFs it also bounds positive
origin-contraction gates and the information available from completed
Appendix D root slots at low precision. Each execution conclusion
specifies which source operations it charges.

\subsection{An interface specifies operations and
information}\label{v8-an-interface-specifies-operations-and-information}

Fix finite sender and receiver domains X and Y. A directed interface
specifies the sender's legal observations, their costs and attained
joint image Z, the receiver's own input, the coefficient ring, and the
permitted operations after communication. An observation is charged when
acquired; reusing a register does not acquire another observation. A
scalar, a jet, a coefficient-extracted jet and a whole word subject to
arbitrary lookup are different interfaces. Overlapping input views must
be declared as additional shared information or handled by conditioning.
A local table address is not the same primitive as a completed
contraction.

Let O be a discrete valuation ring with uniformizer pi. For readability
the equations below write these as \(\mathcal O,\pi\). Let
\(R_p=\mathcal O/\pi^p\), with \(p\ge1\). Choose a finite family of
scalar functions on X, or scalar coefficient readouts of functions,
whose O-linear span is \(\mathcal E\subseteq\mathcal O^X\). The family
must be the one actually generated by the permitted gates before
division. For scalar polynomial decoders it is the evaluation algebra
\(\mathcal O[u_1,\ldots,u_h]\). For indivisible jet decoders it is the
module generated by their legal terminal coefficient readouts. Both are
closed under the receiver's permitted integral linear recombinations.

At a common denominator budget \(L\ge0\), define

\[
 \mathcal C_{L,p}(\mathcal E)=
 \big((\pi^{-L}\mathcal E\cap\mathcal O^X)+\pi^p\mathcal O^X\big)
                /\pi^p\mathcal O^X. \tag{U1}
\]

This is an attainable-column module, not just a number. A target matrix
is realizable precisely when every receiver column belongs to (U1),
provided the receiver may choose any finite permitted polynomial for
that column. This criterion makes no assertion that such a polynomial is
cheap to find or evaluate. The polynomial DAG semantics use one
indeterminate per distinct message and survive arbitrary fan-out and
overlapping products. Intermediate joint tests and non-polynomial lifts
require a new closure.

\subsection{The exact target budget}\label{v8-the-exact-target-budget}

Write a Smith presentation in an ambient integral basis as

\[
 \mathcal E=\bigoplus_{i=1}^{r}\pi^{s_i}\mathcal O e_i
 \ \subseteq\ \bigoplus_{i=1}^{n}\mathcal O e_i,
 \qquad n=|X|,\quad s_i\ge0. \tag{U2}
\]

Such a presentation exists because O is a DVR and X is finite. This is
ordinary Smith theory over O; no principal-ideal property is assumed of
a double-truncated jet ring. Let \(v_{iy}\in R_p\) be the coordinates of
the desired column y in this ambient basis. Define the truncated order
\(\nu_p(0)=p\), and \(\nu_p(v)=t\) when
\(v\in\pi^tR_p\setminus\pi^{t+1}R_p\).

\textbf{Theorem U (exact observable-module budget).} An integral target
matrix with those coordinates has a permitted reconstruction with common
loss L if and only if

\[
 v_{iy}=0\quad(i>r),\qquad
 \nu_p(v_{iy})\ge\min\{p,(s_i-L)_+\}\quad(i\le r). \tag{U3}
\]

If any outside coordinate is nonzero, no finite loss suffices. Otherwise
the exact minimum loss is

\[
 L_{\min}=\max\bigl(\{0\}\cup
     \{(s_i-\nu_p(v_{iy}))_+:i\le r,\ \nu_p(v_{iy})<p\}\bigr).
 \tag{U4}
\]

The attainable module itself is

\[
 \mathcal C_{L,p}(\mathcal E)\cong
 \bigoplus_{i=1}^{r}
 \mathcal O/\pi^{[p-(s_i-L)_+]_+}\mathcal O. \tag{U5}
\]

Its maximum attainable matrix separation rank, with arbitrarily many
receiver columns, is exactly \(\#\{i\le r:s_i<L+p\}\). For a specified
target, evaluating (U3) or (U4) requires the ambient change of basis as
well as the multiset of Smith exponents.

\textbf{Proof.} Scale (U2) by \(\pi^{-L}\) and intersect with the
ambient integral lattice. This replaces each exponent by \((s_i-L)_+\).
Reduction modulo \(\pi^p\) gives (U5), and coordinate membership gives
(U3). A nonzero target coordinate of order t less than p survives
exactly when \((s_i-L)_+\le t\). Taking the maximum of these necessary
and sufficient lower bounds on L proves (U4). The same argument
constructs a realizing numerator for each column; closure under integral
linear combinations and the definition of E supply the legal
pre-division expression. Thus the value of (U4) is basis independent,
though its coordinate computation uses a Smith basis.

For the rank assertion, selecting each nonzero cyclic generator as one
column gives a matrix whose image needs exactly the indicated number of
generators. A submodule of an r-generated module over a finite chain
ring needs at most r generators: pull it back along a surjection from a
free O-module of rank r. Conversely an integral matrix of fraction-field
rank r has an integral factorization of width r over the DVR. These
facts give both the upper bound and the matching width witness.
\(\square\)

This exact membership invariant is stronger than output matrix rank. For
example, extracting bit j from an ordinary dyadic integer encoding has a
single target column and hence rank at most one, but requires loss
\(2^j-1\). The proof uses the evaluation lattice on \(0,\ldots,2^j\),
whose Newton directions have exponents \(v_2(i!)\); the target's last
binomial coordinate is odd. Conversely \(\binom{x}{2^j}\bmod2\) is
precisely that bit. An arbitrary digit test performs this operation in
polynomial bit time outside the polynomial decoder interface.
\hyperref[v8-section-nonlinear-budget]{Complete scalar proof}

\subsection{Concrete complete
instances}\label{v8-concrete-complete-instances}

For h scalar messages in public classes \(c_i+\pi^{v_i}\mathcal O\), set
\(w_q(n)=\sum_{a\ge1}\lfloor n/q^a\rfloor\), where q is the residue
field size. Newton interpolation gives the exact universal width

\[
 H=\#\{\boldsymbol n\in\mathbb N^h:
       \sum_i(v_i n_i+w_q(n_i))<L+p\}. \tag{U6}
\]

The arbitrary actual-image version uses (U2) and can be strictly
smaller. At residue precision, for unrestricted messages,
\(H_q(h,L)\le q^h\binom{h+L}{h}\), with equality when L\textless q.
These factorial weights are classical; the theorem adds the prescribed
decoder, output loss, target membership and source-interface accounting.

For whole jets in \(\mathcal O[z]/z^k\), integral polynomial gates and
one terminal scalar coefficient readout, the coefficient features have
multinomial orders \(\mu(\alpha,j)\). If \(q>L+p+k-2\), the exact
capacity is \(q^h\#\{(\alpha,j):\mu(\alpha,j)<L+p\}\). For every q, an
integral Newton change of basis on the digit variables, followed by
ordinary DVR Smith reduction of a finite weighted coefficient matrix,
gives the exact replacement. This is a finite mathematical
characterization. Both the column cutoff and the degree bound \(P+k-2\)
of the universal matrix depend on its raw precision \(P=L+p\). When
increasing \(L\) at fixed output precision, recompute with the new
cutoff or use the exact untruncated evaluation module on the prescribed
finite image. The construction supplies no large-instance time bound.
\hyperref[v8-section-jet-polynomial-capacity]{Jet formula} ·
\hyperref[v8-section-jet-newton-smith]{All-field Smith characterization}

Arbitrary decoding of \emph{linear} observations has a different
complete answer. On a full module, a decoder of T from E exists
precisely when \(\ker E\subseteq\ker T\). Over a finite commutative
Frobenius ring this is equivalent to row-module inclusion and to a
linear decoder. Unequal precisions in \((\mathcal O/\pi^P)[z]/z^K\) are
handled by multiplying observation rows by \(\pi^{P-P_j}z^{K-K_j}\),
with the analogous scaling for target rows. An exact deterministic
adaptive linear-query tree, whose only input access is through its
queries, has a zero execution whose fixed query list already passes this
criterion. On a structured promise \(S\), the exact fiber test checks
pairs in \(S\), or equivalently \((S-S)\cap\ker E\subseteq\ker T\) when
both maps extend linearly. \hyperref[v8-section-duality-typeset]{Duality
and adaptive query proof}

\subsection{How the invariant becomes a
barrier}\label{v8-how-the-invariant-becomes-a-barrier}

On the v7 common matching restriction the final residue matrix has rank
\(2^m\). Fix the graph-only layout menu before selecting its single
matching and orientation, under the simultaneous matching hypotheses. A
block whose cut has imbalance \(\delta_b\), hazardous overlap
\(\ell_b\), and an observable capacity bound \(H_b\) contributes at most
\(2^{m-\delta_b+\ell_b}H_b\). Thus every such fixed nonlinear hybrid
reconstruction satisfies

\[
 1\le\sum_b\min\{1,2^{-\delta_b+\ell_b}H_b\}. \tag{U7}
\]

\(H_b\) must bound a fraction-field representative on every rectangular
fiber used to condition the overlap and minority-side bits. This lets
the sum of the block representatives retain the correct integral residue
even when individual blocks have denominators. Joint guards, raw inputs
and additional divisions require their own dependence or loss analysis.
If one trace covers every distinct designated output at rate at most
\(W\) per charged operation, the number of those outputs divided by
\(W\) lower-bounds its time. A fused evaluation must justify that
coverage for its own executed operations.

The source audit supplies two more concrete interfaces. A fixed
cross-product with a MAIN operand obtained by explicit
private-assignment enumeration retains a sharp native work barrier, even
without assigning which edges each operand computes. Separately, the
complete D-slot menu at one coordinate and low precision factors through
fewer row bits than the matching target requires, ruling out
\emph{arbitrary} decoding of that entire menu at that precision. The
printed reconstruction uses higher precision losses and more than one
coordinate. \hyperref[v8-section-native-escape]{Exact native hypotheses and
bounds}

\subsection{A complementary actual gate
theorem}\label{v8-a-complementary-actual-gate-theorem}

For bounded-degree expanders the constructed occurrence-split CNFs have
the exact count slice \(3\sum_{I\in\mathrm{IS}(G)}\prod_{v\in I}s_v\).
Its positive origin contraction needs \(2^{\Omega(N)}\) binary scalar
gates, even with arbitrary product overlap and DAG reuse. The proof
counts actual gate contexts; it does not materialize polynomial
monomials. It extends to positive local atoms, small fixed subcubes, and
adaptive traces whose entire control footprint fixes only a sufficiently
small linear number of switches. Modular execution requires a positive
terminal lift below its modulus. The result uses and credits classical
monotone counting and expander-cover techniques.

The source's stated signed upper bound, taken as an input theorem, gives
a subexponential procedure outside this positive execution contract. The
ordinary printed readouts use signed differences or modular lifts.
\hyperref[v8-section-stronger-family]{Full CNFs and gate proof}

The shortened cubic CNF puts three bounds on the same full counting
function. At any fixed certified D coordinate, one matching restriction
has physical parity matrix \(I_{2^m}\). It witnesses both the sparse
MAIN cross-product rank bound and the one-coordinate low-precision
decoder obstruction. Its exact count on the matching alone has a linear
positive circuit. The exponential positive gate theorem requires
correctness on the full availability cube or one of its permitted
low-codimension subcubes. \hyperref[v8-section-common-native-family]{Common
source inputs and matching}

On a choice of these cubic expanders, the exact physical count takes at
least \(2^{(N-10)/12}\) distinct values. Across all root coordinates,
every completed D slot at precision \(p_i\le a+1-i\) is determined by a
joint encoding of \(O(s^{2\varepsilon}\log^2s)=o(N)\) bits. Even an
arbitrary decoder then needs \(\Omega(N)\) additional bits of
input-dependent information to recover the full count. This proof
compares the number of possible whole menus with the number of possible
counts, so it requires no common cut across those coordinates.
\hyperref[v8-section-pooled-menu-barrier]{All-coordinate exact-count proof}

\subsection{What is still not
determined}\label{v8-what-is-still-not-determined}

A large observable loss need not require long execution if the program
keeps sparse or factored values. A large rank need not require many
gates. An arbitrary scalar word predicate can inspect digits unavailable
to a low-loss polynomial. Two raw clause table entries can supply a
switch and its complement; positive arithmetic on such two-rail inputs
simulates every Boolean circuit with constant gate overhead. Restricting
only the word ``native'' leaves that simulation intact.

Consequently a full signed native hybrid barrier still needs a
source-derived data-flow or operation invariant that survives its actual
reconstruction, or a costly-acquisition theorem for the completed
registers it genuinely requests. No such theorem for every optimized
source-preserving signed execution is proved here. The results above
give exact interface tests, larger proven submodels, and explicit escape
mechanisms rather than an unconditional SAT or P-versus-NP conclusion.

\subsection{Attribution}\label{v8-attribution}

Smith normal form, Newton/P-ordering factorials, finite-local-ring
residue selectors and shifted powers, Frobenius-ring duality, balanced
monotone circuit decompositions, induced-matching rank arguments, and
the use of rooted caterpillar ratios and pairs of terminal counts all
precede this project. The candidate contributions are the specific
\nolinkurl{decoder/loss/readout} synthesis, exact small-field jet
evaluation profile, native source reductions and precision obstructions,
and their joint coverage and escape analysis. They remain subject to
external expert and broader literature review. The public evidence
catalogue identifies the available historical records.
\href{https://doi.org/10.1016/j.tcs.2017.09.002}{Bulyovszky and Horváth}
· \href{https://arxiv.org/html/2604.19094}{Chan, Heilman and Panova}

\section{Exact scalar polynomial-decoder
capacity}\label{v8-section-nonlinear-budget}

On a balanced cut of the genuine pinned cubic-expander family, a one-way
polynomial recombination over the Appendix D coefficient ring needs
\(\Omega(N/\log N)\) distinct scalar messages when its denominator loss
is polynomial in \(N\). This width is attainable with zero loss and
polynomial work when the sender may design its messages. Each message
may be reused without limit. The underlying scalar capacity is exact at
every precision, including for an arbitrary prescribed finite set of
attainable message tuples.

\subsection{The interface}\label{v8-the-interface}

Let \(\mathcal O\) be a discrete valuation ring with uniformizer
\(\pi\), fraction field \(K\), and residue field \(\kappa\) of
cardinality \(q\). Write \(R_p=\mathcal O/\pi^p\mathcal O\) for
\(p\ge1\). On a finite Cartesian set of sender and receiver inputs
\(X\times Y\), all sender-dependent data transmitted to the receiver
consist of \(h\) distinct scalar values \(u_i(x)\in\mathcal O\). The
receiver may use arbitrary functions of \(y\), constants, additions,
subtractions, multiplications and arbitrary reuse in a finite arithmetic
DAG. The pertinent condition on its lifted output is

\[
 F(x,y)=\pi^{-L}P_y(u_1(x),\ldots,u_h(x)),\qquad
 P_y\in\mathcal O[U_1,\ldots,U_h],\quad L\in\mathbb N.                 \tag{V8-S1}
\]

Equality modulo \(\pi^p\) in (V8-S1), meaning a difference in
\(\pi^p\mathcal O\), suffices. Receiver coefficients may differ
arbitrarily with \(y\). In particular, an addition or multiplication may
join two registers that both depend on all \(h\) messages. The same
symbol \(U_i\) replaces every reference to a shared sender value.
Induction in topological order proves (V8-S1) for an integral polynomial
suffix with a final fixed denominator \(\pi^L\). A fixed-scalar division
by a unit is integral. For intermediate fixed denominators, one
sufficient syntactic calculation assigns loss zero to integral inputs,
takes the maximum at addition or subtraction, adds the two losses at
multiplication, and adds \(a\) at a division by \(\pi^a\). A fixed scale
\(c\in K\) may use \(\max(0,L_{\rm in}-v(c))\) for \(c\ne0\). Common
factors can improve this bound if a polynomial identity proves the
cancellation. The loss is a bound on a denominator exponent, not the bit
length of that exponent.

Finite input registers can be lifted separately on the sender's own
view. If every scalar argument and lifted polynomial coefficient is
known modulo \(\pi^{L+p}\), integral polynomial evaluation preserves
this congruence and division by \(\pi^L\) gives the required output
precision. An implementation with heterogeneous precisions, intermediate
rounding, or an increase in precision can instead prove (V8-S1) at its
actual interface. A joint zero test, address lookup, digit extraction,
or canonical lift followed by extra arithmetic needs its own analysis.
Decisions and canonical lifts made entirely within one private view may
be absorbed into that view's functions. Denominator powers on separate
summands can differ.

Some input-dependent operations admit polynomial implementations at the
actual precision. In any finite ring of interest here, on a register
guaranteed to be a unit, \(u^{-1}=u^{|R^\times|-1}\). Over \(R_P\) an
exponent \(E\ge P\) divisible by the order of the unit group gives the
exact zero-or-one indicator of nonzero residue by \(u^E\). Hence an
equality test of residues \(u\bmod\pi=a\bmod\pi\) has the integral
polynomial representative \(1-(u-a)^E\). These observations also justify
such gates inside an integral polynomial numerator computed at precision
\(P\). They do not give arbitrary equality testing at all \(P\) digits.

Over a commutative ring \(R\), the \textbf{separation rank} on
\(X\mid Y\) is the least \(r\) such that
\(M(x,y)=\sum_{j=1}^r a_j(x)b_j(y)\) over \(R\), equivalently the inner
width in \(M=UV\). Over fields this is ordinary matrix rank. Over
\(\mathbb Z/4\), the matrix \(2I_2\) has zero determinant and zero
residue matrix, yet its image needs two generators and its separation
rank is two. The theorem below gives this width over \(R_p\) and,
separately, a low-rank representative over \(K\).

\subsection{The exact scalar capacity}\label{v8-the-exact-scalar-capacity}

Allow a fixed public congruence \(u_i(x)\in c_i+\pi^{v_i}\mathcal O\),
where \(c_i\in\mathcal O\) and \(v_i\) are nonnegative integers. Set
\(c_i=v_i=0\) for an unrestricted scalar. Define

\[
 w_q(n)=\sum_{j\ge1}\left\lfloor n/q^j\right\rfloor,\qquad
 s(\boldsymbol n)=\sum_{i=1}^h\big(v_i n_i+w_q(n_i)\big),\qquad
 H=\#\{\boldsymbol n\in\mathbb N^h:s(\boldsymbol n)<L+p\}.          \tag{V8-S2}
\]

\textbf{Theorem 1.} Every matrix satisfying (V8-S1) has a \(K\)-valued
representative \(\widetilde F\) of ordinary rank at most \(H\) such that
\(F-\widetilde F\) has all entries in \(\pi^p\mathcal O\). This
assertion even permits nonintegral \(F\). If \(F\) is pointwise
integral, its matrix modulo \(\pi^p\) has separation rank at most \(H\)
over \(R_p\). For arbitrary sender and receiver functions, both the
universal finite-ring width and the minimum possible \(K\)-rank bound on
integral representatives are sharp. In particular the exact largest
ordinary rank of the residue matrix at \(p=1\) equals \(H\). The case
\(h=0\) has \(H=1\), corresponding to arbitrary receiver-only output.

\textbf{Proof of the upper bound.} Choose representatives
\(d_0=0,d_1,\ldots,d_{q-1}\) of the residue field. For the base-\(q\)
expansion \(n=\sum_j b_j q^j\) set \(a_n=\sum_j d_{b_j}\pi^j\). Two such
points agree modulo \(\pi^s\) precisely when their integer indices agree
modulo \(q^s\). The first \(q^s\) points exhaust all residue classes
modulo \(\pi^s\). Therefore among \(a_0,\ldots,a_{n-1}\) at least
\(\lfloor n/q^s\rfloor\) points agree with \emph{any} \(u\in\mathcal O\)
modulo \(\pi^s\). For \(u=a_n\) the count is exactly
\(\lfloor n/q^s\rfloor\). Counting valuation one level at a time proves

\[
 v\!\left(\prod_{j<n}(u-a_j)\right)\ge w_q(n),\qquad
 v\!\left(\prod_{j<n}(a_n-a_j)\right)=w_q(n).
\]

Infinity at a root is allowed. For coordinate \(i\) use the monic
polynomial \(\Phi_{i,n}(U)=\prod_{j<n}(U-c_i-\pi^{v_i}a_j)\). At an
admissible \(u_i\), it has valuation at least \(v_i n+w_q(n)\), with
equality at \(c_i+\pi^{v_i}a_n\). Products
\(\Phi_{\boldsymbol n}=\prod_i\Phi_{i,n_i}\) form an integral basis of
the polynomial ring as a module, by the univariate monic triangular
change of basis in each coordinate. In the expansion of \(P_y\), discard
all terms with \(s(\boldsymbol n)\ge L+p\). Their total, divided by
\(\pi^L\), lies in \(\pi^p\mathcal O\) at every input. Each surviving
term is a function of \(x\) times a function of \(y\) over \(K\), so
their sum has rank at most \(H\).

For integral \(F\), that sum is itself integral. An integral matrix of
\(K\)-rank \(r\) factors through \(\mathcal O^r\) by ordinary DVR Smith
form. Reducing the factors gives separation rank at most \(r\) over
\(R_p\). Individual terms of the sum may have denominators; their
separate reductions are unnecessary. The same argument gives an often
smaller bound if an independently known downward-closed set contains the
monomial support of every \(P_y\). The triangular Newton change of basis
only lowers coordinate exponents, so intersect that set with the region
in (V8-S2).

\textbf{Sharpness.} Use the finite lower set
\(I=\{\boldsymbol n:s(\boldsymbol n)<L+p\}\) as both the set of row
indices and the set of column indices. Let the row \(\boldsymbol j\)
supply \(u_i=c_i+\pi^{v_i}a_{j_i}\). For each column \(\boldsymbol n\),
take the polynomial function

\[
 F_{\boldsymbol n}(u)
 =\pi^{(s(\boldsymbol n)-L)_+}\,
                \Phi_{\boldsymbol n}(u)/\pi^{s(\boldsymbol n)}.
                                                                    \tag{V8-S3}
\]

Every numerator in \(\pi^L F_{\boldsymbol n}\) is an integral
polynomial. Before multiplication by the leading power in (V8-S3), the
evaluation matrix is integral, vanishes unless \(j_i\ge n_i\) for every
\(i\), and has a unit in each diagonal entry. Order the lower set by a
linear extension of the componentwise order. Its normalized matrix is
triangular and invertible over \(\mathcal O\). Thus the matrix of (V8-S3)
is integrally row-equivalent to the diagonal with entries
\(\pi^{(s(\boldsymbol n)-L)_+}\), all nonzero in \(R_p\). Its
\(R_p\)-image is a direct sum of \(|I|\) nonzero cyclic modules and
needs \(|I|\) generators by reduction of that image modulo its own
maximal-ideal multiple. Any submodule of an \(r\)-generated
\(R_p\)-module has at most \(r\) generators. To see this, take a
surjection from \(\mathcal O^r\) onto the module and pull the submodule
back to a free submodule of rank at most \(r\) over the DVR. The image
of a width-\(r\) matrix factorization therefore needs at most \(r\)
generators. This proves optimality over \(R_p\). The reduction of any
integral lift of \(K\)-rank \(r\) has width at most \(r\), giving the
claimed sharpness for lifts too. At \(p=1\) every diagonal entry in (V8-S3)
is a unit. \(\square\)

\textbf{Capacity on the actual message image.} Let
\(Z\subset\mathcal O^h\) be the nonempty finite set of distinct sender
tuples. Form the submodule \(\mathcal E_Z\subset\mathcal O^Z\) of
evaluations of integral polynomials at those tuples. This has full rank
over \(K\): for \(z\ne z'\) choose a differing coordinate, and multiply
the corresponding affine \(K\)-polynomials to interpolate each
individual point indicator. In an integral basis \(e_i\) of
\(\mathcal O^Z\), write
\(\mathcal E_Z=\bigoplus_i \pi^{s_i}\mathcal O e_i\) by the DVR Smith
theorem. The attainable integral columns at loss \(L\), reduced modulo
\(\pi^p\), form the module

\[
 \big((\pi^{-L}\mathcal E_Z\cap\mathcal O^Z)+\pi^p\mathcal O^Z\big)/
                   \pi^p\mathcal O^Z
 \ \cong\ \bigoplus_i \mathcal O/\pi^{[p-(s_i-L)_+]_+}\mathcal O.       \tag{V8-S4}
\]

Terms of length zero are zero modules. Intersection, reduction and the
choice of a realizing polynomial for each generator prove (V8-S4). Its
exact maximum separation rank with arbitrarily many receiver columns is
\(\#\{i:s_i<L+p\}\), by the same generator argument as in (V8-S3). In
finite data this algebra can be generated without enumerating circuit
paths. For each coordinate reduce polynomials by the monic product over
its distinct attained scalar values, then compute the Smith data of the
resulting finite evaluation matrix. This can be expensive and is only a
certificate calculation. Repeated, correlated or identical messages can
reduce the attainable rank. In particular \(u_2=u_1\) gives the
evaluation algebra of a single message.

For unrestricted scalars and \(p=1\) write \(H_q(h,L)\) for (V8-S2). Then

\[
 H_q(h,0)=q^h,\qquad
 H_q(h,L)\le q^h {h+L\choose h}\le [q(L+1)]^h.                      \tag{V8-S5}
\]

Write \(n_i=q t_i+r_i\), \(0\le r_i<q\). Since \(w_q(n_i)\ge t_i\), only
weak compositions with \(\sum_i t_i\le L\) can contribute. If \(L<q\),
every such \(n_i<q^2\) and \(w_q(n_i)=t_i\), giving \textbf{equality in
the first inequality of (V8-S5)}. For \(p>1\) use \(L+p-1\) in place of
\(L\). For \(q=2,h=1\) a loss of seven gives exactly ten independent
residue columns, indexed by \(0,\ldots,9\), while the scalar has 256
possible residues modulo \(2^8\). The width can still be attained with a
suitable polynomial numerator and one common division.

\subsection{Overlapping scopes and a common hard
output}\label{v8-overlapping-scopes-and-a-common-hard-output}

On the common v7 restriction, the \(m\) independent matched pairs carry
row bits \(A\) and column bits \(B\). After all other signs are set to
one, the target is \(D_\lor^{\otimes m}\) with
\(D_\lor=\left(\begin{smallmatrix}0&1\\1&1\end{smallmatrix}\right)\).
Its determinant is a unit. Consider a finite sum of blocks representing
the target modulo \(\pi\), where each block uses a source cut
\(P\mid Q\) of the free signs. The sender may read \(P\cup O_Q\) with
\(O_Q\subseteq Q\), and the receiver may read \(Q\cup O_P\) with
\(O_P\subseteq P\). After this private reading the receiver obtains only
the sender's \(h\) distinct scalars as additional inputs and has a
polynomial decoder (V8-S1). All other communication or routing used by the
block must be included in its views, coefficient functions, or messages.
Blocks may have different cuts, overlaps, losses and shared registers.

For one block, transpose the roles of matrix rows and columns in its
proof if necessary so that \(\delta=|A\cap P|-|B\cap P|\ge0\). Its two
\emph{computation} views keep their original direction. Let

\[
 \ell=|A\cap O_P|+|B\cap O_Q|,\qquad
 W=(B\cap P)\cup(A\cap Q)\cup(A\cap O_P)\cup(B\cap O_Q).
\]

The sets forming \(W\) are disjoint and \(|W|=m-\delta+\ell\). On a
fiber of these bits, the sender reads only free row bits and the
receiver only free column bits. Theorem 1 gives a congruent field
representative of rank at most \(H_q(h,L)\). Extend it by zero off this
fiber. The row and column masks preserve its rank. Sum the
representatives over the fibers, then over the blocks on the \emph{same}
oriented restriction. Their sum is integral with the original invertible
residue. This proves

\[
             1\ \le\ \sum_{b}
                  \min\{1,\ 2^{-\delta_b+\ell_b}H_q(h_b,L_b)\}.         \tag{V8-S6}
\]

The cap by one is valid when a block's fraction-field lift has size
\(2^m\) in each dimension. Even if the blocks have nonintegral outputs,
their congruent low-rank representatives can be added first. If sender
tuples are constrained to known congruence classes, use (V8-S2).
Conditioning works for correlated overlap data as long as the resulting
fibers are rectangles with the specified private views; counting raw
shared bits always gives the displayed estimate. For an adaptive route,
its guard must belong to an augmented block of the claimed kind, or have
a separate rank bound. An arbitrary joint test cannot enter the list of
private coefficients for free.
\PaperSixEvidence{evidence/research_v7/theory/hybrid_budget.md}{Common genuine output
and geometry} ·
\PaperSixEvidence{evidence/research_v6/overlap/final_separator.md}{Conditioning on
overlaps}

\textbf{Many possible cuts.} If there are \(t\ge1\) blocks, a common
margin \(\delta_b\ge d\), and uniform caps \(\ell,h,L\), (V8-S5)--(V8-S6) give

\[
 \log_2 t+\ell+h\log_2 q+\log_2{L+h\choose h}\ \ge\ d.               \tag{V8-S7}
\]

Use the route-specific \(d_M,d_D\) in (V8-S6) when available. On the
certified D layouts at the published schedule
\(d_D=\Omega(N/\sqrt{\log N})\). With \(\log q,\log(L+1)=O(\log N)\),
keeping \(\ell+h\log_2(q(L+1))=o(d_D)\) forces \(t\ge2^{(1-o(1))d_D}\).
In a fixed evaluator required to materialize all distinct final blocks
at rate at most \(W_0\) per paid unit of work, this gives
\(T\ge t/W_0\). Reuse inside a block is free in this count. If an
execution merges or bypasses the blocks, this materialization condition
must be demonstrated for its actual trace. No enumeration of monomials
or Newton indices appears in the cost argument.

\textbf{One fixed balanced source cut.} There is a stronger statement if
the entire procedure has a single one-way interface. Let \(G\) be one of
the simple cubic expanders and let \(P\mid Q\) be any fixed
\(\eta\)-balanced cut of its pin signs. The whole family has
\(C_G(\sigma)=\prod_{vw\in E}(\sigma_v\lor\sigma_w)\), including on the
four-clone normalization. Edge expansion gives at least \(\gamma\eta N\)
crossing edges. A greedy induced matching in the crossing edges keeps
\(cN\) of them for a constant \(c>0\) depending only on \(\eta,\gamma\):
choosing one edge discards crossing edges with an endpoint at distance
at most one from either chosen endpoint, a bounded number in a cubic
graph. Set all other signs to one and orient every chosen edge from
\(P\) to \(Q\). The target has rank \(2^{cN}\) up to integer rounding on
this cut. For \(h\) directed messages, zero overlap and correct residue,

\[
             q^h {L+h\choose h}\ \ge\ 2^{cN}.
                                                                    \tag{V8-S8}
\]

If \(q,L+1\le N^{O(1)}\), (V8-S8) implies \(h=\Omega(N/\log N)\). With at
most \(\ell\) hazardous overlap bits the right side can be divided by
\(2^\ell\), by the same fiber argument. Each cut in this one-block
statement gets its own matching. Equation (V8-S6) supplies the distinct
argument needed for sums whose blocks choose different cuts.

\textbf{Sharp width on the genuine family.} The source D residue field
has \(q=2^e\), \(e=4\lceil\log_2 s\rceil\), where \(s=\Theta(N)\) on
these fixed-density graphs. An arbitrary one-way redesign can pack \(e\)
sender signs per field element using an \(\mathbb F_2\)-basis
\(\theta_1,\ldots,\theta_e\). Lift that basis within \(\mathcal O\). Let
\(\beta_j\) be the corresponding trace-dual basis in \(\kappa\). For
\(u=\sum_j b_j\theta_j\), the receiver obtains each \(b_j\) in the
residue field using the integral polynomial lift of

\[
       \operatorname{Tr}_{\kappa/\mathbb F_2}(\beta_j u)
             =\sum_{a=0}^{e-1}(\beta_j u)^{2^a}.
\]

The dual basis exists because finite fields are separable. Fixed
coefficient lifts, additions and repeated squaring evaluate these
polynomials correctly modulo two. A residue answer can use precision one
throughout. For the full D precision \(P_0=n+1\), the sender supplies
integral lifts at that precision. Raise each trace value to an exponent
\(E\ge P_0\) divisible by \((q-1)q^{P_0-1}=|R_{P_0}^{\times}|\). Its
value modulo \(2^{P_0}\) is exactly the bit, since any unit has \(E\)-th
power one and a multiple of two has \(E\)-th power zero. One may take
\(E=P_0(q-1)q^{P_0-1}\), whose binary length is \(O(eP_0+\log P_0)\).
The receiver evaluates all edge clauses using \(x+y-xy\) for Boolean OR
and multiplies them. The source's final residue representative gives the
exact count, zero or one. Thus \(\lceil |P|/e\rceil=O(N/\log N)\)
messages suffice with \(L=0\) and polynomial bit work even for the
complete output on this family. Together with (V8-S8) this gives
\(\Theta(N/\log N)\) directed scalar width for each fixed balanced D cut
in this redesign class. The packing functions are chosen freely at this
interface. No equality with a specific printed child contraction is
used.

\textbf{Source precision.} At the D root with coefficient nilpotence
index one, the fixed interpolation among distinct residue elements uses
only unit denominators. In the subsequent tapered Hermite readout, the
frozen proof gives \(v_2(E_t(0))\ge-6h_{\rm H}D_t\), an integral
truncated reciprocal, and an evaluation loss at most \(4M(a+1)\). A
common loss bound from native scalar evaluations to the residue is
consequently

\[
       L_{\rm H}=6h_{\rm H}L_{\rm grid}+4M(a+1),\qquad L_{\rm grid}=2^a .
\]

Here \(h_{\rm H}=\lceil3M/u\rceil\), \(M=O(s)\),
\(L_{\rm grid}\le s^{1/1000}\), and \(a=O(\log s)\). Thus
\(L_{\rm H}<q\) for all sufficiently large instances and the closed form
in (V8-S5) applies to any interface with a denominator budget at most
\(L_{\rm H}\). Native samples have shell-dependent precisions. The
source proves each individual term's error bound before summing; a
common denominator exponent does not imply that every source sample was
supplied to that common precision. Local node data may be lifted within
their own view as prescribed by that proof. Additional division in a
redesign needs its own budget. Graph-only layouts describe a choice of
scopes; establishing a narrow directed interface in a concrete
implementation requires tracing its actual sender-to-receiver
dependencies. The sharp upper construction allows arbitrary local
functions on \(P\).
\href{https://github.com/openai/math/blob/main/preprints/The-Exponential-Time-Hypothesis-is-False-October-25-2026/build/quantitative-contraction.tex\#L10-L108}{Frozen
arithmetic and packing parameters, lines 10--108} ·
\href{https://github.com/openai/math/blob/main/preprints/The-Exponential-Time-Hypothesis-is-False-October-25-2026/build/quantitative-contraction.tex\#L244-L350}{Frozen
precision estimate, lines 244--350}

\subsection{Escapes and antecedents}\label{v8-escapes-and-antecedents}

One arbitrary lookup can defeat the denominator hierarchy. Encode an
\(m\)-bit sender input as its unsigned integer in \(\mathbb Z/2^m\). A
receiver permitted arbitrary equality testing against its \(m\)-bit
integer has the identity matrix of rank \(2^m\) after one message, with
\(O(m)\) bit work. For an integral polynomial receiver at loss zero and
\(q=2\), (V8-S5) instead gives exact maximum residue rank two from one such
message. The first procedure inspects higher digits through a joint
predicate. The finite-field trace construction above uses all
\emph{residue} coordinates of a large field; it is an allowed polynomial
receiver and attains the source order of width.

The precise loss needed for this packed-integer equality matrix is
\textbf{\(2^m-1-m\)} for \(m\ge1\), when the sender supplies the
ordinary integer representative \(x=0,\ldots,2^m-1\). In this
consecutive grid, the columns \(\prod_{i<j}(x-i)\) for \(0\le j<2^m\)
generate the evaluation lattice because the degree-\(2^m\) product
vanishes on it. Dividing each column by the \(2\)-part of \(j!\) gives
an integral triangular matrix with unit diagonal. Thus (V8-S4) has
exponents \(v_2(j!)\). To produce all columns of the identity modulo two
requires every exponent \(\le L\), and this condition also suffices
since the evaluation module then surjects onto the full residue space.
The largest exponent is \(v_2((2^m-1)!)=2^m-1-m\) by the floor sum. More
generally a single block of positive width \(h\) and hard rank \(2^d\)
needs \(L+1\ge 2^{d/h}/q\) by (V8-S5). In a dense implementation that
actually writes an entire numerator at precision \(L+1\) in the source
coordinate representation, that write comprises \(e(L+1)\) bits. A
factored algorithm, a sparse integer, or a different execution of a
polynomial identity requires an independent cost analysis. The formal
size of \(L\) alone supplies no per-run work lower bound.

The actual-image test (V8-S4) can detect a loss even for a \emph{single
column}, whose matrix rank is only one. The exact minimum loss to
extract the bit of weight \(2^j\) of a nonnegative integer \(x\), using
this integer as the sole scalar message, is \(2^j-1\). On the grid
\(0,\ldots,2^j\) the target column is zero except for its last entry. In
the basis \(\binom{x}{i}\), the evaluation matrix is unit triangular
with diagonal one. Hence the coordinate of the target in its last basis
column is one modulo two. That direction first survives when
\(L\ge v_2((2^j)!)=2^j-1\). Conversely \(\binom{x}{2^j}\bmod 2\) is the
requested bit for every \(x\ge0\): in characteristic two, factor
\((1+t)^x\) into the factors \(1+t^{2^a}\) corresponding to the one bits
of \(x\) and read the coefficient of \(t^{2^j}\). The binomial
polynomial has precisely the permitted denominator. The proof depends on
the specified scalar encoding; the trace packing gives a different,
zero-loss encoding.

The arithmetic factors in (V8-S2) are classical. Bhargava introduced
\(P\)-orderings and generalized factorials for polynomial functions on
subsets of Dedekind rings. Li and Sha prove a multivariate canonical
representation in Theorem 3.2 and counting formulas over residue rings
in Theorems 3.9 and 3.11. Their Theorem 3.11 contains the same sum of
floor weights. The use here is an exact scalar matrix capacity after a
bounded denominator and a terminal reduction, together with the directed
circuit and genuine source-cut consequences. A second exact
specialization for scalar readout of indivisible jet messages is in the
companion note.
\href{https://doi.org/10.1515/crll.1997.490.101}{Bhargava, 1997} ·
\href{https://arxiv.org/abs/1704.04965}{Li and Sha, 2019} ·
\hyperref[v8-section-jet-polynomial-capacity]{Jet readout calculation}

The checker reproduced five scalar Smith profiles from ordinary,
unnormalized monomial matrices over quotient rings at primes 2, 3 and 5.
It exhaustively generated receiver columns subject to divisibility on
four points in a shared residue fiber, tested duplicated register
values, and verified packed equality at one to six sender bits. Running
it requires Python 3 and NumPy.
\PaperSixEvidence{evidence/research_v8/theory/check_nonlinear_budget.py}{Reproducible checker} ·
\PaperSixEvidence{evidence/research_v8/theory/nonlinear_budget_checks.json}{Saved finite checks}

\section{Closed jet capacities at large residue
size}\label{v8-section-jet-polynomial-capacity}

Sending a whole jet need not give the receiver arbitrary access to its
coefficient vector. When the receiver can add and multiply jets and
apply a single terminal coefficient readout, the output matrix has a
smaller capacity than an arbitrary lookup on jet words. There is an
exact formula at positive finite precision when the residue field
exceeds the active polynomial degree. This includes the field and
polynomial delay sizes of the printed D root. The theorem describes
redesigned jet interfaces when they have the stated gates.

\subsection{The finite-precision
formula}\label{v8-the-finite-precision-formula}

Let \(\mathcal O\) be a DVR with uniformizer \(\pi\), fraction field
\(K\), finite residue field \(\kappa\) of cardinality \(q\) and
characteristic \(p_0\). Write \(A=\mathcal O[z]/z^k\), \(k\ge1\), and
\(R_p=\mathcal O/\pi^p\), \(p\ge1\). There are \(h\) distinct sender
messages \(U_i(x)\in A\). For each receiver input \(y\), let
\(C_y\in A[T_1,\ldots,T_h]\). Allow arbitrary products and repeated
references to the same message. With common loss \(L\ge0\), the tested
scalar is

\[
                  F(x,y)=\pi^{-L}[z^{k-1}]C_y(U(x)).
                                                                    \tag{V8-J1}
\]

The same result holds for any fixed scalar linear functional over
\(\mathcal O\) whose coefficients are integral as an upper bound: every
such functional can be expressed by multiplying by a fixed jet and
reading the top coefficient. The sharpness below concerns the top
readout. A single coefficient readout of order \(j<k\) reduces to the
same theorem at jet length \(j+1\). Intermediate extraction and
reinsertion of coefficients constitute a different gate set. All
formulas are about ranks on a sender/receiver cut; one may insert their
capacities into the overlap and multi-cut argument of the
\hyperref[v8-section-nonlinear-budget]{scalar theorem}.

Put \(P=L+p\) and \(D=P+k-2\). For \(\alpha\in\mathbb N^h\),
\(0\le j<k\), define

\[
 \mu(\alpha,j)=
 \min_{\substack{b_{ir}\ge0,\ \sum_{r=0}^{k-1} b_{ir}=\alpha_i\\
                         \sum_{i,r} r b_{ir}=j}}
 \left(\sum_i b_{i0}
      +\sum_i v_\pi\!\binom{\alpha_i}{b_{i0},\ldots,b_{i,k-1}}\right).
                                                                    \tag{V8-J2}
\]

A minimum over the empty set is infinite. Integer multinomials are
mapped to \(\mathcal O\), with \(v_\pi(0)=\infty\), so the definition
applies in either characteristic. Set

\[
              J_{q,\mathcal O}(h,k,L,p)
                =q^h\#\{(\alpha,j):0\le j<k,\ \mu(\alpha,j)<P\}.
                                                                    \tag{V8-J3}
\]

The set is finite because \(\mu(\alpha,j)\ge(|\alpha|-j)_+\), hence any
active \(\alpha\) has \(|\alpha|\le D\).

\textbf{Theorem 2.} There is a \(K\)-valued representative congruent to
(V8-J1) modulo \(\pi^p\) with ordinary matrix rank at most
\(J_{q,\mathcal O}\). For integral \(F\), the matrix over \(R_p\) has
separation rank at most \(J_{q,\mathcal O}\). If \(q>D\), the bound is
exact over all finite choices of sender inputs and receiver polynomials,
both for separation rank over \(R_p\) and as a bound on the least
\(K\)-rank of an integral representative. For \(p=1\) it gives the exact
maximal residue-field matrix rank. No bound on circuit degree, number of
multiplication gates, or fan-out is required.

\textbf{Upper bound.} Fix lifts \(a\in\mathcal O\) of each residue.
Partition the sender inputs by the constant residues \(a_i\) of all
\(h\) jets. In such a fiber write

\[
     U_i=a_i+\eta_i,\qquad
     \eta_i=\pi X_{i0}+\sum_{r=1}^{k-1}X_{ir}z^r,\quad X_{ir}\in\mathcal O.
\]

Expand \(C_y\) about the scalar vector \(a\) using the ordinary
multivariable translation, so all coefficients remain in \(A\). For each
monomial \(\eta^\alpha\) and each jet index \(j\), put
\(f_{\alpha j}=[z^j]\eta^\alpha\). In the \(X_{ir}\) its coefficient on
monomial \(\prod X_{ir}^{b_{ir}}\) is exactly
\(\pi^{\sum_i b_{i0}}\prod_i{\alpha_i\choose b_{i0},\ldots,b_{i,k-1}}\)
for the constraints in (V8-J2). Distinct arrays yield distinct monomials.
Hence every coefficient of \(f_{\alpha j}\) is divisible by
\(\pi^{\mu(\alpha,j)}\). After division by \(\pi^L\), terms with
\(\mu\ge P\) contribute zero modulo \(\pi^p\). The surviving features on
each fiber are indexed by (V8-J3). Each depends only on the sender; its
multiplier only on the receiver and the fixed fiber. Extend each feature
by zero off that fiber. The sum over all fibers has the stated
field-rank bound and the required congruence. For integral \(F\) use a
DVR factorization of this integral representative and then reduce its
factors.

\textbf{Attainment when \(q>D\).} For each active pair write
\(g_{\alpha j}=\pi^{-\mu(\alpha,j)} f_{\alpha j}\) as a coordinate
polynomial. It has a nonzero coefficient in the residue field. The
monomials of two different \(g_{\alpha j}\) have disjoint supports
because the row sums of their exponents recover \(\alpha\) and the
\(z\)-weighted sum recovers \(j\). Each variable degree is at most
\(|\alpha|\le D<q\). Evaluation on the entire grid \(\kappa^{hk}\) is
injective on polynomials having degree less than \(q\) in each variable.
This follows by repeated use of the univariate root bound. The
evaluations of the normalized \(g_{\alpha j}\) are therefore linearly
independent modulo \(\pi\) on each fiber. Choose the \(X_{ir}\) to be
the fixed residue representatives and select as many rows as there are
active features supporting a unit minor. This produces integral test
jets.

It remains to verify that one polynomial receiver per feature can select
that \emph{whole fiber}. For each residue representative \(a\), Hermite
CRT in \(\mathcal O[T]\) gives \(E_a(T)\) congruent to one modulo
\((T-a)^{P+k-1}\) and to zero modulo \((T-a')^{P+k-1}\) at every other
representative \(a'\). The ideals are pairwise comaximal since their
centers differ by units. If \(T\) is a jet with constant residue \(a'\),
then \(T-a'\in(\pi,z)\), and \((\pi,z)^{P+k-1}\subseteq\pi^P A\): every
surviving monomial has \(z\)-degree at most \(k-1\). The product
\(E_{\boldsymbol a}(U)=\prod_i E_{a_i}(U_i)\) is consequently a fiber
indicator modulo \(\pi^P A\). To select feature \((\alpha,j)\) for that
fiber, use

\[
 \pi^{-\min(\mu(\alpha,j),L)}
 [z^{k-1}]\,
 z^{k-1-j}E_{\boldsymbol a}(U)
              \prod_i(U_i-a_i)^{\alpha_i}.
                                                                    \tag{V8-J4}
\]

This is a decoder of the form (V8-J1), with an integral numerator after
multiplication by \(\pi^L\). On the chosen fiber (V8-J4) agrees modulo
\(\pi^p\) with \(\pi^{(\mu(\alpha,j)-L)_+}g_{\alpha j}\); on every other
fiber it is zero modulo \(\pi^p\). In particular (V8-J4) is everywhere
integral. Assemble the square submatrices on the chosen rows of all
fibers against these receiver columns. Over \(R_p\) the resulting matrix
is a block-diagonal invertible multiplier times the diagonal with
entries \(\pi^{(\mu(\alpha,j)-L)_+}\) in each fiber. All diagonal
entries survive modulo \(\pi^p\). Its image needs \(J_{q,\mathcal O}\)
generators, forcing the claimed inner width by the submodule property
over a chain ring. Reducing a factorization of any integral lift gives
the matching lower bound on its \(K\)-rank. \(\square\)

Unlike an arbitrary nonlinear function of a finite-ring element, the
decoder in (V8-J1) can only produce the indicated coordinate polynomials
before division. A coarser bound that avoids computing multinomial
valuations is

\[
 q^h\!\left[1+\sum_{d=1}^{P+k-2}{h+d-1\choose d}
                  \big(k-\max(0,d-P+1)\big)\right]\quad(h\ge1).       \tag{V8-J5}
\]

Indeed there are \({h+d-1\choose d}\) choices of \(\alpha\) of degree
\(d\) and \(j\ge\max(0,d-P+1)\) is necessary for survival. For
\(\alpha=0\) only \(j=0\) occurs. At \(k=1\), (V8-J2) reads
\(\mu(\alpha,0)=|\alpha|\); (V8-J3) becomes \(q^h{h+P-1\choose h}\),
exactly the scalar formula when \(P-1<q\). For a purely scalar interface
with larger \(P\), the factorial theorem supplies the full answer.

\subsection{An explicit expression at residue
precision}\label{v8-an-explicit-expression-at-residue-precision}

Take \(L=0,p=1\). For \(k\le q\), denote
\(g(\alpha)=p_0^{\min_{i:\alpha_i>0}v_{p_0}(\alpha_i)}\) for nonzero
\(\alpha\). The exact maximal residue rank in (V8-J3) is

\[
 q^h\left[1+\sum_{\substack{0<|\alpha|<k}}
       \left(1+\left\lfloor\frac{k-1-|\alpha|}{g(\alpha)}\right\rfloor\right)
                                                         \right].    \tag{V8-J6}
\]

After reduction, \(\eta_i=\sum_{r\ge1} X_{ir}z^r\). For \(\alpha\ne0\),
the product \(\eta^\alpha\) is a \(g(\alpha)\)-th power, starts in
degree \(|\alpha|\), and has coefficients only when \(j-|\alpha|\) is a
nonnegative multiple of \(g(\alpha)\). For the converse, take
\(j=|\alpha|+g(\alpha)t<k\). The leading coefficient at \(t=0\) is
nonzero. For \(t>0\), choose an \(i\) for which \(\alpha_i/g(\alpha)\)
is prime to \(p_0\). In \(\eta^{\alpha/g(\alpha)}\), select a single
\(X_{i,1+t}z^{1+t}\) and the lowest term in every remaining factor. This
gives a distinct monomial with coefficient \(\alpha_i/g(\alpha)\ne0\) in
the residue field; \(1+t<k\) follows from \(j<k\). Its Frobenius power
realizes degree \(j\). This proves (V8-J6).

For example, one jet over \(\mathbb F_4[z]/z^4\) gives rank exactly
\(4(1+3+1+1)=24\) for a top readout, while an unrestricted decoder of
the jet word could reach \(4^4=256\). For \(q<k\) the upper bound can be
strict because polynomial functions on residue digits obey \(X^q=X\).
With \(h=1,q=2,k=3,L=0,p=1\), each fiber has features
\(1,X_1,X_2,X_1^2\). Their span on binary digits has dimension three,
and the fiber selectors realize all of it. Thus the exact rank is six
instead of the eight from the uncorrected count. The hypothesis on field
size plays an essential role in the equality proof.

An independent finite calculation forms raw columns
\([z^j]\prod U_i^{n_i}\) across the indicated sender domain. It cuts
each \(n_i\) off at \(q(P+k-1)\), the degree of the monic null
polynomial \(\prod_{a}(T-a)^{P+k-1}\) over the finite jet ring. Smith
elimination on these raw evaluation matrices tests the
elementary-divisor profiles without building Taylor features or Hermite
selectors. For two length-two jets over \(\mathbb Z/9\), it finds 27
unit directions and 45 directions divisible by three across all 6,561
inputs. For one length-three jet over the quadratic unramified dyadic
ring at precision two, all 4,096 inputs give 16 unit and eight
first-order directions. The checker also verifies the rank-six
small-field exception and marks its one test on a proper
higher-coefficient digit grid.
\PaperSixEvidence{evidence/research_v8/theory/check_nonlinear_budget.py}{Reproducible checker} ·
\PaperSixEvidence{evidence/research_v8/theory/nonlinear_budget_checks.json}{Saved finite checks}

Classical work gives normal forms for polynomial functions over finite
residue rings and analyzes dual numbers by coordinate derivatives. The
calculation here projects to one scalar coefficient, permits several
reused jet messages, tracks an output denominator exponent, and gives a
rank witness. A redesigned MAIN calculation over the source field
\(\mathbb Q\) requires an integral DVR model and a verified loss bound
before the mixed-precision formula applies.
\href{https://arxiv.org/abs/1704.04965}{Li and Sha, Theorems 3.2, 3.9
and 3.11} · \href{https://arxiv.org/abs/1910.00238}{Al-Ezeh, Al-Maktry
and Frisch on dual numbers}

\section{Exact projected jet capacity over every residue
field}\label{v8-section-jet-newton-smith}

The projected Taylor coefficients of a jet polynomial have an exact
finite evaluation matrix over a discrete valuation ring. Its Smith
exponents give the optimal width for every residue-field size, including
when different Taylor coefficients define the same residue function. The
matrix uses the degree bound and the number of scalar jet coordinates.
It does not require a row for every possible finite-precision jet word.

\subsection{The matrix and its
interpretation}\label{v8-the-matrix-and-its-interpretation}

Let \(\mathcal O\) be a DVR with uniformizer \(\pi\), fraction field
\(K\) and residue field \(\kappa\) of size \(q\). Let
\(A=\mathcal O[z]/z^k\) for \(k\ge1\). A sender supplies \(h\ge0\) jets
\(U_i(x)\in A\). Each receiver input \(y\) determines an integral
polynomial \(C_y\in A[T_1,\ldots,T_h]\). For \(L\ge0\) and \(p\ge1\),
the required scalar output is congruent modulo \(\pi^p\mathcal O\) to

\[
                 \pi^{-L}[z^{k-1}] C_y(U_1(x),\ldots,U_h(x)).
                 \tag{V8-A1}
\]

Write \(P=L+p\), \(D=P+k-2\) and \(b=hk\). Fix representatives of the
constant residues of the jets. On one such residue fiber use independent
digit variables \(X_{ir}\) and write

\[
 U_i=a_i+\eta_i,\qquad
 \eta_i=\pi X_{i0}+\sum_{r=1}^{k-1}X_{ir}z^r,\qquad
 f_{\alpha j}(X)=[z^j]\prod_i\eta_i^{\alpha_i}.
\]

All the \(X_{ir}\) range independently over \(\mathcal O\) on the whole
fiber. The polynomial \(f_{\alpha j}\) is independent of the chosen
residues \(a_i\). Denote by \(\mu(\alpha,j)\) the minimum valuation of
its ordinary monomial coefficients, with \(\mu=\infty\) for the zero
polynomial. Take the column indices

\[
 \mathcal F=\{(\alpha,j):0\le j<k,\ \mu(\alpha,j)<P\}
\]

and the row indices \(I=\{\nu\in\mathbb N^b:|\nu|\le D\}\). To form an
explicit integer-valued basis, fix residue representatives
\(d_0=0,d_1,\ldots,d_{q-1}\). If \(n=\sum_s n_s q^s\), put
\(a(n)=\sum_s d_{n_s}\pi^s\). For each \(\nu\in I\) define

\[
 B_\nu(X)=\prod_{e=1}^b\prod_{r<\nu_e}(X_e-a(r)),\qquad
 t_\nu=\sum_{e=1}^b\sum_{u\ge1}\lfloor\nu_e/q^u\rfloor.
\]

Coordinate \(e\) here enumerates the pairs \((i,r)\). Expand each
feature in the \textbf{monic} Newton basis, \[
 f_{\alpha j}=\sum_{\nu\in I} C_{\nu,(\alpha,j)} B_\nu.
\] The matrix for the theorem is \[
                 M_{\nu,(\alpha,j)}=\pi^{t_\nu} C_{\nu,(\alpha,j)}.
                 \tag{V8-A2}
\]

Let \(d_1,\ldots,d_s\) be its finite DVR Smith exponents, with
multiplicity. Exponents at least \(P\) may be discarded. For
\(R_p=\mathcal O/\pi^p\), separation rank means the least \(r\) such
that the output matrix factors through \(R_p^r\).

\textbf{Theorem.} For any finite sender and receiver input sets, (V8-A1)
has a congruent \(K\)-valued matrix of ordinary rank at most

\[
                        q^h\,\#\{i:d_i<P\}.                       \tag{V8-A3}
\]

The same number is the exact largest separation rank of an integral
output over \(R_p\), and the exact largest minimum \(K\)-rank of an
integral lift. On the full jet domain and on a finite witnessing grid,
the attainable integral output columns modulo \(\pi^p\) form a module
isomorphic to

\[
 \bigoplus_{\boldsymbol a\in\kappa^h}
 \ \bigoplus_{i:\,d_i<P}\mathcal O/\pi^{\,p-(d_i-L)_+}\mathcal O .
                                                                    \tag{V8-A4}
\]

For a prescribed proper sender image, the rank bound in (V8-A3) still
holds. Its actual attainable module can depend on additional
coincidences there.

\subsection{Proof}\label{v8-proof}

\textbf{1. Active features.} Translate each \(T_i\) by \(a_i\). The
coefficient of a translated jet monomial is an arbitrary element of
\(A\). Its contribution to the top readout is an integral scalar
combination of the \(f_{\alpha j}\). An ordinary monomial of
\(f_{\alpha j}\) uses \(|\alpha|\) factors, at most \(j\) with positive
\(z\)-degree. Its coefficient is therefore divisible by
\(\pi^{(|\alpha|-j)_+}\). A feature with \(|\alpha|>D\) vanishes
coefficientwise modulo \(\pi^P\), as does every omitted feature with
\(\mu\ge P\). After division by \(\pi^L\), they vanish modulo \(\pi^p\).
Every remaining feature has total degree at most \(D\) in its digit
variables. The monic coordinatewise changes of basis only lower
exponents, which proves that the indicated expansions use precisely an
allowed finite set of Newton basis vectors. The argument permits
multinomial coefficients that are zero in equal characteristic.

\textbf{2. Evaluation and divisibility.} For a coordinate
\(x\in\mathcal O\), the number of \(a(r)\), \(0\le r<n\), congruent to
\(x\) modulo \(\pi^u\) is at least \(\lfloor n/q^u\rfloor\). Equality
holds for \(x=a(n)\), by the ordinary base-\(q\) digit ordering. Summing
over \(u\) proves that \(G_\nu=B_\nu/\pi^{t_\nu}\) takes integral values
everywhere. Evaluate these \(G_\nu\) on the grid with points
\(x(\lambda)=(a(\lambda_1),\ldots,a(\lambda_b))\), \(\lambda\in I\). The
entry for column \(\nu\) vanishes unless \(\lambda_e\ge\nu_e\) for every
\(e\), and its value at \(\lambda=\nu\) is a unit. The square grid
evaluation matrix \(V\) is triangular with unit diagonal entries in any
linear extension of the coordinate order, hence belongs to
\({\rm GL}_{|I|}(\mathcal O)\). Every evaluation row elsewhere is an
integral linear combination of grid rows, obtained by multiplying its
\(G_\nu\)-row by \(V^{-1}\). In particular a linear combination of the
\(G_\nu\) is divisible pointwise by \(\pi^m\) exactly when its Newton
coefficient vector in this normalized basis is divisible by \(\pi^m\).
This holds over the whole fiber.

The grid evaluation matrix of the active Taylor features equals \(VM\).
Thus (V8-A2) gives their \textbf{evaluation} Smith exponents. This step
also proves invariance of all exponents below \(P\) under changes of the
digit representatives, since they describe the actual fiber evaluation
module at precision \(P\).

\textbf{3. Globally realizable fiber selections.} Let \(n=P+k-1\). For
each residue representative \(a\), the ideals \((T-a')^n\) for the
different representatives \(a'\) are pairwise comaximal in
\(\mathcal O[T]\). The Chinese remainder theorem supplies \(E_a(T)\)
equal to one modulo \((T-a)^n\) and zero modulo the other powers. At a
jet of constant residue \(a'\) one has \(T-a'\in(\pi,z)\). Every
surviving term of \((\pi,z)^n\) has \(z\)-degree at most \(k-1\) and
\(\pi\)-valuation at least \(P\). Hence \(E_a(T)\) is the residue
indicator modulo \(\pi^P A\). Their product over the \(h\) coordinates
selects the desired whole fiber. Multiplying this product by
\(z^{k-1-j}\prod_i(T_i-a_i)^{\alpha_i}\) realizes any one active feature
there and is zero on all other fibers modulo \(\pi^P\). Integral linear
combinations realize each fiber column independently. This proves
equality of the global raw column module modulo \(\pi^P\) with the
product of the separate raw modules. Selectors only need to be exact to
this precision.

\textbf{4. Denominators and rank.} In an integral basis of the
normalized Newton coefficient module \(\mathcal O^I\), the raw image of
\(M\) has nonzero coordinate ideals \(\pi^{d_i}\mathcal O\). Thus the
numerators divisible by \(\pi^L\), after division and reduction modulo
\(\pi^p\), have ideals \(\pi^{(d_i-L)_+}R_p\). A direction with
\(d_i\ge P\) disappears. Altering any raw numerator by \(\pi^P\) does
not affect its divisibility by \(\pi^L\) or its divided value modulo
\(\pi^p\). The grid divisibility equivalence and fiber selection
therefore prove (V8-A4), including the case where integrality requires
cancellation between active Taylor features. This use of saturation is
important when the normalized Taylor monomials are dependent as residue
functions.

For an arbitrary finite set of rows and a possibly nonintegral output,
take each fiber's DVR Smith expression and discard its terms with
\(d_i\ge P\) before division. Their evaluations, divided by \(\pi^L\),
belong to \(\pi^p\mathcal O\) at every point. Extend each retained
function by zero off its fiber. The resulting \(K\)-matrix has at most
the number of separated terms in (V8-A3). Pointwise integrality of the
target makes the reconstructed sum integral on the given rows, even if
individual terms have denominators. An integral matrix of \(K\)-rank
\(r\) has a factorization through \(\mathcal O^r\) by DVR Smith form, so
reduction gives width at most \(r\).

For attainment choose all the grid rows on all residues and columns
realizing generators of (V8-A4). Each displayed nonzero summand needs one
generator, as seen by quotienting the module by its own
\(\pi\)-multiple. Over \(R_p\), every submodule of a module generated by
\(r\) elements needs at most \(r\) elements. Indeed, pull the submodule
back under a surjection from \(\mathcal O^r\) and use the DVR submodule
theorem. The image of any matrix factoring through \(R_p^r\)
consequently needs at most \(r\) generators. This gives the matching
lower bound on width. Every integral lift of \(K\)-rank \(r\) reduces to
a factorization of width at most \(r\), so the same witness gives
sharpness for lifts. The empty product convention covers \(h=0\).
\(\square\)

\subsection{Specializations}\label{v8-specializations}

If \(q>D\), then all \(t_\nu=0\). In the ordinary coefficient basis
different \(f_{\alpha j}\) have disjoint monomial supports, since the
exponents in any monomial recover both the row sums \(\alpha_i\) and its
\(z\)-degree \(j\). Dividing each column by its individual
\(\pi^{\mu(\alpha,j)}\) gives columns independent modulo \(\pi\).
Changing from monomials to the monic Newton basis is invertible over
\(\mathcal O\). Hence the active Smith exponents are precisely the
individual minima \(\mu(\alpha,j)\). This recovers the earlier closed
expression, which avoids a Smith calculation at large residue size.
\hyperref[v8-section-jet-polynomial-capacity]{Large-field formula and
residue specialization}

If \(L=0,\ p=1\), an alternative calculation works entirely over
\(\kappa\). In every active \(f_{\alpha j}\) put \(\pi=0\) and reduce it
modulo \(X_e^q-X_e\) in every digit variable. The resulting polynomials
have degree below \(q\) in each variable. Their coefficient-column rank
is exactly the rank of their evaluations, by interpolation on the digit
grid. Multiplication by \(q^h\) gives the universal residue rank. Terms
becoming equal must be added in \(\kappa\) before taking rank.

For a finite prescribed image \(Z\subset A^h\), use the ordinary DVR
column module of all functions
\(([z^j]\prod_i U_i^{\alpha_i})_{U\in Z}\), for \(0\le j<k\). If \(S_i\)
is the set of the exact constant coefficients of the attained \(U_i\),
the monic polynomial \(\prod_{c\in S_i}(T_i-c)^k\) vanishes identically
at coordinate \(i\) on \(Z\). Reduction in each variable therefore gives
finitely many columns. Their DVR exponents \(s_i\) replace the
collection of all per-fiber \(d_i\) in (V8-A3)--(V8-A4), omitting any free
ambient direction with zero raw image. This follows by intersecting the
finite evaluation image after division with \(\mathcal O^Z\) and
reducing. Unlike scalar evaluations at distinct points, these projected
jet evaluations can have deficient fraction-field rank. For three jets
\(U=zb\) with distinct \(b\) and \(k=2\), every such scalar output is
affine in \(b\).

\subsection{A positive-loss example by
hand}\label{v8-a-positive-loss-example-by-hand}

Take \(\mathcal O=\mathbb Z_2,\ h=1,\ k=3,\ L=1,\ p=2\). The raw
precision is \(P=3\). On either residue fiber put \(\eta=2x+yz+wz^2\),
and write \(f_{nj}=[z^j]\eta^n\). The coefficient-minimum test retains
nine functions per fiber. Modulo eight,

\[
 f_{20}=2f_{10},\qquad f_{31}=f_{21}=2f_{32}.
\]

The first congruence uses \(x^2\equiv x\bmod2\). The next two use
\(y^2\equiv y\bmod2\) as well. The surviving span is generated by

\[
        1,\quad y,\quad w,\quad 2x,\quad
        y^2+4xw-y,\quad 12x^2w+6xy^2 .
\]

The final three functions have a pointwise factor of two. Divide those
three by two and reduce all six functions modulo two. The result is \[
       1,\quad \bar y,\quad\bar w,\quad\bar x,\quad
       \binom{y}{2}\bmod2,\quad\bar x\bar y .
\] The penultimate function is the second binary digit of \(y\). That
digit varies independently of \(\bar x,\bar y,\bar w\), and the other
five functions are distinct squarefree monomials of these three bits.
Thus the six normalized evaluation columns have a unit minor. Their raw
exponents modulo eight are \(0,0,0,1,1,1\) per fiber. After the
permitted division each fiber supplies six free \(\mathbb Z/4\)
directions. The exact width is \textbf{12}, compared with the
coefficient-minimum upper bound \textbf{18}. The pointwise even
expression \(y^2-y\) explains why an analysis of ordinary coefficients
alone misses cancellations at small residue size.

A separate raw computation enumerates all 512 jets over \(\mathbb Z/8\)
in this example. It forms the 30 columns consisting of all coefficients
of \(T^n\), for \(0\le n<10\). The monic polynomial
\(\prod_{a=0}^1(T-a)^{P+k-1}\) vanishes modulo \(\pi^P\) on these jets,
so these columns span all raw outputs. Their measured orders are six
zeros and six ones. The same program compares the raw and
weighted-Newton matrices on additional complete input domains, including
a binary length-five jet with \(L=p=1\), where the exact width is \(22\)
and the coefficient bound is \(28\). For integer digit nodes it uses
\(X^b=\sum_j {b\brace j}(X)_j\) in the monic basis and the row
multiplier \(j!\), a unit times \(\pi^{w_q(j)}\) at prime \(q\).
\PaperSixEvidence{evidence/research_v8/checks/jet_newton_smith.py}{Checker} ·
\PaperSixEvidence{evidence/research_v8/checks/jet_newton_smith.json}{Exact outputs}

The product Newton basis and factorial ideals are classical. Bhargava
developed \(P\)-orderings and polynomial-function normal forms, and Li
and Sha gave multivariate residue-ring expansions and counting formulas.
Al-Ezeh, Al-Maktry and Frisch analyze polynomial functions and
derivatives on dual numbers. Here these ingredients apply to the
submodule formed by a single jet-coefficient readout with several
polynomial arguments and bounded division. The calculation determines a
matrix capacity. Its construction alone gives no runtime bound or
priority claim.
\href{https://doi.org/10.1515/crll.1997.490.101}{Bhargava} ·
\href{https://doi.org/10.1142/S1793042119500854}{Li and Sha} ·
\href{https://doi.org/10.1515/ms-2021-0039}{Al-Ezeh, Al-Maktry and
Frisch}

\section{Arbitrary decoding on a full linear
domain}\label{v8-section-duality-typeset}

The polynomial receiver has a complementary unrestricted result when the
target and observations are linear on the whole input module. Let
\(E:M\to U\) and \(T:M\to V\) be module homomorphisms. An arbitrary set
function \(D:U\to V\) satisfying \(D(E x)=T x\) for every \(x\) exists
if and only if \(\ker E\subseteq\ker T\). Necessity compares a kernel
vector with zero. Conversely, define \(D(E x)=T x\) on the image. If
\(E x=E y\), then \(x-y\in\ker E\) and \(T x=T y\), so the definition is
single valued. Extend by zero outside the image. Its restriction to the
image is linear. A linear extension to all of \(U\) needs a separate
hypothesis.

\subsection{Finite Frobenius duality}\label{v8-finite-frobenius-duality}

Let \(R\) be a finite commutative ring with a generating additive
character \(\chi\), meaning that its kernel contains no nonzero ideal.
For \(E\in R^{m\times n}\) and \(T\in R^{s\times n}\), the following
conditions are equivalent: an arbitrary decoder recovers \(Tx\) from
\(Ex\) on all \(R^n\); \(\ker E\subseteq\ker T\); every row of \(T\)
belongs to the row module of \(E\); and \(T=CE\) for some
\(C\in R^{s\times m}\).

For a submodule \(S\subseteq R^n\), define \[
 S^\perp=\{x:\chi(\sum_i a_i x_i)=1\text{ for all }a\in S\}.
\] The pairing is nondegenerate: a nonzero coordinate generates a
nonzero ideal, and some element of that ideal has character different
from one. A translation by a vector with nontrivial character shows that
a sum of characters over a finite group is zero unless the character is
trivial. Summing \(\chi(a\cdot x)\) over \(a\in S\) and \(x\in R^n\) in
the two orders gives \(|S||S^\perp|=|R|^n\). Inclusion and cardinality
then give \(S^{\perp\perp}=S\). For \(S=\operatorname{Row}(E)\),
orthogonality to all multiples of a row means the ideal generated by its
dot product lies in \(\ker\chi\). Hence \(S^\perp=\ker E\). Kernel
inclusion therefore puts every target row in \(S^{\perp\perp}=S\),
giving its row coefficients in \(C\). The remaining implications are
immediate. No Smith decomposition over \(R\) is used.

\subsection{Double truncation and heterogeneous
precision}\label{v8-double-truncation-and-heterogeneous-precision}

For \(R=(\mathbb Z/p^P)[z]/z^K\), the character \[
 \chi(f)=\exp\left(2\pi i\,[z^{K-1}]f/p^P\right)
\] is generating. Multiply a nonzero coefficient \(f_j\) by
\(z^{K-1-j}\) to place it in the top coefficient. For an unramified
coefficient extension use its trace before this character. The trace
pairing is perfect: modulo the prime it is the nondegenerate
finite-field trace pairing, so its determinant is a unit at every finite
precision. For completeness the finite-field trace is a nonzero
polynomial \(X+X^p+\cdots+X^{p^{e-1}}\) of degree less than \(p^e\).
Multiplication by a nonzero field element is bijective, proving the
nondegeneracy. The source's localized unramified order and its
completion have the same finite quotients; the finite-ring argument
applies to both.

Suppose row \(E_jx\) is observed modulo \((\pi^{P_j},z^{K_j})\), while
row \(T_ix\) is needed modulo \((\pi^{p_i},z^{b_i})\), with positive
exponents bounded by \(P,K\). Set \[
 \widehat E_j=\pi^{P-P_j}z^{K-K_j}E_j,\qquad
 \widehat T_i=\pi^{P-p_i}z^{K-b_i}T_i.
\] The first multiplier has kernel exactly \((\pi^{P_j},z^{K_j})\):
below degree \(K_j\) its coefficients must be divisible by
\(\pi^{P_j}\), and all higher degrees are annihilated. It embeds the
observation quotient into \(R\) as a module, without being claimed to be
a unital ring embedding. The target embeds in the same way. Thus
arbitrary recovery is equivalent to \[
 \operatorname{Row}_R(\widehat T)
       \subseteq\operatorname{Row}_R(\widehat E). \tag{D1}
\] If coefficients range in a smaller base-ring module, flatten their
coordinates and use that domain. One cannot add arbitrary jet-valued
inputs to a correctness promise.

\subsection{Adaptive acquisition}\label{v8-adaptive-acquisition}

Fix a catalog of linear queries with actual precision and nonnegative
acquisition costs. A deterministic algorithm may choose later queries
from the answers, cache results and use arbitrary exact decoding. All
access to the unknown input must be included in this catalog. Run it on
zero. Every input in the common kernel of the queries on that path gives
the same zero answers and follows the same path; correctness implies its
target is zero. The fixed zero-path list therefore already satisfies the
kernel certificate, and (D1) when applicable. The worst-case acquisition
cost is at least the minimum cost of a fixed catalog list satisfying
that certificate. A repeated acquisition can be charged once; a joint
query requires a catalogued implementation cost. The result is an
acquisition statement and does not make a decoder's execution free in a
time upper bound.

Over a field, to recover \(f(a)\) for every polynomial of degree at most
\(d\) by jets at points \(t\ne a\), let \(k_t\) be the number of
consecutive Taylor coefficients observed at \(t\) on the zero path. If
\(\sum_t k_t\le d\), then \(g(X)=\prod_t(X-t)^{k_t}\) gives all those
zero jets but has \(g(a)\ne0\). Thus \(\sum_t k_t\ge d+1\). Taylor
coefficients, equivalently Hasse derivatives, avoid
characteristic-dependent factorial issues. At finite precision the row
criterion must replace this bare degree argument because nearby residues
can already determine a target.

\subsection{Necessary boundaries}\label{v8-necessary-boundaries}

In \(R=\mathbb F_2[x,y]/(x^2,xy,y^2)\), observe \(xa\) and request
\(ya\). Their kernels both equal \((x,y)\), so a set decoder sends \(0\)
to \(0\) and \(x\) to \(y\). But no scalar \(c\) has \(cx=y\); there is
no linear decoder on all of \(R\). This ring does not satisfy the
Frobenius hypothesis.

On the rational promise \(v=u^2\), observing \(u\) recovers \(v\) by
squaring even though no linear map from \(\mathbb Q\) gives all those
values. The observation and target are restrictions of linear maps on
\(\mathbb Q^2\), but the promise is not the full module. More generally,
on a set \(X\) the exact collision test is
\((X-X)\cap\ker E\subseteq\ker T\). Neither nonlinear native CNF images
nor their differences may be replaced by a larger linear space without a
proof. Generating-character duality is classical; its use here specifies
the decoder and precision interface.

\section{Physical switched CNFs and positive
computation}\label{v8-section-stronger-family}

The optional vertices below give bounded-occurrence CNFs for both source
algorithms. Contracting each vertex's equality cycle and availability
clause yields local tables with entries \(0,1,s_v\). Computing the exact
count from those entries by a fixed nonnegative sum-product circuit
requires exponentially many scalar gates on bounded-degree expanders.
The theorem permits overlapping inputs to products and unrestricted
fan-out. On another restriction of the same clause scopes, the decision
bit and count each have a full-rank dyadic matching matrix. The raw
availability clause also exposes \(1-s_v\) at a different table address.
The choice of local interface affects which lower bound applies.

\subsection{F1. The common formula}\label{v8-f1.-the-common-formula}

Let \(G=(V,E)\) be a simple graph with \(N=|V|\), fixed edge
orientations, and minimum degree at least one. Optionally choose a
collection \(\mathcal D\) of vertex-disjoint triples of distinct
vertices, with a fixed order within each triple. Write
\(L_1(u)=u,\ L_0(u)=\neg u\). The input data consist of independent
availability bits \(s_v\), edge-request bits \(t_e\), and triple-request
bits \(r_D\). The following use \emph{logical} names to describe the
occurrences to be connected. First allocate a different physical Boolean
variable for every displayed occurrence, including two displayed
occurrences with the same logical name. Only then insert the literal
signs. Connect all occurrences of each logical name by a directed
implication cycle, using every copy exactly once on each side of the
implications.

\begin{itemize}
\tightlist
\item
  For vertex \(v\) use the switch clause \(\neg x_v\vee L_{s_v}(x_v)\),
  with its two distinct physical variables.
\item
  For edge \(e=u\to v\) use the conflict clause \(\neg x_u\vee\neg x_v\)
  and the request clause \(x_u\vee x_v\vee L_{t_e}(x_u)\).
\item
  For \(D=(i,j,k)\in\mathcal D\), introduce a logical wire \(y_D\). Use
  \((\neg x_i\vee y_D),\,(\neg x_j\vee y_D),\,(x_i\vee x_j\vee\neg y_D)\),
  and \((y_D\vee x_k\vee L_{r_D}(x_k))\).
\item
  Include one detached fixed clause \(b\vee c\) on two fresh, un-cloned
  variables. It supplies an odd factor of three.
\end{itemize}

Within each cycle make non-neighbors of every pair of copies of that
logical name belonging to the same original clause. Such pairs form a
matching on the slots. The cycles of the full construction have length
at least four. With one forbidden pair, put other slots in both
intervening arcs. With \(j\ge2\) forbidden pairs, the order
\(a_1,\ldots,a_j,b_1,\ldots,b_j\), with any unpaired slots inserted
anywhere, separates every pair \(a_i,b_i\). For zero pairs any order
works. Thus an equality clause cannot be contained in an original
clause, including an inactive width-three request. Original clauses have
mutually disjoint physical occurrences. Distinct equality clauses have
distinct supports. Consequently there are no units, tautologies,
duplicate clauses, repeated literals within a clause, or subsumptions
for any choice of data. In particular, normalization leaves the entire
clause list in place. Each private copy occurs in its original clause
and in exactly two implication clauses. The sparsifier's only possible
thresholds on nonempty proper subsets eventually exceed three, so the
formula also reaches a single unchanged sparsification leaf. The general
rational parity front end can still create sign-dependent equations in
later steps.
\href{https://github.com/openai/math/blob/main/preprints/The-Exponential-Time-Hypothesis-is-False-October-25-2026/build/main_full_version.tex\#L3647}{MAIN
lines 1465--1485, 3576--3603, 3647--3675}

\textbf{Semantics.} Implication around a directed cycle of Boolean
values forces them all equal. Conversely each logical choice extends
uniquely on its cycle. Each \(y_D\) is uniquely the disjunction of
\(x_i,x_j\). Consequently the exact count is three times the number of
sets \(I\) satisfying \[
 I\in{\rm IS}(G),\qquad I\subseteq\{v:s_v=1\},\qquad
 I\cap e\ne\varnothing\ (t_e=1),\qquad I\cap D\ne\varnothing\ (r_D=1).
 \tag{F1}
\] Here \({\rm IS}(G)\) is the family of independent sets. Setting an
edge or triple request to zero makes its original clause automatically
true after cycle equality. These identities prove (F1) on every input.
Every satisfiable instance has at least three models. The decision bit
varies, including on a single fixed topology with at least one edge.

Put \(K=2N+5|E|+10|\mathcal D|\). There are exactly \(K+2\) declared
variables, \(K\) equality clauses and \(N+2|E|+4|\mathcal D|+1\) other
clauses. The copy cycles account for \(2K\) incidences, so \[
 M_{\rm MAIN}=3K+2,\qquad M_{\rm D}=2M_{\rm MAIN}.
 \tag{F2}
\] For D every private degree-three variable needs one trivalent
equality tensor and a degree-one variable needs one unary all-ones
tensor. There are no internal edges in those attachments. Each of the
original clause incidences gives one pre-subdivision wire, and
subdivision doubles the wire count. Every variable occurs. The D
construction contracts to (F1) at precision \(P_0=K+3\). The root is the
source's sparse clause branch for all sufficiently large bounded-degree
instances. MAIN's direct factor sum is the same exact count. This
applies to the two direct entry points and to the combined program's
sparse dispatch.
\href{https://github.com/openai/math/blob/main/preprints/The-Exponential-Time-Hypothesis-is-False-October-25-2026/build/main_full_version.tex\#L1675}{MAIN
lines 1675--1704} ·
\href{https://github.com/openai/math/blob/main/preprints/The-Exponential-Time-Hypothesis-is-False-October-25-2026/build/quantitative-reduction.tex\#L65}{D
reduction lines 65--97, 362--369}

A smaller count-only version drops every edge request and every triple
block. Use this construction on graphs of minimum degree at least two.
Each vertex then has at least four copies, so the two switch copies can
be made nonadjacent. Its number of cycle copies is \(K_0=2N+2|E|\). On
cubic graphs these are precisely five copies per vertex, with the three
graph-edge ports and two switch ports each at their own copy. It has
\(5N+2\) declared variables, \(15N/2+1\) clauses, \(15N+2\) MAIN
incidences, and \(30N+4\) physical D wires. Formula (F1) then equals
three times the number of independent sets of the enabled induced
subgraph. The version with requests has both this counting slice and the
decision slice below.

\subsection{F2. The terminal restriction and its source
labels}\label{v8-f2.-the-terminal-restriction-and-its-source-labels}

For graph-only source layouts, choose the induced matching
\(\mathcal M\) and its row/column orientation by the geometric matching
argument before fixing any requests. Turn off every triple request,
activate the edge requests precisely for \(\mathcal M\), and disable all
availability switches outside its endpoints. Each edge
\(a_ib_i\in\mathcal M\) must supply exactly one selected endpoint. Since
the matching is induced, the choices on different edges are independent.
The resulting count and decision are \[
 3\prod_{i=1}^{m}(s_{a_i}+s_{b_i}),\qquad
 \prod_{i=1}^{m}(s_{a_i}\lor s_{b_i}).
 \tag{F3}
\] The per-edge matrices, with input order \(0,1\), are respectively
\(\bigl(\begin{smallmatrix}0&1\\1&2\end{smallmatrix}\bigr)\) and
\(\bigl(\begin{smallmatrix}0&1\\1&1\end{smallmatrix}\bigr)\). Both have
determinant \(-1\). The unscaled tensor products are unimodular. The
physical decision matrix has rank \(2^m\) over every field. The physical
count matrix includes the scalar three, so its rank is \(2^m\) over
fields of characteristic different from three and zero in characteristic
three. Both physical matrices are invertible over
\(\mathbb Z/2^p\mathbb Z\) for every \(p\ge1\). For the count-only slice
with every request off and every outside vertex disabled the answer is
\(3\prod_i(1+s_{a_i}+s_{b_i})\). The transfer block
\(\bigl(\begin{smallmatrix}1&2\\2&3\end{smallmatrix}\bigr)\) is
unimodular, giving the same dependence on the characteristic for the
physical count. The decision on this latter slice is always true. The
source computes exact counts. An arbitrary zero test applied to a count
is a separate, potentially full-rank operation. Equation (F3) directly
tests a circuit whose represented terminal object is the decision bit.

De Colnet and Mengel's Lemma 14 already derives an exponential
value-matrix rank from an induced matching for a graph-indexed
Boolean-input function. Here the three displayed blocks come from
counting or deciding the physical switched CNF. Their unimodularity also
gives the required dyadic statements.
\href{https://proceedings.kr.org/2021/20/kr2021-0020-de-colnet-et-al.pdf}{de
Colnet and Mengel, Lemma 14}

Each \emph{availability} switch has exactly two clause ports sharing a
factor label in MAIN and two physical old-side ports sharing a vertex
label in D. Request bits have already been fixed on the matching
restriction, so their owning factors contribute no extra dependence. For
MAIN, summing the certified cell upper bound \(2M_{\rm MAIN}/r_M^{2Q}\)
over variable labels and unused coordinates gives availability-label
density at most \(M_{\rm MAIN}/N\) relative to product uniform. In D
write \(D_{\rm arr}\) for the full number of ordered label cells. The
actual source bound per cell is
\(M_{\rm D}/D_{\rm arr}+\sqrt{2M_{\rm D}\Delta}\). Since
\(D_{\rm arr}\le s^{1/1000},\ \Delta=\lceil s^{9/10}\rceil,\ M_{\rm D}=\Theta(s)\),
it is \((1+o(1))M_{\rm D}/D_{\rm arr}\) uniformly in the cell. Summing
over the opposite endpoint and unused coordinates and dividing by the
two ports per labelled switch gives density at most
\((1+o(1))M_{\rm D}/(2N)=(1+o(1))M_{\rm MAIN}/N\).
\href{https://github.com/openai/math/blob/main/preprints/The-Exponential-Time-Hypothesis-is-False-October-25-2026/build/main_full_version.tex\#L2423}{MAIN
lines 2423--2458} ·
\href{https://github.com/openai/math/blob/main/preprints/The-Exponential-Time-Hypothesis-is-False-October-25-2026/build/quantitative-contraction.tex\#L132}{D
contraction lines 73--109, 132--155}

For cubic \(G\) and \(\mathcal D=\varnothing\),
\(M_{\rm MAIN}/N=57/2+2/N\) with edge requests and \(15+2/N\) without
them. For graphs of maximum degree four and \(|\mathcal D|\le N/3\) the
corresponding bound with all requests is \(46+2/N\). Thus at the printed
\(r_D=64\) each D symbol has mass below \(3/4\) for all sufficiently
large sizes, while \(r_M\to\infty\). At each used coordinate, take a
symbol of mass in \([1/4,3/4]\) if one exists, or accumulate smaller
symbols until reaching \(1/4\). This supplies a quarter-balanced union.
On cubic expander graphs the v7 simultaneous-cut proof, under its stated
size and menu conditions, therefore applies directly to these
\textbf{actual signed binary factors} and to the count in (F3). Its
geometric matching is chosen using only \(G\) and the graph-only source
labels of the \(s_v\). The terminal restriction to
\(t_e=1_{e\in\mathcal M}\) is made afterward on that same graph.
Arbitrary dependence of a layout or a circuit on the request values
before this choice would change the quantifier order.
\PaperSixEvidence{evidence/research_v7/theory/hybrid_budget.md}{Weighted matching and
guarded rank theorem}

\subsection{F3. A positive scalar circuit theorem with overlapping
products}\label{v8-f3.-a-positive-scalar-circuit-theorem-with-overlapping-products}

Assume \(G\) has maximum degree \(\Delta\), minimum degree at least one,
and fixed edge expansion \(h>0\), meaning
\(|E(U,V\setminus U)|\ge h\min(|U|,N-|U|)\) for all \(U\subseteq V\).
Turn off the requests and consider the Boolean availability word. The
count is \[
 F(s)=3\sum_{I\in{\rm IS}(G)}\prod_{v\in I}s_v.
 \tag{F4}
\] Consider any acyclic scalar circuit with binary addition and
multiplication over nonnegative reals. Each input atom must have a
multilinear polynomial on the Boolean cube with nonnegative
coefficients; every one of its monomials has fewer than \(\ell\)
distinct variables. This allows arbitrary nonnegative constants and
arbitrary positive-Möbius local preprocessors on fewer than \(\ell\)
variables. Two inputs to a product may use the same variables. Any gate
may have arbitrarily many consumers. The target need agree with the
circuit only on Boolean inputs.

\textbf{Theorem.} There are constants \(\alpha,\beta,c>0\), depending
only on \(h,\Delta\), such that for all sufficiently large \(N\) and
\(\ell=\lceil\alpha N\rceil\), any such circuit computing (F4) needs at
least \(2^{cN}\) distinct binary gates. The same bound holds on every
subcube obtained by fixing any \(b\le\beta N\) switches to any Boolean
values, with the same atom condition on the free variables.

\emph{Proof.} Reduce exponents using \(s_v^2=s_v\). At each gate retain
the family of subsets with positive multilinear coefficient. Uniqueness
of a Boolean multilinear representative follows by finite differences.
Addition takes the union of the two families. Multiplication takes all
pairwise unions of a set from each input. Positivity prevents loss of
any such union. After fixing switches, write \(W\subseteq V\) for the
free vertices. The target family is exactly \({\rm IS}(G[W])\). An
independent free set appears already with no chosen fixed vertices, and
every remaining term arises by deleting fixed vertices from an
independent set. These arguments also apply to any positive prefactor in
(F4).

For a gate \(g\), unroll the circuit just to reason about possible
monomials. In a derivation choose one nonzero summand at each sum, both
operands at each product, and one monomial per input atom. Fix an
occurrence of \(g\) in such a completed derivation and remove the
subtree below that occurrence. Call the union of the other leaf
monomials a context. Let \(S_g\) be the union of all sets produced at
\(g\) and \(T_g\) the union of all its contexts, over every such
occurrence. For every set produced at \(g\) and every context, inserting
its derivation gives a valid root monomial. Other occurrences of a
shared gate simply retain the monomials they used in the derivation.
Hence \textbf{there is no graph edge between \(S_g\) and \(T_g\)}, even
when the two sets overlap. Any edge between them would appear in one of
the resulting root monomials, contradicting its independence.

Let \(\delta=(1+2^\Delta)^{-1}\), take \(\alpha=\delta/32\), and let
\(b\le N/4\). For a uniformly random independent set \(I\) of \(G[W]\),
each vertex is included with probability at least \(\delta\). Among the
sets omitting \(v\), erase its occupied neighbors and insert \(v\). The
map has at most \(2^\Delta\) preimages per set containing \(v\). This
remains true after any sequence of vertex exclusions, which simply
deletes those vertices. Thus \(\mathbb E|I|\ge \delta|W|\ge\delta N/2\),
and \(\Pr[|I|\ge4\alpha N]\ge\delta/2-4\alpha\ge\delta/4\), since
\(|I|\le N\).

For all large \(N\), \(4\alpha N\ge3\ell\). In any derivation of a root
set of size at least \(3\ell\), take a lowest occurrence producing at
least \(\ell\) vertices. It is a binary product: leaves have support
below \(\ell\), and a selected child of a sum has the same monomial. Its
set \(A\) has size at most \(2(\ell-1)\), since it is a union of its two
child sets. Its context \(C\) consequently contains at least \(\ell\)
vertices outside \(A\). In particular \(|S_g|,|T_g|\ge\ell\).

For this gate each vertex of \(T_g\) has a neighbor in
\(V\setminus S_g\), using minimum degree in the original graph and
absence of cross edges. Therefore
\(|V\setminus S_g|\ge |T_g|/\Delta\ge\ell/\Delta\). Expansion across
\(S_g\) gives at least \(h\ell/\Delta\) crossing edges. Their other
endpoints all belong to \(V\setminus(S_g\cup T_g)\). Bounded degree then
gives \[
 |W\setminus(S_g\cup T_g)|\ \ge\ h\ell/\Delta^2-b.
\] Choose \(\beta=\min(1/4,h\alpha/(2\Delta^2))\). Every root set
represented through this gate omits a fixed set of at least
\(h\alpha N/(2\Delta^2)\) free vertices. Under the uniform
independent-set distribution its probability is at most
\((1-\delta)^{h\alpha N/(2\Delta^2)}\), by conditioning on each
successive exclusion and applying the preceding insertion bound. The
balanced gates cover total probability at least \(\delta/4\). A union
bound now requires at least
\((\delta/4)(1-\delta)^{-h\alpha N/(2\Delta^2)}\) \textbf{distinct
binary gates}. Decreasing \(c>0\) absorbs the constant. The proof used
derivations only to classify root support; the objects counted are
actual circuit gates. \(\blacksquare\)

Jerrum and Snir's parse-tree content argument underlies the use of one
gate with its possible contexts. Jukna's Lemma 8 for Boolean-input
monotone counting circuits gives a balanced cover by pairwise unions of
support sets, already permitting overlap. Bova, Capelli, Mengel and
Slivovsky's Corollary 2 gives the precise exclusion probability used
here by complementing their vertex covers. Korhonen's tropical MWIS
Lemma 10 rules out a graph edge across two product supports and also
rules out a shared vertex because its weight would be counted twice. For
Boolean exponents, shared vertices survive. The estimate in the proof
handles them in the original expanding graph and subtracts the fixed
vertices in the subcube. The resulting application covers
positive-coefficient atoms whose individual monomials have degree below
\(\ell\). De Colnet and Mengel's arithmetic-circuit framework permits
literal leaves of both polarities and uses ``positive'' for a
nonnegative output function. Here the sign condition concerns the Möbius
coefficients of every input atom.
\href{https://snir.cs.illinois.edu/listed/J7.pdf}{Jerrum and Snir,
Theorem 3.2} ·
\href{https://web.vu.lt/mif/s.jukna/ftp/counting-web.pdf}{Jukna, Lemma
8} · \href{https://arxiv.org/pdf/1411.1995}{Bova et al., Corollary 2 and
Theorem 4} ·
\href{https://doi.org/10.4230/LIPIcs.ICALP.2021.87}{Korhonen, Lemma 10
and Theorem 1} ·
\href{https://proceedings.kr.org/2021/20/kr2021-0020-de-colnet-et-al.pdf}{de
Colnet and Mengel, §§2.1--2.2}

\subsection{F4. Source execution and precision
contracts}\label{v8-f4.-source-execution-and-precision-contracts}

\begin{enumerate}
\def\labelenumi{\arabic{enumi}.}
\tightlist
\item
  On the count slice, contract each logical vertex's implication cycle
  together with its availability clause. Each copy is the same bit on
  any nonzero assignment. The resulting local tensor is zero on unequal
  external bits, one for the common bit zero, and \(s_v\) for the common
  bit one. In the small formula these are all external graph ports. In
  the version with requests, one may first delete the inactive request
  factors at the logical level using their equality certificates and
  eliminate the uniquely determined unused OR gate wires. This costs
  \(O(N+|E|+|\mathcal D|)\) symbol or Boolean operations for bounded
  degree. The remaining conflict tables have constant zero/one entries.
  In Appendix D the subdivision equalities give the same identification
  at \(T=0\). Hence (F4) is precisely the origin contraction of explicit
  \(0,1,s_v\) tensors in both sources. It is not an independently
  supplied counting oracle.
\item
  A fixed straight-line scalar evaluation of those tensors by
  nonnegative sum-product operations pays one unit per binary gate it
  actually evaluates. Its gates may be reused, reordered, and shared
  between intermediate products with any overlap. Any preprocessing that
  supplies a legal local positive atom is included, even when the lower
  bound grants its value as a free leaf. On this contract the circuit
  theorem lower-bounds the scalar work. If the charged execution of cost
  \(T\) covers that circuit and each paid unit covers at most \(w\)
  distinct binary gates, then it covers at most \(wT\) gates. Thus
  \(T\ge2^{cN}/w\). A compressed primitive requires its own scalar
  implementation or a proved delivery bound.
\item
  Fix a deterministic evaluation with a fixed initial control state. A
  selector may ask adaptive questions about arbitrary functions of local
  input groups, interleaved with arithmetic and cache accesses. Record
  operation and operand identities and every choice of address, shortcut
  and termination. For the \(i\)-th control choice on input \(s\),
  suppose every other input that follows the same control prefix and
  agrees with \(s\) on a specified set \(Q_i\) makes the same next
  choice. This prefix records choices and identities. Numerical operands
  may vary. With \(Q=\bigcup_i Q_i\), induction gives the same complete
  trace throughout the subcube \(s'_Q=s_Q\). On that subcube the
  executed arithmetic is one correct positive DAG, including distinct
  producers for reused values. Each incoming cache entry must satisfy
  the atom hypothesis or have a charged positive producer DAG valid on
  the same subcube, with the relevant producer control choices covered
  as well. If \(|Q|\le\beta N\), F3 requires \(2^{cN}\) distinct gates
  in the resulting DAG. An execution with fewer gates therefore has no
  such certificate on \(\beta N\) or fewer coordinates, independently of
  the number of plans. This is a dependence statement. One global oracle
  test may have footprint \(N\) in one call. A bound on literal reads
  requires an access model that charges those reads.
\item
  A fixed straight-line circuit may perform every addition and
  multiplication modulo \(2^p\). Give its atoms nonnegative integral
  positive-Möbius lifts and let \(\widetilde C\) be the output of the
  same circuit over the integers. Assume correctness modulo \(2^p\) and
  \(\widetilde C(\mathbf1)<2^p\), where \(\mathbf1\) sets all free
  switches to one on a restricted subcube. Every \(\widetilde C(s)\) is
  between zero and \(\widetilde C(\mathbf1)\). When the true counts fit
  this precision, their congruence gives exact agreement with the
  integer lift everywhere, so F3 applies. Every ordinary
  variable-elimination or explicit tensor-contraction calculation on the
  positive origin tensors that is valid over nonnegative integers has
  the true integral terminal lift by distributivity. The source working
  precision \(P_0=K+3\) suffices for \(3\cdot2^N\). Intermediate
  residues may be represented canonically. For example, \(1+(2^p-1)s\)
  returns the Boolean complement modulo \(2^p\) using nonnegative
  integer constants. Its prescribed integer lift at \(s=1\) equals
  \(2^p\), demonstrating the role of the terminal bound.
\end{enumerate}

The restriction to positive Möbius coefficients after local contraction
matters. Raw complementary switch entries, a signed Taylor difference,
and unrestricted reduction with negative residues can restore Boolean
computation. MAIN actually uses signed differences and rational
interpolation. D actually uses signed deformed tensors, higher precision
and reconstruction coefficients. Consequently F3 measures the cost of
\textbf{staying in this explicit native origin submodel}; taking the
source's claimed bit upper bound as an input theorem, it gives an
exponential separation between positive origin computation and the
source's signed subexponential counting algorithms. The bound does not
transfer to those signed algorithms or to an unrestricted interleaving
of their primitives. \hyperref[v8-section-native-closure]{Local-table
simulation} ·
\href{https://github.com/openai/math/blob/main/preprints/The-Exponential-Time-Hypothesis-is-False-October-25-2026/build/main_full_version.tex\#L1945}{MAIN
lines 1945--2014} ·
\href{https://github.com/openai/math/blob/main/preprints/The-Exponential-Time-Hypothesis-is-False-October-25-2026/build/quantitative-contraction.tex\#L173}{D
contraction lines 173--244, 558--619}

\subsection{F5. Uniform embeddings on expanding ambient
graphs}\label{v8-f5.-uniform-embeddings-on-expanding-ambient-graphs}

The fixed-topology restrictions prove unconditional statements under
execution contracts. The following separate, uniform reductions explain
the information retained by this family when the input topology is also
variable.

Take a graph \(H\) on \(n_0\) vertices, padded if needed by disabled
isolated vertices to an even order \(n\). Build a disjoint cubic
expander \(B\) on \(n\) vertices and match each vertex of \(H\) and its
padding to a different vertex in \(B\). Call the resulting \(2n\)-vertex
graph \(G\). Alon's Theorem 1.3 constructs a suitable cubic \(B\) in
deterministic polynomial time on all sufficiently large even \(n\) by
taking \(\epsilon<3-2\sqrt2\). Its spectral bound gives edge expansion
\(h_B>0\). Indeed, for \(U\) of size at most \(n/2\), the quadratic form
of the graph Laplacian on \(1_U-|U|/n\) gives
\(|\partial U|\ge(3-\lambda)|U|(1-|U|/n)\). Constant-size exceptional
inputs can be handled directly or padded.
\href{https://web.math.princeton.edu/~nalon/PDFS/expexpg.pdf}{Alon,
Theorem 1.3}

For any cut of size \(z\le n\) in \(G\), identify its old-side indices
\(A\) and its backbone indices \(B'\). Put \(t=|A\triangle B'|\). The
perfect matching gives \(t\) crossing edges. The intersection has size
\((z-t)/2\), which is at most both \(|B'|\) and \(n-|B'|\), since
\(z\le n\). The backbone gives at least \(h_B(z-t)/2\) further crossing
edges. Thus \(G\) has edge expansion at least \(\min(1,h_B/2)\),
independent of \(H\). Disabling the backbone and the padding leaves
precisely the choices in \(H\).

\begin{itemize}
\tightlist
\item
  For \textbf{counting}, take \(H\) cubic, enable all its vertices, and
  deactivate all requests. A cubic graph already has even order, so
  sufficiently large instances need no padding and the ambient graph is
  four-regular. Handle the finitely many small orders by direct
  computation. The exact answer is \(3|{\rm IS}(H)|\), so one evaluation
  followed by exact division by three returns the desired
  independent-set count. Greenhill proves cubic independent-set counting
  \(\#\mathrm P\)-complete under polynomial-time counting reductions.
  Both the count-only and the full source formula implement this
  embedding.
  \href{https://web.maths.unsw.edu.au/~csg/papers/complexity.pdf}{Greenhill,
  Section 3.2}
\item
  For \textbf{decision}, start with the usual width-three SAT instances
  whose clauses use three distinct variables. Replace each variable of
  occurrence at least two by distinct occurrences connected with the
  source implication cycle. The resulting clauses have widths two and
  three, each variable has at most three occurrences, and satisfiability
  is preserved. Give every resulting literal occurrence a vertex. Make
  each clause a clique on its two or three occurrences and add an edge
  between opposite literals of the same variable in different clauses.
  This simple graph \(H\) has maximum degree at most four. A clause
  vertex has at most two same-clause neighbors and at most two opposite
  occurrences. Activate the request on each two-clause's clique edge,
  add and activate a disjoint triple request for each three-clause's
  clique, and leave all other requests off. Enable the original vertices
  and disable the padding and backbone. A satisfying assignment supplies
  a true literal per clause, giving an independent set that meets every
  request. Conversely such a set chooses exactly one literal per clause
  and never both polarities of any variable. Extend these consistent
  choices to an assignment satisfying the whole formula. This is a
  polynomial-time decision reduction. For an input with \(n\) literal
  occurrences, the construction uses \(O(n)\) vertices, edges and output
  literal occurrences. With consecutive binary variable identifiers the
  explicit clause list occupies \(O(n\log(n+2))\) bits. Multiple choices of true
  literals can arise from one satisfying assignment.
  \href{https://github.com/openai/math/blob/main/preprints/The-Exponential-Time-Hypothesis-is-False-October-25-2026/build/main_full_version.tex\#L1465}{Occurrence-copy
  construction in MAIN}
\end{itemize}

The latter graphs have maximum degree at most five, but their average
density keeps the printed packing balanced. The old half contributes at
most \(2n\) edges, the backbone \(3n/2\), and the perfect matching
\(n\). Hence \(|E|/N\le9/4\) for \(N=2n\), and the disjoint requested
triples on the old half satisfy \(|\mathcal D|/N\le1/6\). Formula (F2)
gives \(M_{\rm MAIN}/N\le 179/4+2/N<45\) eventually. Thus the retained
source symbol masses satisfy the bound of F2 even here.

The simultaneous matching proof for cubic graphs also works with any
fixed maximum degree \(d\). The convexity and flow parts use only edge
expansion. In its thinning step an edge conflicts with at most \(2d^2\)
other edges whose endpoints share or neighbor one of its ends. Marking
with probability one half and retaining an isolated marked edge has
probability at least \(2^{-(2d^2+1)}\). Signed indicators for two edges
with disjoint conflict neighborhoods use disjoint independent random
bits, so each summand has only \(O_d(1)\) dependent partners. Since
\(|E|\le dN/2\), the variance is \(O_d(N)\), as required in the original
Chebyshev union bound. The constants may depend on \(d\). This covers
the bounded-degree expanding topologies just constructed under the same
graph-only layout and menu conditions. The reductions themselves entail
no unconditional superpolynomial time bound. They also give no such
conclusion from any one encoded instance for a topology-dependent
program with arbitrary advice.

\textbf{Finite checks.} A generic exact \(\#\)DPLL evaluator checks the
actual occurrence-split formulas, without being told their implication
cycles. Across 456 \nolinkurl{request/availability} words on four small
graphs, the physical count agrees with (F1) and every normalized clause
list has the stated antichain, degrees and sizes. The three matching
matrices in F2 have computed binary ranks 2, 4, and 8 at one, two, and
three edges. A separate six-vertex exercise verifies all 729 subcube
support identities, all 729 vertex-exclusion probabilities, and every
pair of vertex sets with no cross edge, including 74 pairs with overlap.
These are finite exercises supporting the written proofs.
\PaperSixEvidence{evidence/research_v8/checks/stronger_family.py}{Reproducible checker} ·
\PaperSixEvidence{evidence/research_v8/checks/stronger_family.json}{Recorded output}

\section{Actual native registers and restricted cross
multiplication}\label{v8-section-native-escape}

The rational MAIN design and the fixed-density Appendix D clause network
accept the common four-clone input. There are \(N\) externally varied
unary pin signs, \(n_M=4N\) private Boolean variables, \(M_M=12N\) MAIN
incidences and \(M_D=24N\) physical D wires. The D size parameter is
\(s=21N/2\). In the triangle replacement graph the internal edges form
disjoint triangles and the exterior edges form a perfect matching. Write
\(U_M\) for the product of the internal edge ORs and \(U_D\) for the
product of the exterior edge ORs, so the full count is \(C_G=U_MU_D\).
The v7 graph-only coordinate ties put each triangle's three pins into
one MAIN group and each exterior pair's two pins into one D group. These
placements satisfy the source certificates. The output of a prescribed
independent-label tie-breaking algorithm can differ. The statements
below distinguish actual register values, the cost of their printed
evaluators, and restricted ways to compute those values.
\PaperSixEvidence{evidence/research_v7/reviews/hybrid_independent.md}{v7 construction
and placement quantifiers} ·
\PaperSixEvidence{evidence/research_v7/source/cost_contracts.md}{v7 source contracts}

\subsection{Actual arithmetic inputs and
registers}\label{v8-actual-arithmetic-inputs-and-registers}

For MAIN take \(B=\mathbb Q[z]/(z^J)\). A group register has graph-fixed
indices \((t,E,U,\vec y,g)\) with \(1\le t\le Q\), \(U\subseteq E\),
\(|E|\le J-1\) and a measurement vector in the printed box. Each
unrestricted measurement fiber must be nonempty to use the
decomposition. Fix a representative in each fiber using the layout
alone. The sum for \(g\) ranges over the fiber \(\mathcal X_{g,E}(y_g)\)
with private bits in \(E\cap S_g\) set to one. For a summation point
\(x_g\), fill the other private groups with their fixed representatives
to obtain \(x^{\rm rep}\). If \(v\) indexes an external pin and \(j(v)\)
is its private pin clone, the literal argument at that unary factor
\(f(v)\) is \[
 \theta_v=\langle P_{f(v),j(v)}(t),x^{\rm rep}\rangle+
 \langle P_{f(v),j(v)}(t+z)-P_{f(v),j(v)}(t),\mathbf1_U\rangle .
\] Its two possible values as a factor are \(\ell_1(\theta_v)=\theta_v\)
and \(\ell_0(\theta_v)=1-\theta_v\). The other clause factors are sign
independent at this fixed private point. All factors assigned to that
group are included. Thus the literal summand in \(\Sigma_g\) is \[
 c_g(x_g)\prod_{v\in V_g}\ell_{\sigma_v}(\theta_v),\qquad c_g(x_g)\in B,
\tag{N1}
\] where \(V_g\) denotes the externally signed unary factors assigned to
\(g\). A completed root register is \(\prod_g\Sigma_g\), and a fixed
rational linear readout of this jet produces a scalar. At the prescribed
parameters \(J\), \(D=r_M^{2Q}\) and the bit lengths of the rational
coefficients are polynomial in \(N\). The row denominator divides
\((R-1)!Q!\). For direct clauses the integer factor data and \(O(N)\)
total polynomial degree imply polynomial bit lengths in each jet
summand, sum and product, by clearing this common denominator to the
total degree and adding at most \(O(N)\) bits for assignment sums.
\href{https://github.com/openai/math/blob/main/preprints/The-Exponential-Time-Hypothesis-is-False-October-25-2026/build/main_full_version.tex\#L1714}{MAIN
rows and measurements, 1714--1763} ·
\href{https://github.com/openai/math/blob/main/preprints/The-Exponential-Time-Hypothesis-is-False-October-25-2026/build/main_full_version.tex\#L1810}{Fibers
and representatives, 1810--1848} ·
\href{https://github.com/openai/math/blob/main/preprints/The-Exponential-Time-Hypothesis-is-False-October-25-2026/build/main_full_version.tex\#L1886}{Corrections
and factor products, 1886--1928 and 1964--1997} ·
\href{https://github.com/openai/math/blob/main/preprints/The-Exponential-Time-Hypothesis-is-False-October-25-2026/build/main_full_version.tex\#L2623}{Coefficient
bounds, 2623--2641 and 3974--4008}

For D a root slot \((t,\zeta,\beta)\) has \(t\ne0\), \(i=v_2(t)<u\), a
prescribed residue-distinct interpolation probe \(\zeta\), and a
nonnegative mode tuple with \(\sum_Z w_Z\beta_Z<p_t\). Each original row
or column vertex goes to the child specified by its symbol in coordinate
\(i\). In child \(g\), every old port of this group becomes one edge to
its packet tensor. The tensor at a packet extracts the prescribed
targeted mode coefficient from the product of the local row forms for
\(g\), and multiplies it by the prescribed powers of its source column
moments. The root slot value is the product of all 64 child
contractions. Omitting an old clause or inserting a new two-ended
predicate changes this definition. The native storage is
\(\mathcal A_{p_t}\), with coefficient degree
\(e=4\lceil\log_2s\rceil\), \(\ell=1\) on this family and
\(p_t=\operatorname{poly}(s)\). Arithmetic in it and entries of the
original bounded-degree tensors cost polynomially many bit operations.
The published general packet-entry routine instead takes
\(\exp(O(L\log s))\) time per entry in its stated upper bound. Immutable
original and newly created entries keep their creation-precision
representatives.
\href{https://github.com/openai/math/blob/main/preprints/The-Exponential-Time-Hypothesis-is-False-October-25-2026/build/quantitative-contraction.tex\#L11}{D
entry and lift conventions, 11--71} ·
\href{https://github.com/openai/math/blob/main/preprints/The-Exponential-Time-Hypothesis-is-False-October-25-2026/build/quantitative-contraction.tex\#L250}{Nodes
and precision, 250--268 and 338--368} ·
\href{https://github.com/openai/math/blob/main/preprints/The-Exponential-Time-Hypothesis-is-False-October-25-2026/build/quantitative-contraction.tex\#L432}{Source
moments, 432--454} ·
\href{https://github.com/openai/math/blob/main/preprints/The-Exponential-Time-Hypothesis-is-False-October-25-2026/build/quantitative-contraction.tex\#L487}{Child
entries, 487--553}

\subsection{The cost of a single completed
call}\label{v8-the-cost-of-a-single-completed-call}

\textbf{Printed acquisition bound.} A feasible fixed MAIN tuple formed
by the printed evaluator visits at least \(2^{4N/r_M}\) private Boolean
vectors before completing the group product. A single fixed D root slot
formed by the printed recursion visits more than \(2^{H/32}\)
assignments in just one of its descendants, where \[
 H=\left\lceil s/\lfloor\sqrt{\lceil\log_2s\rceil}\rfloor\right\rceil.
\] Consequently neither printed evaluator supplies the corresponding
completed factor at polynomial cost, for any feasible specialization of
the indices on growing instances.

For MAIN the second scan for every group tests its entire Boolean cube,
and the largest variable group has size at least \(4N/r_M\). At the
printed parameters the base-two logarithm of the bound is
\(\Omega(N(\log\log N)^2/\sqrt{\log N})\). An unrestricted measurement
tuple rejected as infeasible has an identically zero contribution. For
D, at each active call of \(M\) edges, the 64 children for any required
slot have edge counts \(m_g\) satisfying \[
 \sum_g m_g=2M,\qquad M/32\le\max_gm_g\le 3M/32.
\tag{N2}
\] Every parent edge has two ports, and each port becomes exactly one
child edge. The upper inequality is the source's packing bound. Starting
at the active root, follow a largest child of the requested slot. At
each further active node select any required slot and a largest child.
Sizes shrink strictly. The first leaf has more than \(H/32\) edges since
its parent had more than \(H\). The printed leaf sums over every
assignment to these edges, giving \(\Omega(N/\sqrt{\log N})\) in the
binary logarithm. This includes the zero mode and a zero probe. The
statement counts the actual unmodified loops. Caching, recognizing a
constant before invoking recursion, and factoring a child contraction
change the charge.
\href{https://github.com/openai/math/blob/main/preprints/The-Exponential-Time-Hypothesis-is-False-October-25-2026/build/main_full_version.tex\#L1850}{Printed
MAIN scan, 1850--1877} ·
\href{https://github.com/openai/math/blob/main/preprints/The-Exponential-Time-Hypothesis-is-False-October-25-2026/build/main_full_version.tex\#L2747}{MAIN
parameters, 2747--2786} ·
\href{https://github.com/openai/math/blob/main/preprints/The-Exponential-Time-Hypothesis-is-False-October-25-2026/build/quantitative-contraction.tex\#L73}{D
thresholds, 73--109} ·
\href{https://github.com/openai/math/blob/main/preprints/The-Exponential-Time-Hypothesis-is-False-October-25-2026/build/quantitative-contraction.tex\#L560}{D
children and literal leaf scan, 487--495, 525--529, 560--589}

These conclusions also control individually requested registers and
leaves under the source collision certificates. In MAIN,
\(\Xi\le 2Qn_M\) implies \[
 |S_{t,g}| \ge n_M/r_M-\sqrt{2Qn_M}
       =(1-o(1))\, n_M/r_M
\] for every group at the published polylogarithmic \(Q,r_M\). Thus even
one feasible individual \(\Sigma_g\) requires \(2^{(1-o(1))4N/r_M}\)
visits if evaluated by the printed local scan. This uses the joint
potential attained by the source or the v7 tied construction; a promise
of an upper group-size bound alone would not give this conclusion.
\href{https://github.com/openai/math/blob/main/preprints/The-Exponential-Time-Hypothesis-is-False-October-25-2026/build/main_full_version.tex\#L2463}{MAIN
joint potential, 2463--2510}

At \textbf{every active D call}, write \(D_{\rm arr}=64^{2a}\). Its
collision certificate gives
\(\lvert X_c-M/D_{\rm arr}\rvert\le\sqrt{2M\Delta}\). Since \(M>H\),
\(D_{\rm arr}\le s^{1/1000}\), and \(\Delta=\lceil s^{9/10}\rceil\), its
relative error is bounded uniformly by \[
 \eta_s=s^{1/1000}\sqrt{2\Delta/H}
         =O(s^{-49/1000}(\log s)^{1/4})=o(1).
\] Marginalizing cell counts over the \(i\)-th symbol on either side
gives \((1-\eta_s)2M/64\le m_g\le(1+\eta_s)2M/64\) for \textbf{every}
child at every active call. Thus every actual leaf under the active root
has more than \((1-\eta_s)H/32\) edges, including a leaf chosen
individually using its indices. Each root child is itself eventually
active. Even freely selecting one of these genuine leaves does not
produce a polynomial-cost call to the published direct-summation
primitive. A shortcut for a particular leaf's \emph{value} can have a
different cost.
\href{https://github.com/openai/math/blob/main/preprints/The-Exponential-Time-Hypothesis-is-False-October-25-2026/build/quantitative-contraction.tex\#L73}{D
uniform active-call parameters, 73--111} ·
\href{https://github.com/openai/math/blob/main/preprints/The-Exponential-Time-Hypothesis-is-False-October-25-2026/build/quantitative-contraction.tex\#L114}{Collision
and marginalization, 114--170}

\subsection{A single cross multiplication with sparse native MAIN
fibers}\label{v8-a-single-cross-multiplication-with-sparse-native-main-fibers}

Let a MAIN group be formed from a graph-fixed list of \(b_g\) private
summands of (N1), with any fixed sign-independent jet weights. The
complete fiber is included as a special case. Across \textbf{any}
row/column split of externally varied signs, a scalar readout of the
full product of these groups has rank at most \[
 \kappa\prod_{g=1}^{r_M} b_g,\qquad
 \kappa=J.
\tag{N3}
\] If the readout involves only orders through \(j\), take
\(\kappa=j+1\). Indeed after choosing one summand per group the product
over \emph{all} pins separates as
\(A(\sigma_{\rm row})B(\sigma_{\rm col})\) in the single algebra \(B\).
The scalar bilinear form \(\lambda(ab)\) has rank at most \(\dim B=J\),
and descends to \(\mathbb Q[z]/z^{j+1}\) in the low-order case.
Summation over the choices proves (N3). The finite bound charges actual
contributing private points and avoids a factor of \(J\) per group.

\textbf{One-coordinate theorem.} Fix a certified D coordinate \(i\) on
the cubic expander family. Consider \(K\) cross products, each a scalar
MAIN register as in (N3) times a full D group product at \(i\).
Different MAIN coordinates, measurements, Taylor pairs, lists and
readouts are allowed. The D group functions may be arbitrary functions
of their own signed unary factors lifted separately into the common
characteristic-zero field. Fixed coefficients and fixed linear
recombinations may be collected or distributed. Suppose the final sum is
integral and has the correct terminal count residue on the full sign
cube. There is \(c>0\) independent of \(N\) such that \[
 \sum_{h=1}^K \kappa_h\prod_{g=1}^{r_M} b_{hg}\ \ge\ 2^{cN}.
\tag{N4}
\] Pure MAIN terms use a constant D product, pure D terms use a constant
MAIN scalar of capacity one, and a fixed scalar adds one. A separately
supplied bypass is charged its actual matrix rank on the chosen
restriction.

The D unary pin marginal has density at most \(48\) times the uniform
symbol law, so no symbol contains more than \(3N/4\) pins. A union of
whole groups with between \(N/4\) and \(3N/4\) pins exists by inspecting
a large atom, or adding atoms of size less than \(N/4\). A cubic
expander of expansion \(\gamma\) has at least \(\gamma N/4\) edges
crossing this cut. Greedily keep a crossing edge and exclude all edges
within line-graph distance two. At most \(1+4+4\cdot3=17\) candidate
edges are removed per choice, so this gives an induced matching of size
at least \(\gamma N/68\). Direct its pairs from the chosen union to the
complement. Fix all other pin signs to one. The target is a Kronecker
power of the binary OR matrix, of rank \(2^m\) in every characteristic.
The assembled integral lift has this full fraction-field rank because
its determinant is nonzero after reduction modulo two. Every D group now
sees only row signs or only column signs, so every D product has rank at
most one. The Hadamard product rank inequality and (N3) give (N4), with
\(c=\gamma/68\). No claim that either factor itself equals the v7
internal or exterior edge product is needed.
\PaperSixEvidence{evidence/research_v7/source/cost_contracts.md}{Common input, density
and output lift} · \PaperSixEvidence{evidence/research_v7/theory/hybrid_budget.md}{The
v7 residue transfer}

For a fixed finite list using \textbf{different} D coordinates, the v7
common matching instead gives a matched target of rank \(2^m\) and every
full D product has rank at most \(2^{m-d_D}\), with
\(d_D=\Omega(N/\sqrt{u})\). The same calculation replaces \(cN\) by
\(d_D\) in (N4). At the source \(u=\Theta(\log N)\).

For an execution bound, take a fixed representation in (N4) with zero
bypass. Collect repeated products with fixed coefficients, and combine
repeated private points in a list. Suppose a single trace covers all its
\(K\) designated products and every designated point in each list, with
at most \(W=\operatorname{poly}(N)\) designated outputs delivered per
unit of work. Shared computation and repeated reads are permitted. If
its total time is \(T\), the coverage of each of these sets gives
\(K,\max_{h,g} b_{hg}\le WT\). An adaptive trace that covers only the
terms on its branch requires its own guarded representation. With
\(\kappa_*=\max_h\kappa_h\), (N4) gives \[
 \log_2 T\ \ge\ \frac{cN-\log_2\kappa_*}{r_M+1}-\log_2 W.
\tag{N5}
\] On the single D coordinate this is
\(\Omega(N(\log\log N)^2/\sqrt{\log N})\). Across the full native D
coordinate menu it is \(\Omega(N(\log\log N)^2/\log N)\). If exactly one
product of linear MAIN and D forms is first distributed symbolically,
its \(K\) can be bounded by the product of the two actual entrance
counts, giving an additional constant additive term in the denominator
in place of individually charging all cross terms. For a gate whose
input fiber is summed by other means, \(b_{hg}\le WT\) must be proved
anew. (N4) remains a statement about the number of formal private points
and does not imply that every arithmetic evaluation lists them.
\PaperSixEvidence{evidence/research_v7/source/cost_contracts.md}{v7 execution contract}

\textbf{Canonical lifts of completed D slots.} For a native slot at
\(t\), let \(\widehat Z\) be the canonical lift of its completed product
modulo \(2^{p_t}\). Let \(\widetilde Z_g\) be separate integral lifts of
its child residues. A scalar multiplier \(b(\sigma)\in\mathcal K\) can
include a MAIN readout and its fixed coefficient. For every input and
every requested terminal precision \(b_0\ge1\) one has \[
 p_t+v_2(b(\sigma))\ge b_0
 \quad\Longrightarrow\quad
 b(\sigma)\bigl(\widehat Z-\textstyle\prod_g\widetilde Z_g\bigr)
                          \in 2^{b_0}\mathcal O.
\] This follows by reduction of the two lifts to the same product. With
\(v_2(0)=+\infty\), the condition includes vanishing multipliers. If it
holds term by term, the assembled integral sum retains its terminal
residue upon making the replacements. This gives (N4) for a genuine
whole-slot lift with those multipliers. The lift of the complete product
itself need not split across the cut.

The ordinary source readouts satisfy the condition on these direct
clause instances for all sufficiently large \(N\), even with
\(b_0=P_0=4N+1\). To check the loss, put \(\delta=(R-1)!Q!\) and
\(d=12N\), the sum of the degrees of the clause factors. Every MAIN jet
coefficient of a complete tuple lies in \(\delta^{-d}\mathbb Z\). The
same holds for graph-fixed private sublists with integer-coefficient jet
weights. Weights in \(\mathbb Z_{(2)}[z]/(z^J)\) give the corresponding
bound \(\delta^{-d}\mathbb Z_{(2)}\). Integer shifts, Boolean
substitutions and truncation preserve the respective bounds. For the
published Hermite origin readout take any MAIN node \(v\in[Q]\) and put
\[
 B_v(T)=\prod_{w\ne v}\left(\frac{T-w}{v-w}\right)^J .
\] Its factor at zero is \(\bigl((-1)^{v-1}\binom Qv\bigr)^J\), an
integer. The coefficient of \(z^k\) in \(1/B_v(v+z)\) has denominator of
two-adic valuation at most \(k\lfloor\log_2 Q\rfloor\), by expanding the
factors \((1+z/(v-w))^{-J}\). Multiplying the given jet by this inverse
through order \(J-1\), evaluating at \(-v\), and multiplying by
\(B_v(0)\) gives the node's exact Hermite contribution. This follows
from the interpolation polynomial having the prescribed jet at \(v\) and
zero jets at the other nodes. The coefficient-extraction readout and the
origin readout therefore have loss, including their MAIN input, at most
\[
 L_M=d\,v_2(\delta)+(J-1)\lfloor\log_2 Q\rfloor=O(N\log N).
\] Here \(J=1+\lfloor dR/Q\rfloor\), \(v_2(\delta)\le R+Q\), and the
published \(R,Q\) are \(O(\log N)\). Möbius signs are integral.
\href{https://github.com/openai/math/blob/main/preprints/The-Exponential-Time-Hypothesis-is-False-October-25-2026/build/main_full_version.tex\#L2623}{MAIN
layout denominator and degree, 1688--1704 and 2623--2641} ·
\href{https://github.com/openai/math/blob/main/preprints/The-Exponential-Time-Hypothesis-is-False-October-25-2026/build/main_full_version.tex\#L1886}{Jet
order, 1886--1902} ·
\href{https://github.com/openai/math/blob/main/preprints/The-Exponential-Time-Hypothesis-is-False-October-25-2026/build/main_full_version.tex\#L4126}{Origin
functional, 4126--4165}

For the D root put \(M=24N\), \(h=\lceil3M/u\rceil\). Each scalar
coefficient of the ordinary D origin readout after distributing its
linear operations has loss at most \(L_D(t)=6hD_t+4M(a+1)\). The source
bounds the loss of \(E_t(0)\) by \(6hD_t\). The inverse local Taylor
series are integral, evaluation at \(-t/\alpha\) loses at most
\(4M(a+1)\), and residue-distinct Lagrange probes have unit
denominators. The root precision is \(p_t=P_0+8hD_t+8M(a+1)\). For every
active node \(i=v_2(t)<u\), \(a=2u\) gives
\(D_t\ge2^{u+1}\ge s^\varepsilon\). Thus
\(hD_t=\Omega(N^{1+\varepsilon}/\log N)\) uniformly, including
\(t=\pm1\), and eventually \[
 p_t-L_D(t)-L_M\ \ge\ P_0.
\] This proves the replacement for a native MAIN scalar multiplied by a
native D slot, including one ordinary D reconstruction coefficient.
Additional weights, divisions, or nonlinear readouts require their own
congruence proof. For example, at precision two the two child residues
\(2,2\) have canonical product zero, while their integer product is
four. Multiplication by \(1/4\) distinguishes those two lifts at parity.
\href{https://github.com/openai/math/blob/main/preprints/The-Exponential-Time-Hypothesis-is-False-October-25-2026/build/quantitative-contraction.tex\#L40}{D
ring and lifts, 40--71} ·
\href{https://github.com/openai/math/blob/main/preprints/The-Exponential-Time-Hypothesis-is-False-October-25-2026/build/quantitative-contraction.tex\#L250}{D
multiplicities, origin weights and probe interpolation, 250--345}

\subsection{The specified triangle and matching
factors}\label{v8-the-specified-triangle-and-matching-factors}

A single native MAIN register with scalar linear readout equal to
\(U_M\) satisfies the sharper local assertion \[
 \kappa\,|\mathcal X_{g,E}(y_g)|\ \ge\ 2^{k_g}
\tag{N6}
\] for every group \(g\), where \(k_g\) is its number of whole
triangles. In each of those triangles keep two signs, put one on each
side and freeze every other sign in the input to one. The target is the
Kronecker power of \(k_g\) OR matrices. All other groups give a fixed
jet multiplier that can be absorbed in the readout. Each private point
has rank at most \(\kappa\) by (N1), as in (N3). The same estimate
applies to any graph-only weighted list of native private points
computing the register. At the designated v7 MAIN coordinate, all
triangles are tied and some \(k_g\ge N/(3r_M)\).

At any other MAIN coordinate a single complete register with this
readout computing \(U_M\) can split at most \(2\log_2\kappa\) triangles
across its groups. Color the groups independently with two equiprobable
colors. Each split triangle meets both colors with probability at least
one half. Some group cut therefore crosses at least half the split
triangles. Keep two oppositely colored signs in each crossing triangle,
and fix its third sign and every remaining sign to one. This again gives
a Kronecker power of OR with one factor per chosen triangle. Every
varied sign belongs to one unary factor, and the fixed measurement
representatives ensure that its sign enters just that factor's assigned
\(\Sigma_g\). The entire completed jet consequently factors as
\(A(\sigma_{\rm row})B(\sigma_{\rm col})\) across this cut. Its single
scalar readout has rank at most \(\kappa\), regardless of the fiber
lengths. At most \(\log_2\kappa\) triangles can cross the chosen cut, as
claimed. Hence some \(k_g\ge (N/3-2\log_2\kappa)/r_M\) whenever such a
register computes \(U_M\), even on an untied layout. Together with (N6)
this excludes a polynomial-size private list in every group of that
register. The assertion uses one completed jet and one scalar linear
readout.

The zero measurement tuple has an exact value. The published interval
satisfies \(6R\le Q\) and \(A_0=\lfloor Q(r-1)/(2r)-(R-1)/2\rfloor>0\)
eventually. The exceptional matrix includes the weight-zero tensor
eigenvector for every orbit of the coordinate swaps. This vector is its
orbit indicator, a product of symmetric one-coordinate basis elements.
After placement, the corresponding per-copy and per-group row counts
ones in that orbit's occupied positions, with positive multiplicities
for repeated private variables. Every private variable occurs. If all
those rows measure zero, every private bit is zero. The all-one
measurement likewise has a unique preimage by taking Boolean
complements. At zero measurements any nonempty \(E\) gives an empty
restricted fiber. If \(E=U=\varnothing\), all input forms are zero and a
positive edge OR forces the complete product to vanish.
\href{https://github.com/openai/math/blob/main/preprints/The-Exponential-Time-Hypothesis-is-False-October-25-2026/build/main_full_version.tex\#L2241}{Source
symmetric eigenvectors and exceptional rows, 2241--2269} ·
\href{https://github.com/openai/math/blob/main/preprints/The-Exponential-Time-Hypothesis-is-False-October-25-2026/build/main_full_version.tex\#L2537}{Embedding
and group restriction, 2537--2591} ·
\href{https://github.com/openai/math/blob/main/preprints/The-Exponential-Time-Hypothesis-is-False-October-25-2026/build/main_full_version.tex\#L2796}{Published
interval, 2796--2805}

At a fixed assignment to all edges of a D root child, its original unary
pin tables select exactly one sign word on its pins. Every other word
gives a zero product. Therefore any sign-independent list of individual
child edge assignments has support size at most the list length, even
with fixed signed weights. If a full D root product has the correct
exterior-edge OR residue, every exterior pair must have both pins in the
same group. Otherwise freeze all signs except the endpoints of a split
pair to one. The product of the individual integral group residues in
\(\mathbb F_{2^e}\) gives a matrix of rank at most one across these two
endpoints, whereas their OR has rank two. This separability condition is
met at the designated v7 tied D coordinate. At any coordinate at which
such a product computes \(U_D\) modulo two, let \(k_g\) be the number of
its group \(g\)'s whole exterior pairs. Freeze the other groups to ones.
All \(3^{k_g}\) satisfying local words must occur in the list for \(g\).
The \(N/2\) co-located pairs imply some \(k_g\ge N/128\). This is a
bound on direct fixed-list expansion on the whole cube. A sign-adaptive
enumeration can first impose the unary bits, and the printed recursion
already uses a different decomposition.

\subsection{What the D zero mode actually
evaluates}\label{v8-what-the-d-zero-mode-actually-evaluates}

Write \(a=2u\) at the root, \(L=2^a\), \(\alpha=2^{a+1}\),
\(S=\{-L,\ldots,L-1\}\), and \(\mathcal C(T)=\prod_{d\in S}(T-d)\). At a
fixed active \(t\), let \(c_t=\mathcal C'(t)/\mathcal C'(0)\), an odd
unit of \(\mathbb Z_{(2)}\). For real row and column slots in one packet
set \(q=y-x\). The source Cauchy formula gives exact polynomial
identities \[
 W_{xy}(t+\alpha z)=
 \begin{cases}
 c_t\displaystyle\prod_{d\ne t}\bigl(1+\alpha z/(t-d)\bigr),&q=t,\\[3pt]
 c_t\,\dfrac{\alpha z}{t-q}
    \displaystyle\prod_{d\ne t,q}\bigl(1+\alpha z/(t-d)\bigr),&q\ne t .
 \end{cases}
\tag{N7}
\] Every omitted or retained \(d\) in this formula belongs to \(S\).
Since \(v_2(t-d)\le a\) for \(d\ne t\), the coefficient of \(z^j\) is
divisible by \(2^j\). At \(z=0\) the surviving coefficient is \(c_t\) on
the real shifts \(x\mapsto x+t\) and zero elsewhere. These pairs are
local in coordinate \(i=v_2(t)\). At any \(\zeta\), reduction modulo two
gives the same partial Boolean wiring. A row without a real matched
column is fixed to zero, while every unmatched column bit is summed
subject to its original tensor. Therefore all zero-mode probes at a
fixed \(t\) yield the same mod-two child and root values.
\href{https://github.com/openai/math/blob/main/preprints/The-Exponential-Time-Hypothesis-is-False-October-25-2026/build/quantitative-contraction.tex\#L173}{D
deformation and integrality, 173--241} ·
\href{https://github.com/openai/math/blob/main/preprints/The-Exponential-Time-Hypothesis-is-False-October-25-2026/build/quantitative-contraction.tex\#L387}{Locality
of the shifts, 387--403}

At a nonzero mode tuple the \textbf{whole product} has divisibility by
\(2^{\sum_Z w_Z\beta_Z}\). Its targeted child coefficient already
acquires each required power of two. A source-only child can survive
even though its full product vanishes modulo two. Here a stronger lower
bound follows directly from the crossing coefficient. In the source
expansion with tile representative \(y_0\), \[
 v_2(t-y_0+x)\le i,\qquad
 v_2\!\left(\frac{\mathcal C(t+\alpha\zeta)}
                  {\mathcal C'(0)(t+\alpha\zeta-y_0+x)}\right)
                 \ge a+1-i.
\] The first fraction has the normalized numerator gain \(a+1\), and the
additional ratio \(2^{i+1}/(t+\alpha\zeta-y_0+x)\) has valuation at
least one. Hence each targeted occurrence of weight parameter \(b\)
supplies at least \(a+1-i+b\) powers of two. The valuation convention
includes zero; at \(\zeta=0\) all targeted coefficients are exactly
zero. In particular a full nonzero-\(\beta\) slot cannot supply \(U_D\)
modulo two, either alone or after multiplication by an integral scalar.
\href{https://github.com/openai/math/blob/main/preprints/The-Exponential-Time-Hypothesis-is-False-October-25-2026/build/quantitative-contraction.tex\#L418}{D,
418--454} ·
\href{https://github.com/openai/math/blob/main/preprints/The-Exponential-Time-Hypothesis-is-False-October-25-2026/build/quantitative-contraction.tex\#L487}{Target
coefficients and source moments, 487--523}

\textbf{An exact obstruction at the unit shift.} If zero is chosen among
the permitted residue-distinct interpolation lifts, no individual
completed D root slot with \(t=-1,\zeta=0\) equals \(U_D\) at its native
precision, for either choice of the bipartite sides and every certified
packet layout. This statement also holds for a direct additional
evaluation at that specific legal probe. It is a statement about values
modulo \(2^{p_t}\), stronger precision than parity alone.

Indeed the symmetric grid gives \(c_{-1}=-1\). Every nonzero mode
vanishes exactly here. For \(\beta=0\), if old vertices are columns,
each subdivision equality row whose two arguments are
\(b_1,b_2\in\{0,-1\}\) contributes \[
 (1-b_1)(1-b_2)+b_1b_2>0
\] after its two row bits are summed. All old column tensors have a
satisfying local port assignment for every pin sign, and use disjoint
column ports. At least one summand is strictly positive for every sign
word. If old vertices are rows, their Boolean multilinear extensions are
evaluated at points with entries \(0,-1\). A positive OR
\(b_1+b_2-b_1b_2\) is nonpositive and is nonzero precisely when some
argument is \(-1\). A positive unary is nonpositive, and a negative
unary \(1-b\) is positive. Every implication \(1-b_1+b_1b_2\) and every
trivalent equality
\(\prod_{j=1}^3(1-b_j)+\prod_{j=1}^3b_j=1-\sum_j b_j+\sum_{j<k}b_jb_k\)
is strictly positive. Thus every nonzero summand has the same sign. The
all-one assignment to the subdivision column equality bits is feasible.
Consequently the support is empty if an OR has no real shifted port, and
otherwise is exactly the subcube requiring all unmatched unary pins to
have sign zero. The exterior matching OR has neither full nor empty nor
proper-subcube support since it is nonconstant and each coordinate takes
both values among satisfying assignments.

In both orientations each factor after summing its local row bits has
absolute value at most seven. There are at most \(M_D\) nonisolated row
vertices and \(2^{M_D}\) column assignments. Thus
\(|F(-1)|\le 2^{M_D}7^{M_D}<2^{4M_D}<2^{p_t}\); reducing at native
precision preserves zero versus nonzero. Before reducing the same
support proof works at every odd \(t\), where \(c_t<0\), on the exact
rational zero-probe value. For small odd displacements it also survives
native reduction: if \(c_t=n_t/d_t\) in lowest terms with \(d_t>0\) odd,
then \(\max(|n_t|,d_t)\le(2L)^{|t|+1}\) by the factorial formula for the
grid derivative. Clearing \(d_t^{M_D}\) in the original deformation
yields an integer of bit length \(O(M_D(1+(|t|+1)(a+1)))\) since there
are \(2^{2M_D}\) terms each a product of \(M_D\) affine factors.
Whenever \((|t|+1)(a+1)=o(L/u)\), this is less than \(p_t\ge 24M_D L/u\)
(odd \(t\) has \(D_t=L\)). The parity of a positive even value can still
coincide with a target zero. High-precision equality and parity
therefore have different implications here.
\href{https://github.com/openai/math/blob/main/preprints/The-Exponential-Time-Hypothesis-is-False-October-25-2026/build/quantitative-contraction.tex\#L175}{D
grid kernel, 175--202} ·
\href{https://github.com/openai/math/blob/main/preprints/The-Exponential-Time-Hypothesis-is-False-October-25-2026/build/quantitative-reduction.tex\#L65}{Root
network orientations and original tensors, reduction 65--92} ·
\href{https://github.com/openai/math/blob/main/preprints/The-Exponential-Time-Hypothesis-is-False-October-25-2026/build/quantitative-contraction.tex\#L250}{Precision,
250--266}

For arbitrary \(t\), the partial wiring also gives cheap necessary tests
for the mod-two zero-mode candidate on a proposed layout. Every pin must
be matched. An unmatched old row pin fixes its sign, and an unmatched
old column pin sums out to one, erasing its sign. Each exterior pair
must have both pins in the same connected component of the partially
rewired network, because freezing other signs to one would otherwise
give rank at most one in the endpoints, against the OR matrix's rank
two. A quick cardinality exclusion is \((3M_D/L)(L-|t|)<N\), since at
most this many real row-column matches are possible and every unary port
must occur. Sufficient and computationally checkable data would be a
layout in which all components have \(O(\log N)\) edges and each local
parity table equals the product of its assigned exterior ORs.
Enumerating assignments and sign words component by component would
verify that certificate in polynomial time. No such layout is assumed.

\subsection{Arbitrary nonlinear decoding of a whole low-precision
coordinate}\label{v8-arbitrary-nonlinear-decoding-of-a-whole-low-precision-coordinate}

Fix one root D coordinate \(i<u\) and put \(D_t=L/2^i\), the number of
grid points of exact valuation \(i\). Give a decoder every complete root
slot at these \(t\), all allowed interpolation probes, and all allowed
mode tuples. The decoder may perform arbitrarily many arithmetic
operations with integral coefficients and then divide the assembled
output by \(2^{p-1}\) under a pointwise divisibility promise. More
generally allow any decoder whose final bit is determined by the input
residues modulo \(2^p\). This includes a polynomial with coefficient
denominators cleared by \(2^{p-1}\) and any odd denominators. Suppose \[
 1\le p\le a+1-i.
\] All nonzero mode tuples vanish modulo \(2^p\). Take the induced
matching of \(m\ge \gamma N/68\) in (N4). For \(\beta=0\) the product of
all children in a group union on the row side is a polynomial
\(P_t(z)\in\mathbb Z_{(2)}[2z]\) before evaluation at the probes. The
original root tensors are zero/one and every local \(W\) has (N7), so
local affine factors, summation over Boolean assignments, and products
of children have this property. Canonical lifts of the source residues
agree with this polynomial at each probe modulo \(2^p\). Its
coefficients are base scalars; the probe \(\zeta\) can still be an
arbitrary extension-field lift. For \(j<p\), sending the coefficient of
\((2z)^j\) modulo \(2^{p-j}\) suffices. Thus all root slot values are
determined by the column signs and at most \[
 B_{\rm msg}=D_t\,p(p+1)/2
\tag{N8}
\] bits of row information. Equal messages give identical rows of the
decoded Boolean target. The matching target has \(2^m\) distinct rows.
Therefore exact decoding requires \(D_t p(p+1)/2\ge m\). At the source
\(D_t\le 2^{2u}\le 2^{2\varepsilon}s^{2\varepsilon}\), because
\(u=\lfloor\varepsilon\lceil\log_2 s\rceil\rfloor\). Here
\(p=O(\log N)\), while \(m=\Omega(N)\), a contradiction for all
sufficiently large \(N\). The proof allows arbitrary nonlinear
postprocessing under the indicated congruence condition. If the decoder
also receives \(b\) bits of arbitrary one-sided information, add \(b\)
to (N8). Actual D Hermite reconstruction uses much larger inverses of
powers of two and combines different coordinates. Its denominator
analysis is essential.
\href{https://github.com/openai/math/blob/main/preprints/The-Exponential-Time-Hypothesis-is-False-October-25-2026/build/quantitative-contraction.tex\#L49}{D
exact lifting, 49--71} ·
\href{https://github.com/openai/math/blob/main/preprints/The-Exponential-Time-Hypothesis-is-False-October-25-2026/build/quantitative-contraction.tex\#L275}{Tapered
Hermite formula and denominator allowance, 275--345} ·
\href{https://github.com/openai/math/blob/main/preprints/The-Exponential-Time-Hypothesis-is-False-October-25-2026/build/quantitative-contraction.tex\#L301}{Coordinate-specific
node count, 301--313}

At parity, one may also expose selected \textbf{individual} child
contractions. Fix a single active \(t\), and in each group allow every
\(\beta\) whose surviving source moments reduce to at most one distinct
Boolean tile form for that child. All positive target exponents give
zero. A source mode of index \(b=0\) reduces to the parity of all column
bits in its tile. Every \(b\ge1\) instead gives the parity of the bits
at odd \(j\). Positive powers of the same form agree on Boolean
assignments. Apart from the 64 zero-mode group values, there are
therefore at most two further child bits per source tile, independent of
\(\zeta\) and the sizes of the positive exponents. A tile has a unique
source group. The number of tiles is at most \(6M_D/D_t\). Thus the
whole menu of these individual child residues requires at most \[
 64+12M_D/D_t
\tag{N9}
\] row bits, a bound that even counts the values wholly on the column
side. The full product menu for the entire coordinate at parity adds at
most \(D_t\) more bits by (N8). At the root \(D_t\ge 2^{u+1}\) and
\(u=\Theta(\log N)\), so both bounds together are \(o(N)\). Arbitrary
nonlinear decoding of these bits cannot produce the matching target. For
several \(t\) in the same coordinate, sum the bound in (N9) over those
nodes. General source-only \(\beta\) can request products of many
distinct tile parities, giving additional local moments. The number of
genuinely distinct moments in that larger interface, and their cost
under compressed exact computation, require further analysis.
\href{https://github.com/openai/math/blob/main/preprints/The-Exponential-Time-Hypothesis-is-False-October-25-2026/build/quantitative-contraction.tex\#L402}{D
tile definitions and multiplicity, 402--447} ·
\href{https://github.com/openai/math/blob/main/preprints/The-Exponential-Time-Hypothesis-is-False-October-25-2026/build/quantitative-contraction.tex\#L487}{Source
powers and child definition, 487--519}

\subsection{Entry, small-network, and complete-register
access}\label{v8-entry-small-network-and-complete-register-access}

There is a polynomial-cost positive construction at the literal
input-entry interface. The original MAIN unary factor returns
\(\Phi_{\rm pin}(1)=\sigma_v\) and \(\Phi_{\rm pin}(0)=1-\sigma_v\). The
original D unary tensor has \(T_{\rm pin}(1)=\sigma_v\). These addresses
are fixed. Group the MAIN reads by the chosen MAIN symbol and the D
reads by the chosen D symbol. Use the public fixed OR polynomial
\(x+y-xy\) within each group, take the prescribed triangle or
exterior-edge product locally, then take the two completed products and
multiply. All intermediate values of the Boolean OR-and-product
evaluation have constant bit length. This costs \(O(N)\) word accesses
and Boolean operations after graph-only preprocessing. If the IDs and
input tables are supplied in an unsorted sequential representation,
emitting a constant number of \(\Theta(\log N)\)-bit records per
incidence, sorting by group and edge in a mergesort, and scanning the
records gives \(O(N\log^2 N)\) elementary multitape bit work, with
ordinary binary encodings. This construction genuinely calls the
original immutable entry/factor descriptions. The two computed products
use additional arithmetic on the entry values and are not thereby
identified with \(\Sigma_g\) or \(Z_{g,t,\zeta,\beta}\).

If fresh small native instances may be submitted, the same edge
predicate is also the exact legal contraction \[
 \sum_{x,y\in\{0,1\}}T_u(x)\,T_v(y)\,\mathrm{OR}(x,y).
\] This is a two-edge bipartite network with the two unary leaves on one
side and the existing public OR table on the other. A constant parameter
size that meets the tensor theorem's fixed constraints, or its finite
direct-summation case, suffices. MAIN's constant-size direct Boolean sum
gives the same value for the two unary factors and one fixed OR. Realize
every edge this way and use the v7 group allocation. There are \(3N/2\)
calls of bounded bit cost apart from reading their original unary
descriptors. Such calls compose the existing local predicates on new
subinstances; none is automatically a descendant produced by the
immutable full-input root schedule.
\href{https://github.com/openai/math/blob/main/preprints/The-Exponential-Time-Hypothesis-is-False-October-25-2026/build/quantitative-reduction.tex\#L15}{D
legal inputs and local table convention, reduction 15--48, 65--94} ·
\href{https://github.com/openai/math/blob/main/preprints/The-Exponential-Time-Hypothesis-is-False-October-25-2026/build/main_full_version.tex\#L1675}{MAIN
direct polynomial interface, 1675--1704}

The remaining native-register question is precise. On an immutable,
certified full-input layout, do there exist nondegenerate MAIN fibers
and D slots or charged combinations of them that equal the two v7
factors (or yield the target through their cross product), and can their
\emph{values} be computed in polynomial time by exploiting internal
factorization or source-compatible arithmetic at the necessary
precision? The printed acquisitions, graph-fixed sparse private
enumerations, singleton/zero measurement specializations, nonzero mode
residues, unit-shift exact register, and single-coordinate low-loss
channel have the bounds above. A proof equating other specified
registers with the desired factors, a fast exact contraction for those
actual registers, or a source-specific obstruction to all of them is
still needed. The local-entry and fresh-small-instance constructions
settle their explicitly broader access models.

\subsection{Exact checks on small source
tensors}\label{v8-exact-checks-on-small-source-tensors}

The accompanying checker compares (N7) coefficient by coefficient with
direct construction of the Cauchy polynomial for \(a=1,2,3\). Across 258
distinct node/displacement choices all 3,720 coefficient positions agree
and have the stated two-divisibility bounds. A separate computation
takes a legal two-variable clause network with two unary pins, one
positive edge OR and the two directed implications between its
variables. Each variable has four ports, realized by two trivalent
equality vertices as prescribed in the fixed-density reduction. With toy
packet length four, 16 deterministic layouts in both orientations and
four sign words each give all 128 predicted zero/nonzero outcomes at
\(t=-1,\zeta=0\).

One explicit packet ordering in the output produces the exact values
\((-421038,\allowbreak 184437,\allowbreak 165807,\allowbreak -72657)\)
on sign words \(00,01,10,11\). Reduction modulo two gives exactly the OR
truth table. This example shows concretely why the
\nolinkurl{integer/native-precision} support obstruction cannot simply
be transferred to parity. The toy exercises the source tensor identities
and the slot-shift rule. It makes no assertion of an active printed root
or an asymptotic 64-symbol packing certificate at this small size.
\PaperSixEvidence{evidence/research_v8/native/check_native_escape.py}{Exact checker} ·
\PaperSixEvidence{evidence/research_v8/native/native_escape_checks.json}{Results and explicit packet
ordering}

\section{One counting family for both native
routes}\label{v8-section-common-native-family}

One fixed cubic expander supplies a genuine counting function for the
positive gate theorem, the enumerated MAIN cross-product theorem, and
the one-coordinate D information theorem. A choice of cubic expanders
also yields an information obstruction for the \textbf{combined menu of
all root coordinates} at the specified low precisions. The combined-menu
obstruction concerns the exact count and allows arbitrary joint
decoding.

\subsection{Source inputs and a common
matching}\label{v8-source-inputs-and-a-common-matching}

Take a simple cubic graph \(G\) of edge expansion \(\gamma>0\) on \(N\)
vertices. Use the shortened count-only physical CNF of F1. Each vertex
has five private copies, two for its availability clause and one for
each of its three graph-edge ports. Directed implications force equality
around its cycle. Each edge forbids choosing both endpoints. The
detached clause on two separate variables has three models. If the
independent availability signs are \(s_v\), the exact count over the
whole cube is

\[
 F_G(s)=3\sum_{I\in{\rm IS}(G)}\prod_{v\in I}s_v.                 \tag{C1}
\]

All switches disabled still gives three models. The physical formula has
\(n=5N+2\) variables, \(15N/2+1\) clauses and \(M_{\rm MAIN}=15N+2\)
incidences. The sparse D reduction gives \(M_{\rm D}=30N+4\) wires,
\(s_{\rm D}=n+\#\mathrm{clauses}=25N/2+3\), coefficient length
\(\ell=1\), and terminal precision \(P_0=n+1=5N+3\). These quantities
stay fixed as switches vary. MAIN uses its direct clause-factor entry
point and D uses its sparse clause branch.
\hyperref[v8-section-stronger-family]{Physical construction and positive
gate theorem} ·
\href{https://github.com/openai/math/blob/main/preprints/The-Exponential-Time-Hypothesis-is-False-October-25-2026/build/main_full_version.tex\#L1675}{MAIN
factor interface} ·
\href{https://github.com/openai/math/blob/main/preprints/The-Exponential-Time-Hypothesis-is-False-October-25-2026/build/quantitative-reduction.tex\#L65}{Sparse
D entry}

Fix a graph-only root D layout with its published collision certificate,
then choose one active coordinate \(i\). Both ports of an availability
clause belong to one old clause vertex and have the same coordinate
symbol. The relative density estimate of F2 bounds each symbol's share
of switches by \((1+o(1))(15+2/N)/64\). A union of groups therefore
holds between \(N/4\) and \(3N/4\) switches for all sufficiently large
sizes. At least \(\gamma N/4\) graph edges cross this cut. Choosing a
crossing edge in a cubic graph removes at most \(1+4+12=17\) candidates
at line-graph distance at most two. The resulting induced matching has
\(m\ge\gamma N/68\) pairs. Orient each pair from the chosen group union
to its complement.

\textbf{Fix the switches outside this matching to zero.} On the
remaining row switches \(x_j\) and column switches \(y_j\), the exact
physical count becomes

\[
                 3\prod_{j=1}^m(1+x_j+y_j).                       \tag{C2}
\]

Each factor has integer value matrix
\(\left(\begin{smallmatrix}1&2\\2&3\end{smallmatrix}\right)\). Modulo
two it is the identity. The entire physical parity matrix is exactly
\(I_{2^m}\), including its odd global factor. Consequently the matrix of
any integral characteristic-zero lift with the correct parity has rank
\(2^m\). This same restriction and row orientation serve both
comparisons below. \hyperref[v8-section-stronger-family]{Availability-label
density and count-only transfer block}

\subsubsection{MAIN crossed with a D
product}\label{v8-main-crossed-with-a-d-product}

Take a fixed sum of \(K\) products of the form in N4. In product \(h\),
MAIN forms each group by summing a graph-fixed list of \(b_{hg}\) native
private points with fixed sign-independent weights. It applies a fixed
linear readout to the single completed jet, of dimension \(\kappa_h=J\),
or \(\kappa_h=j+1\) when it reads only through coefficient \(j\). The D
operand is a product of group-local scalar functions at the chosen
coordinate, lifted separately to the common characteristic-zero field.
Fixed scalar weights are permitted. Require the assembled sum to be
integral and have the parity of (C1) at every input.

For one specific MAIN summand, the replacement arguments
\(\theta_1,\theta_2\) at the two physical switch ports are graph-fixed
jets once the private point, correction and representatives are fixed.
The switch factor is \[
      1-\theta_1\bigl(1-L_{s_v}(\theta_2)\bigr),
      \qquad L_1(w)=w,\quad L_0(w)=1-w .
\] The same \(s_v\) occurs in exactly this one factor. The factor's two
ports are owned by its factor group regardless of where their private
variable labels lie. All other clause factors are fixed at that point.
Choose one summand in each group. The resulting complete jet product
separates as a product of one row jet and one column jet across the
matching. A scalar linear readout of this product is a bilinear form of
rank at most \(\kappa_h\). Summing over private choices bounds its rank
by \(\kappa_h\prod_g b_{hg}\). Every D child at this coordinate owns the
tables of its original vertices, so its switches all belong to one side
of the chosen cut. A product of their separate lifts has matrix rank at
most one there. Multiplication entrywise and then addition over \(h\)
gives

\[
                    \sum_{h=1}^{K}\kappa_h\prod_g b_{hg}
                         \ \ge\ 2^m.                                \tag{C3}
\]

This proof includes no uncharged bypass. For an extra summand, its
actual rank on this same restriction adds to the left side. Repeated
factors may be collected before counting. If \textbf{one charged
execution} covers the \(K\) designated products and all designated
points in each list at rate at most \(W\) per paid unit, then
\(K,\max_{h,g}b_{hg}\le WT\). Thus
\(\log_2 T\ge (m-\log_2\kappa_*)/(r_M+1)-\log_2 W\), where
\(\kappa_*=\max_h\kappa_h\). The logical size of a sum computed by a
different contraction supplies no delivery certificate.
\hyperref[v8-section-native-escape]{Native rank and trace argument N3--N5}
·
\href{https://github.com/openai/math/blob/main/preprints/The-Exponential-Time-Hypothesis-is-False-October-25-2026/build/main_full_version.tex\#L1810}{Fixed
MAIN private representatives} ·
\href{https://github.com/openai/math/blob/main/preprints/The-Exponential-Time-Hypothesis-is-False-October-25-2026/build/main_full_version.tex\#L1964}{Source
group ownership}

The canonical lift of a completed D slot can replace the product of
separate lifts under the same error calculation as N4. With
\(M=M_{\rm D}\), \(h_H=\lceil3M/u\rceil\), \(a=2u\),
\(L_{\rm grid}=2^a\) and \(D_i=L_{\rm grid}/2^i\), the two lifts agree
modulo \(2^{p_t}\), where \(p_t=P_0+8h_HD_i+8M(a+1)\). For the ordinary
D readout the coefficient loses at most \(6h_HD_i+4M(a+1)\) dyadic
digits. For the ordinary MAIN readout on the direct binary clauses, put
\(d=M_{\rm MAIN}\) and \(\delta=(R-1)!Q!\). The MAIN tuple and its
origin readout lose at most
\(d\,v_2(\delta)+(J-1)\lfloor\log_2 Q\rfloor=O(N\log N)\), including
graph-fixed sublists with dyadically integral jet weights. Since
\(D_i\ge 2^{u+1}\ge s_{\rm D}^{\varepsilon}\) at every active root
level, \(2h_HD_i+4M(a+1)\) eventually dominates this MAIN loss.
Consequently the ordinary cross multiplier preserves equality through
\(P_0=5N+3\) when making the replacement. Additional denominators
require verifying \(p_t+v_2(b(s))\ge P_0\) for their actual multiplier
\(b(s)\) separately. \hyperref[v8-section-native-escape]{Source lifting and
both readout estimates}

\subsubsection{An arbitrary decoder at the same
coordinate}\label{v8-an-arbitrary-decoder-at-the-same-coordinate}

Give the decoder every \textbf{completed} D root slot at that
coordinate, for every active grid point \(t\) of valuation \(i\),
allowed mode tuple and interpolation probe. Suppose the input to the
decoder from a slot is its residue modulo \(2^{p_i}\), where
\(1\le p_i\le a+1-i\). Every nonzero-mode complete product vanishes at
that precision. At least one targeted coefficient in such a product
carries a factor of \(2^{a+1-i}\), and all other factors are integral.
The zero-mode kernel entries \(W_{xy}(t+\alpha z)\) belong to
\(\mathbb Z_{(2)}[2z]\) by N7. All original clause and equality tables
on this family are binary. On (C2), the product of the row-side children
for a fixed \(t\) therefore belongs to \(\mathbb Z_{(2)}[2z]\) before
evaluating at the probe. Its coefficients on \((2z)^j\), \(0\le j<p_i\),
are needed only modulo \(2^{p_i-j}\). Sending those base coefficients
suffices even for extension-ring probes. The receiver computes its own
side locally.

The reduction uses one row message of \[
                       D_i\,p_i(p_i+1)/2
\] bits for all the slots and probes together. Correct parity on (C2)
would require at least \(m\) bits since its \(2^m\) rows are distinct.
The displayed bit count is \(o(N)\) at the source schedule and
\(m=\Omega(N)\). Thus every decoder of that single complete
low-precision menu fails on some availability input. The conclusion
permits adaptive reads within the supplied menu.
\hyperref[v8-section-native-escape]{Kernel, targeted divisibility and
fixed-coordinate proof N7--N8} ·
\href{https://github.com/openai/math/blob/main/preprints/The-Exponential-Time-Hypothesis-is-False-October-25-2026/build/quantitative-contraction.tex\#L402}{Source
crossing-mode expansion}

The positive origin circuit for (C1), correct throughout the full cube,
independently requires \(2^{\Omega(N)}\) distinct binary gates by F3. On
(C2) alone, a positive circuit with \(O(m)\) gates computes the exact
count. The graph, layouts and terminal function can be common, and (C3)
and the one-coordinate decoder test can share the very same matching.
The gate lower bound and the interface obstructions quantify over
different programs and domains of correctness. Combining them does not
force the existence of a signed execution that materializes the
enumerated points. \hyperref[v8-section-stronger-family]{Positive scalar,
control and precision contracts F3--F4}

The same five-copy family also admits a topology on which arbitrary
joint decoding fails at low precision across all root coordinates for
the full exact count.
\hyperref[v8-section-pooled-menu-barrier]{All-coordinate theorem and proof}

\section{The entire low-precision root
menu}\label{v8-section-pooled-menu-barrier}

On a single fixed family of simple cubic expander topologies, the exact
number of models takes exponentially many distinct values as the
availability switches vary. A joint decoder given every completed
Appendix D root slot, across all active coordinates, probes and modes at
the low precisions specified below, receives only sublinearly many bits
of information about the switches. This gives an exact-count obstruction
for arbitrary joint postprocessing.

The physical input is the shortened five-copy CNF. Its terminal function
(C1) is three times the number of independent sets of the enabled
induced subgraph. For a cubic graph of order \(N\), its sparse D root
has \(M_{\rm D}=30N+4\) wires, size \(s_{\rm D}=25N/2+3\), coefficient
length one and terminal precision \(P_0=5N+3\). At the active root put
\(u=\lfloor\varepsilon\lceil\log_2s_{\rm D}\rceil\rfloor\), \(a=2u\),
\(L_{\rm grid}=2^a\), with the printed value
\(\varepsilon=1/(128\cdot1024)\). There are \(L_{\rm grid}/2^i\) active
grid points at valuation \(0\le i<u\).
\hyperref[v8-section-common-native-family]{Physical construction and
matching restriction (C2)} ·
\href{https://github.com/openai/math/blob/main/preprints/The-Exponential-Time-Hypothesis-is-False-October-25-2026/build/quantitative-contraction.tex\#L73}{Root
parameters}

\subsection{Encoding the complete pooled
menu}\label{v8-encoding-the-complete-pooled-menu}

Fix a certified graph-only root layout of a binary clause network with \(\ell=1\) before the availability input is chosen. At each active centre \(t\) of valuation \(i<u\), keep the original \(p_t\)-legal mode tuples and all \(M(2L_{\rm grid}-1)+1\) prescribed, residue-distinct binary-coordinate probes, where \(M\) is the root wire count. The observer receives the completed child products reduced to \(1\le p_i\le a+1-i\), without rebuilding a shorter menu, receiving individual child values, or receiving their entries. The interpolation-point count is not the Hermite truncation order \(q_t\). The graph, original addresses and fixed parameters are public.
\PaperSixEvidence{evidence/research_work/sources/quantitative-contraction.tex}{Archived source: precision and Hermite order, lines 250--267; probes, lines 338--346}. At each active \(t\), the whole zero-mode product, before probe
evaluation, is a polynomial in \(2z\) over \(\mathbb Z_{(2)}\). All
original tensors are binary, the local kernels satisfy N7, and each
defining zero-mode child contraction uses finite sums and products of
these entries. Reducing each entry at a supplied probe commutes with the
child's sums and products. Multiplying the resulting child residues
therefore gives exactly this particular completed slot modulo
\(2^{p_i}\). Every nonzero mode in a complete product has a targeted
coefficient with dyadic valuation at least \(a+1-i\), and all its other
factors are integral. It vanishes modulo \(2^{p_i}\). For the zero mode,
the coefficients on \((2z)^j\) are needed only modulo \(2^{p_i-j}\) for
\(0\le j<p_i\). These are base-ring bits even at an extension-valued
probe. \hyperref[v8-section-native-escape]{Kernel N7 and targeted
divisibility}

Thus the entire global menu of one input is determined by a binary
string of length

\[
       B=\sum_{i=0}^{u-1}\frac{L_{\rm grid}}{2^i}\,
                       \frac{p_i(p_i+1)}2
          < 2L_{\rm grid}\binom{a+2}{2}
          =O\bigl(s_{\rm D}^{2\varepsilon}\log^2 s_{\rm D}\bigr).  \tag{C4}
\]

The code in (C4) consists of coefficients of the \textbf{whole product}
at each node. It need not be efficiently computed from the input. It
upper-bounds the number of distinct whole menus by \(2^B\). It supplies
no private row message for a common cut. A proof about a single-bit
target would require information beyond its two possible output values.

\subsection{A cubic topology carrying exponentially many distinct exact
counts}\label{v8-a-cubic-topology-carrying-exponentially-many-distinct-exact-counts}

Chan, Heilman and Panova use rooted independent-set counts to encode
binary choices in caterpillar ratios and obtain many distinct scalar
counts from pairs of terminal counts. Their proof of Theorem 1.3 gives
the rooted operations in Lemma 3.1 and Corollary 3.2, and applies the
square-root counting step in equations (3.3)--(3.5). Section 6.6
identifies the caterpillars. Here the switch choices and both terminal
evaluations are realized within one expander host and its original
physical CNF. \href{https://arxiv.org/html/2604.19094}{Chan, Heilman and
Panova, §§3 and 6.6}

Take a caterpillar with spine \(v_0,\ldots,v_k\) and a private leaf
\(w_j\) adjacent to each \(v_j\), for \(k\ge1\). Its degree deficits
relative to three total \(d=2k+4\). Each leaf has two, and each spine
endpoint has one. Take a disjoint simple cubic expander \(H\) on \(d\)
vertices, which exists with uniform expansion \(\gamma_0>0\) at all
sufficiently large even orders. Join the \(d\) deficits to distinct
vertices of \(H\), then split each resulting degree-four core vertex
into two adjacent vertices, assigning two external edges to each. The
graph \(G_k\) is simple, cubic, and has \(N=2(k+1)+2d=6k+10\) vertices.
The only edges between caterpillar vertices are the stated tree edges.
\hyperref[v8-section-stronger-family]{Exact-order expander construction
used in F5}

To check expansion, partition \(G_k\) into clusters indexed by the
vertices of \(H\). Start each with the connected pair replacing its core
vertex. Assign \(w_j\) and \(v_j\) to the cluster at one of the two core
spokes of \(w_j\). Distinct leaves use distinct core vertices, so every
cluster is connected and has at most four vertices. Each original edge
of \(H\) still joins its corresponding two clusters. For a cut with
smaller side of size \(x\), let \(f,e,b\) count full, empty and partial
clusters on one side. Then \(\min(f,e)\ge x/4-b\). At least \(b\) edges
cross inside partial clusters. At least \(\gamma_0\min(f,e)-3b\)
original edges go from a full to an empty cluster, since a partial core
index has three original incident edges. Hence the number of edges
crossing the cut is at least \[
       \max\{b,\ \gamma_0 x/4-(\gamma_0+3)b\}
                       \ \ge\ \frac{\gamma_0 x}{4(\gamma_0+4)} .
\] This proves the uniform expansion needed in F3 and in the matching
argument above.

Disable all split-core switches, enable \(v_0,\ldots,v_{k-1}\) and
\(w_0\), and let the switch of \(w_j\) be \(s_j\in\{0,1\}\) for
\(1\le j\le k\). The only other switch needed below is \(v_k\). For a
prefix ending at \(v_j\) when its inclusion is permitted, write
\(A_j,B_j\) for the independent-set counts with \(v_j\) absent and
present, respectively. Then \[
       (A_0,B_0)=(2,1),\qquad
       A_j=(1+s_j)(A_{j-1}+B_{j-1}),\qquad B_j=A_{j-1}.           \tag{C5}
\] An absent \(v_j\) leaves \(w_j\) freely selectable when enabled. A
present \(v_j\) forces both \(w_j\) and \(v_{j-1}\) absent. All counts
in (C5) are positive. Put \(r_j=A_j/B_j\). Since \(r_0=2\) and
\(r_{j-1}>1\), the formula \[
                r_j=(1+s_j)(1+1/r_{j-1})
\] puts \(r_j\) in \((1,2)\) for \(s_j=0\), and in \((2,4)\) for
\(s_j=1\). Once \(s_j\) is determined, \(r_{j-1}=(1+s_j)/(r_j-1-s_j)\).
Therefore the \(2^k\) leaf words yield \(2^k\) distinct ratios \(r_k\),
and hence \(2^k\) distinct pairs \((A_k,A_k+B_k)\).

Disabling \(v_k\) gives the physical count \(3A_k\). Enabling it gives
\(3(A_k+B_k)\). Both are answers of the same fixed topology. Applying
the pair-to-scalar argument to its attained count set \(\mathcal S_k\)
gives \(|\mathcal S_k|^2\ge 2^k\). All these counts are less than
\(2^{P_0}\). \href{https://arxiv.org/html/2604.19094\#S3.E5}{Earlier
pair-to-scalar argument, equation (3.5)}

It follows that \[
              \log_2|\mathcal S_k|\ \ge\ k/2=(N-10)/12.            \tag{C6}
\]

\textbf{Pooled-coordinate consequence.} On the shortened CNF over
\(G_k\), an arbitrary decoder whose only input-dependent information is
the combined complete-slot menu in (C4) fails to recover the full exact
count on some availability word for sufficiently large \(k\). On the
fixed native layout, suppose any additional input-dependent information
has at most \(2^b\) possible states for each fixed menu, including the
supplied state of an input-dependent cache or control transcript.
Correctness on all inputs then requires \[
                         b+B\ \ge\ (N-10)/12 .
\] Indeed a fixed topology and layout give at most \(2^{B+b}\) possible
inputs to the final decoder. This accounting includes a fixed-length
\(b\)-bit string, or at most \(b\) adaptive binary answers with all
addressing and termination determined by the previous observations. An
independently conveyed stopping time or string length is part of the
side information. Equation (C6) requires at least \(|\mathcal S_k|\)
distinct decoder inputs. Since \(s_{\rm D}=\Theta(N)\) and the printed
\(\varepsilon=1/(128\cdot1024)\), (C4) is \(o(N)\). Adaptive selection
from the same menu cannot distinguish two inputs having identical entire
menus and side information.

The exact-count argument uses a different restriction from (C2), on the
same \(G_k\) and full function (C1). It gives an information
requirement. Charging acquisition, or allowing individual child
registers, the larger native precisions, additional deformations, or
registers from MAIN changes the proposed execution model. The printed
reconstruction uses its high-precision root data.
\href{https://github.com/openai/math/blob/main/preprints/The-Exponential-Time-Hypothesis-is-False-October-25-2026/build/quantitative-contraction.tex\#L250}{Source
precision and root schedule} ·
\href{https://github.com/openai/math/blob/main/preprints/The-Exponential-Time-Hypothesis-is-False-October-25-2026/build/quantitative-contraction.tex\#L487}{Definition
of the completed native slots}

Direct enumeration of the independent sets checked both terminal choices
for every leaf word through \(k=7\). At \(k=7\) it found 128 distinct
ordered pairs and 146 distinct scalar counts. A generic exact CNF solver
also checked all 28 terminal instances through \(k=3\), after
constructing all five-copy clauses on finite cubic completions. The
finite test uses cycle-and-chord cores. Uniform asymptotic expansion in
the theorem uses the stated expander cores and the cut proof.
\PaperSixEvidence{evidence/research_v8/checks/common_family_counts.py}{Checker} ·
\PaperSixEvidence{evidence/research_v8/checks/common_family_counts.json}{Saved finite results}

\section{The larger raw-entry closure and its
escape}\label{v8-section-native-closure}

An optimization model must specify which source values may be named and
which identities its replacement must preserve. Two operations in the
actual source expose the distinction. The direct rational factor
interface accepts a clause polynomial evaluated on an argument tuple.
The sparse Appendix D entry stores the clause satisfaction table on its
ports. Those operations are local. The printed programs also prescribe
how the local values enter the recovered answer. Access to the first
operations alone does not enforce the latter condition.
\href{https://github.com/openai/math/blob/main/preprints/The-Exponential-Time-Hypothesis-is-False-October-25-2026/build/main_full_version.tex\#L1667}{MAIN
lines 1667--1686} ·
\href{https://github.com/openai/math/blob/main/preprints/The-Exponential-Time-Hypothesis-is-False-October-25-2026/build/quantitative-reduction.tex\#L65}{D
reduction lines 65--97}

\textbf{Proposition C1 (local-table universality).} Fix an input
topology with independent switch bits \(s_1,\ldots,s_N\). Suppose each
switch has two graph-known legal local evaluations returning
respectively \(s_i\) and \(1-s_i\). Allow scalar addition and
multiplication, constants zero and one, reuse of registers, and free
wiring between previously computed registers. Let correctness be
required only of the final Boolean result. Every fan-in-two Boolean
circuit over AND and NOT on these switches has an implementation in this
model of size \(O(N+S)\), where \(S\) is its Boolean gate count. This
holds over the nonnegative integers and after reduction at \textbf{any
fixed working precision} \(2^p\), \(p\ge1\).

\emph{Proof.} Maintain a pair \(a,\bar a\) with values zero or one
summing to one, using the two evaluations for an input gate. Negation
swaps the two entries. From pairs for \(a,b\), the output pair for AND
is \[
       \bigl(ab,\ \bar a+a\bar b\bigr).
\] The summands in the second component cannot both be one. Thus both
outputs remain bits, even for correlated inputs or two uses of the same
register. A gate costs at most two multiplications and an addition. This
establishes the claim over the integers. Its gate values reduce to
exactly the same bits modulo \(2^p\). A fan-out edge just reads the
stored pair. If signed constants, subtraction, or a Boolean zero test
are available, one input rail suffices through \(1-a\) or the predicate
\(a=0\). Over \(\mathbb F_2\), \(1+a\) also suffices. \(\blacksquare\)

The two evaluations are actual original clause entries. For a unary pin
they are obtained at arguments one and zero. For the pin-free switch
\((\neg x\vee L_s(y))\), where \(L_1(y)=y\) and \(L_0(y)=\neg y\), the
argument tuples \((1,1)\) and \((1,0)\) give precisely \(s,1-s\).
Different occurrences in this clause are distinct ports prior to
contraction. On the pin-free source family, the second tuple is locally
a legal table address, even though equality makes it inconsistent with
the contracted vertex gadget. Whether an optimizer can address that
entry after contracting the gadget is a material interface choice.
\hyperref[v8-section-stronger-family]{Pin-free construction and positive
contraction}

The simulation extends to a nonnegative integer answer of \(b\) bits if
the output circuit supplies those \(b\) bits and the arithmetic
precision is at least \(b\). Form \(\sum_{j<b}2^j a_j\) at \(O(b)\)
additional word operations. At fixed precision an AND gate has constant
bit cost. At polynomially bounded working precision each simulated gate
costs polynomially many bit operations even if a full word must be
touched. No count oracle or public truth table of the final problem
occurs in the simulation.

\textbf{Corollary for the uncontracted CNF tables, conditional on the
source theorem.} Assume the supplied fixed-density counting algorithm
runs in its stated deterministic \(2^{o(N)}\) bit time on the switched
CNFs of Section F1, which have linear size at fixed graph degree. On
each fixed such topology and the inactive-request slice, the exact count
has a \(2^{o(N)}\)-size straight-line circuit over the nonnegative
integers with uncontracted clause entries as inputs. To see this,
hardwire the topology and fixed signs. A \(T\)-step bit machine can be
simulated by a Boolean circuit of size \(\operatorname{poly}(T,N)\) by
composing its bounded-configuration transition circuit for \(T\) steps.
This size is \(2^{o(N)}\) here. Two legal switch entries supply each
pair of input literals for C1. Apply C1 and assemble the \(O(N)\) output
bits by repeated doubling and addition. Thus the source upper bound
itself excludes a \(2^{\Omega(N)}\) lower bound for this output-only
raw-table circuit model. The exponential result in Section F3 starts
with the contracted tables and their nonnegative Möbius coefficients.
\hyperref[v8-section-stronger-family]{Switched CNF and Section F3} ·
\href{https://github.com/openai/math/blob/main/preprints/The-Exponential-Time-Hypothesis-is-False-October-25-2026/build/quantitative-reduction.tex\#L86}{Source
fixed-density bit bound}

\textbf{Converse for bounded-word machines.} If each supposed primitive
on \(B\) stored bits has a Boolean circuit of size
\(\operatorname{poly}(B)\), a fixed source-style straight-line program
with \(T\) such operations has a Boolean circuit of size
\(T\operatorname{poly}(B)\), plus its ordinary input/output wiring.
Compose those circuits in topological order and preserve fan-out. For
data-dependent indexing into \(m\) available registers one can pay
\(O(mB)\) per access by multiplexing; a bounded deterministic run can be
unrolled using its full explicitly bounded configuration. Exact rational
arithmetic requires a bit-size bound, and a contraction delivered in one
black-box call requires an implementation cost. This converse makes no
unit-cost assumption for either of them.

\textbf{Implication for the scientific question.} On any fixed topology
obeying C1, a gate lower bound for \emph{all} such replacements of the
final output is also a general Boolean-circuit lower bound for that
particular input-indexed function (up to the stated simulation cost).
This identifies a model boundary without assuming a complexity
conjecture. A superpolynomial bound for a Boolean decision function in
NP would in particular separate that function from polynomial-size
circuits. For a counting target the implication concerns circuits for
the answer bits. The decision consequence needs its own decision target.
A lower bound on uniform running time has its own uniform quantifiers.
Source restrictions such as preservation of intermediate contractions,
integral one-sided messages, specific nonlinear interfaces, or positive
formal lifts give different mathematical questions.

The older pinned graph count has the explicit Boolean circuit
\(\prod_{uv\in E}(1-(1-s_u)(1-s_v))\), or the same AND/OR DAG in C1. It
uses \(O(N+|E|)\) scalar operations and admits word-sized Boolean
intermediates. The simulation gives no fast evaluator for arbitrary
counting targets. It proves that the named set of local operations and a
final correctness check, by themselves, do not isolate the interpolation
algorithms. \PaperSixEvidence{evidence/research_v7/source/cost_contracts.md}{Original
target and terminal contracts}


\part{Actual native observations: reconstruction and acquisition}
We now specialize from modules of possible observations to literal probes
of the source algorithms. We distinguish separate registers from joint
evaluation, and mathematical decode correctness from charged operations.
On the four-clone family with subdivision equalities on the row side,
J9 compares two precisions of the completed-slot interface at $t=L-1$.
The paper's erasure argument covers every certified layout in that
orientation, whereas the upper bound uses a chosen layout. The checked
upper semantic certificate recovers all pins from a printed $10M+1$-probe
prefix at $10M$ bits for each singleton moment under its seating and
source guards. The separate lower certificate projects unchanged
high-precision modes of the seated clone file to the low receiver, up to
the stated $M/6-3M/L$ bound under its mask and preservation guards.
This is not one guard-free Lean theorem. The $N(10M+1)$ bound counts
request occurrences, not the cost of executing each request. It is
multiplied by 64 if all group receipts are acquired separately.
The jointly compiled first-coefficient readout is another interface;
the five-copy v8 witness and the guarded J2/J6 orientation remain separate.
The source statement names and hypothesis scopes are retained. The v8
scalar, large-residue jet and all-residue Newton equation tags carry an
explicit V8 prefix; the v9 contradiction assumption and layer identity
are labelled Q1 and J6-layer rather than the theorem names J1 and J2.

\section{Overview and principal
findings}\label{v9-overview-and-principal-findings}

The principal result of this revision is an executable native escape on
the common pinned CNF family used in the preceding rank-budget work. It
closes an acquisition gap in the earlier multiplication example. We
construct a legal, deterministic Appendix D layout, independent of every
pin sign, on which actual source-child moments recover all signs using
polynomial bit work. A second result recovers those signs from
polynomially many actual completed root slots at one node. Its required
first-coefficient readouts also admit a polynomial bit evaluator,
although separately evaluating every sampled completed slot has not been
shown polynomial. Both the old triangle factor and the old exterior
factor can then be computed from the recovered signs. They are obtained
through D information; we do not assert that one genuine MAIN register
equals the first factor.

This positive result accompanies a nonlinear adaptive obstruction on
another explicitly specified interface. MAIN diagonal pins restrict the
support of each genuine jet. A decoder that receives only those jets,
together with a guarded low-precision completed D menu, must make at
least \(2^{\Omega(N/r^{Q+1})}\) queries even with arbitrary arithmetic,
branching and caching between replies. The new reconstruction changes
the D orientation, its certified layout and the information available to
the decoder. It therefore supplies an escape from a proposed broader
obstruction, while leaving the restricted theorem intact.

\subsection{What is fixed, what varies, and what is
charged}\label{v9-what-is-fixed-what-varies-and-what-is-charged}

The frozen source has two distinct mechanisms. MAIN uses rational
polynomial rows and shared measurements to form local Taylor sums.
Appendix D deforms a bipartite Boolean tensor network using packetwise
Cauchy matrices, evaluates in a finite-precision unramified extension
ring, expands cross-group dependence into source moments and targeted
mode coefficients, and multiplies child contractions. The main
references in this revision are the archived files
\nolinkurl{main_full_version.tex},
\nolinkurl{quantitative-reduction.tex} and
\nolinkurl{quantitative-contraction.tex}. The public evidence catalogue
records historical paths, hashes and availability. We use the source's
printed asymptotic parameter schedule; finite checks at smaller
parameters validate exact identities but do not instantiate that
schedule.

A \textbf{child value} is one actual group contraction at a fixed
address. A \textbf{completed slot} is the product of all children for a
fixed node, probe and mode tuple. The \textbf{sum of modes} is a third
observable. The distinctions matter: a completed product may vanish
while its source child contains a useful moment. A \textbf{compiled
readout} is a specified linear combination or coefficient of the native
expressions evaluated through an algebraically equivalent joint
algorithm. It need not materialize the individual values in that
combination.

Our algorithms charge ordinary bit operations. Their integers, rational
numerators and denominators, extension coordinates and addresses all
have polynomial length in the physical input. A claim of polynomial
query count alone does not include acquisition cost. Our hardness
statements classify a uniformly prescribed evaluation problem under
standard counting reductions; they are not unconditional superpolynomial
lower bounds. A chosen certified layout is not the output of an
unspecified deterministic source tie-breaking rule. When a layout
carries the reduction's input incidence data, it is an explicit part of
the evaluation problem.

\subsection{Main results at a glance}\label{v9-main-results-at-a-glance}

\begin{OfflineResearchTable}
\OfflineResearchHeaderRule
\noindent%
\begin{minipage}[t]{0.31500\OfflineResearchTableWidth}\raggedright
Result
\end{minipage}\hfill%
\begin{minipage}[t]{0.31500\OfflineResearchTableWidth}\raggedright
Proven statement
\end{minipage}\hfill%
\begin{minipage}[t]{0.31500\OfflineResearchTableWidth}\raggedright
Essential scope
\end{minipage}%
\OfflineResearchRow
\nobreak
\OfflineResearchHeaderRule
\nobreak
\noindent%
\begin{minipage}[t]{0.31500\OfflineResearchTableWidth}\raggedright
Native moment reconstruction, J7
\end{minipage}\hfill%
\begin{minipage}[t]{0.31500\OfflineResearchTableWidth}\raggedright
At most \(N+64\) actual child values
recover all \(N\) pin signs, with polynomial acquisition and decoding.
\end{minipage}\hfill%
\begin{minipage}[t]{0.31500\OfflineResearchTableWidth}\raggedright
One graph-only certified layout; subdivision rows; a permitted probe set
containing zero.
\end{minipage}%
\OfflineResearchRow
\noindent%
\begin{minipage}[t]{0.31500\OfflineResearchTableWidth}\raggedright
Completed native reconstruction, J8--J9
\end{minipage}\hfill%
\begin{minipage}[t]{0.31500\OfflineResearchTableWidth}\raggedright
At most \(N(10M+1)\) completed
values at one node, read at \(10M\) bits, recover all signs. Their
required first coefficients have a polynomial joint evaluator.
\end{minipage}\hfill%
\begin{minipage}[t]{0.31500\OfflineResearchTableWidth}\raggedright
Any
prescribed unit-separated probe set of the source size; no per-sample
polynomial acquisition claim.
\end{minipage}%
\OfflineResearchRow
\noindent%
\begin{minipage}[t]{0.31500\OfflineResearchTableWidth}\raggedright
Sharp precision order, J9
\end{minipage}\hfill%
\begin{minipage}[t]{0.31500\OfflineResearchTableWidth}\raggedright
The same orientation erases all completed
information at that node through \(M/6-3M/L\) bits on every certified
layout; one layout permits recovery at \(10M\) bits.
\end{minipage}\hfill%
\begin{minipage}[t]{0.31500\OfflineResearchTableWidth}\raggedright
Single node and
completed-slot interface; gives order \(\Theta(M)\), not an optimal
constant.
\end{minipage}%
\OfflineResearchRow
\noindent%
\begin{minipage}[t]{0.31500\OfflineResearchTableWidth}\raggedright
Adaptive mixed query bound, J1--J2
\end{minipage}\hfill%
\begin{minipage}[t]{0.31500\OfflineResearchTableWidth}\raggedright
A real exact execution incurs
\(A2^{-k}+B2^{-k_0}\ge1\), with \(k=\Omega(N/r^Q)\),
\(k_0=\Omega(N/r^{Q+1})\); bounded-error extension included.
\end{minipage}\hfill%
\begin{minipage}[t]{0.31500\OfflineResearchTableWidth}\raggedright
Only the
declared MAIN jet queries and the erased D channel carry input
information.
\end{minipage}%
\OfflineResearchRow
\noindent%
\begin{minipage}[t]{0.31500\OfflineResearchTableWidth}\raggedright
Guarded native precision, P5--P11
\end{minipage}\hfill%
\begin{minipage}[t]{0.31500\OfflineResearchTableWidth}\raggedright
Every completed slot vanishes
through \((1-o(1))M/(2L)\) bits on a constructed ordering; a natural
random ordering gives \((1-o(1))M/(2L^2)\) bits with high probability.
\end{minipage}\hfill%
\begin{minipage}[t]{0.31500\OfflineResearchTableWidth}\raggedright
Subdivision columns; printed padded load; all signs at the fixed
order.
\end{minipage}%
\OfflineResearchRow
\noindent%
\begin{minipage}[t]{0.31500\OfflineResearchTableWidth}\raggedright
First-extra-digit witness, P10
\end{minipage}\hfill%
\begin{minipage}[t]{0.31500\OfflineResearchTableWidth}\raggedright
Increasing the old precision threshold
by one can expose \(\Omega(M)\) independent signs in the actual
completed menu.
\end{minipage}\hfill%
\begin{minipage}[t]{0.31500\OfflineResearchTableWidth}\raggedright
General binary tensors; the origin count is
constant.
\end{minipage}%
\OfflineResearchRow
\noindent%
\begin{minipage}[t]{0.31500\OfflineResearchTableWidth}\raggedright
Exact acquisition classification, A1--A4
\end{minipage}\hfill%
\begin{minipage}[t]{0.31500\OfflineResearchTableWidth}\raggedright
A prescribed actual register
can encode a counting-hard problem, including on read-once literal CNFs
with a supplied certified layout.
\end{minipage}\hfill%
\begin{minipage}[t]{0.31500\OfflineResearchTableWidth}\raggedright
Hardness of that requested value,
not of the easy original count or every execution.
\end{minipage}%
\OfflineResearchRow
\OfflineResearchHeaderRule
\end{OfflineResearchTable}

Here \(M=24N\) on the common four-clone family, \(r=64\) for D,
\(L=2^{2u}\), and \(u=\lfloor\varepsilon\lceil\log_2s\rceil\rfloor\)
with \(\varepsilon=1/131072\). MAIN uses its own alphabet \(r\) and word
length \(Q\); its \(r^{2Q}\) is at most a constant times \(N^{1/128}\).
The formulas using MAIN and D alphabets are local to their respective
sections. All asymptotic assertions concern sufficiently large inputs
under these fixed constants. The resulting finite threshold can be very
large.

\subsection{Why the positive result answers the earlier acquisition
question}\label{v9-why-the-positive-result-answers-the-earlier-acquisition-question}

The earlier multiplication example factored the count as \(U_MU_D\)
using group functions of pin signs. That was a statement about available
representations. It did not prove that those group functions could be
obtained from the original statistics. The new construction retains
every original CNF tensor and port. It places nonpin wires at the
endpoints of each packet and ensures that every pin belongs to a source
tile. At the extreme node, the zero-probe wiring then has a matching
structure. Actual source moments become power sums of normalized bit
marginals. Pins have deterministic marginals equal to their signs, while
the nonpin contributions are public. Invertible rational Vandermonde
systems recover the signs.

For completed products the corresponding nonzero modes vanish at the
zero probe. Their first derivative need not vanish. The targeted
equality derivative and the Cauchy coefficient derivative have the same
fixed sign for every term, giving a nonzero public multiplier of the
source moment. Finite-precision interpolation recovers this coefficient
from actual probes, including when zero is absent. Alternatively, its
exact finite derivative formula gives a polynomial-time compiled
readout. The proof accounts for its bit sizes and does not charge an
exponentially expanded expression as an execution.

There is still no general SAT algorithm improvement here: the witness
family itself is easy when its raw pins are read. The contribution is a
source-faithful acquisition and execution result showing that this
family cannot certify the proposed stronger nonlinear barrier on the
newly covered native interface.

\subsection{Relationship to the preceding
theory}\label{v9-relationship-to-the-preceding-theory}

The v7/v8 results retain their stated matrix, observable-module,
finite-precision, and representation hypotheses. This revision does not
replace those theorems by a runtime conclusion. Instead it tests the
information they expose against actual native value acquisition. The
diagonal-support theorem gives an unconditional query obstruction
without restricting the decoder to linear or positive computation. The
moment algorithm shows that a neighboring native interface can defeat
it. The acquisition reductions show why asking for a particular
large-precision value may itself be difficult, while a final-answer
program may avoid that value.

We preserve the distinction between exact counting and SAT decision. On
the pinned family the count is zero or one, so a query theorem for that
interface also distinguishes satisfiability there. The counting-hard
register reductions do not turn into a SAT decision lower bound. Nothing
in this revision proves \(\mathrm P\ne\mathrm{NP}\), invalidates the
source's upper bound, or excludes arbitrary optimized signed hybrid
computation.

\section{Synthesis and exact interfaces}\label{v9-section-synthesis}

This chapter separates three questions about a native interpolation
interface: which distinctions its values retain, whether those values
can be obtained economically, and whether a final-answer program must
use that interface. They require different proofs. Throughout, a public
layout is fixed before the tensor entries or pin signs vary. A layout
supplied as part of a uniform evaluation problem is explicitly
additional input.

The v9 outcome includes an executable escape on the earlier common
pinned family. The native-moment construction, proved separately, should
therefore take precedence over an attempt to declare a universal hybrid
lower bound. The guarded and diagonal obstructions remain valid on their
stated interfaces. They do not exclude the new access channel or its
layout. \hyperref[v9-section-moment-escape]{Native moments} ·
\hyperref[v9-section-joint-algorithms]{Restricted queries} ·
\hyperref[v9-section-precision-frontier]{Precision}

\subsection{The common physical input}\label{v9-the-common-physical-input}

Let a simple cubic graph be obtained by replacing every vertex of a
simple cubic base graph by a triangle. Its triangle edges partition into
disjoint triples and its exterior edges are a perfect matching. Write
\(N\) for its number of vertices. At vertex \(v\) put a unary pin whose
sign \(\sigma_v\) fixes its Boolean value, and put a positive
two-literal OR on every graph edge. Replace the four occurrences of each
variable by four private clones tied by a directed implication cycle.
Each clone occurs three times; the total number of private variables is
\(4N\) and the number of literal incidences is \(12N\). Every satisfying
assignment is forced by its pins. Consequently the exact count is

\[
 C_G(\sigma)=U_M(\sigma)U_D(\sigma),\qquad
 U_M=\prod_{\{v,w\}\text{ internal}}(\sigma_v\lor\sigma_w),\quad
 U_D=\prod_{\{v,w\}\text{ exterior}}(\sigma_v\lor\sigma_w).
\]

The sparse tensor compiler replaces each literal incidence by a path
through a two-port equality. Its \(M=24N\) physical wires join old
clause/clone vertices to subdivision equalities. Either side may be
chosen as rows in the bipartite contraction theorem. This choice affects
native values and is recorded in each result. These graphs were useful
for simultaneous rank lower bounds in v7. The count on the full pin cube
also has an elementary linear scan once the raw signs are available.
None of our claims treats this family as an unrestricted hard SAT
family. \PaperSixEvidence{evidence/research_v7/reviews/hybrid_independent.md}{Original
common construction and expansion}

\subsection{A nonlinear adaptive obstruction on a specified native
interface}\label{v9-a-nonlinear-adaptive-obstruction-on-a-specified-native-interface}

The following elementary principle is independent of arithmetic circuits
and remains valid for arbitrary internal computation.

\textbf{Support-cover theorem.} Suppose each public pattern \(w\) in a
nonempty finite set \(W\) has two legal inputs \(x(w,0),x(w,1)\) whose
required answers differ. A native query \(a\) has a publicly designated
response \(0_a\) and a support set \(S_a\subseteq W\): outside \(S_a\),
both inputs give exactly \(0_a\). Initial state and query choices may
depend only on public information and earlier replies. A deterministic
exact decoder has a real execution satisfying

\[
                  \sum_{a\text{ requested}}|S_a|\ \ge\ |W|.       \tag{S1}
\]

The request that reaches the bound is charged before its answer.
Repeated cached answers give no new coverage. Queries whose answer is
identically public may be inserted for free.

\textbf{Proof.} Follow the designated responses. Before the accumulated
support sizes reach \(|W|\), at least one pattern lies outside the union
of all supports requested so far. Both corresponding witnesses follow
this real prefix. It cannot halt, since it has the same transcript and
two different correct outputs. If the next request first reaches the
bound, a witness realizing the preceding prefix pays for that request.
If this never happens, the designated path remains real at every finite
prefix and an exact finite worst-case decoder cannot terminate. This
proves (S1) without charging an unreachable suffix.

For a randomized decoder capped at \(T\) actual requests, with
\(|S_a|\le c\), choose \(w\) uniformly and its witness bit fairly. For
any fixed random seed, truncate the abstract tree to depth \(T\). It
agrees with the actual computation on every legal input. Its designated
path covers at most \(cT\) patterns. Outside this union both equally
likely answers have the same transcript. The seed's success probability
is at most \(1/2+cT/(2|W|)\). Average over the random seed. Worst-case
error at most \(\epsilon<1/2\) requires \(cT\ge(1-2\epsilon)|W|\). This
is a direct finite-distribution argument; the standard adversary method
is not a new contribution.
\PaperSixEvidence{evidence/research_v9/formal/DiagonalQueryTree.lean}{Actual-tree Lean proof}

In the MAIN construction, equal full endpoint words give a literal
singleton row in the exceptional measurement basis and an identically
constant interpolation row. At each requested measurement address, a
diagonal pin must therefore agree with one prescribed sign.
Deterministic, source-certified word choices give an independent set of
\(k=\Omega(N/r^Q)\) such pins, with at least \(k_0=\Omega(N/r^{Q+1})\)
in every node/group. Reserve an unrelated triangle. On each selected
word, either all remaining signs are one, or two adjacent reserved signs
are zero. These are the two witnesses: \(U_D=1\) on both and \(U_M=C_G\)
changes from one to zero.

For a completed MAIN jet \(|S_a|\le1\); for an individual group jet
\(|S_a|\le2^{k-k_0}\). The first is a whole jet at a single fixed
address, not a table of all addresses. Thus an exact computation incurs

\[
                   A2^{-k}+B2^{-k_0}\ge1,                       \tag{S2}
\]

where \(A,B\) count actual completed and individual group requests. A
source-certified D ordering with subdivision columns makes every
completed D slot zero through \(q=(1-o(1))M/(2L)\) bits. Freely
supplying that whole menu does not change (S2), so it is a genuine
mixed-interface theorem on one family and one pair of fixed layouts.
Rational operations, digit tests, comparisons, caching and arbitrary
nonlinear computation on the answers are allowed. The D side does not
include arbitrary individual child values or all printed precision
levels. \hyperref[v9-section-joint-algorithms]{Detailed construction and
proof}

\textbf{Optimal size of this diagonal mechanism.} Write \(K=r^Q\). The
source incidence certificate gives at most \(2M_M/K^2\) incidences in
any full ordered cell, where \(M_M=12N\). Exactly \(K\) cells are
diagonal, and exactly \(K/r\) are diagonal with a specified symbol at a
specified coordinate. Hence, on every layout meeting this cell bound,

\[
 \#\{\text{all diagonal pins}\}\le\frac{2M_M}{K},\qquad
 \#\{\text{diagonal pins in }(t,g)\}\le\frac{2M_M}{rK}.          \tag{S3}
\]

Sum the per-cell bound over these cells; each pin contributes one of the
counted incidences. The construction attains both scales to constant
factors. More diagonal pins alone cannot improve the resulting exponent
beyond these scales. This says nothing about other lower-bound methods.

\textbf{A precise input-channel escape.} On the two-witness promise used
above, one raw reserved pin entry reveals which of the two correct
answers is required. On the entire pin cube, reading all \(N\) signs and
scanning the \(3N/2\) graph edges computes \(C_G,U_M,U_D\). Thus (S2)
cannot survive unrestricted input-dependent preprocessing. The native
moment construction goes further: it obtains all signs through genuine
child statistics on a different legal layout, with a polynomial bit-cost
evaluator for those statistics. This eliminates the earlier
arbitrary-group-function acquisition gap for that family and access
channel.

\subsection{Erasure, denominators, and
computation}\label{v9-erasure-denominators-and-computation}

Let \(R=\mathbb Z_{(2)}\). Its elements have odd denominators and its
residue field is \(\mathbb F_2\). The native expansion has

\[
 Q_{t,\beta}(z)=2^v z^k\sum_{j\ge0}c_j(2z)^j,
 \quad v=\sum(\delta_i+b)\beta_{\tau,b},\quad k=|\beta|.
\]

If \(K\) original equality columns have no real shifted port at this
node, a coefficient of residual degree \(j\) can involve at most \(k+j\)
of their independent common bits. The sum of a Boolean monomial
supported on \(h\) bits over \(K\) bits is exactly \(2^{K-h}\). It
follows that

\[
 2^{(K-k-j)_+}\mid c_j,
 \qquad Q_{t,\beta}(\zeta)\in 2^{v+(K-k)_+}R_{\rm ext}.        \tag{S4}
\]

The statements use actual assignment sums, including all source moments.
They permit cancellation and do not apply to an isolated source child.
Since \(v\ge k\), all completed values vanish through \(K\) bits. The
full native proof also bounds the joint coefficient alphabet, rather
than counting the same information repeatedly at different extension
probes. \hyperref[v9-section-precision-frontier]{Precision theorem} ·
\PaperSixEvidence{evidence/research_v9/formal/BooleanSummation.lean}{Boolean-sum Lean proof}

One can connect this to denominator budgets without inferring running
time. Suppose a finite vector of observations \(Q(x)\in R^m\) takes
congruent values modulo \(2^q\) on two inputs, and a rational polynomial
\(H\) reconstructs different output parities. Let \(\tau\ge0\) be a
common pole bound, so every coefficient of \(2^\tau H\) lies in \(R\).
Then

\[
                            \tau\ge q.                        \tag{S5}
\]

Indeed every integral polynomial preserves congruence modulo \(2^q\). If
\(\tau<q\), then \(H(Q(x))-H(Q(x'))\in2^{q-\tau}R\subseteq2R\), contrary
to the parities. Equivalently, subtract each pair of equal-degree
monomials and telescope one coordinate at a time. The same argument
covers an affine decoder with a public constant term. It is independent
of the polynomial degree, but depends on the integral input coordinates
and the stated common coefficient bound. An arbitrary digit-test decoder
is not represented by this assertion. A factored cancellation, division
or fused evaluator can avoid materializing a large intermediate; (S5) by
itself prices no arithmetic gates.

On guarded physical layouts, (S4) yields large erasure precision. On
another general binary tensor interface, the very first extra digit
above \(\delta_i\) already reveals \(\Omega(M)\) independent clause
signs. On the common pin family, the matched-layout moment readout
recovers signs at linear precision, while its completed outer-node menu
is identically zero at sufficiently small linear precision. These
statements concern different input classes or different layers of the
same interface. They are compatible, and rule out a universal
extrapolation of the old logarithmic-precision estimate.

\subsection{An exact value can be hard without being
necessary}\label{v9-an-exact-value-can-be-hard-without-being-necessary}

The signed mesh transfer at \(t=-1\) is

\[
 K_{Y,X}=l_Y(-X)=\begin{pmatrix}1&2\\0&-1\end{pmatrix}.
\]

Pairing two such ports with one auxiliary pin realizes every Boolean
function:

\[
 \sum_{y,z\in\{0,1\}}b(yz)l_y(-x)l_z(-1)=b(x).
\]

This identity holds over every commutative ring, lifts to arbitrary row
tables and then to a joint network with overlapping variable scopes. On
general binary tensors, a fixed graph and packing can therefore have a
prescribed native register that is counting-hard as its input table bits
vary. On literal sparse-CNF inputs, a different reduction uses a
\emph{supplied} certified layout to encode a planar bipartite edge-cover
instance. Its actual register equals

\[
       \#\operatorname{EdgeCover}(H)\prod_c(1+2^{d_c}),
       \qquad d_c\in\{0,1,2\}.
\]

The positive result fits the printed precision and is recovered before
any even factor is divided out. The original formula has disjoint
clauses and an immediately computable count. The reduction proves
uniform complexity of requesting the specified value, not an obligation
to request it, a lower bound on its source-chosen layout, or a lower
bound on the formula's count. The full acquisition theorem is
conditional on polynomial-time original entries and explicitly lists the
distinctions between one completed mode, a child, and the sum of all
modes. \hyperref[v9-section-acquisition-cost]{Acquisition proofs} ·
\PaperSixEvidence{evidence/research_v9/formal/PhysicalAcquisition.lean}{Finite network Lean identity}

This gives the operational conclusion of v9. A representation lower
bound, an observation obstruction, and a difficult uniformly prescribed
value are separate mathematical statements. The new native
reconstruction shows that some previously plausible information barriers
are bypassed within an actual child interface. Establishing a full
runtime barrier still requires a restriction on legal information
acquisition and a proof that every permitted exact execution pays for
it. We have not established that result for the source's optimized
full-precision mixed computation.

\section{An executable native moment
escape}\label{v9-section-moment-escape}

On the same four-clone physical CNF of a triangle-replacement cubic
expander, a deterministically constructible, sign-independent certified
Appendix D layout admits two exact recoveries. First, at most \(N\)
specified individual source-child singleton moments plus 64 baseline
children recover all \(N\) pins, with polynomial acquisition
\textbf{and} decoding. Second, polynomially many specified
high-precision \textbf{completed root slots} recover all pins with
polynomial postprocessing. Their required first-coefficient linear
combinations can be evaluated jointly in polynomial time from the input.
Separately materializing each completed slot at every nonzero probe is a
different task. All original CNF tensors and all original ports stay in
the input. Once the signs are recovered, either v7 factor, or their
product, follows by its edge list. Both results use the ordinary sparse
clause conversion, with subdivision equalities as rows and old variable
and clause vertices as columns.
\href{https://github.com/openai/math/blob/main/preprints/The-Exponential-Time-Hypothesis-is-False-October-25-2026/build/quantitative-reduction.tex\#L65}{Source
sparse clause conversion} ·
\href{https://github.com/openai/math/blob/main/preprints/The-Exponential-Time-Hypothesis-is-False-October-25-2026/build/quantitative-contraction.tex\#L487}{Genuine
group and root register definitions} ·
\PaperSixEvidence{evidence/research_v7/reviews/hybrid_independent.md}{v7 family}

\subsection{One graph-only certified
layout}\label{v9-one-graph-only-certified-layout}

Fix the retained coordinate \(i=0\), with \(r=64\), \(M=24N\) physical
wires, \(a=2u\), \(L=2^a\) and \(\mathcal D=r^{2a}=L^{12}\). Each unary
pin has just one physical wire, to a subdivision whose other wire ends
at the trivalent equality column for that pin's private clone. Draw the
symbol \(A_v\) of the old pin column uniformly, using one common symbol
on the two pins of each exterior matching edge as prescribed in v7.
Conditional on \(A_v\), give the pin subdivision row a symbol \(S_v\)
uniform on the \(r-1\) other symbols. Give the private clone column
symbol \(B_v\) with probability \(2/r\) of \(B_v=S_v\) and otherwise
uniformly among the other symbols. Conditional samples use independent
private coins for different pins. All other coordinate-zero symbols and
every vertex symbol in every other coordinate are independent uniform
symbols. The cited old vertices, private clones and subdivisions are
distinct except for the deliberate exterior symbol sharing. This is a
distribution on graph labels alone.

Every individual symbol is uniform. The pin wire \((S_v,A_v)\) never has
equal symbols, so it has an actual mode tile at \(i=0\) regardless of
its eventual packet seat. On the paired wire \((S_v,B_v)\) the equality
probability is \(2/r\). For a fixed equal ordered pair, the two wires'
combined probability is \(0+2/r^2\). For a fixed unequal pair it is \[
          \frac1{r(r-1)}+\frac{r-2}{r^2(r-1)}=\frac2{r^2}.
\] All other wires join independently uniform endpoints in this
coordinate: in particular, the other two incident wires of each affected
clone have independently labeled subdivision rows. The other coordinates
give uniform full ordered-cell marginals, independently of this one.
Thus every full cell has expected total load \(M/\mathcal D\). Its
expected number of pin wires is at most \(\frac{r}{r-1}N/\mathcal D\).

Let \(\Phi,\Phi_{\mathrm{pin}}\) be the respective sums of squared
deviations of total and pin-only loads from their actual cell means.
Their combined expectation is \(O(M)\). Indeed one may group each pair
of exterior-pin gadgets together with its shared symbol and independent
private offsets as a single block of random data. A block meets only
constantly many wires by bounded degree. All other independent vertex
words also meet only constantly many wires. Two physical wires with
disjoint underlying blocks have independent cell indicators. For a wire
let \(I_e\) be its cell-indicator vector and \(p_e=\mathbb E I_e\) its
own mean. The centered pair contribution is
\(\Pr[\mathrm{cell}(e)=\mathrm{cell}(f)]-\langle p_e,p_f\rangle\);
independent pairs contribute zero even when their individual laws are
biased. In the sum of these covariances over cells, a dependent ordered
pair contributes at most one; there are \(O(M)\) such pairs. This
reasoning applies to the pin-only sublist as well. Derandomize the
combined potential by conditional expectation, fixing symbols and the
independent offsets in turn. For example implement
\(S_v=A_v+d_v\bmod r\) with uniform nonzero offset, and
\(B_v=S_v+e_v\bmod r\) with the stated offset law. A pair of wires
involves only a constant number of primitive words and offsets, each
with polynomial domain size and exact polynomial-bit probabilities. This
makes each fixing a polynomial calculation.

The resulting \(\Phi=O(M)\) satisfies the printed \(\Phi\le 2M\Delta\)
eventually. All cells have load \((1+o(1))M/\mathcal D\), with common
padded load \(H_{\rm pad}=(1+o(1))M/\mathcal D\), since
\(M/\mathcal D\gg\sqrt M\) and \(L=o(M/\mathcal D)\). Each cell contains
at least \((1-\frac{r}{24(r-1)}-o(1))M/\mathcal D\) \textbf{nonpin}
wires. Apply the bounded-degree matching construction of J4 just to
these nonpin wires. Independently mark fair coins, retain edges whose
mark is the only mark at their endpoints, and fix the marks by the sum
of the cellwise squared deviations of retained counts from their means.
The survival probability is at least \(2^{-(2d_0-1)}\) at physical
maximum degree \(d_0\), and the sum of variances is \(O_{d_0}(M)\) by
line-graph distance-two independence. We obtain one matching with
\(\Omega(H_{\rm pad})\) nonpin wires in \emph{every} ordered cell, by a
deterministic graph-only polynomial algorithm.
\hyperref[v9-section-joint-algorithms]{Matching construction J4} ·
\href{https://github.com/openai/math/blob/main/preprints/The-Exponential-Time-Hypothesis-is-False-October-25-2026/build/quantitative-contraction.tex\#L114}{Packing
and padded cells, lines 114--171}

Distribute the orientation masks equally often in each unordered packet
type. Each full ordered cell then has \(H_{\rm pad}/L\) seats at each of
positions zero and \(L-1\), with a consistent orientation in each packet
as in P9. Fill all these seats from distinct wires of the matching in
their own cells, then fill every other position with the remaining real
or padding wires in its cell. There are enough seats and matching wires
eventually. Every endpoint slot is real, no pin is in either endpoint
slot, and no old row or column vertex has more than one matched port at
\(t=L-1\). Every pin remains in an unequal-symbol packet. The
construction retains the designated exterior-pair pin ties and uses the
printed physical input size and \(P_0\). It chooses actual legal source
parameters and an admissible order, with no assertion about the original
author's unspecified tie-breaks.
\hyperref[v9-section-precision-frontier]{Endpoint seating proof P9}

\subsection{J7. Polynomial acquisition from individual source
children}\label{v9-j7-polynomial-acquisition-from-individual-source-children}

Write \(t=L-1,\ \zeta=0,\ c=c_t=-\binom{2L-1}{L}\). The true kernel is
the partial shift from zero to the last position of each packet with
coefficient \(c\). Each subdivision row has zero or one matched port.
Its multilinear equality contributes one when unmatched and \(1-cX_y\)
when matched to column bit \(X_y\). Original column tensors are
therefore independent in the partial contraction. Define the local
nonnegative integer weight of an original column \(v\) by multiplying
its table by these unary factors on matched ports. Write \(z_v>0\) for
its local weight sum and \(m_y\) for the local mean of port bit \(X_y\)
in this weighted table. Positivity follows from \(c<0\) and the
existence of a satisfying local assignment to every original column. A
pin column is unmatched, has \(z_v=1\), and has \(m_y=\sigma_v\). Every
nonpin \(z_v,m_y\) is graph-only and computable on a bounded table. For
each group \(g\), the complete zero-mode \textbf{child} therefore has \[
                           Z_g=\prod_{v\in g\ {\rm column}}z_v>0.
\]

Fix a real source tile \(\tau\) containing \(d\ge1\) pin ports, with
source group \(g\). Its real columns have indices \(j_y=(y-\tau)/2\)
within this packet, distinct nonnegative integers below \(L/2\). For
\(b=0,\ldots,d-1\) use the genuine tuple with a single exponent one on
this tile's mode \(b\). It is legal: \(b+1\le L/2<p_t\). In the source
child no positive target exponent occurs. Its immutable packet tensor
multiplies by exactly \(\sum_y j_y^bX_y\), with \(0^0=1\), and the rest
of this child uses the ordinary zero-probe wirings. Its actual value,
before reduction, is consequently \[
                    S_{\tau b}=Z_g\sum_{y\in\tau} j_y^b m_y .              \tag{J7}
\] This follows by the ordinary finite sum over all the child's original
column assignments after summing the original row tensors. It allows
\(d=1\) and padding holes. The \emph{completed product} of this same
nonzero-mode tuple at zero vanishes because its other, targeted child
has a zero mode coefficient there.

All inputs to (J7) and its value have polynomial-bit exact evaluators.
One may sum each old bounded table, form \(Z_g\), and compute the left
side as
\(\sum_y j_y^b(\sum_{X_v} X_y\,{\rm wt}_v)\prod_{v'\ne v}z_{v'}\). This
computes an actual nonnegative integer without division. With at most
\(k_t\le3M/L\) matched ports and \(|c|<2^{2L}\), each
\(Z_g\le2^M(1+|c|)^{k_t}\) has \(O(M)\) bits. Each moment gains at most
\((b+1)\log_2(L+1)+O(1)=O(L\log L)=o(M)\) bits. The real source
precision \(p_t\ge24ML/u\) thus gives the integer exactly. The
coefficients \(j_y^b\), products and all requested addresses have
polynomial bit length. Baselines can be acquired at 64 child addresses,
or computed from the public weights.

The decoder divides \(S_{\tau b}\) by \(Z_g\) over the rationals and
subtracts the public contributions of all nonpin ports. For the \(d\)
distinct remaining indices it obtains the power sums with unknowns
\(\sigma_v\) at orders \(0,\ldots,d-1\). Writing the pin positions as
\(v_1,\ldots,v_d\) with indices \(j_{v_h}\), the Vandermonde determinant
\(\prod_{h<k}(j_{v_k}-j_{v_h})\) is nonzero. Its possible even factors
cause no loss because the actual integers were first recovered at native
precision. Exact linear algebra takes polynomial bit time by determinant
bounds. Summing \(d\) over the tiles gives \(N\). This establishes
polynomial acquisition and decoding from \(N\) genuine full-input
individual child moments. It uses an allowed probe set containing zero.
With a fixed set omitting zero, the source-only child polynomial can be
interpolated at zero from its prescribed evaluations; that simulation
alone does not establish polynomial acquisition of each nonzero-probe
child.

\subsection{J8. Polynomially many completed full-input
probes}\label{v9-j8-polynomially-many-completed-full-input-probes}

Use the very same graph-only layout and odd node. For each chosen
\((\tau,b)\), let \(Q_{\tau b}(z)\) denote the \textbf{whole completed
product} for the singleton tuple as a function of the local probe. Write
\(g'\ne g\) for its target. Its first coefficient in \(z\) equals a
positive, graph-only rational multiple of the power sum in (J7): \[
             [z]Q_{\tau b}(z)=D_{\tau b}\sum_{y\in\tau} j_y^b m_y,
             \qquad D_{\tau b}>0 .                                     \tag{J8}
\] Here the scalar is distinct from the grid-size parameter. To prove
the equality, fix a target row port \(x\) of this packet crossing to
\(\tau\). At \(z=0\) its target coefficient is zero. Its first
derivative is \[
          A'_x(0)=c\,\alpha\, 2^b/(t-\tau+x)^{b+1}<0.
\] Indeed at \(i=0\) a crossing row has the same parity as \(\tau\). Its
denominator is an odd positive integer because
\(t=L-1,\ \tau\in\{0,1\}\), and the root size grows. At least one real
target row exists: the pin wire in this tile also has its real row at
the identical position, with an unequal symbol and the required parity.

At first order every term uses exactly one such target row and all other
local forms at zero. For the subdivision owning this row port, replace
the baseline equality by its partial derivative in that leg, \(-1+2w\),
where \(w=0\) or \(c X\) is the other leg's baseline argument. Leave the
remaining subdivision rows unchanged. Summing this modified target child
gives an integer \(T_x<0\). The replacing factor is strictly negative,
other row factors are strictly positive, and each old column has at
least one satisfying local assignment. All matching ports are nonpins,
including any occurrence in \(w\). Target-child unary pins only sum to
one. Thus \(T_x\) is graph-only. The modification removes a unary weight
when the selected leg was the matched one, replaces a unary weight when
the other leg was matched, or gives \(-1\) when neither leg was matched.
Consequently it preserves vertex-wise computability at bounded degree.
Since the source and target are distinct groups, differentiation of
their product at the point where the targeted coefficient vanishes
yields precisely \[
                D_{\tau b}=
                   \left(\prod_{h\ne g'} Z_h\right)
                   \sum_{x\ {\rm target}} A'_x(0)\,T_x>0 .
\] The sum is over the actual eligible real row ports in this packet.
Distinct eligible ports of the same subdivision are distinct allocations
and have just the indicated one-port derivative. This derivation
preserves the native child sums and their assembly. No cancellation
assumption replaces the sign proof.
\hyperref[v9-section-precision-frontier]{Exact native target coefficient
and finite allocations, P1--P3} ·
\href{https://github.com/openai/math/blob/main/preprints/The-Exponential-Time-Hypothesis-is-False-October-25-2026/build/quantitative-contraction.tex\#L418}{Source
mode formula and child ownership, lines 418--442 and 487--523}

\textbf{Obtaining the coefficient from actual slots.} P3 gives
\(Q_{\tau b}(z)=2^{a+1+b}z\sum_{j\ge0}c_j(2z)^j\) with
\(c_j\in\mathbb Z_{(2)}\). Modulo \(2^{p_t}\) it therefore agrees at
every integral probe with a polynomial of degree at most \(p_t\). At
this node \(p_t=O(ML/u)\), with \(u\to\infty\), so
\(p_t+1\le M(2L-1)+1\) eventually. The printed set supplies that many
binary-coordinate lifts with pairwise unit differences. Ring Lagrange
interpolation at any \(p_t+1\) of them gives the true first coefficient
modulo \(2^{p_t}\). This step does not require that zero belong to the
prescribed set. For each tile, fix its anchor from the public layout;
the anchor and moment order determine the singleton mode. Fix the same
\(p_t+1\) probe indices before the signs are chosen. The resulting list of
\(N(p_t+1)\) requests therefore works for every sign word. Its length
bounds the number of distinct completed-slot addresses used. A source ring element of degree
\(e=4\lceil\log_2s\rceil\) and precision \(p_t\) and all the linear
interpolation weights have polynomial bit representation.
\href{https://github.com/openai/math/blob/main/preprints/The-Exponential-Time-Hypothesis-is-False-October-25-2026/build/quantitative-contraction.tex\#L338}{Printed
root interpolation and precision, lines 338--367}

\textbf{Exact recovery and joint cost.} Let
\(U_{\tau b}=\prod_x(t-\tau+x)^{b+1}\) over the eligible real rows in
this one packet. Each factor is odd and positive. In the summand
indexed by \(x\), the kernel slope \(A'_x(0)\) has denominator dividing
\((t-\tau+x)^{b+1}\). The derivative of the Boolean equality row is
\(-1+2w\), an integer at its baseline argument. The source moment and
the remaining baseline factors are integral on each assignment.
Consequently \(U_{\tau b}\) clears the denominator of every summand
and of their sum. Also
\(\log_2 U_{\tau b}\le L(b+1)\log_2(2L)\). There are at most \(L\) row
choices and \(2^M\) original column assignments. The source moment is at
most \(L^{b+1}\), the remaining row factors have product at most
\((1+|c|)^{k_t}\), the distinguished derivative has absolute value at
most \(1+2|c|\), and \(|A'_x(0)|\le |c|\alpha 2^b\). Hence \[
  \log_2\!\left(1+\left|U_{\tau b}[z]Q_{\tau b}\right|\right)
          =O(M+L^2\log L)=O(M)=o(p_t).
\] Here \(M=\Theta(s)\) and \(L=O(s^{2\varepsilon})\) with the printed
\(\varepsilon\). Multiplying by the public odd denominator in the native
ring and taking the centered representative thus recovers the exact
rational coefficient. Dividing by \(D_{\tau b}\), subtracting
graph-known nonpin terms, and using the same Vandermonde matrices
recovers all pins from the \textbf{completed} queries. No restriction to
a witness subcube is needed.

The postprocessor and the number of actual completed registers in J8 are
polynomial. Moreover the specific source-preserving linear readouts
\([z]Q_{\tau b}\) have a polynomial \emph{joint} acquisition algorithm:
evaluate the proved derivative formula on the actual input via bounded
vertex marginals. The readouts agree with the Lagrange combinations
modulo the requested precision. A program required to return every
nonzero-probe slot value to a caller cannot charge this shortcut as if
those separate outputs had been produced. J7 supplies a different input
interface with polynomial cost for every requested individual child.
Both J7 and J8 concern the pinned full-cube family, whose count also has
a direct edge scan. They give no general speedup for the printed
counting algorithm or for the multi-valued availability family. J2 uses
the opposite D orientation and a bounded-precision completed channel,
where these recoveries are unavailable.
\hyperref[v9-section-joint-algorithms]{J2 and its query model}

\subsection{J9. The single-node precision order is
linear}\label{v9-j9-the-single-node-precision-order-is-linear}

On the J7--J8 layout it suffices to supply every completed evaluation
used in J8 at the public precision \(p=10M\). Indeed \(|c|<2^{2L}\) is
an integer and \(1+|c|\le2^{2L}\). The explicit bounds above, including
the positive odd denominator \(U_{\tau b}\), give \[
 \log_2\bigl(1+|U_{\tau b}[z]Q_{\tau b}|\bigr)
       \le M+2L k_t+O(L^2\log L)
       \le 7M+o(M)<9M
\] eventually. The term \(M\) bounds the Boolean column sum; the costs
of the one distinguished derivative, its public scalar, the number of
possible target rows, the source moment and its odd denominator are all
in the displayed \(O(L^2\log L)\). Thus centered lifting at \(10M\)
recovers the signed integer. P3 modulo \(2^{10M}\) has degree at most
\(10M\). There are \(M(2L-1)+1\ge10M+1\) legal probe lifts once
\(L\ge6\), and the source precision \(p_t\ge24ML/u>10M\) eventually.
Choose just \(10M+1\) of the actual probes for each chosen singleton
tuple and reduce the native residues after acquisition. The
deterministic full-cube sign decoder uses at most \textbf{\(N(10M+1)\)}
completed root addresses at this one fixed node. The interpolation and
subsequent integer arithmetic still have polynomial bit cost.

Conversely, orient any certified physical layout on the same four-clone
input with subdivisions as rows and old vertices as columns. At
\(t=L-1\) there are only \(k_t\le3M/L\) shifted column ports. There are
\(4N=M/6\) distinct trivalent equality columns, one per clone. Each
equality with a shifted port uses at least one of those \(k_t\) ports.
Thus at least \[
                 K\ \ge\ M/6-3M/L
\] independent equality bits have no incoming real partial shift. The
unused-equality theorem P7 applies to original equality columns of every
arity, including these trivalent columns. Its whole-slot valuation is at
least \(v+(K-|\beta|)_+\ge K\). Every complete slot at this node, at all
modes and all legal probes, is consequently identical to zero through
\(K\) bits on \emph{every} sign word. No decoder with only this menu can
even distinguish two sign words. This proof does not use a specially
guarded order. In this orientation it holds for every certified physical
packing. Together with the constructive upper bound, the sufficient
precision order for reconstruction of all signs by the complete
\textbf{single-outer-node} interface is \(\Theta(M)\) when the certified
layout can be chosen. Its exact leading constant, a corresponding
statement at every other mesh point, and polynomial-time acquisition of
each individually specified nonzero-probe completed value are not
implied. \hyperref[v9-section-precision-frontier]{Unused original equality
columns, P5--P7}

\subsection{Exact checks}\label{v9-exact-checks}

On a literal 288-wire four-clone example, twelve genuine source-child
moments recovered all 4,096 possible sign words and all 163 vertex
covers. Nonpin columns with mean one-half shared queried tiles with the
pins. A separate eight-wire physical source checked both
first-coefficient readouts on all four sign words by direct residue-ring
interpolation. Each of two 41-point menus used binary-coordinate lifts
in a degree-six unramified ring at 40 bits, one menu containing zero and
one omitting it. The two computed public factors were \(209216/3\) and
\(1895936/45\). Ten of the actual target-slot evaluations were also
checked by a separate direct enumeration of all original Boolean row
assignments. These examples test the identities on finite packet arrays.
Printed scheduling and asymptotic certification use the preceding
proofs. \PaperSixEvidence{evidence/research_v9/checks/joint_algorithms_moment_escape.py}{Independent
direct moment and probe checker} ·
\PaperSixEvidence{evidence/research_v9/checks/joint_algorithms_moment_escape.json}{Saved finite
output}

\section{Precision and whole-menu
information}\label{v9-section-precision-frontier}

One additional bit above the old vanishing threshold can expose linearly
many independent original clause signs on a fixed certified binary
tensor network. On the v8 cubic counting family, choose new subdivision
equalities as columns. Given \textbf{any} certified vertex labels, one
can deterministically pack the physical wires so that \textbf{every}
completed mode at \textbf{every} retained node and probe is zero through
\((1-o(1))M/(2L)\) bits. The choice places original subdivision
equalities in specific packet positions. The original tensors and packet
formulas determine every child in both constructions.

The universal coefficient bound below also applies to layouts with
neither special feature. It counts shared base coefficients across all
extension-valued probes. A separate bound applies to the extreme grid
nodes on \emph{every} certified packing of the physical cubic family.

\subsection{The source slot and its first possible
digit}\label{v9-the-source-slot-and-its-first-possible-digit}

Use \(R=\mathbb Z_{(2)}\) and binary immutable original tensors with
\(\ell=1\). At the root, \(a=2u\), \(L=2^a\), \(\alpha=2L\), and
\(S=\{-L,\ldots,L-1\}\). At a retained \(t\), let \(i=v_2(t)<u\),
\(D_i=L/2^i\), \(B_i=2^{i+1}\), and \(\delta_i=a+1-i\). For a tile with
representative \(\tau\in[0,B_i)\), the source moment of index \(b\) is
\(L_{\tau,b}(X)=\sum_{y=\tau+B_i j}j^b X_y\) over real columns in that
packet. Each mode belongs to one packet, tile, index, source and target.
The source sums the products over mode tuples of printed weight
\(\sum(b+1)\beta_{\tau,b}<p_t\). Here each product is examined
separately, at a requested \(1\le p\le p_t\).
\href{https://github.com/openai/math/blob/main/preprints/The-Exponential-Time-Hypothesis-is-False-October-25-2026/build/quantitative-contraction.tex\#L11}{Coefficient
ring and lifts, lines 11--71} ·
\href{https://github.com/openai/math/blob/main/preprints/The-Exponential-Time-Hypothesis-is-False-October-25-2026/build/quantitative-contraction.tex\#L387}{Source
budget and modes, lines 387--455} ·
\href{https://github.com/openai/math/blob/main/preprints/The-Exponential-Time-Hypothesis-is-False-October-25-2026/build/quantitative-contraction.tex\#L487}{Child
entries and product identity, lines 487--523}

Write \(c_t=\mathcal C'(t)/\mathcal C'(0)\) and
\(N_t(z)=\prod_{d\ne t}(1+\alpha z/(t-d))\) with \(d\in S\). The former
is an odd unit and the latter belongs to \(1+(2z)R[2z]\). For a crossing
row \(x\) and tile \(\tau\), put \(d=t-\tau+x\) and \(h=v_2(d)\le i\).
The exact rational coefficient of that targeted mode in the source
formula is

\[
 A_{x,\tau,b}(z)=
 c_t\frac{\alpha B_i^b}{d^{b+1}}z\,N_t(z)
      (1+\alpha z/d)^{-b-1}.
 \tag{P1}
\]

The power series after the leading scalar and \(z\) is a unit in
\(R[[2z]]\) with constant term one. At every nonzero integral probe
\(\zeta\), including in the unramified extension, the \textbf{exact}
order of (P1) is

\[
 v_2(A_{x,\tau,b}(\zeta))
   =\delta_i+b+(b+1)(i-h)+v_2(\zeta).
 \tag{P2}
\]

At zero the coefficient is zero. To evaluate the power series at a probe
modulo a finite power of two one needs only a finite truncation.
Equivalently, the factors in its displayed rational expression have unit
values. These observations avoid any completeness assumption on \(R\).
The uniform lower bound \(\delta_i+b\) is attained by the row in residue
\(\tau\bmod B_i\) when present and the probe is a unit. Indeed the row
and tile have the same bit \(i\) when they cross. Within these crossing
positions, \(h=i\) precisely when \(x\equiv\tau\bmod B_i\). All
normalized constant units reduce to one in the base field. The local
kernels lie in \(R[2z]\) and at residue precision implement the partial
shift \(x\mapsto x+t\) on real matched pairs.
\href{https://github.com/openai/math/blob/main/preprints/The-Exponential-Time-Hypothesis-is-False-October-25-2026/build/quantitative-contraction.tex\#L175}{Cauchy
kernel, lines 175--231} ·
\href{https://github.com/openai/math/blob/main/preprints/The-Exponential-Time-Hypothesis-is-False-October-25-2026/build/quantitative-contraction.tex\#L402}{Crossing
expansion, lines 402--442}

For mode tuple \(\beta\), let \(k=|\beta|\),
\(v=v_i(\beta)=\sum_{\tau,b}(\delta_i+b)\beta_{\tau,b}\), with packet
indices implicit. Denote the completed slot before probe evaluation by
\(Q_{t,\beta}(z)\). This is a rational function regular in the integral
local coordinates. It has the form

\[
 Q_{t,\beta}(z)=2^v z^k\sum_{j\ge0}c_{t,\beta,j}(2z)^j,
 \qquad c_{t,\beta,j}\in R.
 \tag{P3}
\]

To see this directly from the children, distribute their product into
independent sums over each group's original assignments. Equivalently
take a single sum over all original row and column port assignments,
weighted by their original table product. For each packet, allocate its
targeted exponents to its distinct real row factors, assigning a row at
most one mode. For a selected row use the targeted coefficient (P1)
times \(l'_{Y_x}=\pm1\). Use its local affine form on other rows.
Multiply by all prescribed \(L_{\tau,b}^{\beta_{\tau,b}}\) on their
source assignments. Summing over these allocations is ordinary
coefficient extraction, with no extra factorial for repeated identical
modes. Pulling out (P1)'s factors proves (P3). Evaluating and reducing
every entry commutes with all finite sums and products, so this
expression gives exactly the native completed residue, even when a child
is acquired recursively at higher internal precision. A completed slot
can cancel or vanish identically. Equation (P2) concerns an individual
kernel coefficient.
\href{https://github.com/openai/math/blob/main/preprints/The-Exponential-Time-Hypothesis-is-False-October-25-2026/build/quantitative-contraction.tex\#L487}{Definition
of each packet tensor, lines 487--509} ·
\href{https://github.com/openai/math/blob/main/preprints/The-Exponential-Time-Hypothesis-is-False-October-25-2026/build/quantitative-contraction.tex\#L544}{Immutable
creation precision, lines 544--553}

There are two elementary refinements to \(v\) when a particular packet
is known. Minimize the sum of the exact row costs
\(\delta_i+b+(b+1)(i-h)\) over feasible allocations to distinct real
rows. To this minimum add \(\sum\beta_{\tau,b}m_{\tau,b}\), where
\(m_{\tau,b}\) is the minimum order of a nonzero weight \(j^b\) on that
real tile. Use \(0^0=1\). An impossible target allocation or a positive
exponent of a zero source form gives an identically zero product. With
the new exponent \(v_*\), (P3) holds with \(v_*\) in place of \(v\).
Moreover, the constant coefficient after dividing by \(2^{v_*}z^k\),
modulo two, has a literal finite test. Sum just the minimum-cost
allocations and the original Boolean assignments. Impose the partial
shift on every unselected row, ignore the selected linear signs, and use
the reduced source forms \(2^{-m_{\tau,b}}L_{\tau,b}\). Reduce the
resulting unweighted original table sum modulo two. A one attains the
order \(v_*+k v_2(\zeta)\) at every nonzero integral probe. A zero gains
at least one further digit. This follows term by term from (P1) and the
common constant term of its unit. The test includes both allocation
multiplicity and cancellation from the old tensors.

\subsection{Finite-precision information for a whole
menu}\label{v9-finite-precision-information-for-a-whole-menu}

Put \(T(n)=n(n+1)/2\) for \(n\ge0\) and \(x_+=\max(x,0)\). For fixed
\(t,\beta\), sending \(c_{t,\beta,j}\bmod 2^{p-v-j}\) whenever
\(0\le j<p-v\) determines \emph{every} value
\(Q_{t,\beta}(\zeta)\bmod 2^p\). This costs at most \(T((p-v)_+)\) base
bits. The fixed, publicly known extension probe simply evaluates the
same shared coefficients. There is no factor from the extension degree
or from the number of probes. Each coefficient refers to the whole slot,
so there is also no multiplicity from the 64 children. This is an
alphabet bound, with no requirement that acquisition compute the
coefficients quickly.

If \(R_{t,b}\) is the number of public packet-tile modes of index \(b\),
a bound on the joint bits at this \(t\) is

\[
 B_t(p)\ \le\ [x^{p-1}]\,
 \frac{1}{(1-x)^3}\prod_{b\ge0}(1-x^{\delta_i+b})^{-R_{t,b}}.
 \tag{P4}
\]

Indeed the product counts multisets by \(v\), and the coefficient from
\((1-x)^{-3}\) at \(p-1-v\) is \(T(p-v)\). Only indices costing less
than \(p\) matter. At the root \(R_{t,b}\le6M/D_i\). Counting all
exponents in (P4) is valid even if the supplied menu uses a subset or
some modes vanish. For arbitrary public precisions, add (P4) over the
actual retained nodes. This sum has \(D_i\) terms at valuation \(i\). At
\(p\le\delta_i\) this reproduces the v8 zero-mode bound. At
\(p=\delta_i+1\) the extra slots cost at most \(R_{t,0}\) bits per node.
All surviving positive exponents are single occurrences of \(b=0\). At
this first digit, a particular slot is \(\bar\zeta\) times a single base
bit after dividing by \(2^{\delta_i}\). Its bit is the XOR of the
original assignment contributions from selecting one row
\(x\equiv\tau\bmod B_i\), with each contribution multiplied by the
parity of its source tile. All unselected rows use the partial shift.
\href{https://github.com/openai/math/blob/main/preprints/The-Exponential-Time-Hypothesis-is-False-October-25-2026/build/quantitative-contraction.tex\#L432}{Source
multiplicity, lines 432--447} ·
\PaperSixEvidence{evidence/research_v8/applications/pooled_menu_barrier.md}{Prior
all-coordinate encoding}

The entire initial band \(p=\delta_i+h\), \(1\le h\le\delta_i\), admits
only single occurrences in nonzero tuples. Formula (P4) specializes to

\[
 T(\delta_i+h)+\sum_{b=0}^{h-1}R_{t,b}T(h-b)
 \ \le\ T(\delta_i+h)+\frac{6M}{D_i}\binom{h+2}{3}.
\]

Thus beyond the old precisions one may elevate \(n_i\) nodes per level
by up to \(h_i\) digits in this band. A sufficient uniform condition for
sublinear pooled information is
\(\sum_i (n_i/D_i)\binom{h_i+2}{3}=o(1)\), for \(M=\Theta(s)\). The
scalar triangular costs over every retained node are
\(O(L\log^2s)=o(M)\). With \(h_i=1\) at a full elevated coordinate the
singleton term is \(O(M)\). The construction below attains that order in
the general binary contraction interface.

\subsection{What unused equality columns
force}\label{v9-what-unused-equality-columns-force}

Suppose \(K=K(t)\) original column equality vertices have no port \(y\)
with a real row \(x=y-t\) in its packet. Binary, trivalent, and
higher-arity common-bit equalities, and a unary \((1,1)\), all qualify.
Require disjoint original vertices, as in the physical sparse
construction. Each contributes one independent common Boolean bit when
its internal equality is imposed. The coefficients in (P3) then obey

\[
 v_2(c_{t,\beta,j})\ \ge\ (K-k-j)_+ .
 \tag{P5}
\]

Expand the actual original assignment sum giving (P3) after imposing
these equalities. Each monomial from the \(k\) source moment factors can
mention at most \(k\) distinct such bits, even if a linear source moment
involves many columns. On every local row the coefficient multiplying
any of the \(K\) bits belongs to \((2z)R[2z]\), since its column has no
surviving shift. Thus a term contributing to power \(j\) of \(2z\) can
mention at most \(j\) further such bits. Normalized target coefficients
depend only on public locations. Each remaining untouched Boolean bit
sums to 2 independently. Other original column tensors use disjoint
ports. This gives (P5). The same reasoning works with the strengthened
allocation-and-source exponent \(v_*\). The Boolean cube summation and
arbitrary-probe vanishing for the resulting finite coefficient expansion
are checked in Lean. \PaperSixEvidence{evidence/research_v9/formal/BooleanSummation_scope.md}{Formal
statements and correspondence}

Consequently a slot costs at most

\[
 \left[T((p-v)_+)-T((K-k)_+)\right]_+
 \tag{P6}
\]

bits for all probes combined. This sums \([p-v-\max(j,K-k)]_+\) over
\(j\ge0\). To count all slots jointly, use the coefficient of
\(y^k x^v\) in \(\prod_b(1-yx^{\delta_i+b})^{-R_{t,b}}\) as the
multiplicity in the sum of (P6). One can instead count actual feasible
tuples with their strengthened \(v_*\). In particular

\[
 Q_{t,\beta}(\zeta)\equiv0\pmod{2^{\,v+(K-k)_+}}
 \quad\text{for every integral probe}.
 \tag{P7}
\]

Since \(v\ge k\), \textbf{every completed slot} is zero through \(K(t)\)
bits. This includes the zero mode. The argument bounds the actual
products and uses their source moments. A source-only child register
need not obey this erasure.

\subsubsection{The outer grid for every instantiation of the printed
padding
construction}\label{v9-the-outer-grid-for-every-instantiation-of-the-printed-padding-construction}

Here a physical layout uses the printed common load: the maximum real
ordered-cell load rounded up to a multiple of \(L\), as specified at
source lines 149--167. Arbitrarily adding more empty packets is not
included. For the v8 shortened physical CNF, \(M=30N+4\), \(s=25N/2+3\)
and \(P_0=5N+3\). If the subdivisions are columns, there are \(E=M/2\)
independent binary column equalities. If the original vertices are
columns, there are \(E=5N+2\) qualifying variable tensors (one trivalent
equality per three-occurrence variable, and two unary \((1,1)\)
tensors). Let \(P\) count the packets. At most \(P(L-|t|)\) real column
ports have a shifted real row. Hence \(K(t)\ge(E-P(L-|t|))_+\). At the
printed collision certificate,

\[
 0\le P L-M\le D_{\rm arr}(\sqrt{2M\Delta}+L)=o(M/L),
 \qquad D_{\rm arr}=r^{2a}=L^{12},\quad r=64 .
 \tag{P8}
\]

Here \(M=\Theta(s)\), \(\Delta=\lceil s^{9/10}\rceil\),
\(L=O(s^{2\varepsilon})\) and \(26\varepsilon<1/20\). Use the source's
common padded load \(H_{\rm pad}\), with \(P=H_{\rm pad}D_{\rm arr}/L\).
For fixed \(0<c<1/6\), every node with \(|t|\ge(5/6+c)L\) therefore has
\(K(t)\ge(c-o(1))M\) in either orientation. If the subdivisions are
columns, \(|t|\ge(1/2+c)L\) suffices. The decoder may receive the full
root menu there through these linear precisions and the v8 precisions
elsewhere. Its total menu information is still \(o(N)\) on every
certified physical layout. The exponentially diverse counts then force
\(\Omega(N)\) additional bits for exact reconstruction.
\href{https://github.com/openai/math/blob/main/preprints/The-Exponential-Time-Hypothesis-is-False-October-25-2026/build/quantitative-contraction.tex\#L132}{Source
packing and padded loads, lines 132--170} ·
\href{https://github.com/openai/math/blob/main/preprints/The-Exponential-Time-Hypothesis-is-False-October-25-2026/build/quantitative-reduction.tex\#L65}{Physical
reduction, lines 65--95} ·
\PaperSixEvidence{evidence/research_v8/applications/pooled_menu_barrier.md}{Physical
family and counts}

\subsubsection{All nodes on a permitted physical
packing}\label{v9-all-nodes-on-a-permitted-physical-packing}

\textbf{Guarded packing theorem.} Fix any of the v8 cubic counting
graphs, take subdivisions as columns, and give the vertices \textbf{any}
graph-only labels satisfying the printed collision certificate. There
exists a deterministic polynomial-time packet ordering preserving those
labels, the common padded load, and every packing condition, such that

\[
 \min_{t\in S\setminus\{0\}}K(t)
     \ \ge\ \frac{M}{2L}-O\!\left(\sqrt{D_{\rm arr}M/L}\right)
     \ =\ (1-o(1))\frac{M}{2L}.
 \tag{P9}
\]

The left side uses the absence of a real shifted port, independently of
whether the node is retained. The proof works with arbitrary original
row tables and binary subdivision columns. It preserves graph-only tie
constraints on the labels, including the v7 exterior-pair ties.

\textbf{Endpoint capacities.} Write \(b_0=H_{\rm pad}/L\), an integer.
For an unordered packet type with \(a'\) unequal coordinate pairs, the
source puts \(H_{\rm pad}2^{a'}/L=b_0 2^{a'}\) packets in that type. Use
each of the \(2^{a'}\) orientation masks in exactly \(b_0\) packets. A
mask fixes which orientation has bit zero at each unequal coordinate and
is used consistently throughout its packet, as required in source lines
157--165. This gives each ordered row/column label cell exactly \(b_0\)
slots at \(y=0\), and exactly \(b_0\) at \(y=L-1\). The two boundary
positions are distinct. Within every packet, each compatible pair label
still has its prescribed \(2^{a-a'}\) positions. Across the type every
cell still has exactly \(H_{\rm pad}\) slots. The property that adding a
number of order \(i\) makes the matched row and column symbols equal
uses their opposite \(i\)-bits, and holds for every mask.

\textbf{Whole equalities in the available positions.} Independently mark
each subdivision equality positive or negative with probability \(1/L\)
per sign, and leave it unmarked otherwise. Each of the \(M/2\) vertices
has two incident ports. For a given sign let \(C\) count marks and
\(X_\gamma\) count their ports in cell \(\gamma\). Then
\(\mathbb EC=M/(2L)\), \(\operatorname{Var}C\le M/(2L)\), and if
\(n_\gamma\) is the cell's real load,

\[
 \mathbb EX_\gamma=n_\gamma/L\le b_0,\qquad
 \operatorname{Var}X_\gamma\le 2n_\gamma/L .
\]

The second bound allows an equality to use the same cell twice, since
its multiplicity \(d\in\{0,1,2\}\) satisfies \(d^2\le2d\). Sum the
squared deviations of \(C\) and all \(X_\gamma\) from these means over
both signs. The expectation is at most \(5M/L\). Fix each mark in turn
to minimize the conditional expectation. Conditional means and variances
of the remaining independent bounded summands are rational and
polynomial-time computable. A zero-probability outcome can be ignored.
For each sign delete a marked equality meeting an overfull cell until
every cell load is at most \(b_0\). The number deleted is at most the
initial \(\sum_\gamma(X_\gamma-b_0)_+\), since each removal reduces this
total by at least one. By Cauchy--Schwarz the total is at most
\(\sqrt{5D_{\rm arr}M/L}\). The initial deviation of \(C\) is at most
\(\sqrt{5M/L}\). Both sets therefore retain the number in (P9), since
\(D_{\rm arr}L/M=o(1)\).

Place each remaining positive port at \(y=0\) in an available slot of
its own ordered cell and each negative port at \(y=L-1\). Match the
other real ports and padding to all remaining positions within each
cell. Each physical edge has the same position at its row and column
ends throughout. For \(t>0\), \(y=0\) has no possible predecessor
\(x=y-t\). For \(t<0\), \(y=L-1\) has none. Both ports of each chosen
equality are protected at the required sign, proving (P9).
\href{https://github.com/openai/math/blob/main/preprints/The-Exponential-Time-Hypothesis-is-False-October-25-2026/build/quantitative-contraction.tex\#L149}{Exact
source construction, lines 149--170}

The leading constant here is optimal \textbf{for certifying erasure by
unsupported equality vertices}. At \(t=1\) every unmatched real column
port is at slot zero or immediately follows a padded position. Hence
\(2K(1)\le P+(PL-M)=(1+o(1))M/L\) by (P8). The last-slot argument gives
the same bound at \(t=-1\). Further arithmetic cancellations can extend
the actual vanishing range.

On the guarded cubic graph the entire public menu therefore has exactly
one state whenever every requested precision is at most the guaranteed
integer in (P9), of order \(s^{1-2\varepsilon}\). In the v8 caterpillar
completion the exact count takes at least \(2^{(N-10)/12}\) distinct
values on the same fixed topology. Any successful joint decoder needs at
least \((N-10)/12\) additional bits in the sense of a conditional
side-information alphabet. Adaptive reads of this menu give the same
answers for all switches. The all-disabled physical switch word has
count three. Equation (P9) uses the constructed graph-only packet order.
The outer-grid conclusion above covers every source packing with its
certificate. The printed order from any particular deterministic
tie-breaking implementation may differ. The published root precisions
\(P_0+8\lceil3M/u\rceil D_i+8M(a+1)\) exceed the range in (P9).
\PaperSixEvidence{evidence/research_v8/applications/pooled_menu_barrier.md}{Range of
exact physical counts} ·
\href{https://github.com/openai/math/blob/main/preprints/The-Exponential-Time-Hypothesis-is-False-October-25-2026/build/quantitative-contraction.tex\#L250}{Printed
Hermite precision, lines 250--268}

\subsubsection{How many mode addresses pass the full-precision valuation
tests}\label{v9-how-many-mode-addresses-pass-the-full-precision-valuation-tests}

On the physical root, at \textbf{every} retained \(t\), at least
\(2^{\Omega(M/\sqrt{u})}\) distinct native exponent tuples pass the
weight, feasible-target-row, nonzero-source-form, and (P7) tests at the
published \(p_t\). This is a count of addresses retained by those tests.
The products could still be zero from their original clause tables or
cancel in the final readout.

Indeed, by (P8) fewer than \(o(P)\) packets have even one padding hole.
The fraction of all packets with unequal symbols in coordinate \(i\) is
\(1-1/r\), since the common padded cell load makes the ordered symbol
pairs uniform. Each full unequal packet offers \(B_i\) real tiles. Their
total number is \(\Theta(PB_i)=\Theta(M/D_i)\), uniformly over \(i<u\).
Set \(d=\lfloor D_i/(4\sqrt u)\rfloor\). On each of these tiles, take an
arbitrary subset of the mode indices \(0\le b<d\), each with exponent
one. Different tiles have disjoint residue classes of potential target
rows inside a packet. Within each tile there are \(D_i/2\) real rows in
the minimum-cost residue, more than enough to allocate its at most \(d\)
targets distinctly. All source forms have the unit-weight position
\(j=1\). Hence \(v_*=v\) for these tuples. Since \(D_i\ge2^{u+1}\) at
the root, \(\delta_i/d=o(1)\) uniformly. Their largest cost is at most

\[
 \frac{6M}{D_i}\left(d\,\delta_i+\frac{d(d-1)}2\right)
 \le \left(\frac{3}{16}+o(1)\right)\frac{MD_i}{u}.
\]

The actual precision is at least \(24 MD_i/u\), and \(K(t)\le M\).
Consequently \(p_t>v+(K-k)_+\) for every one of these tuples. Their
independent Boolean subset choices give \(\Theta(M/\sqrt u)\) address
bits. The paper's partition upper bound in the other direction puts
\emph{all} printed addresses at \(2^{O(M/\sqrt u)}\). The explicit
subset proves the lower bound on filtered addresses. Counting distinct
observed values or operations in a fused evaluator requires further
evidence.
\href{https://github.com/openai/math/blob/main/preprints/The-Exponential-Time-Hypothesis-is-False-October-25-2026/build/quantitative-contraction.tex\#L149}{Native
packing, lines 149--170} ·
\href{https://github.com/openai/math/blob/main/preprints/The-Exponential-Time-Hypothesis-is-False-October-25-2026/build/quantitative-contraction.tex\#L250}{Printed
precision, lines 250--265} ·
\href{https://github.com/openai/math/blob/main/preprints/The-Exponential-Time-Hypothesis-is-False-October-25-2026/build/quantitative-contraction.tex\#L449}{Printed
address count, lines 449--482}

\subsection{A sharp first-extra-digit escape in the general contraction
interface}\label{v9-a-sharp-first-extra-digit-escape-in-the-general-contraction-interface}

\textbf{Single-coordinate capacity witness.} At any chosen active \(i\),
a fixed certified graph and layout of bounded-degree original binary
tensors have at least

\[
       \frac{1-1/r}{3}\left(1-\frac{4}{D_i}\right)M
       \tag{P10}
\]

independently recoverable original clause-sign bits in their completed
\(b=0\) singleton slots modulo \(2^{\delta_i+1}\). This remains true if
only one fixed nonzero-residue probe is available. The printed set
contains such a probe because it has more than one residue-distinct
element. This construction is at the paper's general Boolean
tensor-contraction entry point. Its column pins differ from the
subdivision columns of the sparse-CNF compiler.
\href{https://github.com/openai/math/blob/main/preprints/The-Exponential-Time-Hypothesis-is-False-October-25-2026/build/quantitative-reduction.tex\#L15}{General
input definition and theorem, lines 15--49} ·
\href{https://github.com/openai/math/blob/main/preprints/The-Exponential-Time-Hypothesis-is-False-October-25-2026/build/quantitative-contraction.tex\#L338}{Public
probe schedule, lines 338--345}

Take arbitrary sufficiently large powers of two \(s\), with the printed
root \(a,u,L\), and choose
\(H_{\rm pad}=3L\lfloor s/(3L D_{\rm arr})\rfloor\). Set
\(M=H_{\rm pad}D_{\rm arr}=(1-o(1))s\). Fully occupy every cell with
exactly \(H_{\rm pad}\) ports, giving the collision potential zero. Each
unordered type supplies \(H_{\rm pad}2^{a'}/L\) complete packets,
divisible by three. Within a type group them in triples \(U,V,W\) with
the \textbf{same} orientation mask throughout the triple. Join
corresponding row positions \(x\) of those three packets at a single
original degree-three vertex. Because their labels match in every
coordinate this defines one legal row vertex. At that vertex place the
literal clause

\[
       R_{s_x}(u,v,w)=\neg u\ \lor\ \neg v\ \lor\ L_{s_x}(w),
       \qquad L_1(w)=w,\quad L_0(w)=1-w .
\]

Use a separate unary original column vertex at every slot. In packet
\(U\) pin its bit to one exactly on \([0,B_i)\). In \(V\) use ones
exactly on \([L-B_i,L)\), and in \(W\) use zeros. These are literal
immutable \(0/1\) tables, as are all row tensors. Each edge occupies its
stated position at both ends. Degree is at most three, descriptions have
size \(O(M\log M)\), \(\ell=1\) and one can choose \(P_0=1\). All root
source parameter inequalities hold eventually. Only triples with unequal
coordinate-\(i\) symbols will be used for data, a fraction \(1-1/r\) of
the wires. The other rows may use fixed choices of their sign.

Consider positive nodes of order \(i\) with \(B_i\le t\le L-B_i\). The
partial shift sends every \(U\) argument to zero. The \(V\) argument is
one exactly on

\[
       I_t=[L-B_i-t,L-t).
\]

The \(W\) argument is zero. After the original Boolean row sum, its
multilinear table is \(1-uv(1-L_{s_x}(w))\). At those partial-shift
arguments its value is one, and its derivative in \(u\), modulo two, is
\(s_x\mathbf1_{I_t}(x)\). This also keeps the unselected rows in
equal-symbol triples at one. Query the singleton of tile \(\tau\) in
\(U\) at mode index zero. Its source moment is exactly one, since that
tile has exactly one pinned one. By (P2), after dividing by
\(2^{\delta_i}\) modulo two only rows \(x\equiv\tau\bmod B_i\)
contribute. There is exactly one such row in \(I_t\). Every other child
factor in this digit has the computed baseline value one. Thus the
\emph{actual completed product}, normalized at its first possible digit,
equals \(\bar\zeta s_x\).

The eligible \(t\)'s run from \(3B_i/2\) to \(L-3B_i/2\) in steps of
\(B_i\). Their intervals are disjoint, cover \(L-2B_i\) positions per
triple, and each of their \(B_i\) tile queries returns a different sign.
There are \((1-1/r)M/(3L)\) eligible triples, proving (P10). At the old
threshold the singleton information disappears. The old
entire-coordinate menu has \(o(M)\) bits while the first new digit here
has \(\Omega(M)\). The original undeformed contraction of this witness
is one for every sign word, since the two pin intervals are disjoint. It
proves a menu-information transition, with no exact-count decoding claim
on the v8 cubic family and no acquisition-time lower bound. The question
whether \emph{every} certified cubic-CNF layout has small information at
this first extra digit is open here.

\subsection{Random packet orderings with the same fixed vertex
labels}\label{v9-random-packet-orderings-with-the-same-fixed-vertex-labels}

There is also an all-node obstruction under a simple graph-only
distribution on permitted orderings. Independently choose, for each
unordered type, a uniform orientation mask and use it in all its
packets. Within each ordered row/column label cell uniformly permute its
real and padding slots, independently. With probability \(1-o(1)\), all
completed slots on the physical subdivision-column network
simultaneously vanish through at least
\((1-o(1))M/(2L^2)=\Omega(s^{1-4\varepsilon})\) digits. This probability
is over the packing only. Once successful, it covers every switch word.

For one sign let \(C\) count subdivision equalities with both ports at
the corresponding endpoint. In a type with \(a'\) unequal coordinates,
each ordered cell has probability \(2^{-a'}\) to be eligible for this
endpoint. When eligible it has \(b_0 2^{a'}\) endpoint positions among
\(H_{\rm pad}\), where \(b_0=H_{\rm pad}/L\). Each port therefore hits
with probability \(1/L\). The two ports of an equality in distinct types
are independent. When in the same type they belong to the same ordered
cell, since their common column label and the unordered pairs determine
their entire row label. Sampling two distinct positions in this cell
gives joint endpoint probability

\[
 2^{-a'}\frac{b_0 2^{a'}(b_0 2^{a'}-1)}
                 {H_{\rm pad}(H_{\rm pad}-1)}
       \ \ge\ \frac{1-1/b_0}{L^2}.
\]

Thus \(\mathbb EC\ge(1-1/b_0)M/(2L^2)\). The random data in different
types are independent. Replacing all the data in type \(j\), which
contains \(P_j\) packets, changes \(C\) by at most \(2P_j\), since this
bounds the symmetric difference of its selected port sets. Revealing
successive types gives a conditional-expectation martingale with each
difference bounded by \(2P_j\). Orthogonality of these differences gives
\(\operatorname{Var}C\le4\sum_j P_j^2\le4H_{\rm pad}P\), since
\(P_j\le H_{\rm pad}\). Relative to \((M/(2L^2))^2\), this is
\(O(L^3/D_{\rm arr})=O(L^{-9})\). Chebyshev's inequality with relative
slack \(L^{-2}\), for both signs, gives

\[
 \Pr\left[\min_{t\in S\setminus\{0\}}K(t)\ge
          (1-1/b_0-L^{-2})\frac{M}{2L^2}\right]
       \ge 1-O(L^{-5}).
 \tag{P11}
\]

Equation (P7) gives the stated consequence for every completed slot. The
distribution in (P11) is fixed here, independently of the original
deterministic implementation.

\subsection{Finite checks}\label{v9-finite-checks}

The checker compares the exact rational kernel in
\(\mathbb Q[U]/(U^2+U+1)\) with (P2), including zero and nonunit probes.
It computes completed singleton products using the original
three-literal truth tables and compares a four-position example with an
independent enumeration of the Boolean assignments in each actual child.
At a 16-position block and \(i=1\), its entire eight-bit input cube
gives 256 different menus from one quadratic-extension probe. It also
extracts coefficients of literal complete slots with equality columns
and one or two target occurrences, verifies (P5), and applies the
conditional-expectation seating algorithm to a finite v8 shortened cubic
CNF. The examples use deliberately small grid and group parameters to
test the identities. The asymptotic source parameter and collision
assertions follow from the preceding proofs.
\PaperSixEvidence{evidence/research_v9/checks/precision_frontier.py}{Runnable checker} ·
\PaperSixEvidence{evidence/research_v9/checks/precision_frontier.json}{Recorded results} ·
\PaperSixEvidence{evidence/research_v9/logs/precision_frontier_review.md}{Verification and layout
quantifiers}

\section{MAIN queries and exact joint
evaluators}\label{v9-section-joint-algorithms}

We use the four-clone pinned triangle-replacement cubic expander of
order \(N\) from v7. A unary pin chooses the common value of its four
private clones, so the count is \(C_G=U_MU_D\), where \(U_M\) multiplies
the three positive ORs in every triangle and \(U_D\) multiplies the ORs
on the exterior perfect matching. Our decoders only receive the
input-dependent registers specified in each theorem. For a requested
register its evaluator may compute the true value by any algorithm. We
also treat a uniform evaluation problem on general fixed-density clause
networks. \PaperSixEvidence{evidence/research_v7/reviews/hybrid_independent.md}{v7
common family and tied layouts} ·
\PaperSixEvidence{evidence/research_v8/native/native_escape.md}{v8 native register
definitions}

\subsection{J1. Diagonal pins in every native MAIN
group}\label{v9-j1-diagonal-pins-in-every-native-main-group}

Call a pin diagonal when the full word \(a_{j(v)}\in[r]^Q\) of its
private variable equals the full word \(d_{f(v)}\) of its unary factor.
The ordered auxiliary cell in its assigned copy is fixed by all \(Q\)
swaps. Every permutation in \(M_0=V_0f_{\mathcal I}(W_0)\) fixes that
cell. The corresponding row is supported just on the cell, where the
mismatch polynomial is identically one. The source formula \[
 P_u(T)=\delta_u(I-M_0)+\sum_y M_0[u,y]\varphi_{u,y}(T)\delta_y
\] therefore gives \(P_u(T)=\delta_u\), even if the sole entry of
\(M_0\) vanishes. After embedding, \(P_{f(v),j(v)}(T)=\delta_{j(v)}\) at
every \(T\), and its Taylor correction vanishes identically.

At the published parameters weight zero is exceptional for all
sufficiently large sizes. The tensor product of the \(Q\)
diagonal-coordinate indicators for this cell is a weight-zero
exceptional row consisting of exactly this cell indicator. The
block-diagonal copy matrix and the injective placement within each
occupied copy turn the row into \(e_{j(v)}^\top\). Group restriction at
\emph{every} node preserves it in the variable's group, which is also
this pin factor's group. Every admissible private point of a feasible
measurement address \(y\) thus has one prescribed bit \(w_v(y)\). At its
corrected argument the pin factor is precisely
\(1_{\{\sigma_v=w_v(y)\}}\), at every order and every \(U\subseteq E\).
Empty restricted fibers give zero. Consequently a group jet \(\Sigma_g\)
vanishes unless the input agrees with one prescribed pattern on all
diagonal pins owned by the group. A completed root product requires
agreement on all diagonal pins at once.
\href{https://github.com/openai/math/blob/main/preprints/The-Exponential-Time-Hypothesis-is-False-October-25-2026/build/main_full_version.tex\#L2241}{Swap
basis and exceptional weights, lines 2241--2272} ·
\href{https://github.com/openai/math/blob/main/preprints/The-Exponential-Time-Hypothesis-is-False-October-25-2026/build/main_full_version.tex\#L2313}{Polynomial
rows, lines 2313--2352} ·
\href{https://github.com/openai/math/blob/main/preprints/The-Exponential-Time-Hypothesis-is-False-October-25-2026/build/main_full_version.tex\#L2533}{Copy
embedding and group measurement, lines 2533--2593} ·
\href{https://github.com/openai/math/blob/main/preprints/The-Exponential-Time-Hypothesis-is-False-October-25-2026/build/main_full_version.tex\#L2796}{Published
interval, lines 2796--2807}

\textbf{Certified abundance at every group.} Reserve a triangle in
advance. Remove its vertices and their exterior mates, and choose a
graph-only maximal independent set \(H\) of the remaining graph. Greedy
choice gives \(|H|\ge (N-6)/4\); \(H\) contains at most one vertex in
each triangle. Under the v7 triangle-tied MAIN distribution, the
pin-factor words on \(H\) are independent and uniform: even the shared
primitive symbol in the designated coordinate is used by at most one
such factor. Their distinct private variable words remain independent
and uniform, independently of those factor words. Put \(K=r^Q\), and for
each \(t,g\) let \(X_{tg}\) count the vertices in \(H\) that are
diagonal and have symbol \(g\) at \(t\). For each fixed \(t,g\),
independently across the candidate vertices the success probability is
\(1/(rK)\), whence \(\mu_0=\mathbb E X_{tg}=|H|/(rK)\) and
\({\rm Var}(X_{tg})\le\mu_0\). The counts for different \(t,g\) can be
dependent.

For \(n=4N,M=12N,b=3,w=2\), let \(\Phi\) be the sum of squared
deviations of ordered incidence loads and \(\Xi\) the sum of squared
group-size deviations. In the v7 collision expansion an incidence has at
most four ordered partners including itself that share an endpoint. The
pin ties introduce at most \(2N\) other dependent pairs. Thus
\(\mathbb E\Phi\le4M+2N\) and \(\mathbb E\Xi\le Qn\). The nonnegative
score \[
 \frac{\Phi}{5M}+\frac{\Xi}{Qn}
   +8\sum_{t,g}\frac{(X_{tg}-\mu_0)^2}{\mu_0^2}
\] has expectation at most \(11/6+8Qr/\mu_0<2\) eventually. Here the
published \(r^{2Q}\le2^{\lceil\log_2(4N)\rceil/128}\) gives
\(K=O(N^{1/256})\), while \(Q,r\) are polylogarithmic. Fix the primitive
symbols successively by conditional expectation. The additional score
terms reduce to single and pair events involving at most four words with
symbol equalities and fixed symbols. Enumerating the \(K^4\) possible
completions of these words, respecting ties and previous fixings,
suffices for each term. Equivalently count connected equality
constraints coordinatewise. There are polynomially many events and
primitive symbols, and the rational denominators have polynomial bit
length. The fixed score is less than two. This gives the original joint
source certificate and \(X_{tg}>\mu_0/2\) simultaneously. No particular
copy-placement tie-break is needed once the certified words have been
found. \PaperSixEvidence{evidence/research_v7/reviews/hybrid_independent.md}{v7
expectation and source certificate} ·
\href{https://github.com/openai/math/blob/main/preprints/The-Exponential-Time-Hypothesis-is-False-October-25-2026/build/main_full_version.tex\#L2747}{Printed
exponents, lines 2747--2787}

Write \(P\) for all diagonal pins in \(H\), \(k=|P|\), and
\(k_0=\min_{t,g}X_{tg}\). Summing at any fixed node gives \[
   k=\Omega(N/r^Q)=\Omega(N^{255/256}),\qquad
   k_0=\Omega(N/r^{Q+1})=\Omega(N^{255/256}/r).
\] In particular the same independent set of sign coordinates works at
every node and group. The conclusion specifies deterministically
constructible certified layouts. It does not specify the labels produced
by the source's own independently distributed conditional-expectation
tie-breaking algorithm.

\subsection{J2. Adaptive query obstruction and a mixed
interface}\label{v9-j2-adaptive-query-obstruction-and-a-mixed-interface}

For every word \(w\) on \(P\), set all other signs to one. Every edge is
covered because \(P\) is independent. Keeping \(w\) fixed, set two
vertices of the reserved triangle to zero. Their common internal edge is
now uncovered. The exterior matching remains covered in both inputs
because both exterior mates were reserved outside \(H\). Thus \(C_G\)
and \(U_M\) each take both values on every pattern \(w\), with \(U_D=1\)
on both witnesses.

A deterministic decoder may adaptively request complete rational root
jets and individual rational group jets and use arbitrary computation on
the replies. Public graph-only preprocessing is free. Follow the branch
giving the zero jet at every such request. Each complete request can
match at most one word on \(P\); each individual group request can match
at most \(2^{k-k_0}\) words by J1. After \(A\) complete and \(B\)
individual requests with \[
                        A\,2^{-k}+B\,2^{-k_0}<1,                 \tag{Q1}
\] some word has not matched a requested pattern. Both witness inputs
for that word reproduce all replies, so this branch cannot terminate
correctly. Every such prefix is reachable by the same reasoning. In a
finite worst-case computation some input must therefore incur a prefix
reaching weight at least one, counting the request before its reply. In
particular at least \(2^k\) queries are necessary if only completed MAIN
jets are supplied, and at least \(2^{k_0}\) even with individual MAIN
groups. Caching, cancellations, comparisons, canonical rational
representations and arbitrary precision in replying to each query do not
affect this counting of actual requests. Any fixed function of a
requested jet can be evaluated on the simulated zero response as well.

For a bounded-error variant, choose \(w\) uniformly and either witness
equiprobably. For each fixed random seed with a worst-case limit \(T\)
on queries, the zero path of depth at most \(T\) covers a fraction at
most \(T/2^{k_0}\) of patterns. If a zero path has already become
unreachable its covered fraction is one and the bound is trivial. Off
the covered set the two outputs remain equiprobable, giving success at
most \(1/2+T/2^{k_0+1}\). Worst-case error at most \(1/3\) thus requires
\(T\ge 2^{k_0}/3\); for complete-only queries use \(k\).

For the same physical CNF and v7 exterior-pair-tied D labels, take old
vertices as rows and subdivision equalities as columns. The
guarded-order theorem in the companion precision analysis supplies a
graph-only slot order for which every completed D root slot at every
node, probe and mode is zero modulo \(2^q\), for
\(q=\Omega(M_{\rm D}/L)\), \(M_{\rm D}=24N\),
\(L=2^{2\lfloor\varepsilon\lceil\log_2 s_{\rm D}\rceil\rfloor}\). Both
layouts can be fixed independently before any signs arrive. Giving all
these D residues freely adds no information to the simulated branch. The
same query bounds therefore hold for a single genuine mixed interface
and the actual count. The published D precisions \(p_t\) are larger. Raw
pin entries, per-input cache state, groups from D that have not been
proved constant, and other layouts require separate analysis. In
particular polynomially many genuine MAIN probes cannot produce the v7
MAIN operand on these layouts. On this two-witness promise a raw read of
either designated triangle pin reveals the target. On the full v7 cube a
raw edge scan computes the count in polynomial time. The obstruction
concerns the native query interface. On another certified layout the
high-precision D channel does recover the entire input, as J7--J8 prove.
\hyperref[v9-section-precision-frontier]{All-node guarded erasure, P7--P9}
· \hyperref[v9-section-moment-escape]{Opposite orientation and
high-precision recovery, J7--J8}

\subsection{J3. A graph-testable first-derivative
decoder}\label{v9-j3-a-graph-testable-first-derivative-decoder}

There is also a sufficient test for a positive result on other layouts.
Use the genuine measurement of the all-one private word. All its fibers
are singletons: weight-zero orbit indicators, one per orbit and copy,
measure the total of the occupied Boolean entries of that orbit, and
this total is at its maximum. Every private variable occurs. Let \(B_t\)
be the product of the sign-independent factors at these base arguments,
\(\theta_{tv}=\langle P_{f(v),j(v)}(t),\mathbf1\rangle\), and
\(d_{tvi}=P'_{f(v),j(v)}(t)[i]\). At a node with \(B_t\ne0\) and all
\(\theta_{tv}\notin\{0,1\}\), the zero-order complete product
\(H_t\ne0\) for every sign word. For \(J\ge2\), query the native completed jet at \(E=U=\{i\}\)
and read both its constant coefficient \(H_t\) and its first coefficient
\(h_{ti}\) from that address. The normalizing value \(H_t\) may depend
on the private signs. Let
\(b_{ti}\) be the corresponding first coefficient of just the public
sign-independent product. The singleton product rule gives \[
 \frac{h_{ti}}{H_t}-\frac{b_{ti}}{B_t}+\sum_v\frac{d_{tvi}}{1-\theta_{tv}}
       =\sum_v \frac{d_{tvi}}{\theta_{tv}(1-\theta_{tv})}\,\sigma_v .
\] For \(L_1(w)=w\) the logarithmic derivative is \(d/\theta\), and for
\(L_0(w)=1-w\) it is \(-d/(1-\theta)\). Full column rank of the
displayed graph-only matrix, stacking any admissible nodes and
directions, recovers every sign by rational Gaussian elimination. Both
v7 factors then follow by their edge definitions. At most \(O(nQ)\)
actual native completed addresses are used. Their \emph{values} have a
polynomial bit algorithm: the all-one measurement certifies the only
private point, so evaluate the source products there. Common row
denominators, the bounded total degree and the source coefficient bounds
give polynomial intermediate bit lengths. Fraction-free elimination
costs polynomial bit work by the determinant size bound. We have no
claim of existence of a suitable asymptotic v7 layout; a single diagonal
pin defeats this admissibility test at every node.
\PaperSixEvidence{evidence/research_v8/native/native_escape.md}{Singleton proof and
coefficient sizes} ·
\href{https://github.com/openai/math/blob/main/preprints/The-Exponential-Time-Hypothesis-is-False-October-25-2026/build/main_full_version.tex\#L1886}{Taylor
addresses, lines 1886--1997}

\subsection{J4. Exact contraction of an outer zero-probe
register}\label{v9-j4-exact-contraction-of-an-outer-zero-probe-register}

At a nonzero active node \(t\), put \(d=L-|t|\). At \(\zeta=0\) the
packet kernel on real ports is the partial injection \(x\mapsto x+t\)
weighted by the odd rational unit \(c_t=\mathcal C'(t)/\mathcal C'(0)\).
All targeted mode coefficients vanish exactly, hence so does every
nonzero-mode \emph{completed} product. For the zero mode all surviving
pairs are local. Multiplying the children is exactly the contraction of
this partial wiring on the full root. The number \(k_t\) of actual pairs
is at most \(d\) per packet and at most \(3Md/L\) in total.
\href{https://github.com/openai/math/blob/main/preprints/The-Exponential-Time-Hypothesis-is-False-October-25-2026/build/quantitative-contraction.tex\#L176}{Cauchy
formula, lines 176--203} ·
\href{https://github.com/openai/math/blob/main/preprints/The-Exponential-Time-Hypothesis-is-False-October-25-2026/build/quantitative-contraction.tex\#L387}{Locality
and crossing coefficients, lines 387--441}

Enumerate just the \(k_t\) matched column bits. For a row vertex sum its
table over its own row bits with factors \(l_Y(c_tX)\) on matched ports
and \(l_Y(0)\) elsewhere, obtaining its multilinear extension at those
arguments. For a column vertex condition its table on matched bits and
sum all its unmatched bits. Different column vertices have disjoint port
bits, and every row contribution depends only on the specified matched
bits. Multiply the vertex values and sum over matched assignments.
Isolated vertices contribute their scalar once. These are finite
distributivity identities at the creation-time lifts reduced modulo
\(2^{p_t}\). For bounded-degree physical clause tensors this takes
\(2^{k_t}\operatorname{poly}(s,p_t)\) bit time and polynomial space,
including the source unit algorithm for \(c_t\). Both legal odd nodes
\(t=\pm(L-1)\) have exact zero-mode evaluators in
\(2^{O(M/L)}\operatorname{poly}(s,p_t)\) bit operations. Here
\(M/L=o(M/\sqrt{\log s})\) on the growing fixed-density families. The
printed evaluator visits \(2^{\Omega(M/\sqrt{\log s})}\) assignments for
each such slot. If zero is included among the legal binary-coordinate
lifts the value is a genuine delivered slot. A fixed probe set omitting
zero gives a native interpolated value, as described below.
\href{https://github.com/openai/math/blob/main/preprints/The-Exponential-Time-Hypothesis-is-False-October-25-2026/build/quantitative-reduction.tex\#L65}{Original
tensors and lifts} ·
\PaperSixEvidence{evidence/research_v8/native/native_escape.md}{Printed scan versus
actual values}

View matched bits as variables and the local quantities above as
factors. Given an elimination order of induced width \(w\) on their
primal graph, variable elimination costs
\(\operatorname{poly}(s,p_t)2^{w+1}\) bit operations for bounded-degree
original tables. Multiply only the factors containing the next bit, sum
over that bit and store the resulting table on its at most \(w\)
remaining neighbors. A verified width \(O(\log s)\) thus gives a
polynomial-time register evaluator. The packing need not have this
width. For general initial degree \(\Delta\) and entry cost
\(E_{\rm tab}\), local tabulation adds a factor
\(\operatorname{poly}(s,\mathcal L)2^{O(\Delta)}E_{\rm tab}\) to the
enumeration bound. This fits
\(2^{O(M/L)}\operatorname{poly}(s+\mathcal L)\) at an active root under
the source bounds: \(\Delta=O(s^{9/10})\),
\(\log E_{\rm tab}\le O(s^{9/10})+O(\log(s+\mathcal L))\), and
\(M/L\ge H/L\gg s^{9/10}\). No individual high-degree marginal is
assumed cheap.

\subsubsection{A polynomial-time, nonconstant genuine register on the v7
full
input}\label{v9-a-polynomial-time-nonconstant-genuine-register-on-the-v7-full-input}

On the same four-clone family there is a legal graph-only packet order
that gives polynomial-time exact evaluation of the complete zero-probe
register at \textbf{every active node in a fixed outer interval}
\(|t|\ge(1-\eta)L\), with \(\eta>0\) absolute. Take the old CNF vertices
as columns and the subdivisions as rows. This order is distinct from
P9's guarded subdivision-column construction. Existing certified vertex
labels, including the exterior-pair ties on old pins, can be retained.

Here are the ordering details. The physical graph has bounded degree
\(d_0\). Fix one designated pin wire \(e_*\) for inclusion in a matching
and remove its closed edge neighborhood. Mark each remaining edge
independently with probability \(1/2\), retaining it when it is the sole
marked edge at both endpoints. Add \(e_*\). Retained wires form one
matching. Each candidate survives with probability at least
\(c=2^{-(2d_0-1)}\), since at most \(2d_0-2\) other edges touch it. In
cell \(\gamma\) with \(n_\gamma\) real wires, the expected retained
count \(\mu_\gamma\) is at least \(c(n_\gamma-2d_0)\). The variance is
\(O_{d_0}(n_\gamma)\), since inclusion indicators are independent at
line-graph distance greater than two, and there are only \(O_{d_0}(1)\)
neighbors within that radius. Derandomize the sum over cells of squared
deviations from \(\mu_\gamma\); its expectation is \(O_{d_0}(M)\). Each
mixed moment of two indicators is determined by the marks on two
constant-size closed edge neighborhoods and can be computed under any
partial fixing in polynomial time. The final matching has at least
\(\mu_\gamma-O_{d_0}(\sqrt M)\) wires in every cell. The certified load
estimate yields \(n_\gamma=(1+o(1))M/D_{\rm arr}\) and
\(H_{\rm pad}=(1+o(1))M/D_{\rm arr}\), with \(M/D_{\rm arr}\gg\sqrt M\).
Thus all cells have at least \(c H_{\rm pad}/2\) matching wires
eventually.

Use each orientation mask equally often within each packet type, as in
P9. Every fixed position has exactly \(H_{\rm pad}/L\) seats of each
ordered cell. Set \(h=\lfloor cL/8\rfloor\). Fill the first and last
\(h\) positions of every packet from distinct matching wires, spending
\(2hH_{\rm pad}/L\le c H_{\rm pad}/4\) seats per cell. Place \(e_*\) in
a last position of its own cell, and fill the other seats with the
remaining wires and padding as usual. At any \(|t|\ge L-h\), every
shifted pair uses two band positions, and every original vertex on
either side has at most one matched port. Taking, for example,
\(\eta=c/16\) gives the stated interval for all large \(L\).

In this orientation an unmatched subdivision row has value one; a row
matched on one leg gives \(1-c_tX\) on its incoming column bit. The
original column tensors are now independent. Each column tensor can be
summed on its own ports in constant many ring operations. The zero-mode
slot at native precision is computable in \(\operatorname{poly}(s,p_t)\)
bit time. If \(S_t\) indexes the matched unary pin columns, its exact
rational value before reduction is \[
                A_t\, (1-c_t)^{\sum_{v\in S_t}\sigma_v},
\] where \(A_t\) is a graph-only product of the other original column
sums. At \(t=L-1\), all positions \(0,L-1\) are real and \(e_*\) makes
\(S_t\) nonempty. Here \(c_t=-\binom{2L-1}{L}\) is a negative integer
with \(|c_t|<2^{2L}\); every old fixed table has at least one satisfying
assignment, so \(A_t>0\). All terms are nonnegative and there are at
most \(2^M\) column assignments, each of weight at most
\((1+|c_t|)^{3M/L}\). Its base-two logarithm after summation is at most
\(M+(2L+1)3M/L=O(M)\), below the printed \(p_t\ge24ML/u\). Thus the
actual residue equals the positive integer. A single full-precision
completed probe with \(\zeta=0\) determines the Hamming weight on the
nonempty graph-known set \(S_t\) by comparison with \(A_t(1-c_t)^j\),
\(0\le j\le|S_t|\). It determines a designated sign on the restriction
where the other \(S_t\) signs are fixed. This is a nonconstant exact
native register with polynomial acquisition and decoding. This
particular statistic does not give the v7 exterior OR product. A
different no-matched-pin order recovers every sign from native moments.
\hyperref[v9-section-moment-escape]{J7--J8} At a prescribed probe set
omitting zero it is the polynomially interpolable value discussed in J6;
efficient separate evaluation of all the required nonzero-\(\zeta\)
probes is not asserted.
\href{https://github.com/openai/math/blob/main/preprints/The-Exponential-Time-Hypothesis-is-False-October-25-2026/build/quantitative-contraction.tex\#L114}{Root
packing and shifts} ·
\href{https://github.com/openai/math/blob/main/preprints/The-Exponential-Time-Hypothesis-is-False-October-25-2026/build/quantitative-reduction.tex\#L65}{Physical
unary and equality tables}

\subsection{J5. Joint mixed child moments at native
precision}\label{v9-j5-joint-mixed-child-moments-at-native-precision}

Fix one child at \(\zeta=0\). Every positive target exponent at this
child gives zero. A source-only tuple multiplies the column weight by
integer powers of the actual source forms
\(S_{\tau b}(X)=\sum_{j\in\tau}j^b X_j\), where \(\tau\) is a real
column tile with local index \(j\). Let \({\cal B}_\tau\) list the
requested mode indices in tile \(\tau\), with \(0^0=1\), and set \[
 H_{\tau b}=\sum_{j\in\tau}j^b,\qquad
 V_g=\prod_\tau\prod_{b\in{\cal B}_\tau}(1+H_{\tau b}) .
\] For each matched column pattern in the child, perform the row
reduction of J4. Accumulate a table on vectors \((S_{\tau b})\) by
processing one original column vertex at a time: enumerate its local
assignments consistent with the matched bits, shift each prior index by
the assignment's vector of contributions, and add the product of its
table weight with the old entry. All vector indices are nonnegative
integers; the weights use the actual ring at \(p_t\). Multiply by the
row factor and add the table over matched patterns. Induction on
processed vertices identifies each entry with the sum of the original
weights having precisely that vector.

This single table answers every mixed tuple in the listed indices by its
dot product with \(\prod_{\tau,b}S_{\tau b}^{\beta_{\tau b}}\), using
modular repeated squaring. It permits arbitrary admissible powers and
arbitrary mixtures across tiles. For bounded-degree original vertices
and fixed \(b\le b_*\), construction costs
\(2^{k_t}V_g\operatorname{poly}(s,p_t)\) bit operations and
\(V_g\operatorname{poly}(s,p_t)\) workspace. A new requested tuple costs
\(V_g\operatorname{poly}(s,p_t)\) bit operations. With
\(D_t=L/2^{v_2(t)}\), there are at most \(D_t/2\) positions in each tile
and at most \(6M/D_t\) relevant tiles. Thus
\(\log_2 V_g=O_{b_*}((M/D_t)\log(D_t+1))\). At either outer odd node the
resulting time \(2^{O_{b_*}((M/L)\log L)}\operatorname{poly}(s,p_t)\) is
asymptotically below the printed fixed-density exponent. Writing a
complete explicit list of tuple answers also costs its output length.

At parity one table covers \emph{every} positive mode index.
\(S_{\tau0}\bmod2\) is the XOR of all tile bits, and
\(S_{\tau b}\bmod2\) for every \(b>0\) is the XOR of odd-\(j\) bits.
Positive powers agree with the bit. Replace the index vector by these at
most \(d_g\le12M/D_t\) features and addition by XOR. The same
distribution recurrence applies. For each feature subset its product
moment is the sum of entries whose vector equals one on that subset. A
Boolean superset transform obtains all these values in \(O(d_g2^{d_g})\)
additions by successively adding entries with a feature present into
those with it absent. The whole compressed parity menu for a child costs
\(2^{k_t+d_g}\operatorname{poly}(s)\) at bounded degree, and
\(2^{O(M/L)}\operatorname{poly}(s)\) at the outer nodes. This includes
simultaneous products of different tile moments. For unrestricted \(b\)
at higher precision the two-feature reduction fails, as
\(2\not\equiv2^2\pmod4\).
\href{https://github.com/openai/math/blob/main/preprints/The-Exponential-Time-Hypothesis-is-False-October-25-2026/build/quantitative-contraction.tex\#L432}{Native
packet tensors, lines 432--454 and 487--523}

\subsection{J6. Recovering a count from actual full-input zero modes of
one physical clause
network}\label{v9-j6-recovering-a-count-from-actual-full-input-zero-modes-of-one-physical-clause-network}

Consider the uniform problem of returning the residue of a specified
completed D root slot when the sparse clause input, its certified
graph-only packing and the legal slot address are given. A polynomial
number of such evaluations on a \emph{single literal physical CNF and a
single immutable packing} suffices to recover \(\#3{\rm SAT}\) by
deterministic polynomial-time Turing reduction. The construction uses a
large fixed polynomial enlargement to reach the source mesh. Its role is
a uniform hardness result for these actual registers. The v7 pinned
inputs, and a particular mandatory packing tie-break, have different
quantifiers.

Start with a 3CNF \(F\). Constant exceptional inputs are immediate.
Remove tautologies and repeated literals, retaining all declared
variables. With a fresh \(z\), guard each remaining clause \(C\) by
\(\neg z\lor C\), and add \(z\lor\neg x_j\) for every declared variable.
The resulting width-four CNF \(H\) is satisfied by the all-zero word and
has \(\#H=1+\#F\): if \(z=0\) all \(x_j=0\), while \(z=1\) gives exactly
the models of \(F\). All declared variables occur for the nontrivial
instance. Form its literal old clause and variable equality trees and
its new subdivision equalities exactly as in the sparse source. Take old
vertices as rows and subdivisions as columns. Write \(m\) for its number
of \emph{physical wires},
\(s_H=\#\mathrm{variables}+\#\mathrm{clauses}\), and \(q_H=m/2\) for the
number of subdivision vertices. We have \(s_H\le m=O(s_H)\): every
variable and clause occurs, twice the number of variable-clause
incidences dominates their combined count, and any internal
equality-tree edges only increase \(m\). Every old row tensor is one at
its all-zero input.

Define the degree-at-most-\(m\) rational polynomial \[
 G_H(w)=\sum_{X,Y}A_H(Y)B_H(X)
                 \prod_{\text{row ports }x} l_{Y_x}(wX_x),
                 \qquad G_H(1)=\#H ,
\] with the original port identified on the right and left. At \(w=1\)
the Boolean identity wiring proves the last equality. We will recover
this polynomial through actual completed probes.
\href{https://github.com/openai/math/blob/main/preprints/The-Exponential-Time-Hypothesis-is-False-October-25-2026/build/quantitative-reduction.tex\#L65}{Source
clause encoding and physical tensors, lines 65--97} ·
\href{https://github.com/openai/math/blob/main/preprints/The-Exponential-Time-Hypothesis-is-False-October-25-2026/build/quantitative-contraction.tex\#L176}{Original
contraction and deformation, lines 176--204}

\textbf{One physical padding and packing.} Put
\(b=\lceil\log_2(m+2)\rceil\), \(\lambda_* =4b/\varepsilon\), \(u=4b\),
\(a=8b\), \(L=2^a\), and \(\mathcal D=64^{2a}=L^{12}\), with
\(\varepsilon=1/(128\cdot1024)\). Let the desired common load be
\(h_*=2^{\lambda_*}/\mathcal D\), an even multiple of \(L\) much larger
than \(Lm\). Use \(L\) disjoint copies of \(H\), giving every vertex of
these physical networks the same constant word \(\rho\). For each
ordered word pair \((A,B)\), add \(h_*/2\) disjoint negative unary
clauses on fresh variables; in cell \((\rho,\rho)\) use only
\((h_*-Lm)/2\). The two original row vertices of each padding gadget are
the negative clause and the source's unary variable table \((1,1)\);
give both word \(A\). Its subdivision equality column has word \(B\).
Thus each padding gadget contributes exactly two physical wires to that
cell. Every cell has load \(h_*\), with collision potential zero. The
total wire count is \(M=2^{\lambda_*}\), and the actual sparse-reduction
size is \[
 s=L s_H+2^{\lambda_*}-Lm \in (2^{\lambda_*-1},2^{\lambda_*}].
\] The sparse clause entrance uses the variable count of the padded
formula plus one for \(P_0\), with \(\ell=1\). Its stopping threshold
uses this padded variable-plus-clause count \(s\). The algebraic
factor in the readout still uses \(q_H=m/2\). Consequently the actual
source parameters at the root have the chosen
\(\lambda=\lceil\log_2s\rceil\), \(u,a,L\). The root is active for large
inputs. Its packing conditions follow from the printed schedule or
directly from
\(M\asymp s,\mathcal D=L^{12}, L=2^{2\varepsilon\lambda}\). The final
clause density is bounded, since the dominating padding has one variable
and one clause per gadget. The entire explicit formula, graph, addresses
and labels have polynomial length in the original input:
\(2^{\lambda_*}=(2^b)^{4/\varepsilon}\). The constant polynomial degree
is very large.

Allocate one full packet to each of the \(m\) planted physical wire
types, putting that wire of copy \(j\) in slot \(j\) for \(0\le j<L\).
Its template has only equal coordinate pairs, allowing this ordering.
The remaining positions of this same template still consist of whole
packets of size \(L\). Fill the other templates with the ordinary source
orientation and copy rules. All real ports of every remaining packet
belong to padding gadgets. This is one fixed source-certified layout for
every query in the reduction.
\href{https://github.com/openai/math/blob/main/preprints/The-Exponential-Time-Hypothesis-is-False-October-25-2026/build/quantitative-contraction.tex\#L114}{Packing
construction, lines 114--171}

\textbf{Values and decoding.} Take the \(m+1\) positive odd values
\(t=1,3,\ldots,2m+1<L\). All are active at coordinate zero. For the zero
mode at \(\zeta=0\), the surviving coefficient is \[
       c_t=(-1)^t \frac{(L+t)!(L-1-t)!}{L!(L-1)!}
           =(-1)^t\prod_{h=1}^t\frac{L+h}{L-h}.
\] It is a dyadic unit, negative for all the selected \(t\), with
strictly increasing magnitude. On each planted packet row layer \(j\)
hits column layer \(j+t\) when that layer exists. A matched layer
therefore contracts to \(G_H(c_t)\). Every one of the \(t\) unmatched
row layers contributes \(A_H(0)=1\). Each of the \(t\) unmatched column
layers contributes \(2^{q_H}\). For a padding subdivision equality
column \(v\), let \(n_v(t)\in\{0,1,2\}\) count its ports reached by
matched negative unary rows. Matched \((1,1)\) rows supply one,
unmatched negative unaries supply one, and a matched negative unary
supplies \(1-c_tX\). Each padding column bit sums independently and
gives \(1+(1-c_t)^{n_v(t)}>0\). Thus the exact rational whose residue is
this completed slot is \[
 Z_t=\left[\prod_{v\ {\rm padding\ column}}\bigl(1+(1-c_t)^{n_v(t)}\bigr)\right]
                 2^{tq_H}\,G_H(c_t)^{\,L-t}.                     \tag{J6-layer}
\] The bracket is nonzero and computable from the public padding packet
order.

Write \(c_t=n_t/d_t\) in lowest terms with \(d_t>0\) odd. The factorial
formula gives \(\max(|n_t|,d_t)\le(2L)^t\). Multiply the native residue
by \(d_t^M\). The desired rational is now an integer of absolute value
at most \(2^{2M}(2(2L)^t)^M\), by the independent full \(X,Y\)
expression with binary original tensors. This bound also covers empty
row or column scopes. For \(t\le2m+1\), its base-two logarithm is at
most \(M(3+t(a+1))\), which is eventually less than \(p_t-1\), since
\(p_t\ge24ML/u\), \(a=8b,u=4b,L\ge(m+2)^8\). The centered integer
representative recovers \(d_t^M Z_t\) exactly. Divide by \(d_t^M\) and
by the public nonzero factors in (J6-layer). Since \(L-t\) is odd, the
remaining rational has a unique real \((L-t)\)-th root, \(G_H(c_t)\),
including in the cases of zero or a negative value. Reduce the fraction
and take the integer roots of its numerator and denominator by exact
binary search. All operations above have polynomial bit complexity in
the enlarged physical input and hence also in the original input.
Ordinary exact univariate interpolation at the distinct \(c_t\) recovers
\(G_H(1)\). Subtract one to obtain \(\#F\).

Zero can be chosen among the legal residue-distinct binary-coordinate
interpolation lifts. If instead some already prescribed probe set omits
zero, take all its \(M(2L-1)+1\) native \textbf{zero-mode} slots at each
needed \(t\). Their common polynomial before probe substitution has
degree at most \(M(2L-1)\): remove crossing terms from each row's linear
kernel, multiply the \(M\) affine factors and sum the original Boolean
assignments. The \(\beta=0\) child products are precisely these
evaluations, since each target coefficient has degree zero in its mode
variables and the source exponent is zero. The source's unit-difference
Lagrange interpolation therefore recovers its value at zero modulo
\(2^{p_t}\) in polynomial time. No nonzero mode is sampled and no
alternative deformation is provided.
\href{https://github.com/openai/math/blob/main/preprints/The-Exponential-Time-Hypothesis-is-False-October-25-2026/build/quantitative-contraction.tex\#L250}{Valuation
and precision, lines 250--267} ·
\href{https://github.com/openai/math/blob/main/preprints/The-Exponential-Time-Hypothesis-is-False-October-25-2026/build/quantitative-contraction.tex\#L338}{Native
probes and interpolation, lines 338--346} ·
\href{https://github.com/openai/math/blob/main/preprints/The-Exponential-Time-Hypothesis-is-False-October-25-2026/build/quantitative-contraction.tex\#L487}{Mode
substitution, lines 487--523}

A uniform polynomial-time evaluator on all the specified legal inputs
would imply \({\rm FP}=\#{\rm P}\) in the usual sense for counting
functions. This interpolation gives a polynomial-query hardness
statement. A different construction encodes edge covers in one fixed
mesh point of a certified physical layout. Neither result implies an
unconditional time lower bound or hardness for the particular v7 pinned
family. \hyperref[v9-section-acquisition-cost]{Single-node physical
construction A4} In this construction the original padded formula is
typically very large. Its direct planted-layer identity makes its
original count easy given a count for \(H\), but supplies no method for
evaluating the unknown \(G_H(c_t)\) without either the native oracle or
another computation.

\subsection{Checks and unresolved
cases}\label{v9-checks-and-unresolved-cases}

The accompanying deterministic checker independently evaluates source
swap-orbit rows, pin-pattern cylinders, first-jet identities, literal
shift sums, mixed-moment tables, signed odd-root extraction and
interpolation. Finite packet and swap examples use small analogues of
the source arrays so their truth tables can be enumerated. There are 336
auxiliary rows, 28 singleton diagonal cells, four decoded first-jet sign
words, 1,351 unions of query cylinders, 84 direct full shift
contractions, 729 mixed moments at 13-bit precision and 32 parity
moments. The literal physical checks include two complete interpolations
at five and nine distinct nodes, which recover three and two models
respectively. Thirteen of the sixteen total sampled values of the
single-layer polynomial are negative. The asymptotic schedule and
certification bounds above are analytic.
\PaperSixEvidence{evidence/research_v9/checks/joint_algorithms.py}{Checker} ·
\PaperSixEvidence{evidence/research_v9/checks/joint_algorithms.json}{Replay output}

The partial matching at zero need not have polynomial width. The
extreme-node evaluator computes an actual full-precision register, while
Hermite origin recovery also needs other nodes and local coefficients. A
root product at a nonzero \(\zeta\) retains all same-group entries of
\(W\), so treating it as the same partial matching loses real terms.
Using the superset transform for characteristic-zero high moments also
loses the information in \(j^b\). A polynomial-time decision at the raw
pin-entry interface still computes the v7 factors, but raw entries do
not obey the support bound of J1. The gate lower bounds in earlier
versions and the time of the printed source evaluator are separate from
these query and uniform-hardness statements.

\section{The cost of acquiring a prescribed
value}\label{v9-section-acquisition-cost}

The index of a completed Appendix D root register specifies a graph, a
packing, a grid point, a probe and a mode tuple. Even its zero-mode
product can encode an exact counting computation. There are binary,
bounded-degree inputs on which \textbf{one specified register} equals
the independent-set count of an arbitrary cubic graph. The physical root
on these inputs is a disjoint union of stars. A second construction
fixes even the packing and graph as a function of the input length. At
the same precision other registers can be inexpensive. The distinction
is between computing a uniformly prescribed value and recovering the
terminal answer by any route.

\subsection{A1. Complexity of a prescribed root
register}\label{v9-a1-complexity-of-a-prescribed-root-register}

Use the original bipartite tensor entry point with explicit
polynomial-time root entries, polynomially bounded precision, and a
polynomial length graph and packing. Include the parameters in the
ordinary explicit size bound, for example by requiring the active source
size parameter to be polynomial in the input length. On an active root
the inequality \(M>H\) already enforces this for an explicitly listed
graph. All promises below use the root parameters of the source. Allow a
certified layout to be supplied, or the particular uniform layout
described below. Fix zero as one of the eligible probes before choosing
the tensor entries. The source allows any sufficiently large set of
distinct binary-coordinate residue lifts.
\href{https://github.com/openai/math/blob/main/preprints/The-Exponential-Time-Hypothesis-is-False-October-25-2026/build/quantitative-contraction.tex\#L11}{Coefficient
arithmetic and probes, lines 11--71} ·
\href{https://github.com/openai/math/blob/main/preprints/The-Exponential-Time-Hypothesis-is-False-October-25-2026/build/quantitative-contraction.tex\#L338}{Grid
and probes, lines 338--345}

\textbf{Upper bound.} Every individual completed root register on this
interface belongs to \(\mathrm{FP}^{\#\mathrm P}\), including signed
coefficients, nonzero modes, extension probes and \(\ell>1\). The same
upper bound holds for the sum over all modes at a fixed \((t,\zeta)\).

For a prescribed mode tuple \(\beta\), guess the independent row and
column Boolean port assignments \(Y,X\). In each row-port affine factor
choose its local affine term, or one of its targeted formal mode terms.
Reject unless the multiplicities of the selected variables are exactly
\(\beta\). According to the sign of that row factor, its local choice
has value \(1-\mathrm{local}\) or \(\mathrm{local}\), and a chosen mode
coefficient has the corresponding sign. For this witness multiply the
chosen factors, all old tensor entries and the source moments
\(L_Z(X)^{\beta_Z}\) in the source ring modulo \(2^{p_t}\). Choices are
indexed once each, even when two choices happen to have equal numeric
values. Mode addresses and the total number of row-port choices have
polynomial bit length. The local kernels and mode coefficients are
polynomial-time computable by the source's unit inversions. The source
moment powers are computed by repeated squaring. Use a fixed-width
encoding of choices and reject invalid choices to give each valid
expansion exactly one witness. Distributivity and the
\nolinkurl{one-source/one-target} definition show that the sum of these
weights is exactly the prescribed product of child contractions. This
expands one mathematical coefficient for a counting reduction. No
inference about the algorithm's actual operation count uses this
expansion.
\href{https://github.com/openai/math/blob/main/preprints/The-Exponential-Time-Hypothesis-is-False-October-25-2026/build/quantitative-contraction.tex\#L206}{Local
kernels, lines 206--242} ·
\href{https://github.com/openai/math/blob/main/preprints/The-Exponential-Time-Hypothesis-is-False-October-25-2026/build/quantitative-contraction.tex\#L402}{Modes
and children, lines 402--523}

For any output coordinate let \(w\) be the least nonnegative residue of
this witness's weight. A nondeterministic polynomial-time machine
additionally guesses a \(p_t\)-bit integer \(d\) and accepts exactly
when \(d<w\). Its accepting-path count, reduced modulo \(2^{p_t}\), is
the coordinate of the root register. This handles negative witnesses and
all carries. There are polynomially many coordinates. For the sum over
the modes use the original deformed product at \(t+\alpha\zeta\)
directly, without a mode choice. Each witness still has polynomial
length. This upper bound requires polynomial-time \textbf{original}
entries. The full contraction theorem permits more expensive original
oracles.
\href{https://github.com/openai/math/blob/main/preprints/The-Exponential-Time-Hypothesis-is-False-October-25-2026/build/quantitative-contraction.tex\#L194}{Root
sum, lines 194--204} ·
\href{https://github.com/openai/math/blob/main/preprints/The-Exponential-Time-Hypothesis-is-False-October-25-2026/build/quantitative-reduction.tex\#L26}{General
entries, lines 26--49}

\textbf{Single-register hardness.} There is a polynomial-time reduction
from cubic independent-set counting to \emph{one} prescribed completed D
root register with \((t,\zeta,\beta)=(-1,0,0)\). All original entries
are binary. The source maximum degree is three and the number of
physical wires is linear in the starting graph size. The prescribed
physical graph has only star components. Counting independent sets of
cubic graphs is \(\#\mathrm P\)-complete under polynomial-time counting
reductions. Therefore a uniform polynomial-time evaluator for these
named registers would imply \(\mathrm{FP}=\#\mathrm P\) in the usual
sense that every \(\#\mathrm P\) function is polynomial-time computable.
The result concerns exact acquisition on this input class.
\href{https://web.maths.unsw.edu.au/~csg/papers/complexity.pdf}{Greenhill,
Theorem 3.1}

At \(t=-1\) the Cauchy entry is zero unless \(y=x-1\) in the same
packet. Its value there is \(\mathcal C'(-1)/\mathcal C'(0)=-1\). The
equality follows by cancelling the two factorial products on the grid
\(-L,\ldots,L-1\). All crossing-mode coefficients are zero at probe
zero, as their numerator is \(\mathcal C(-1)=0\) and the crossing
denominator is nonzero. Hence the actual zero-mode completed product
equals the whole deformed evaluation. The row spin \(Y\) facing a
matched column spin \(X\) carries the matrix with row indices \(Y=0,1\)
and column indices \(X=0,1\) \[
             K_{Y,X}=l_Y(-X)=
             \begin{pmatrix}1&2\\0&-1\end{pmatrix}.
\tag{A1}
\] Unmatched rows have argument zero.
\href{https://github.com/openai/math/blob/main/preprints/The-Exponential-Time-Hypothesis-is-False-October-25-2026/build/quantitative-contraction.tex\#L176}{Cauchy
polynomial, lines 176--204} ·
\href{https://github.com/openai/math/blob/main/preprints/The-Exponential-Time-Hypothesis-is-False-October-25-2026/build/quantitative-contraction.tex\#L387}{Crossing
geometry, lines 387--442}

Take a cubic graph \(H\). Introduce one simulated binary NAND row per
edge and one ternary equality column per vertex. The simulated wire for
each incidence should put a row port opposite a prescribed column port.
Replace each such wire by two separate physical wires. At even position
\(y\) put the prescribed column and a new unary all-ones dummy row. At
position \(x=y+1\) put the prescribed row and a new column pinned to
zero. Every remaining wire has a new all-ones row and a new zero-pinned
column. The next paragraph constructs these slots legally. An all-ones
row sums to one for any incoming argument since \(l_0+l_1=1\). Each
dummy column has precisely one allowed assignment. Each original NAND
row facing \(X_i,X_j\) evaluates to its multilinear extension at
\((-X_i,-X_j)\), which is \(1-X_iX_j\). Each old equality column chooses
one bit for its vertex. The integer before reduction is exactly
\(|\mathrm{IS}(H)|\). Taking \(\ell=1, P_0=|V(H)|+1\) ensures this
integer fits both the root and register precisions. At the origin every
physical NAND sees two pinned zeros. The terminal value there is
\(2^{|V(H)|}\). There is no count oracle among the leaf tables.

If the published choice of distinct probes omits zero, the polynomial
\(Q_{-1,0}(z)\) for the \emph{single zero mode} is still recoverable by
querying this same mode at the published probes. Its local affine
factors have coefficients in \(\mathcal O[z]\) and degree at most
\(2L-1\). The complete product before evaluating the probe thus has
degree at most \(M(2L-1)\). Interpolation over the residue ring at the
published \(M(2L-1)+1\) points recovers its constant term in polynomial
additional work. The one-call reduction needs zero as a selected probe.
This alternative uses polynomially many genuine zero-mode calls.

\textbf{Packing lemma.} Any bounded-degree simulated bipartite network
of \(m\) wires admits the even/odd arrangement above in an active source
root on \(O(m)\) physical wires. All ordered label-pair classes can be
made exactly equally loaded. The topology and the ordering depend only
on the simulated incidence graph.

Take \(r=64\) and \(s=16r^2m\) for sufficiently large \(m\), with
\(a=2u,L=2^a,P=r^{2a}\) from the printed root schedule. Give each
simulated vertex the same symbol \(g\) in coordinate zero. For the other
\(a-1\) coordinates, assign independent uniform words to simulated
vertices. With \(P'=r^{2(a-1)}\), the source pair-indicator argument
yields \[
       \sum_c (m_c-m/P')^2\le 2m\Delta_*,\qquad
       m_c\le m/P'+\sqrt{2m\Delta_*}\le2m/P'
\tag{A2}
\] by deterministic conditional expectation, where \(\Delta_*\) is the
fixed maximum simulated degree. The last inequality holds eventually
because \(P'\le s^{1/1000}\). In particular, a common class contains no
more than \(2m/P'\) simulated incidences. Choose \[
       h=L\left\lfloor\frac{s}{PL}\right\rfloor,\qquad M=Ph .
\] Then \(M=(1-o(1))s\) and \(h\ge s/(2P)\ge8m/P'\) eventually. Allocate
\(h\) physical positions in every ordered full class using the source's
ordinary template packing. In every class used by a simulated wire, the
two coordinate-zero symbols coincide. Its \(h\) positions therefore
occur in \(h/2\) adjacent even/odd pairs. Changing bit zero preserves
every higher label and the equal-symbol coordinate is free. Assign a
distinct pair to each incidence of that class. For its physical even
wire give the dummy row the simulated row label, and for its physical
odd wire give the dummy column the simulated column label. Fill all
other positions with fresh isolated unary pairs with exactly the
required labels. Every physical ordered class now has \(h\) real wires.
Its collision potential is zero and the padded maximum rounded to a
multiple of \(L\) is precisely \(h\). The printed root inequalities hold
for large \(s\), since \(M\sim s\), \(P,L\le s^{1/1000}\),
\(\Delta=\lceil s^{9/10}\rceil\), and root degree is bounded. The lists
and conditional expectations have polynomial length and construction
time. A fixed finite range can be counted directly in the reduction, or
use constant padding.
\href{https://github.com/openai/math/blob/main/preprints/The-Exponential-Time-Hypothesis-is-False-October-25-2026/build/quantitative-contraction.tex\#L73}{Packing
and balance, lines 73--170}

\subsection{A2. One packing fixed before the input
switches}\label{v9-a2-one-packing-fixed-before-the-input-switches}

The same shift simulates an arbitrary Boolean constraint \(f\) on \(d\)
incident column spins using only a single binary-valued physical row of
degree \(2d\) and \(d\) additional simulated unary columns pinned to
one. Write its old row table as \[
       G_f(Y,Z)= f(Y_1Z_1,\ldots,Y_dZ_d).
\tag{A3}
\] Pair each \(Y_j\) with the intended source bit \(X_j\) and each
\(Z_j\) with a pinned-one bit, using the packing lemma. For any
one-coordinate function \(b\) and \(x\in\{0,1\}\), \[
        \sum_{y,z\in\{0,1\}} b(yz)l_y(-x)l_z(-1)=b(x).
\tag{A4}
\] When \(x=0\), only \(y=0\) remains and the two \(z\)-weights sum to
one. When \(x=1\), the \(y=0\) contribution is \(2b(0)\) and the \(y=1\)
contribution is \(-2b(0)+b(1)\). Applying this identity in each
coordinate proves that the contracted row gives exactly \(f(X)\). All
tables in (A3) take values zero and one, even though the internal
weights in (A4) have signs.

For an explicit fixed-graph family, list all pairs of \(n\) labelled
vertices. Create one Boolean occurrence of each vertex for every pair
containing it. In a public order join its occurrences by successive
equality constraints along a path. The single value for a vertex extends
uniquely. The row for each pair is \(1-x_ix_j\) if it is an edge and one
otherwise. An input graph changes one table bit per potential pair.
Simulate each equality row by (A3), using four ports and two fixed
auxiliary columns. A pair row can use its original two-port NAND or
all-ones table directly by (A1). Each occurrence in a column has at most
three ports and uses a single Boolean equality or a unary all-ones
table. The simulated graph, source packing and prescribed slot are
therefore uniformly fixed given \(n\). The input table bits alone encode
\(H\). There are \(O(n^2)\) physical wires, maximum degree four, and all
original entries are binary. At the origin the pair and equality rows on
zero arguments all equal one. Every nontrivial physical column belongs
to a star with all-ones leaves, so its sum is public. Thus the original
terminal count here is independent of the input graph. The shifted value
remains \(|\mathrm{IS}(H)|\). In particular the hardness survives
polynomial-time graph-only preprocessing, including polynomially many
values computed within that time. Unbounded preprocessing that leaves a
polynomial-length advice string gives a different, nonuniform complexity
statement.

The general construction also gives a parsimonious one-evaluation
reduction directly from \(\#3\mathrm{SAT}\). Replace repeated
occurrences of each logical variable by copies constrained along a path
of equalities, so each occurrence is in at most three constraints. Use
the Boolean table of each clause or equality in (A3). This preserves one
extension per assignment. Free variables with no occurrences can be
handled by multiplying by their known power of two, or by an independent
equality-column switch with one tautological constraint for a strictly
count-preserving encoding. The total number of wires before source
packing is linear in the number of literal occurrences plus declared
variables, and maximum old degree is six. Taking \(P_0=n+1\) suffices.
This uses the general source tensor entry point. The additional old
tables and the un-subdivided star topology distinguish it from the
frozen direct clause-to-network encoding.
\href{https://github.com/openai/math/blob/main/preprints/The-Exponential-Time-Hypothesis-is-False-October-25-2026/build/quantitative-reduction.tex\#L65}{Literal
physical encoding, lines 65--97}

\textbf{Aggregating all modes on literal clause inputs.} For any
graph-only packing returned by the source algorithm, a different uniform
oracle returns \(F(t+\alpha\zeta)\), the \emph{sum over all modes}, at a
prescribed active root node and supplied probe. These evaluations are
counting-hard on the ordinary sparse clause branch under polynomial-time
Turing reductions, even with the source's chosen layout and probe
ordering. An empty clause gives count zero immediately, and an empty
clause list gives its known free-variable count. Otherwise delete
tautological clauses and remove unused declared variables, retaining
their known power-of-two factor separately. Apply the frozen clause
conversion to this normalized 3-CNF; when all remaining variables occur
and clauses are nonempty, its literal incidence count is at least both
the variable and clause counts, so the resulting root is active at all
sufficiently large sizes. It uses \(s=n+m\), bounded-degree binary
tables and \(P_0=n+1\). The number of active \(t\)'s and of probes per
node is polynomial. Query each evaluation at its printed \(p_t\), use
source Lagrange interpolation and tapered Hermite recovery, then restore
the factor for unused variables. This gives the exact input count in
polynomial oracle time. For finitely many source sizes handled by direct
summation the reduction can directly return the count. Together with the
upper bound this classifies uniform exact evaluation relative to
\(\#\mathrm P\). One sum over all modes is a different oracle from one
mode-product query. The direct sparse-CNF reduction need not make any
individual summand hard.
\href{https://github.com/openai/math/blob/main/preprints/The-Exponential-Time-Hypothesis-is-False-October-25-2026/build/quantitative-reduction.tex\#L65}{Sparse
CNF reduction} ·
\href{https://github.com/openai/math/blob/main/preprints/The-Exponential-Time-Hypothesis-is-False-October-25-2026/build/quantitative-contraction.tex\#L250}{Tapered
Hermite and interpolation}

\subsection{A3. What state counting prices in an adaptive
execution}\label{v9-a3-what-state-counting-prices-in-an-adaptive-execution}

Let \(F\) be a target on a finite promise and \(A=|F(\mathcal X)|\).
Suppose a free completed menu has at most \(2^B\) states, and the
incoming input-dependent state has at most \(2^b\) states given that
menu. On a fixed menu and state, future addresses, stopping and the
final answer are deterministic functions of the previous query replies.
Signed arithmetic, division with its prescribed lift, arbitrary digit
tests and cross-route caching between replies are allowed. If a query at
a history \(h\) costs \(c_h>0\), and its number of possible \emph{new
reply values conditional on that history} is at most \(2^{\rho c_h}\),
then \[
             \max_{x\in\mathcal X}\sum_{h\in\mathrm{trace}(x)}c_h
             \ge \frac{\log_2 A-B-b}{\rho},\qquad \rho>0.
\tag{A5}
\] The conditional bound is relative to inputs which reach \(h\) with
this free state. If precision is refined, only the conditional number of
compatible refinements counts, irrespective of how elsewhere-produced
low residues were cached. Arbitrarily many zero-cost deterministic
operations may be contracted into edges of the query tree. A separately
conveyed input-dependent address, timing choice or cached payload is an
acquisition. For every fixed pair of initial states build the answer
tree. Assign a child the parent mass divided by the number of children.
Every terminal mass is at least \(2^{-\rho T}\) if its trace has cost at
most \(T\). Terminal masses sum to one. Thus at most \(2^{\rho T}\)
outputs can arise per initial state, and counting the initial states
proves (A5). Equivalently one may use the product of the branching
factors along a path and omit any artificial integral rounding of their
logarithms. Genuine operation costs must independently justify the
conditional bound.

On the v8 cubic caterpillar host, \(\log_2 A\ge(N-10)/12\) and the low
completed menu has \(B=o(N)\). Acquiring at most \(w\) novel Boolean
answers per charged operation, with all control determined as above,
gives a worst-case cost at least \(((N-10)/12-B-b)/w\). This prices
actual bit queries in the stated model. A high native word can carry
many coordinate bits in a single word call. A high-precision zero or a
copy can cost no novel input bits. Query cardinality alone cannot assign
a superpolynomial price to each high word. In particular, on precisely
the restricted family proving the distinct-output bound, the leaf bits
\(s_j\) can each be read from one fixed original binary clause entry.
The two terminal answers are computed by iterating \[
      (A_j,B_j)=((1+s_j)(A_{j-1}+B_{j-1}),A_{j-1}),\quad (A_0,B_0)=(2,1)
\] and returning \(3A_k\) or \(3(A_k+B_k)\) according to the final
switch. This costs \(O(k)\) input-bit reads and \(O(k^2)\) ordinary bit
operations, since the two integers have \(O(k)\) bits. The
matching-order read bound is tight up to constants on this witness if
arithmetic is free. It leaves the rest of the fixed host's availability
cube and the cost of D-only acquisition open.
\PaperSixEvidence{evidence/research_v8/applications/pooled_menu_barrier.md}{v8 counts
and promise} ·
\PaperSixEvidence{evidence/research_v8/models/native_closure.md}{Physical switch table
addresses}

The source's high precision, the actual number of input-dependent
states, and the cost of one requested value thus have separate
witnesses. On the general uniform problem A1--A2 give counting hardness
for a specific quantity. They neither force an optimizer to request that
quantity nor preclude a fast algorithm on particular inputs. In
particular, general \(\#\mathrm P\)-hardness gives no unconditional
superpolynomial lower bound, and the source's frozen reconstruction
retains its stated subexponential upper bound.

\subsection{A4. A single fixed mesh point on the literal sparse clause
branch}\label{v9-a4-a-single-fixed-mesh-point-on-the-literal-sparse-clause-branch}

The uniform single-register problem is also counting-hard on the frozen
clause-to-tensor conversion itself. The original CNF has disjoint
variable sets in its clauses, width at most four, and at most one
occurrence per variable. Its count is immediate at the origin. A
supplied certified layout and just the specified register \((-1,0,0)\)
suffice to recover the number of edge covers of an input planar
bipartite graph with degree two on one side and three on the other. The
reduction uses a zero probe and polynomially many padding clauses. Here
the choice of packing encodes the desired incidences. No claim concerns
the output of the source's particular label-selection or in-class
ordering rule on this formula.
\href{https://pierre.senellart.com/publications/amarilli2017conjunctive.pdf}{Edge-cover
hardness, Amarilli, Monet and Senellart, Theorem D.1} ·
\href{https://github.com/openai/math/blob/main/preprints/The-Exponential-Time-Hypothesis-is-False-October-25-2026/build/quantitative-reduction.tex\#L65}{Literal
construction with unary variable leaves and subdivision equalities}

Let \(H=(V,E)\) be such a graph, and let \(o\) count its vertices of odd
degree. The number \(o\) is even by the handshaking identity. Think of a
bit \(X_e=0\) as selecting the edge \(e\) for the cover. For each vertex
of degree two create an original clause with two \textbf{negative}
literals. For each degree-three vertex create an original clause with
four negative literals, its fourth port intended for an auxiliary bit
\(X_v\), and one positive unary clause intended to pin this same bit.
Give \textbf{every} literal occurrence so far its own fresh CNF
variable. With \(k=|E|+o\) there are \(I=2k\) hard row ports, two for
each intended edge or auxiliary bit. Add \(k\) negative unary clauses on
distinct new variables, one per intended bit. In the literal sparse
construction put old vertices on the row side, including each
occurrence's unary \((1,1)\) variable table. Every original incidence is
subdivided by a distinct two-port equality column. Designate the
subdivision of each of the \(k\) extra negative unary clauses as the
column for its intended bit. Its two physical wires end at its negative
clause and its all-ones variable row. Every fresh variable occurs once,
as required.

In a constant equal-label class, give each hard row port one even/odd
pair. Its hard row wire goes in the odd position and a \emph{different},
designated column wire goes in the preceding even position. Use the two
different physical wires of the same designated equality for its two
hard occurrences. They are distinct from all the hard row wires. Each
intended row now sees exactly its intended column bit with coefficient
\(-1\) at the prescribed mesh point. The other original endpoint on an
odd wire is a subdivision column of the hard occurrence, and the other
original endpoint on an even wire is a unary row. No activated hard row
or positive unary is left unmatched.

For an all-negative clause of even width \(d\), its multilinear table at
these arguments equals \(1-\prod_{j=1}^d(-X_j)=1-\prod_j X_j\). Each
positive unary becomes \(-X_v\), so it forces \(X_v=1\) at the price of
a sign. There are \(o\) such signs and their product is one. The hard
row factors therefore enforce that every vertex has at least one
incident bit zero. They sum on designated columns to
\(\#\mathrm{EdgeCover}(H)\).

Every old row outside this hard set is unary. An all-ones row always
contributes one. An unmatched negative unary also contributes one, and a
matched negative unary contributes \(1+X\) on its incoming column bit.
Every designated equality has both ports already used to feed hard rows,
so these other contributions concern only the other subdivision columns.
For each such column \(c\), compute \(d_c\in\{0,1,2\}\), the number of
its ports feeding matched negative unary rows, from the public order.
Summing its own bit gives \(1+2^{d_c}\). Consequently the \emph{single
actual completed zero-mode product} before reduction equals the positive
integer \[
    \#\mathrm{EdgeCover}(H)\prod_{\text{other subdivision columns }c}
                                   \bigl(1+2^{d_c}\bigr).
\tag{A6}
\] No unknown normalizing factor remains. This integer is at most
\(2^{|E|}5^{M/2}<2^{2M}\) on the constructed network of \(M\) physical
wires. The inequality \(|E|\le M/2\) suffices. The published \(p_t\) at
the active root exceeds \(2M\), so the least nonnegative residue is the
actual integer. Compute the public product and divide exactly, including
all its powers of two. This is integer postprocessing \emph{after} exact
recovery. At the origin all physical clauses use distinct variables, and
their independent counts multiply to \(\prod_C (2^{|C|}-1)\).

Here are the size and certification details. The unpadded literal
construction has \(M_0=2(I+k)=6k\) wires,
\(s_0=\#\mathrm{variables}+\#\mathrm{clauses}\le M_0\). Write
\(\delta=M_0-s_0\). Choose \[
    \lambda=2\lceil\log_2(M_0+2)\rceil,\quad
    M=2^\lambda,\quad u=\lfloor\varepsilon\lambda\rfloor,\quad
    a=2u,\quad L=2^a,\quad P=64^{2a},\quad h=M/P.
\] For sufficiently large input \(u\ge1\). The common load \(h\) is an
even multiple of \(L\), since \(13a<\lambda\), and
\(h\ge M^{1-24\varepsilon}>M_0\) eventually. Use the same full word on
every planted vertex, so the planted component belongs to one ordered
diagonal class. Fill this class by adding \((h-M_0)/2\) further negative
unary clauses, each on its own unused variable. For every other ordered
word pair \((A,B)\), add \(h/2\) such clauses, labeling both old unary
rows \(A\) and their subdivision column \(B\). Every added clause
supplies two wires and increases \(s\) by two. All cells now have
exactly \(h\) real wires, and \[
       s=M-\delta\in(2^{\lambda-1},2^\lambda].
\] Thus the actual CNF parameters give exactly the chosen
\(\lambda,u,a,L\). The ordinary template packing has zero collision
potential, no unoccupied slots, \(M\asymp s\) and an active root at all
sufficiently large sizes. The diagonal class has \(h/2\) disjoint
even/odd pairs, enough for the \(I\) prescribed pairs since \(h\ge6k\).
Fill other positions of that class arbitrarily with its remaining wires.
Other classes follow the source orientation rules. The formula, order
and rational/integer decoding work have polynomial size, with
\(M=O(M_0^2)\). Any fixed finite exception is handled directly. The
one-node reduction requires a native probe set containing zero. As in
A1, interpolating just that node's zero-mode polynomial yields the same
reduction with polynomially many actual zero-mode calls if the
prescribed probe set omits zero.

A4 supplies physical sparse-branch one-node hardness for a certified
order. The printed algorithm is free to choose another legal order, and
the origin on these particular formulas is inexpensive. Thus (A6) cannot
give a time lower bound for evaluating their origin or for a full signed
hybrid.

\PaperSixEvidence{evidence/research_v9/formal/ShiftedAcquisition_scope.md}{Formal algebra and scope} ·
\PaperSixEvidence{evidence/research_v9/checks/acquisition_cost.py}{Finite checker} ·
\PaperSixEvidence{evidence/research_v9/checks/acquisition_cost.json}{Recorded output} ·
\PaperSixEvidence{evidence/research_v9/reviews/cost_audit.md}{Independent technical review}

\section{Current theorem coverage and release evidence}\label{paper6-current-coverage}

The final R119 replay confirms the J2 probability endpoints in the
\hyperref[r47-v9-review-consequences]{source-review supplement} and the J6
source-precision count and clock in the
\hyperref[sec:retained-physical-execution]{retained-execution appendix}.
The theorem index gives their exact declarations and hypotheses.
The table retains the proof boundaries recorded at v8 checkpoint 1
and its native-menu bridge, and at v9 R17, R18, R20 and R21.

\begin{OfflineResearchTable}
\OfflineResearchHeaderRule
\noindent%
\begin{minipage}[t]{0.14074\OfflineResearchTableWidth}\raggedright
Identifier
\end{minipage}\hfill%
\begin{minipage}[t]{0.31165\OfflineResearchTableWidth}\raggedright
Written scope
\end{minipage}\hfill%
\begin{minipage}[t]{0.49261\OfflineResearchTableWidth}\raggedright
Verified boundary at the historical checkpoint
\end{minipage}%
\OfflineResearchRow
\nobreak
\OfflineResearchHeaderRule
\nobreak
\noindent%
\begin{minipage}[t]{0.14074\OfflineResearchTableWidth}\raggedright
v8 U / V8-S / V8-J
\end{minipage}\hfill%
\begin{minipage}[t]{0.31165\OfflineResearchTableWidth}\raggedright
Actual evaluation modules, scalar/common-cut bounds and large-residue jet capacity at the stated field-size and loss thresholds.
\end{minipage}\hfill%
\begin{minipage}[t]{0.49261\OfflineResearchTableWidth}\raggedright
The preserved finite algebra and exact-decoder components support these statements; this checkpoint does not claim new global replay of the full written DVR and large-residue theories.
\end{minipage}%
\OfflineResearchRow
\noindent%
\begin{minipage}[t]{0.14074\OfflineResearchTableWidth}\raggedright
v8 V8-A1--A4
\end{minipage}\hfill%
\begin{minipage}[t]{0.31165\OfflineResearchTableWidth}\raggedright
Exact projected jet matrix and Smith capacities for each actual DVR and input domain.
\end{minipage}\hfill%
\begin{minipage}[t]{0.49261\OfflineResearchTableWidth}\raggedright
The v8 proof checkpoint identifies the literal CRT image, active Newton matrix, upper bound, sharp full-domain width/rank witnesses and saturation. A proper sender image need not attain the universal width; equal residue cardinalities alone do not identify matrices over different DVRs.
\end{minipage}%
\OfflineResearchRow
\noindent%
\begin{minipage}[t]{0.14074\OfflineResearchTableWidth}\raggedright
v8 F / N / C1--C3
\end{minipage}\hfill%
\begin{minipage}[t]{0.31165\OfflineResearchTableWidth}\raggedright
Positive finite computation, source-qualified cross multiplication and one-coordinate information bounds.
\end{minipage}\hfill%
\begin{minipage}[t]{0.49261\OfflineResearchTableWidth}\raggedright
Finite gates, decoder kernels and selected native identities keep their stated coefficient, semantic and acquisition assumptions. This is not an unrestricted signed algorithm bound; the full F1 request construction is not the C6 no-request specialization.
\end{minipage}%
\OfflineResearchRow
\noindent%
\begin{minipage}[t]{0.14074\OfflineResearchTableWidth}\raggedright
v8 C4--C6
\end{minipage}\hfill%
\begin{minipage}[t]{0.31165\OfflineResearchTableWidth}\raggedright
Whole completed menu after low-bit projection, and exponentially many exact outputs on the same five-copy CNF topology.
\end{minipage}\hfill%
\begin{minipage}[t]{0.49261\OfflineResearchTableWidth}\raggedright
\path{V8PooledMenu.actual_pooled_decoder} connects the actual tape to the exact scalar range and \(B+b\ge(N-10)/12\); \path{no_menu_only_decoder} uses \(12B+10<N\) at a valid active root. The root layout and prescribed probes are supplied and legal before the switches; only whole nonzero-mode products erase. The checked split-core construction takes a cubic expanding core; the general existence theorem and acquisition costs remain outside these certificates.
\end{minipage}%
\OfflineResearchRow
\noindent%
\begin{minipage}[t]{0.14074\OfflineResearchTableWidth}\raggedright
v9 S / J1--J5
\end{minipage}\hfill%
\begin{minipage}[t]{0.31165\OfflineResearchTableWidth}\raggedright
Finite execution support, diagonal MAIN and specified mixed-query obstruction, tests and joint local evaluators.
\end{minipage}\hfill%
\begin{minipage}[t]{0.49261\OfflineResearchTableWidth}\raggedright
Exact and bounded-error query results require their own actual-path contracts and erased D interface; an arbitrary layout, high-precision D data or raw entries are not included. A favorable derivative-test layout is a separate sufficient condition.
\end{minipage}%
\OfflineResearchRow
\noindent%
\begin{minipage}[t]{0.14074\OfflineResearchTableWidth}\raggedright
v9 J6
\end{minipage}\hfill%
\begin{minipage}[t]{0.31165\OfflineResearchTableWidth}\raggedright
Uniform multi-node reduction through one planted CNF and one supplied packing, with the outer root mode \(\beta=0\).
\end{minipage}\hfill%
\begin{minipage}[t]{0.49261\OfflineResearchTableWidth}\raggedright
\path{J6R15SignedReadNative.printedFallback_raw_count} connects genuine recursive receipts in the prescribed nonzero-probe fallback, centered integer lift, odd-root extraction and final \(\#F\) under its source guards. This top-level certificate is not an unguarded uniform implementation of both paper branches or an oracle bit-time theorem. A4 is a separate single-mesh edge-cover reduction.
\end{minipage}%
\OfflineResearchRow
\noindent%
\begin{minipage}[t]{0.14074\OfflineResearchTableWidth}\raggedright
v9 J7--J9: recovery
\end{minipage}\hfill%
\begin{minipage}[t]{0.31165\OfflineResearchTableWidth}\raggedright
Seated four-clone source, subdivision equalities as rows, \(t=L-1\); actual singleton moments and signed first-coefficient recovery.
\end{minipage}\hfill%
\begin{minipage}[t]{0.49261\OfflineResearchTableWidth}\raggedright
R17 supplies guarded all-pin recovery. R21's \path{cubic_graph_implication_public_J9} fixes one graph-only seat layout, singleton choices and a \(10M+1\)-probe prefix before the signs vary. The decoder uses only the \(10M\)-bit receiver reductions of the original \(p_t\)-bit recursive products at those same addresses; it does not rerun the source at \(10M\) bits. Separate collection uses \(N(10M+1)\) request occurrences, or 64 times as many individually requested group receipts. This is not a bound on raw-slot acquisition cost.
\end{minipage}%
\OfflineResearchRow
\noindent%
\begin{minipage}[t]{0.14074\OfflineResearchTableWidth}\raggedright
v9 P1--P11 / J9: erasure
\end{minipage}\hfill%
\begin{minipage}[t]{0.31165\OfflineResearchTableWidth}\raggedright
Coefficient and receiving-ring budgets, independent-bit erasure and packing conditions; J9's single-node universal paper lower side.
\end{minipage}\hfill%
\begin{minipage}[t]{0.49261\OfflineResearchTableWidth}\raggedright
In R21 the entire common low menu, including zero mode, vanishes at \(q=2N=M/12\). For every sign and address this is a projection of the same original product and file, not a second run. No decoder given only that full menu recovers all signs. The generic theorem \path{printed_product_low_erases} also gives erasure at \(q\le p_t\), \(q\le M/6-3M/L\) on the legal oriented packing with the literal four-clone certificate. P10 is a separate binary-tensor witness; P11 retains its stated random law and simultaneous guards.
\end{minipage}%
\OfflineResearchRow
\noindent%
\begin{minipage}[t]{0.14074\OfflineResearchTableWidth}\raggedright
v9 J8/J9: compiled joint readout
\end{minipage}\hfill%
\begin{minipage}[t]{0.31165\OfflineResearchTableWidth}\raggedright
Joint first-coefficient evaluator on the certified source addresses, distinct from giving each completed register.
\end{minipage}\hfill%
\begin{minipage}[t]{0.49261\OfflineResearchTableWidth}\raggedright
The R18 signed equality \path{compiles_literal_recursive_readouts} is per address, not merely equal decoded pins. The four-clone scope gives at most \(16N(65+2M)M\) old-table assignment visits in the specified plan. The mathematical polynomial-bit implementation uses these factors and bounded width; the audit does not verify a full bit machine, raw-register cost, or preprocessing and decoder traffic.
\end{minipage}%
\OfflineResearchRow
\noindent%
\begin{minipage}[t]{0.14074\OfflineResearchTableWidth}\raggedright
v9 A1--A4 and R17 dispatch
\end{minipage}\hfill%
\begin{minipage}[t]{0.31165\OfflineResearchTableWidth}\raggedright
Specific register-acquisition and fixed-host hardness, conditional prices and supplied-layout edge-cover readout.
\end{minipage}\hfill%
\begin{minipage}[t]{0.49261\OfflineResearchTableWidth}\raggedright
Each acquisition theorem keeps its own original-entry and charged-event hypotheses. The R17 weight-tape theorem prices construction and dispatch given the layout/precision data; it does not price writing all tensor entries, recursive evaluation, leaves, or final interpolation.
\end{minipage}%
\OfflineResearchRow
\OfflineResearchHeaderRule
\end{OfflineResearchTable}

J9 scope and provenance. R21 applies to the actual triangle-replacement of a simple cubic graph. It uses the directed copy implication (false/source, true/successor), positive OR on original graph edges, and equality on each of the four trivalent clone-variable columns per pin. The R20 joined equality-cycle instance is historical, not a pointwise model of the original implication tensors. This was a mismatch in the Lean instantiation; the original written J9 argument is unchanged. Its generic four-clone erasure and receiver lemmas are used by R21.

The R21 joined theorem keeps the printed RootBounds for \(s=21N/2\), \(C=3\), a monic positive-degree modulus with irreducible mod-2 reduction, positive local truncation order \(\ell\) and creation precision, and its constructed graph-only packing at \(t=L-1\). This is neither arbitrary-layout high recovery nor a theorem at all centers; it does not forbid side information, price a bit machine, or imply a counting or SAT time bound.

\PaperSixEvidence{evidence/checkpoints/v8/Paper6_v8_Proof_Checkpoint_1.zip}{v8 proof checkpoint 1}\par
\PaperSixEvidence{evidence/checkpoints/v8/Paper6_v8_Native_Menu_Bridge.zip}{v8 native-menu and decoder bridge}\par
\PaperSixEvidence{evidence/checkpoints/v9/r17/Beyond_Rank_Budgets_v9_Lean_Round17.zip}{v9 R17 proof checkpoint}\par
\PaperSixEvidence{evidence/checkpoints/v9/r17/J9R17AllPinsEvidence.md}{v9 R17 all-pin evidence}\par
\PaperSixEvidence{evidence/checkpoints/v9/r18/J9R18CompiledEvidence.md}{v9 R18 compiled-branch evidence}\par
\PaperSixEvidence{evidence/checkpoints/v9/r18/audit/summary.json}{v9 R18 strict replay record}\par
\PaperSixEvidence{evidence/research_work/sources/source_manifest.json}{archived original source and hashes}\par
\PaperSixEvidence{evidence/README.md}{historical evidence inventory and preservation limits}\par
\PaperSixEvidence{evidence/checkpoints/v9/r20/Beyond_Rank_Budgets_v9_Lean_Round20.zip}{R20 historical equality-cycle package and reusable proofs}\par
\PaperSixEvidence{evidence/checkpoints/v9/r20/J9R20SameFileEvidence.md}{R20 proof scope}\par
\PaperSixEvidence{evidence/checkpoints/v9/r21/Paper6_J9_Literal_Implication_Patch.zip}{R21 literal-implication Lean patch}\par
\PaperSixEvidence{evidence/checkpoints/v9/r21/J9R21ImplicationEvidence.md}{R21 statement and assumptions}\par
\PaperSixEvidence{evidence/checkpoints/v9/r21/P6_J9_INDEPENDENT.md}{Independent R20 source-model review}\par
\PaperSixEvidence{evidence/checkpoints/v9/r21/audit/summary.json}{R21 strict Lean receipt}\par
\PaperSixEvidence{evidence/J9_OVERLAY_MANIFEST.json}{J9 supplemental evidence hashes}\par

\subsection{Featured theorem index}\label{paper6-featured-index}

This index supplements the current coverage table. The inherited table retains
its own checkpoint cutoff; the later v8 receipts below do not certify a wider
acquisition model.

\begin{OfflineResearchTable}
\OfflineResearchHeaderRule
\noindent%
\begin{minipage}[t]{0.16085\OfflineResearchTableWidth}\raggedright
Original result
\end{minipage}\hfill%
\begin{minipage}[t]{0.49261\OfflineResearchTableWidth}\raggedright
Verified interface and evidence
\end{minipage}\hfill%
\begin{minipage}[t]{0.29154\OfflineResearchTableWidth}\raggedright
Limits
\end{minipage}%
\OfflineResearchRow
\nobreak
\OfflineResearchHeaderRule
\nobreak
\noindent%
\begin{minipage}[t]{0.16085\OfflineResearchTableWidth}\raggedright
v9 J9
\end{minipage}\hfill%
\begin{minipage}[t]{0.49261\OfflineResearchTableWidth}\raggedright
One graph-fixed layout and public address set. Prescribed high-precision
observations recover all pin signs; projection of the same original registers makes the
entire low menu zero. \path{cubic_graph_implication_public_J9} on
\path{ImplicationPublicWorld}. \PaperSixEvidence{evidence/checkpoints/v9/r21/audit/summary.json}{R21
replay}; \PaperSixEvidence{evidence/checkpoints/v9/r21/J9R21ImplicationEvidence.md}{literal
implication scope}.
\end{minipage}\hfill%
\begin{minipage}[t]{0.29154\OfflineResearchTableWidth}\raggedright
Two modules and 13 checks. Keeps the directed implication tables and
the printed field, layout and size guards. Request count is distinct from acquisition
cost.
\end{minipage}%
\OfflineResearchRow
\noindent%
\begin{minipage}[t]{0.16085\OfflineResearchTableWidth}\raggedright
v8 C3
\end{minipage}\hfill%
\begin{minipage}[t]{0.49261\OfflineResearchTableWidth}\raggedright
MAIN factors can separate within a group; the Appendix D children use a
group-union cut. See \path{C3_literal_factor_cut} and \path{C3_source_readout_rank};
Checkpoint 8 extends the fixed-list readout and child-lift precision.
\PaperSixEvidence{evidence/showcase/Paper6_v8_Proof_Checkpoint_8.zip}{Checkpoint 8};
\PaperSixEvidence{evidence/showcase/C3_fixed_readout.md}{readout scope}.
\end{minipage}\hfill%
\begin{minipage}[t]{0.29154\OfflineResearchTableWidth}\raggedright
Five-copy source and
correct terminal parity. This is not the J9 family or a full global D
polynomial-compatibility or adaptive-cost theorem.
\end{minipage}%
\OfflineResearchRow
\noindent%
\begin{minipage}[t]{0.16085\OfflineResearchTableWidth}\raggedright
v8 N2, D path
\end{minipage}\hfill%
\begin{minipage}[t]{0.49261\OfflineResearchTableWidth}\raggedright
At a stable active physical root, keep the prescribed legal centre, probe
and mode. Choose a largest group there and continue along legal largest-group successors
to an actual all-bits leaf. Its edge count exceeds one thirty-second of the printed
cutoff. See \path{physical_terminal_at_native_path}.
\PaperSixEvidence{evidence/showcase/N2_D_scope.md}{Source scope};
\PaperSixEvidence{evidence/showcase/N2_D_audit_index.json}{four strict receipts}.
\end{minipage}\hfill%
\begin{minipage}[t]{0.29154\OfflineResearchTableWidth}\raggedright
Four modules and
14 checks. A largest-choice path, not every individual leaf, MAIN acquisition or full
bit-time. The printed source guards remain.
\end{minipage}%
\OfflineResearchRow
\noindent%
\begin{minipage}[t]{0.16085\OfflineResearchTableWidth}\raggedright
v8 N1, N2, N4, N6
\end{minipage}\hfill%
\begin{minipage}[t]{0.49261\OfflineResearchTableWidth}\raggedright
The original statements and versioned index remain in the paper. The
preceding additions concern specific verified source and precision interfaces.
\end{minipage}\hfill%
\begin{minipage}[t]{0.29154\OfflineResearchTableWidth}\raggedright
Do not
infer complete original-family coverage from a local D source lemma or from C3. Geometry,
leaf path, global source compatibility and timing remain distinct.
\end{minipage}%
\OfflineResearchRow
\OfflineResearchHeaderRule
\end{OfflineResearchTable}

\subsection*{Later source proof map}\label{paper6-late-coverage}
These receipts supplement the earlier coverage tables. They do not change
any theorem or its assumptions.

\textbf{v8 N2, Appendix D.} Checkpoint 21 keeps the actual requested
centre, probe and mode, follows a largest-group path to the original
all-bits leaf, and checks the printed visit exponent. Checkpoint 22
preserves the collision certificate on the physical labels selected by
the original full-64 label run. The largest-route proof does not cover
every selected leaf, and the label proof takes the formatted inputs for
each call as given. See the \PaperSixEvidence{evidence/closeout/v8/n2/scope.md}{path receipt}
and \PaperSixEvidence{evidence/closeout/v8/labels/scope.md}{label receipt}.

\textbf{v8 N8, four-copy old row.} Checkpoint 22 verifies the printed
single-coordinate obstruction on the original fixed packing and all
its legal probes and mode products, after projection to the stipulated
low precision. The actual OR matching target forces the stated bit
budget even for nonlinear decoding with one-sided information. It is
distinct from the five-copy pooled result and does not price acquisition.
See the \PaperSixEvidence{evidence/closeout/v8/n8/scope.md}{N8 proof receipt}.

\textbf{v9 J6, original gateway and ancestor reads.} R35 returns all
64 genuine group results at each prescribed nonzero outer probe and
feeds their product into the exact count readout under the existing
input and routing guards. R36 verifies actual ancestral origins and
paid page reads with the old creation precision, before the current
precision lift. Neither receipt creates all descendant records or
closes original whole-program bit-time or space.
See the \PaperSixEvidence{evidence/closeout/v9/r35/scope.md}{gateway receipt},
\PaperSixEvidence{evidence/closeout/v9/r36/round.md}{ancestor and read receipts}, and
\PaperSixEvidence{evidence/closeout/CURRENT_PROOF_CROSSWALK_zh.md}{detailed source crosswalk}.

\appendix
\section*{Appendices: source-snapshot evidence and assessments}
\addcontentsline{toc}{section}{Appendices: source-snapshot evidence and assessments}
The original scopes, contribution audits, review notes and
reproducibility records are kept under their versioned headings.
Every statement of formal status and every verification count below
describes its original snapshot, not the
\hyperref[paper6-current-coverage]{current proof evidence}.
The two sets of theorem identifiers still refer to different constructions.

\section{V8 archive: Scope and remaining source obligations}\label{v8-section-coverage}

Every row refers to an actual theorem with its stated input domain. No
row alone establishes a lower bound for every optimized implementation
of the complete signed interpolation program.

\begin{OfflineResearchTable}
\OfflineResearchHeaderRule
\noindent%
\begin{minipage}[t]{0.17955\OfflineResearchTableWidth}\raggedright
Interface
\end{minipage}\hfill%
\begin{minipage}[t]{0.22680\OfflineResearchTableWidth}\raggedright
Proved conclusion
\end{minipage}\hfill%
\begin{minipage}[t]{0.23625\OfflineResearchTableWidth}\raggedright
Gates and sharing covered
\end{minipage}\hfill%
\begin{minipage}[t]{0.30240\OfflineResearchTableWidth}\raggedright
Remaining requirement or escape
\end{minipage}%
\OfflineResearchRow
\nobreak
\OfflineResearchHeaderRule
\nobreak
\noindent%
\begin{minipage}[t]{0.17955\OfflineResearchTableWidth}\raggedright
Polynomial scalar receiver, one directed interface
\end{minipage}\hfill%
\begin{minipage}[t]{0.22680\OfflineResearchTableWidth}\raggedright
Exact
finite-precision capacity and target-specific minimum loss
\end{minipage}\hfill%
\begin{minipage}[t]{0.23625\OfflineResearchTableWidth}\raggedright
Arbitrary
integral polynomial DAG, repeated messages, multiplication of
overlapping dependencies
\end{minipage}\hfill%
\begin{minipage}[t]{0.30240\OfflineResearchTableWidth}\raggedright
Actual sender image, all shared information
and denominator losses must be recorded; arbitrary higher-digit tests
change the closure
\end{minipage}%
\OfflineResearchRow
\noindent%
\begin{minipage}[t]{0.17955\OfflineResearchTableWidth}\raggedright
Whole-jet polynomial receiver, one scalar readout
\end{minipage}\hfill%
\begin{minipage}[t]{0.22680\OfflineResearchTableWidth}\raggedright
Exact all-field
Newton--Smith capacity; explicit closed form for large residue fields
\end{minipage}\hfill%
\begin{minipage}[t]{0.23625\OfflineResearchTableWidth}\raggedright
Arbitrary jet multiplication and reuse before the final readout
\end{minipage}\hfill%
\begin{minipage}[t]{0.30240\OfflineResearchTableWidth}\raggedright
Intermediate coefficient extraction and reinsertion are different gates;
finite matrix cutoff depends on allowed loss and output precision
\end{minipage}%
\OfflineResearchRow
\noindent%
\begin{minipage}[t]{0.17955\OfflineResearchTableWidth}\raggedright
Linear observations and linear target on a full module
\end{minipage}\hfill%
\begin{minipage}[t]{0.22680\OfflineResearchTableWidth}\raggedright
Arbitrary
decoding iff kernel inclusion; finite Frobenius row criterion; adaptive
zero-path acquisition bound
\end{minipage}\hfill%
\begin{minipage}[t]{0.23625\OfflineResearchTableWidth}\raggedright
Arbitrary nonlinear postprocessing and
deterministic adaptive query selection
\end{minipage}\hfill%
\begin{minipage}[t]{0.30240\OfflineResearchTableWidth}\raggedright
Native CNF coefficient tuples
need not form a module; a structured promise uses its actual
differences
\end{minipage}%
\OfflineResearchRow
\noindent%
\begin{minipage}[t]{0.17955\OfflineResearchTableWidth}\raggedright
Fixed nonlinear hybrid block sum on one common matching restriction
\end{minipage}\hfill%
\begin{minipage}[t]{0.22680\OfflineResearchTableWidth}\raggedright
Sum of overlap-adjusted observable capacities must cover full target
rank
\end{minipage}\hfill%
\begin{minipage}[t]{0.23625\OfflineResearchTableWidth}\raggedright
Per-block nonlinear polynomial computation and shared messages
\end{minipage}\hfill%
\begin{minipage}[t]{0.30240\OfflineResearchTableWidth}\raggedright
A runtime bound requires a trace delivering the charged blocks; merging,
joint guards and bypasses must be charged explicitly
\end{minipage}%
\OfflineResearchRow
\noindent%
\begin{minipage}[t]{0.17955\OfflineResearchTableWidth}\raggedright
One D coordinate, sparse native MAIN private sums
\end{minipage}\hfill%
\begin{minipage}[t]{0.22680\OfflineResearchTableWidth}\raggedright
Exponential
private-point capacity and a conditional time bound with the published
MAIN exponent
\end{minipage}\hfill%
\begin{minipage}[t]{0.23625\OfflineResearchTableWidth}\raggedright
Unprescribed division of graph edges between operands,
fixed recombination, ordinary canonical D slot lifts
\end{minipage}\hfill%
\begin{minipage}[t]{0.30240\OfflineResearchTableWidth}\raggedright
Private points
and represented products must be covered by the actual trace; arbitrary
compressed full-register evaluation remains open
\end{minipage}%
\OfflineResearchRow
\noindent%
\begin{minipage}[t]{0.17955\OfflineResearchTableWidth}\raggedright
All completed D slots at one coordinate, low precision
\end{minipage}\hfill%
\begin{minipage}[t]{0.22680\OfflineResearchTableWidth}\raggedright
The entire menu
cannot recover the matching target, even with an arbitrary decoder
\end{minipage}\hfill%
\begin{minipage}[t]{0.23625\OfflineResearchTableWidth}\raggedright
All
allowed probes and mode products, unrestricted processing of the
available residues
\end{minipage}\hfill%
\begin{minipage}[t]{0.30240\OfflineResearchTableWidth}\raggedright
Published Hermite reconstruction uses greater
precision losses and several coordinates; general individual-child mixed
moments are not covered
\end{minipage}%
\OfflineResearchRow
\noindent%
\begin{minipage}[t]{0.17955\OfflineResearchTableWidth}\raggedright
All completed low-precision D root slots across all coordinates
\end{minipage}\hfill%
\begin{minipage}[t]{0.22680\OfflineResearchTableWidth}\raggedright
An
exact-count decoder needs a linear number of additional information bits
on a fixed cubic expanding family
\end{minipage}\hfill%
\begin{minipage}[t]{0.23625\OfflineResearchTableWidth}\raggedright
Entire pooled menu, all probes and
modes, arbitrary joint decoding and adaptive reads of that menu
\end{minipage}\hfill%
\begin{minipage}[t]{0.30240\OfflineResearchTableWidth}\raggedright
This
is a global output-range argument; full native precision, individual
children, raw input access and extra MAIN data require separate
accounting
\end{minipage}%
\OfflineResearchRow
\noindent%
\begin{minipage}[t]{0.17955\OfflineResearchTableWidth}\raggedright
Positive origin contraction of the switched CNF
\end{minipage}\hfill%
\begin{minipage}[t]{0.22680\OfflineResearchTableWidth}\raggedright
Exponentially many
actual scalar gates on an expander count slice
\end{minipage}\hfill%
\begin{minipage}[t]{0.23625\OfflineResearchTableWidth}\raggedright
Overlapping products,
DAG reuse, bounded-support positive local atoms; specified
small-footprint adaptive traces
\end{minipage}\hfill%
\begin{minipage}[t]{0.30240\OfflineResearchTableWidth}\raggedright
Complementary raw literal entries,
signed interpolation and modular identities without a bounded positive
lift bypass this model
\end{minipage}%
\OfflineResearchRow
\noindent%
\begin{minipage}[t]{0.17955\OfflineResearchTableWidth}\raggedright
Raw clause-table entries with free arithmetic wiring
\end{minipage}\hfill%
\begin{minipage}[t]{0.22680\OfflineResearchTableWidth}\raggedright
Explicit
simulation of Boolean circuits with constant overhead
\end{minipage}\hfill%
\begin{minipage}[t]{0.23625\OfflineResearchTableWidth}\raggedright
Local table
evaluations, positive two-rail arithmetic, finite precision and reuse
\end{minipage}\hfill%
\begin{minipage}[t]{0.30240\OfflineResearchTableWidth}\raggedright
This enlarged model does not isolate the interpolation method; a
universal barrier would need a further source-derived restriction
\end{minipage}%
\OfflineResearchRow
\OfflineResearchHeaderRule
\end{OfflineResearchTable}

\subsection{Inputs and quantifiers}\label{v8-inputs-and-quantifiers}

The earlier four-clone pinned graph family has count
\(\prod_{uv\in E}(s_u\lor s_v)\). It supplies common full-rank
restrictions and native-register tests, but the full count has a
linear-size Boolean circuit. The new physical switched CNF has a count
slice \(3\sum_{I\in\mathrm{IS}(G)}\prod_{v\in I}s_v\), and a separate
request-controlled decision slice. Its gate lower bound uses the whole
count slice on an expander, not the matching restriction alone.

A common shortened five-copy cubic family now supports the native MAIN/D
cross-product and coordinate-menu statements as well as the positive
gate theorem. Its matching parity matrix is the identity. The
pooled-menu theorem uses another restriction of this same full counting
function, retaining exponentially many exact answer values. It is not a
single-bit or SAT-decision lower bound.

A fixed common matching is chosen after graph-only layouts and before
the request signs. Input-dependent layouts that see the request signs
first require a new argument. A single-cut width theorem may choose a
matching for that cut; a sum of different-cut blocks must use the
simultaneous restriction theorem. The positive support theorem is a gate
statement on its specified input representation and does not prove a
signed mixed runtime lower bound on the same instances.

\subsection{Source and precision}\label{v8-source-and-precision}

Source extraction used the frozen TeX files in the historical v8
archive. Its \nolinkurl{SOURCE_MANIFEST.json} recorded the four source
hashes. Public availability is listed in the evidence catalogue. The original algorithm is not an
executable extracted by this Lean project. Direct MAIN and sparse D
entry points are covered at the indicated interfaces; later
data-dependent parity work in the general MAIN front end is not silently
treated as a graph-only layout.

For finite precision, equality of whole-slot residues permits
replacement by separated child lifts only after bounding the loss of
every multiplier. The native audit proves that inequality for ordinary
MAIN/D weights and records a counterexample for arbitrary lossy weights.
A statement about parity does not follow from a statement about integer
support: the saved small tensor example has nonzero integer values on
all four inputs and the OR truth table after reduction modulo two.

\subsection{Status}\label{v8-status}

The deliverable proves complete theorems for the interfaces above. It
does not prove that the OpenAI manuscript is false, that its full
technique cannot yield a polynomial-time algorithm, or that SAT lacks
polynomial algorithms. The exact counting, decision, acquisition,
representation, word-operation and bit-operation conclusions remain
separately identified.

\section{V8 archive: Prior work and contribution assessment}\label{v8-section-contribution-audit}

This comparison uses primary papers and their exact statements. It does
not certify worldwide priority. The v7 audit, including Nazarov,
Dias--Lima and Tulsyan, belongs to the historical record covered
by the public evidence catalogue.

\begin{OfflineResearchTable}
\OfflineResearchHeaderRule
\noindent%
\begin{minipage}[t]{0.23625\OfflineResearchTableWidth}\raggedright
Prior tool or result
\end{minipage}\hfill%
\begin{minipage}[t]{0.34020\OfflineResearchTableWidth}\raggedright
Source
\end{minipage}\hfill%
\begin{minipage}[t]{0.36855\OfflineResearchTableWidth}\raggedright
Consequence for this paper
\end{minipage}%
\OfflineResearchRow
\nobreak
\OfflineResearchHeaderRule
\nobreak
\noindent%
\begin{minipage}[t]{0.23625\OfflineResearchTableWidth}\raggedright
P-orderings and factorial ideals
\end{minipage}\hfill%
\begin{minipage}[t]{0.34020\OfflineResearchTableWidth}\raggedright
Bhargava, \emph{P-orderings and
polynomial functions on arbitrary subsets of Dedekind rings} (1997),
\href{https://doi.org/10.1515/crll.1997.490.101}{DOI}
\end{minipage}\hfill%
\begin{minipage}[t]{0.36855\OfflineResearchTableWidth}\raggedright
The scalar floor
weights and Newton normalizations are classical.
\end{minipage}%
\OfflineResearchRow
\noindent%
\begin{minipage}[t]{0.23625\OfflineResearchTableWidth}\raggedright
Multivariable polynomial functions over residue rings
\end{minipage}\hfill%
\begin{minipage}[t]{0.34020\OfflineResearchTableWidth}\raggedright
Li--Sha,
Theorems 3.2, 3.9, 3.11, \href{https://arxiv.org/abs/1704.04965}{paper}
\end{minipage}\hfill%
\begin{minipage}[t]{0.36855\OfflineResearchTableWidth}\raggedright
Product normal forms and their counting formulas precede the capacity
calculation.
\end{minipage}%
\OfflineResearchRow
\noindent%
\begin{minipage}[t]{0.23625\OfflineResearchTableWidth}\raggedright
Polynomial functions on dual numbers
\end{minipage}\hfill%
\begin{minipage}[t]{0.34020\OfflineResearchTableWidth}\raggedright
Al-Ezeh--Al-Maktry--Frisch,
Theorem 3.5, \href{https://arxiv.org/abs/1910.00238}{paper}
\end{minipage}\hfill%
\begin{minipage}[t]{0.36855\OfflineResearchTableWidth}\raggedright
The
one-variable length-two \nolinkurl{coefficient/derivative} model is a
close predecessor. The paper studies several jet arguments, terminal
readouts, arbitrary lengths and bounded division.
\end{minipage}%
\OfflineResearchRow
\noindent%
\begin{minipage}[t]{0.23625\OfflineResearchTableWidth}\raggedright
Polynomial functions on finite local rings
\end{minipage}\hfill%
\begin{minipage}[t]{0.34020\OfflineResearchTableWidth}\raggedright
Bulyovszky--Horváth,
Theorem 1 and Lemma 3,
\href{https://doi.org/10.1016/j.tcs.2017.09.002}{DOI}
\end{minipage}\hfill%
\begin{minipage}[t]{0.36855\OfflineResearchTableWidth}\raggedright
Shifted powers
on residue classes and selector polynomials already handle arbitrary
finite local rings. The projected readout and bounded division are the
particular application here.
\end{minipage}%
\OfflineResearchRow
\noindent%
\begin{minipage}[t]{0.23625\OfflineResearchTableWidth}\raggedright
Factorial-weighted fixed divisors
\end{minipage}\hfill%
\begin{minipage}[t]{0.34020\OfflineResearchTableWidth}\raggedright
Prasad--Rajkumar--Reddy, Proposition
3.3, Lemma 4.1 and Theorem 4.2,
\href{https://arxiv.org/abs/1803.05780}{paper}
\end{minipage}\hfill%
\begin{minipage}[t]{0.36855\OfflineResearchTableWidth}\raggedright
The integrality test
behind the weighted Newton matrix is an established fixed-divisor
mechanism; Smith reduction is standard.
\end{minipage}%
\OfflineResearchRow
\noindent%
\begin{minipage}[t]{0.23625\OfflineResearchTableWidth}\raggedright
Generating characters and finite-ring duality
\end{minipage}\hfill%
\begin{minipage}[t]{0.34020\OfflineResearchTableWidth}\raggedright
Wood (1999),
\href{https://doi.org/10.1353/ajm.1999.0024}{DOI};
Dougherty--Kör--Leroy,
\href{https://leroy.perso.math.cnrs.fr/Articles/chi10.pdf}{author paper}
\end{minipage}\hfill%
\begin{minipage}[t]{0.36855\OfflineResearchTableWidth}\raggedright
The double-annihilator identity and Frobenius ring condition are
established algebra. Unequal precision and acquisition interpretations
retain explicit domain hypotheses.
\end{minipage}%
\OfflineResearchRow
\noindent%
\begin{minipage}[t]{0.23625\OfflineResearchTableWidth}\raggedright
Gate content and complementary contexts
\end{minipage}\hfill%
\begin{minipage}[t]{0.34020\OfflineResearchTableWidth}\raggedright
Jerrum--Snir, Theorem 3.2,
\href{https://snir.cs.illinois.edu/listed/J7.pdf}{author paper}
\end{minipage}\hfill%
\begin{minipage}[t]{0.36855\OfflineResearchTableWidth}\raggedright
The
support replacement argument has a classical antecedent.
\end{minipage}%
\OfflineResearchRow
\noindent%
\begin{minipage}[t]{0.23625\OfflineResearchTableWidth}\raggedright
Balanced Boolean-domain support covers
\end{minipage}\hfill%
\begin{minipage}[t]{0.34020\OfflineResearchTableWidth}\raggedright
Jukna, Lemma 8,
\href{https://web.vu.lt/mif/s.jukna/ftp/counting-web.pdf}{author paper}
\end{minipage}\hfill%
\begin{minipage}[t]{0.36855\OfflineResearchTableWidth}\raggedright
Correctness only on zero-one inputs and overlapping repeated variables
are not new model ideas.
\end{minipage}%
\OfflineResearchRow
\noindent%
\begin{minipage}[t]{0.23625\OfflineResearchTableWidth}\raggedright
Expander cover and avoidance probability
\end{minipage}\hfill%
\begin{minipage}[t]{0.34020\OfflineResearchTableWidth}\raggedright
Bova--Capelli--Mengel--Slivovsky, Corollary 2 and Theorem 4,
\href{https://arxiv.org/abs/1411.1995}{paper}
\end{minipage}\hfill%
\begin{minipage}[t]{0.36855\OfflineResearchTableWidth}\raggedright
The independent-set
exclusion bound is the complement of their stated vertex-cover bound.
\end{minipage}%
\OfflineResearchRow
\noindent%
\begin{minipage}[t]{0.23625\OfflineResearchTableWidth}\raggedright
Independent-set product geometry
\end{minipage}\hfill%
\begin{minipage}[t]{0.34020\OfflineResearchTableWidth}\raggedright
Korhonen, Lemma 10 and Theorem 1,
\href{https://doi.org/10.4230/LIPIcs.ICALP.2021.87}{proceedings}
\end{minipage}\hfill%
\begin{minipage}[t]{0.36855\OfflineResearchTableWidth}\raggedright
The
no-cross-edge observation is known; Booleanized support permits overlap
that tropical repeated weights exclude.
\end{minipage}%
\OfflineResearchRow
\noindent%
\begin{minipage}[t]{0.23625\OfflineResearchTableWidth}\raggedright
Arithmetic value matrices and induced matchings
\end{minipage}\hfill%
\begin{minipage}[t]{0.34020\OfflineResearchTableWidth}\raggedright
De Colnet--Mengel,
Theorem 2 and Lemma 14,
\href{https://proceedings.kr.org/2021/20/kr2021-0020-de-colnet-et-al.pdf}{proceedings}
\end{minipage}\hfill%
\begin{minipage}[t]{0.36855\OfflineResearchTableWidth}\raggedright
The matching rank method is established. Here the physical clause
transfer matrices and finite-precision interfaces are calculated
explicitly.
\end{minipage}%
\OfflineResearchRow
\noindent%
\begin{minipage}[t]{0.23625\OfflineResearchTableWidth}\raggedright
Caterpillar ratios and pair-to-scalar count range
\end{minipage}\hfill%
\begin{minipage}[t]{0.34020\OfflineResearchTableWidth}\raggedright
Chan--Heilman--Panova, Theorem 1.3 and its proof,
\href{https://arxiv.org/html/2604.19094}{paper}
\end{minipage}\hfill%
\begin{minipage}[t]{0.36855\OfflineResearchTableWidth}\raggedright
The encoding and the
square-root counting step are prior work. A qualitative fixed-host range
result also follows from their constructions. The present source
application uses a particular switched host inside a cubic expander and
the joint low-precision root menu.
\end{minipage}%
\OfflineResearchRow
\OfflineResearchHeaderRule
\end{OfflineResearchTable}

The following are the concrete candidate contributions to assess
further:

\begin{enumerate}
\def\labelenumi{\arabic{enumi}.}
\tightlist
\item
  The exact \nolinkurl{loss/precision/attained-image} decoder module and
  its target-specific minimum loss, including deficient jet images.
\item
  The all-residue-field projected jet capacity, computed by a finite
  weighted Newton matrix and ordinary DVR Smith form, with explicit
  small-field separation from the coefficient-minimum count.
\item
  The physical switched CNF shared by the native rank, low-precision and
  positive-circuit applications; source ownership, canonical-lift
  compatibility and precision thresholds are part of each application.
\item
  The pooled-coordinate exact-count obstruction, using a cubic expanding
  topology with exponentially many count values and a global encoding of
  the complete low-precision probe menu.
\item
  The positive contraction theorem with its real scalar-DAG, subcube,
  local-input, finite-precision and execution boundaries, and its Lean
  finite gate-count proof.
\end{enumerate}

None of these descriptions claims a general circuit lower-bound method
or an unrestricted obstruction to the source algorithm. In particular
the raw complementary-input simulation is classical Boolean arithmetic
in a source-specific legal-entry form. Its purpose is to locate a
necessary restriction, not to claim a new universal gate set.

The independent review files give fuller theorem-by-theorem comparisons.
The \hyperref[v8-section-final-targeted-audit]{targeted comparison}
identifies the closest all-field and fixed-host predecessors, including
overlap beyond the table's short descriptions. The new pooled-menu
construction and the full synthesis need outside expert review and a
broader specialized search before any novelty claim is used in a
submission. The mathematical arguments and exact checks are provided
independently of such a priority claim.

\section{V8 archive: Targeted comparisons for jet capacity and the pooled menu}\label{v8-section-final-targeted-audit}

The jet calculation specializes established interpolation and module
theory to a particular scalar output. The closest graph paper already
contains the binary continued-fraction construction and the passage from
pairs of terminal counts to individual counts. The source-specific
result here is the bound on how much the entire menu of completed
Appendix D products can reveal on a fixed family of physical inputs,
under the stated reductions in precision. The comparisons below identify
exact overlaps and deductions from earlier work. They support no claim
of worldwide priority.

\subsection{Projected polynomial functions on
jets}\label{v8-projected-polynomial-functions-on-jets}

Write \(P=L+p\), and let \(A_P=(\mathcal O/\pi^P)[z]/(z^k)\). A receiver
uses a polynomial in \(h\) whole jets with jet-ring coefficients. It
reads one scalar coefficient, divides by \(\pi^L\) on inputs where that
scalar is divisible, and keeps \(p\) output digits. This distinction
matters because attaining a scalar function and attaining a specified
whole jet function impose different equations.

\begin{OfflineResearchTable}
\OfflineResearchHeaderRule
\noindent%
\begin{minipage}[t]{0.27405\OfflineResearchTableWidth}\raggedright
Primary result
\end{minipage}\hfill%
\begin{minipage}[t]{0.67095\OfflineResearchTableWidth}\raggedright
Overlap and hypotheses
\end{minipage}%
\OfflineResearchRow
\nobreak
\OfflineResearchHeaderRule
\nobreak
\noindent%
\begin{minipage}[t]{0.27405\OfflineResearchTableWidth}\raggedright
Bulyovszky and Horváth, \emph{Polynomial functions over finite
commutative rings}, Theorem 1, Lemma 3, Theorem 2 and §§4.2--4.3
\end{minipage}\hfill%
\begin{minipage}[t]{0.67095\OfflineResearchTableWidth}\raggedright
For
every finite local ring with maximal ideal \(M\) of nilpotency index
\(t\), one-argument ring-valued polynomial functions are generated by
translates of \(1_M(x)x^j\) for \(0\le j<t\). Their proof separates
residue fibers. They also solve the table interpolation system and
handle nonprincipal maximal ideals by expanding over a coefficient ring.
This applies directly to \(A_P\) for every \(k,P\). Here \(M=(\pi,z)\)
has index \(P+k-1\). Products of the same indicators and multivariate
expansion give the \(h\)-argument generating statement immediately, with
total degree at most \(P+k-2\) on each fiber. Taking top coefficients of
the ring multiples of these generators already gives the one-argument
raw scalar span.
\href{https://riemann.math.unideb.hu/~ghorvath/publications/BuHo.pdf}{Author
manuscript} ·
\href{https://doi.org/10.1016/j.tcs.2017.09.002}{Journal}
\end{minipage}%
\OfflineResearchRow
\noindent%
\begin{minipage}[t]{0.27405\OfflineResearchTableWidth}\raggedright
Al-Ezeh, Al-Maktry and Frisch, Lemma 2.7, Theorem 3.5 and Corollary 3.7
\end{minipage}\hfill%
\begin{minipage}[t]{0.67095\OfflineResearchTableWidth}\raggedright
For one argument and \(z^2=0\), the top component of
\((f_0+zf_1)(a+zb)\) is exactly \(bf'_0(a)+f_1(a)\). Thus the raw
projected feature description in this case follows by taking a component
of their formula. Their equality test concerns the whole function and
also tests its constant component.
\href{https://arxiv.org/pdf/1910.00238}{Paper}
\end{minipage}%
\OfflineResearchRow
\noindent%
\begin{minipage}[t]{0.27405\OfflineResearchTableWidth}\raggedright
Al-Maktry, Theorems 3.4, 5.17 and 6.1, and Elliott, Theorems 1.1 and
2.10
\end{minipage}\hfill%
\begin{minipage}[t]{0.67095\OfflineResearchTableWidth}\raggedright
Al-Maktry treats one argument in an extension with several
infinitesimal directions, all pairwise products zero, and gives a
linear-system recognition test. Elliott treats a different extension
with \(\epsilon_i^2=0\) and products of distinct directions retained. He
identifies integer-valued polynomials and derivative conditions allowing
coefficients in a total quotient ring. These are substantive precedents
for nilpotents, readout coordinates and divisibility. The parameter
counting infinitesimal directions in either work has a different meaning
from the number of jet arguments here.
\href{https://arxiv.org/pdf/2002.01304}{Al-Maktry} ·
\href{https://arxiv.org/pdf/1608.00171}{Elliott}
\end{minipage}%
\OfflineResearchRow
\noindent%
\begin{minipage}[t]{0.27405\OfflineResearchTableWidth}\raggedright
Bhargava, \emph{P-orderings and polynomial functions}, and Li and Sha,
Theorems 3.2, 3.9 and 3.11
\end{minipage}\hfill%
\begin{minipage}[t]{0.67095\OfflineResearchTableWidth}\raggedright
The product Newton basis and factorial
weights for ordinary polynomial functions on residue rings are
classical. Moreover, at \textbf{raw} precision \(P=1\), taking the
Dedekind domain \(\kappa[z]\) and the ideal \((z^k)\) already gives a
multivariate normal form for whole-jet-valued polynomial functions of
arbitrary jet length. Projection of its generators gives an immediate
finite computation of the scalar image. For \(P,k>1\) the ideals
\((\pi)\) and \((z)\) in \(A_P\) are incomparable, so the chain-ring
argument applied to the whole double-truncated ring would require a new
justification.
\href{https://doi.org/10.1515/crll.1997.490.101}{Bhargava} ·
\href{https://arxiv.org/pdf/1704.04965}{Li and Sha}
\end{minipage}%
\OfflineResearchRow
\noindent%
\begin{minipage}[t]{0.27405\OfflineResearchTableWidth}\raggedright
Prasad, Rajkumar and Reddy, Proposition 3.3, Lemma 4.1 and Theorem 4.2
\end{minipage}\hfill%
\begin{minipage}[t]{0.67095\OfflineResearchTableWidth}\raggedright
Proposition 3.3 supplies finite tests for bounded-degree integer-valued
polynomials under its nonvanishing determinant hypotheses. These hold on
the full DVR digit grid. On a Cartesian product, Lemma 4.1 and Theorem
4.2 give the triangular evaluation test and express the fixed divisor of
a primitive polynomial as the ideal generated by its factorial-weighted
Newton coefficients. Factoring out the common content over a DVR covers
any nonzero polynomial. Applied to every integral linear combination of
the projected Taylor features, this gives the weighted divisibility test
underlying (V8-A2). \href{https://arxiv.org/pdf/1803.05780}{Paper}
\end{minipage}%
\OfflineResearchRow
\OfflineResearchHeaderRule
\end{OfflineResearchTable}

Higher jet substitution itself also has a general treatment. Vojta's
Lemma 1.7 and Proposition 5.1 identify the coefficient laws of a
polynomial in several jet arguments over any base ring via
Hasse--Schmidt derivations. The object there is the universal
coefficient algebra, whereas (V8-A2) identifies functions after evaluation
on a specific finite-residue domain.
\href{https://math.berkeley.edu/~vojta/jets.pdf}{Author manuscript}

\textbf{Contribution position.} Write the allowed projected features in
the monic Newton basis as columns of \(C\). The classical fixed-divisor
test reads their integral linear combinations through
\(M=\operatorname{diag}(\pi^{t_\nu})C\). A raw Smith direction of
exponent \(d\) supplies the ideal \(\pi^{(d-L)_+}(\mathcal O/\pi^p)\)
after division and reduction. This last deduction, target membership in
Smith coordinates and counting nonzero cyclic summands are elementary
Smith theory. The manuscript specifies the particular feature columns,
computes them on a degree-bounded grid, proves attainment for the
indicated decoder and exhibits small-field cancellations, including
width 12 against a coefficient-minimum bound of 18. These are an
explicit specialization and examples. The formula gives a finite matrix
characterization for each fixed raw precision. For a proper sender image
the exact answer can be smaller.
\hyperref[v8-section-jet-newton-smith]{All-field theorem and examples}

\subsection{Exact counts on one host and the complete low-precision
menu}\label{v8-exact-counts-on-one-host-and-the-complete-low-precision-menu}

\begin{enumerate}
\def\labelenumi{\arabic{enumi}.}
\item
  \textbf{Rooted tree counts and two terminal values.} Law's equations
  (1)--(2) give the standard root-absent and total independent-set
  recurrences. Equation (C5) follows by adding a spine vertex and
  optionally one leaf. Chan, Heilman and Panova's Lemma 3.1 and
  Corollary 3.2 encode binary words by the ratio of root-present to
  root-absent counts, using the update \(r\mapsto 1/(d+r)\),
  \(d\in\{1,2\}\). In the proof of Theorem 1.3, equations (3.3)--(3.5)
  explicitly use pairs of terminal scalar counts and the square-root
  pigeonhole argument. Section 6.6 identifies the trees as caterpillars.
  This is a direct methodological predecessor of (C5)--(C6).
  \href{https://www.combinatorics.org/ojs/index.php/eljc/article/download/v17i1n18/pdf/}{Law,
  §1} · \href{https://arxiv.org/html/2604.19094v1\#S3}{Chan, Heilman and
  Panova, §3}
\item
  \textbf{A fixed subcubic host follows from that predecessor.} This
  deduction uses the actual gadgets in the proof of Lemma 3.1. Start
  from a singleton. A digit 1 adds a spine edge. A digit 2 adds a spine
  edge and a leaf at each end of that edge, because the auxiliary marked
  tree is a rooted edge and is glued at both ends of the extension.
  Among \(2m\) steps fix all odd digits to 1 and allow the even digits
  independent choices. The digit-2 edges are disjoint. All \(2^m\) trees
  hence occur as induced subgraphs of one spine with one optional leaf
  at each vertex. The first ratio is \(1/2\). At every subsequent step
  \(d=\lfloor1/r_{\rm new}\rfloor\) and \(r_{\rm old}=1/r_{\rm new}-d\),
  so the ratios distinguish the words. If \(a,b\) are the counts with
  the final root present and absent, its two availability settings give
  \(b,a+b\). Their pair determines \(a/b\), giving at least \(2^{m/2}\)
  scalar counts on that one host. Thus even the qualitative
  fixed-subcubic-host range needs only a short specialization of the
  earlier proof. The present single-leaf update uses every spine step
  for an independent bit and yields the explicit \((N-10)/12\) bound
  after its expanding cubic completion. Uniform expansion concerns the
  ambient host while the enabled induced subgraph can be a caterpillar.
  \href{https://arxiv.org/html/2604.19094v1\#S3}{Predecessor's gadget
  proof} · \hyperref[v8-section-pooled-menu-barrier]{Cubic completion and
  (C6)}
\item
  \textbf{What the ETH comparisons actually estimate.} Dias and Lima's
  Definition 5.1 assumes rational Hermite recovery for every polynomial
  through a specified degree. Section 5.1 explicitly explains why
  tapered modular recovery in Appendix D falls outside that model.
  Theorems 8.1 and 8.3 charge specified clause-product evaluation costs.
  Their Proposition 4.1 and Appendix B.1 \emph{do} examine Appendix D's
  prescribed higher root precision and estimate the certified
  running-time bound from mode tuple enumeration. Counting tuple indices
  at that precision does not control how many different menus of reduced
  slot \emph{values} varying inputs can produce.
  \href{https://arxiv.org/html/2611.00517v1}{Dias and Lima, §§4--5, 8
  and Appendix B.1}
\item
  \textbf{Other linear-functional and residue barriers.} Tulsyan's
  Theorem 3.7 bounds substitution-algebra dimension for a linear readout
  of a product representing the permanent. Proposition 3.9 bounds the
  number of CNF variables in a bounded-weight sum of counts representing
  the permanent by using its value on the all-ones matrix. Proposition
  5.2 treats the main-text measurement template, and Theorem 6.4 treats
  residues of restrictions of a set family, including Hamiltonian-cycle
  queries. These specify different inputs and observables. None of these
  statements supplies the joint value-range estimate for Appendix D's
  completed products.
  \href{https://hexagonmath.org/pdf/2611.00149v1}{Tulsyan, §§3.5, 5, 6
  and 11}
\end{enumerate}

\textbf{Contribution position.} The frozen source supplies the valuation
lower bound \(a+1-i\) for a crossing coefficient and the expansion into
full products over child contractions. The manuscript derives a
simultaneous encoding for \emph{every completed product} of one public,
graph-only root layout across all active nodes, modes and probes at
public precisions \(p_i\le a+1-i\). Nonzero mode tuples vanish there. At
each node the one remaining zero-mode product is a polynomial in \(2z\)
over the dyadic base ring, so \(p_i(p_i+1)/2\) bits encode its residues
at every probe. This produces (C4) and uses more information about the
observed slots than the number of their addresses. On the switched CNF
over the cubic host, the count range then forces \(b+B\ge(N-10)/12\) for
any exact decoder having at most \(2^b\) extra states per menu. This
final inequality is the ordinary counting principle. The substantive
source application is the joint encoding with its actual multiplication,
base ring, access rules and precision thresholds. The reconstruction in
the source uses higher precision. The inequality concerns exact-count
information and provides no time lower bound for acquiring it.
\href{https://github.com/openai/math/blob/main/preprints/The-Exponential-Time-Hypothesis-is-False-October-25-2026/build/quantitative-contraction.tex\#L385-L523}{Source,
mode expansion and completed products} ·
\href{https://github.com/openai/math/blob/main/preprints/The-Exponential-Time-Hypothesis-is-False-October-25-2026/build/quantitative-contraction.tex\#L338}{Precisions
used for recovery} · \hyperref[v8-section-pooled-menu-barrier]{Pooled
theorem}

\section{V8 archive: Formal verification and reproduction}\label{v8-section-formal-status}

\subsection{Completed clean replay}\label{v8-completed-clean-replay}

The historical v8 combined replay used Lean 4.24.0 and Mathlib commit
\nolinkurl{f897ebcf72cd16f89ab4577d0c826cd14afaafc7}. It compiled all
project modules into a fresh temporary directory, with warnings treated
as errors. The 46 preserved v6/v7 modules and their hashes were checked
before adding the seven v8 modules. The result is \textbf{53 successful
modules and 404 checked declaration-dependency reports}. The seven new
modules account for 97 reports; the previous 46 account for 307. These
numbers count selected \texttt{\#print\ axioms} declarations, not
distinct mathematical results. The authoritative machine-readable
records are \PaperSixEvidence{evidence/research_v8/formal/replay_results_v8.json}{replay results}, the
\PaperSixEvidence{evidence/research_v8/formal/replay_all.py}{replay script} and the individual files
in \texttt{replay\_logs/}.

No project declaration uses \texttt{sorry}, \texttt{admit}, a
project-local \texttt{axiom}, or \texttt{native\_decide}. Every reported
dependency belongs to Lean's usual \texttt{propext},
\texttt{Classical.choice}, and \texttt{Quot.sound} set. This is a
kernel-checked classical formalization; it is not an axiom-free
foundational development. The replay checks the pinned Mathlib checkout
identity and the locations of directly imported Mathlib compiled
objects. It does not hash all Mathlib source files or rebuild Mathlib
from its foundations. Historical project source hashes remain unchanged.

\subsection{What the new kernel proofs
establish}\label{v8-what-the-new-kernel-proofs-establish}

\begin{OfflineResearchTable}
\OfflineResearchHeaderRule
\noindent%
\begin{minipage}[t]{0.23625\OfflineResearchTableWidth}\raggedright
Module
\end{minipage}\hfill%
\begin{minipage}[t]{0.08505\OfflineResearchTableWidth}\raggedright
Checked reports
\end{minipage}\hfill%
\begin{minipage}[t]{0.31185\OfflineResearchTableWidth}\raggedright
Formal mathematical content
\end{minipage}\hfill%
\begin{minipage}[t]{0.31185\OfflineResearchTableWidth}\raggedright
Written identification still required
\end{minipage}%
\OfflineResearchRow
\nobreak
\OfflineResearchHeaderRule
\nobreak
\noindent%
\begin{minipage}[t]{0.23625\OfflineResearchTableWidth}\raggedright
\PaperSixEvidence{evidence/research_v8/formal/ArbitraryDecoder.lean}{\nolinkurl{ArbitraryDecoder.lean}}
\end{minipage}\hfill%
\begin{minipage}[t]{0.08505\OfflineResearchTableWidth}\raggedright
8
\end{minipage}\hfill%
\begin{minipage}[t]{0.31185\OfflineResearchTableWidth}\raggedright
Arbitrary set decoding of linear observations is equivalent to
kernel containment; the decoder on the image is linear. An actual
recursive adaptive query tree has a same-zero-path certificate.
\end{minipage}\hfill%
\begin{minipage}[t]{0.31185\OfflineResearchTableWidth}\raggedright
Ambient extension over general finite Frobenius rings,
precision-rectangle embeddings and the full-domain hypothesis for any
native application.
\end{minipage}%
\OfflineResearchRow
\noindent%
\begin{minipage}[t]{0.23625\OfflineResearchTableWidth}\raggedright
\PaperSixEvidence{evidence/research_v8/formal/NonlinearProgram.lean}{\nolinkurl{NonlinearProgram.lean}}
\end{minipage}\hfill%
\begin{minipage}[t]{0.08505\OfflineResearchTableWidth}\raggedright
5
\end{minipage}\hfill%
\begin{minipage}[t]{0.31185\OfflineResearchTableWidth}\raggedright
A topologically indexed register DAG with reuse and overlapping
multiplication has a multivariate-polynomial semantics with one variable
per acquired message. The Boolean two-rail AND identities hold.
\end{minipage}\hfill%
\begin{minipage}[t]{0.31185\OfflineResearchTableWidth}\raggedright
A
data-flow proof for a source program containing divisions, guards,
precision changes or nonlocal accesses; polynomial degree alone is not a
gate count.
\end{minipage}%
\OfflineResearchRow
\noindent%
\begin{minipage}[t]{0.23625\OfflineResearchTableWidth}\raggedright
\PaperSixEvidence{evidence/research_v8/formal/FiniteMessage.lean}{\nolinkurl{FiniteMessage.lean}}
\end{minipage}\hfill%
\begin{minipage}[t]{0.08505\OfflineResearchTableWidth}\raggedright
5
\end{minipage}\hfill%
\begin{minipage}[t]{0.31185\OfflineResearchTableWidth}\raggedright
An arbitrary decoder through a finite directed message has an explicit
indicator factorization; its field rank is bounded by the alphabet size,
yielding the stated bit bounds.
\end{minipage}\hfill%
\begin{minipage}[t]{0.31185\OfflineResearchTableWidth}\raggedright
The native message factorization and
full-rank target premise. The pooled-menu application uses range
cardinality, not this matrix-rank wrapper.
\end{minipage}%
\OfflineResearchRow
\noindent%
\begin{minipage}[t]{0.23625\OfflineResearchTableWidth}\raggedright
\PaperSixEvidence{evidence/research_v8/formal/NewtonCapacity.lean}{\nolinkurl{NewtonCapacity.lean}}
\end{minipage}\hfill%
\begin{minipage}[t]{0.08505\OfflineResearchTableWidth}\raggedright
19
\end{minipage}\hfill%
\begin{minipage}[t]{0.31185\OfflineResearchTableWidth}\raggedright
Product falling-factorial and binomial grids, unit-triangular
evaluation, ideal tests, weighted-matrix factorization, normalized node
grids, and diagonal precision/loss arithmetic.
\end{minipage}\hfill%
\begin{minipage}[t]{0.31185\OfflineResearchTableWidth}\raggedright
General DVR factorial
ideals, Smith-module decomposition and the entire projected-jet/CRT
capacity theorem. These have written proofs but are not claimed as fully
formalized.
\end{minipage}%
\OfflineResearchRow
\noindent%
\begin{minipage}[t]{0.23625\OfflineResearchTableWidth}\raggedright
\PaperSixEvidence{evidence/research_v8/formal/PositiveSupport.lean}{\nolinkurl{PositiveSupport.lean}}
\end{minipage}\hfill%
\begin{minipage}[t]{0.08505\OfflineResearchTableWidth}\raggedright
29
\end{minipage}\hfill%
\begin{minipage}[t]{0.31185\OfflineResearchTableWidth}\raggedright
For a graph satisfying the degree, expansion and finite parameter
inequalities, and input atoms whose every specialized positive monomial
has size below the stated threshold, scalar agreement on a switch
subcube forces the finite lower bound on distinct multiplication
registers. Shared registers, overlapping products and fixed switches are
allowed. A modular version checks the terminal-lift contract.
\end{minipage}\hfill%
\begin{minipage}[t]{0.31185\OfflineResearchTableWidth}\raggedright
The
physical CNF bijection, expander families, rounding to asymptotics and
any adaptive execution-cost contract.
\end{minipage}%
\OfflineResearchRow
\noindent%
\begin{minipage}[t]{0.23625\OfflineResearchTableWidth}\raggedright
\PaperSixEvidence{evidence/research_v8/formal/TruncatedProbeMenu.lean}{\nolinkurl{TruncatedProbeMenu.lean}}
\end{minipage}\hfill%
\begin{minipage}[t]{0.08505\OfflineResearchTableWidth}\raggedright
17
\end{minipage}\hfill%
\begin{minipage}[t]{0.31185\OfflineResearchTableWidth}\raggedright
Actual odd-denominator dyadic coefficients reduce into the
prescribed rings; weighted coefficient messages simultaneously
reconstruct every probe. The finite alphabet and matching-row bounds
hold, including auxiliary bits and both matching blocks.
\end{minipage}\hfill%
\begin{minipage}[t]{0.31185\OfflineResearchTableWidth}\raggedright
Source Cauchy
deformation, vanishing of nonzero root modes, the global all-coordinate
encoding and its asymptotic parameter substitution.
\end{minipage}%
\OfflineResearchRow
\noindent%
\begin{minipage}[t]{0.23625\OfflineResearchTableWidth}\raggedright
\PaperSixEvidence{evidence/research_v8/formal/CaterpillarRange.lean}{\nolinkurl{CaterpillarRange.lean}}
\end{minipage}\hfill%
\begin{minipage}[t]{0.08505\OfflineResearchTableWidth}\raggedright
14
\end{minipage}\hfill%
\begin{minipage}[t]{0.31185\OfflineResearchTableWidth}\raggedright
Explicit independent selections obey the recurrence; the output
is three times their count. Endpoint ratios and output pairs encode the
switches, and every exact finite-transcript decoder satisfies the range
and bit budgets.
\end{minipage}\hfill%
\begin{minipage}[t]{0.31185\OfflineResearchTableWidth}\raggedright
The cubic expander completion, five-copy CNF assembly
and identification with the entire native menu.
\end{minipage}%
\OfflineResearchRow
\OfflineResearchHeaderRule
\end{OfflineResearchTable}

The exact finite circuit statement, its constants and semantic scope are
in \PaperSixEvidence{evidence/research_v8/formal/PositiveSupport_scope.md}{PositiveSupport scope}. The
independent-set objects and remaining source steps are described in
\PaperSixEvidence{evidence/research_v8/formal/CaterpillarRange_scope.md}{CaterpillarRange scope}. The
\hyperref[v8-section-theorem-index]{theorem index} maps paper statements to
formal declarations and written proofs.
\nolinkurl{DECLARATION_INDEX.json} lists all 404 audited declarations
and their module-level dependency unions. Individual declaration reports
are in the corresponding \texttt{replay\_logs/} files.

In particular, the kernel does not merely accept a cover or a large-gate
premise for the positive theorem: it derives the support contexts,
separation, independent-set avoidance estimate and finite
multiplication-register bound from its stated circuit and graph
hypotheses. Conversely, successful compilation does not establish that
the source's signed interpolation satisfies positivity, or that its
full-precision catalog fits a low-precision message contract.

\subsection{Independent finite checks}\label{v8-independent-finite-checks}

All ten Python checkers were rerun in separate temporary directories
containing their source dependencies but no expected result files.
Assertions were enabled. Their outputs reproduce all 26,335 bytes of the
saved JSON records exactly. The recorded environment is CPython 3.12.14
with NumPy 2.3.5. The native fixture uses the declared fixed random seed
451. See \PaperSixEvidence{evidence/research_v8/logs/finite_replay.json}{finite replay}, its
\PaperSixEvidence{evidence/research_v8/logs/finite_replay.log}{text log}, and the
\PaperSixEvidence{evidence/research_v8/checks/replay_all_checks.py}{portable driver}.

These checks cover concrete dyadic jet matrices, finite-ring decoders,
positive support counts, native deformation coefficients, actual small
CNFs and the caterpillar completion. They are corroboration and
regression checks. They do not replace the general mathematical proofs
or prove a running-time lower bound outside the stated interfaces.

\subsection{Reproduction boundary}\label{v8-reproduction-boundary}

The historical v8 archive recorded project Lean sources, source hashes,
replay programs, Python fixtures, logs and paper sources. Its external
Lean and Mathlib distributions were pinned. The
\PaperSixEvidence{evidence/research_v8/reviews/reproduction_release.md}{release reproduction
review} describes that environment and its commands. The
\PaperSixCatalogue{public evidence catalogue} identifies the
records available for this edition.
Compilation in a new project output directory has already succeeded in
the pinned local installation. Acquiring the external pinned
distributions on a different machine is a separate prerequisite.

The historical v8 formalization reported partial coverage of the
general-DVR observable-module theorem, the arbitrary-residue-field jet
capacity and the source-to-model reduction. Its report did not include
a kernel verification of a signed hybrid execution lower bound. The last item is also mathematically
open in this work.

\section{V8 archive: Theorem index}\label{v8-section-theorem-index}

Names below refer to the original v8 companion proof notes. Their scalar S, large-residue J and all-residue A equation families are displayed here as V8-S, V8-J and V8-A, respectively. Statement labels U, F, N and C are unchanged. The rows are archival, not the current status.

\begin{OfflineResearchTable}
\OfflineResearchHeaderRule
\noindent%
\begin{minipage}[t]{0.16065\OfflineResearchTableWidth}\raggedright
Statement
\end{minipage}\hfill%
\begin{minipage}[t]{0.31185\OfflineResearchTableWidth}\raggedright
Result and hypotheses
\end{minipage}\hfill%
\begin{minipage}[t]{0.17955\OfflineResearchTableWidth}\raggedright
Proof and review
\end{minipage}\hfill%
\begin{minipage}[t]{0.29295\OfflineResearchTableWidth}\raggedright
Kernel scope
\end{minipage}%
\OfflineResearchRow
\nobreak
\OfflineResearchHeaderRule
\nobreak
\noindent%
\begin{minipage}[t]{0.16065\OfflineResearchTableWidth}\raggedright
U, U1--U5
\end{minipage}\hfill%
\begin{minipage}[t]{0.31185\OfflineResearchTableWidth}\raggedright
The actual evaluation module, bounded saturation and output
quotient give exact target membership, target minimum denominator loss
and maximal separation width in Smith coordinates.
\end{minipage}\hfill%
\begin{minipage}[t]{0.17955\OfflineResearchTableWidth}\raggedright
\hyperref[v8-section-framework]{Framework};
\PaperSixEvidence{evidence/research_v8/reviews/final_scope.md}{final scope audit}.
\end{minipage}\hfill%
\begin{minipage}[t]{0.29295\OfflineResearchTableWidth}\raggedright
Diagonal-coordinate and minimum-loss arithmetic in
\texttt{NewtonCapacity}; full DVR theorem written.
\end{minipage}%
\OfflineResearchRow
\noindent%
\begin{minipage}[t]{0.16065\OfflineResearchTableWidth}\raggedright
Scalar Theorem 1
\end{minipage}\hfill%
\begin{minipage}[t]{0.31185\OfflineResearchTableWidth}\raggedright
Exact universal polynomial-decoder capacity for
integral scalar messages, arbitrary polynomial DAG reuse and prescribed
\nolinkurl{denominator/output} budgets.
\end{minipage}\hfill%
\begin{minipage}[t]{0.17955\OfflineResearchTableWidth}\raggedright
\hyperref[v8-section-nonlinear-budget]{Scalar capacity};
\PaperSixEvidence{evidence/research_v8/reviews/algebra_independent.md}{algebra review}.
\end{minipage}\hfill%
\begin{minipage}[t]{0.29295\OfflineResearchTableWidth}\raggedright
DAG normal
form in \texttt{NonlinearProgram}; normalized-grid ingredients in
\texttt{NewtonCapacity}.
\end{minipage}%
\OfflineResearchRow
\noindent%
\begin{minipage}[t]{0.16065\OfflineResearchTableWidth}\raggedright
Scalar common-cut budget
\end{minipage}\hfill%
\begin{minipage}[t]{0.31185\OfflineResearchTableWidth}\raggedright
Capacities add over a common target
restriction, with leakage and scope-overlap penalties; freely designed
encodings show the one-way order can be attained.
\end{minipage}\hfill%
\begin{minipage}[t]{0.17955\OfflineResearchTableWidth}\raggedright
\hyperref[v8-section-nonlinear-budget]{Common-cut proof and escape}.
\end{minipage}\hfill%
\begin{minipage}[t]{0.29295\OfflineResearchTableWidth}\raggedright
Uses
preserved rank modules; application written.
\end{minipage}%
\OfflineResearchRow
\noindent%
\begin{minipage}[t]{0.16065\OfflineResearchTableWidth}\raggedright
Large-residue jet capacity
\end{minipage}\hfill%
\begin{minipage}[t]{0.31185\OfflineResearchTableWidth}\raggedright
Finite projected coefficient features give
a universal bound, with equality under the specified residue-size
condition.
\end{minipage}\hfill%
\begin{minipage}[t]{0.17955\OfflineResearchTableWidth}\raggedright
\hyperref[v8-section-jet-polynomial-capacity]{Jet theorem};
\PaperSixEvidence{evidence/research_v8/reviews/algebra_independent.md}{algebra review}.
\end{minipage}\hfill%
\begin{minipage}[t]{0.29295\OfflineResearchTableWidth}\raggedright
General jet
theorem written.
\end{minipage}%
\OfflineResearchRow
\noindent%
\begin{minipage}[t]{0.16065\OfflineResearchTableWidth}\raggedright
All-residue jet theorem
\end{minipage}\hfill%
\begin{minipage}[t]{0.31185\OfflineResearchTableWidth}\raggedright
A factorial-weighted Newton matrix and its
Smith exponents give exact projected jet capacity over every finite
residue field.
\end{minipage}\hfill%
\begin{minipage}[t]{0.17955\OfflineResearchTableWidth}\raggedright
\hyperref[v8-section-jet-newton-smith]{Newton--Smith
proof}; \PaperSixEvidence{evidence/research_v8/reviews/jet_extension_crosscheck.md}{independent
crosscheck}.
\end{minipage}\hfill%
\begin{minipage}[t]{0.29295\OfflineResearchTableWidth}\raggedright
Evaluation grids and matrix/precision ingredients in
\texttt{NewtonCapacity}; CRT and full DVR identification written.
\end{minipage}%
\OfflineResearchRow
\noindent%
\begin{minipage}[t]{0.16065\OfflineResearchTableWidth}\raggedright
Small-field examples
\end{minipage}\hfill%
\begin{minipage}[t]{0.31185\OfflineResearchTableWidth}\raggedright
At positive denominator loss a dyadic example has
exact width 12 while the coefficient bound is 18; arbitrary digit
decoding can escape polynomial capacities.
\end{minipage}\hfill%
\begin{minipage}[t]{0.17955\OfflineResearchTableWidth}\raggedright
\hyperref[v8-section-jet-newton-smith]{Hand calculation};
\hyperref[v8-section-nonlinear-budget]{scalar examples}.
\end{minipage}\hfill%
\begin{minipage}[t]{0.29295\OfflineResearchTableWidth}\raggedright
Exact finite
Python checks; not a general Lean jet theorem.
\end{minipage}%
\OfflineResearchRow
\noindent%
\begin{minipage}[t]{0.16065\OfflineResearchTableWidth}\raggedright
Arbitrary linear-domain decoder
\end{minipage}\hfill%
\begin{minipage}[t]{0.31185\OfflineResearchTableWidth}\raggedright
Kernel containment is necessary and
sufficient for set decoding; finite Frobenius duality gives ambient
linear extension; non-Frobenius and nonlinear-domain examples delimit
it.
\end{minipage}\hfill%
\begin{minipage}[t]{0.17955\OfflineResearchTableWidth}\raggedright
\hyperref[v8-section-duality-typeset]{Duality and counterexamples};
\PaperSixEvidence{evidence/research_v8/reviews/algebra_independent.md}{review}.
\end{minipage}\hfill%
\begin{minipage}[t]{0.29295\OfflineResearchTableWidth}\raggedright
\nolinkurl{arbitrary_decoder_iff}, \nolinkurl{imageDecoder_correct};
full Frobenius extension written.
\end{minipage}%
\OfflineResearchRow
\noindent%
\begin{minipage}[t]{0.16065\OfflineResearchTableWidth}\raggedright
Adaptive zero-path certificate
\end{minipage}\hfill%
\begin{minipage}[t]{0.31185\OfflineResearchTableWidth}\raggedright
On the full module, the common kernel
of the zero-path queries is contained in the target kernel. Target
row-span follows in the stated field or finite commutative Frobenius
matrix setting.
\end{minipage}\hfill%
\begin{minipage}[t]{0.17955\OfflineResearchTableWidth}\raggedright
\hyperref[v8-section-duality-typeset]{Duality proof}.
\end{minipage}\hfill%
\begin{minipage}[t]{0.29295\OfflineResearchTableWidth}\raggedright
Actual query-tree \texttt{same\_zero\_path} and
\nolinkurl{zero_path_certificate} prove kernel containment.
\end{minipage}%
\OfflineResearchRow
\noindent%
\begin{minipage}[t]{0.16065\OfflineResearchTableWidth}\raggedright
F1--F2
\end{minipage}\hfill%
\begin{minipage}[t]{0.31185\OfflineResearchTableWidth}\raggedright
Genuine bounded-occurrence switched CNFs realize
independent-set counts; matching restrictions give the specified
count/decision matrices with the characteristic qualification.
\end{minipage}\hfill%
\begin{minipage}[t]{0.17955\OfflineResearchTableWidth}\raggedright
\hyperref[v8-section-stronger-family]{Physical construction};
\PaperSixEvidence{evidence/research_v8/reviews/positive_independent.md}{independent review}.
\end{minipage}\hfill%
\begin{minipage}[t]{0.29295\OfflineResearchTableWidth}\raggedright
The new
physical assembly is written and checked on finite CNFs; preserved
pinned-family modules are a different construction.
\end{minipage}%
\OfflineResearchRow
\noindent%
\begin{minipage}[t]{0.16065\OfflineResearchTableWidth}\raggedright
F3--F4
\end{minipage}\hfill%
\begin{minipage}[t]{0.31185\OfflineResearchTableWidth}\raggedright
On a bounded-degree expanding graph, positive-coefficient
circuits satisfying the input monomial-size and finite parameter bounds
need exponentially many distinct multiplication registers on the count
cube and sufficiently large subcubes. Reuse and overlap are allowed;
adaptive traces and modular output require explicit contracts.
\end{minipage}\hfill%
\begin{minipage}[t]{0.17955\OfflineResearchTableWidth}\raggedright
\hyperref[v8-section-stronger-family]{Positive circuit proof};
\PaperSixEvidence{evidence/research_v8/reviews/positive_formal_semantics.md}{formal semantics audit}.
\end{minipage}\hfill%
\begin{minipage}[t]{0.29295\OfflineResearchTableWidth}\raggedright
\nolinkurl{finite_product_gate_bound},
\nolinkurl{finite_subcube_gate_bound},
\nolinkurl{finite_modular_gate_bound} derive the finite gate bounds.
Rounding and source execution are written.
\end{minipage}%
\OfflineResearchRow
\noindent%
\begin{minipage}[t]{0.16065\OfflineResearchTableWidth}\raggedright
F5
\end{minipage}\hfill%
\begin{minipage}[t]{0.31185\OfflineResearchTableWidth}\raggedright
A uniform expanding ambient topology contains the stated enabled
counting instances. Complexity reductions are distinguished from
unconditional time lower bounds.
\end{minipage}\hfill%
\begin{minipage}[t]{0.17955\OfflineResearchTableWidth}\raggedright
\hyperref[v8-section-stronger-family]{Embedding}.
\end{minipage}\hfill%
\begin{minipage}[t]{0.29295\OfflineResearchTableWidth}\raggedright
Written.
\end{minipage}%
\OfflineResearchRow
\noindent%
\begin{minipage}[t]{0.16065\OfflineResearchTableWidth}\raggedright
N1--N5
\end{minipage}\hfill%
\begin{minipage}[t]{0.31185\OfflineResearchTableWidth}\raggedright
Native group semantics, printed acquisition cost, sparse
cross-product budgets, loss of canonical lifts and their stated
cost-transfer conditions.
\end{minipage}\hfill%
\begin{minipage}[t]{0.17955\OfflineResearchTableWidth}\raggedright
\hyperref[v8-section-native-escape]{Native
registers}; \PaperSixEvidence{evidence/research_v8/reviews/native_independent.md}{native audit}.
\end{minipage}\hfill%
\begin{minipage}[t]{0.29295\OfflineResearchTableWidth}\raggedright
Written source extraction, with exact fixtures. No unrestricted
compressed-acquisition theorem.
\end{minipage}%
\OfflineResearchRow
\noindent%
\begin{minipage}[t]{0.16065\OfflineResearchTableWidth}\raggedright
N6--N9
\end{minipage}\hfill%
\begin{minipage}[t]{0.31185\OfflineResearchTableWidth}\raggedright
Tests of the earlier multiplication escape; obstructions for
one-coordinate low-precision root menus and the stated limited child
menu; raw-entry access escapes.
\end{minipage}\hfill%
\begin{minipage}[t]{0.17955\OfflineResearchTableWidth}\raggedright
\hyperref[v8-section-native-escape]{Native tests}.
\end{minipage}\hfill%
\begin{minipage}[t]{0.29295\OfflineResearchTableWidth}\raggedright
Finite-message and
weighted coefficient components; full native identification written.
\end{minipage}%
\OfflineResearchRow
\noindent%
\begin{minipage}[t]{0.16065\OfflineResearchTableWidth}\raggedright
C1--C3
\end{minipage}\hfill%
\begin{minipage}[t]{0.31185\OfflineResearchTableWidth}\raggedright
The same five-copy counting family supports positive-circuit,
cross-product and one-coordinate arbitrary-decoder statements.
\end{minipage}\hfill%
\begin{minipage}[t]{0.17955\OfflineResearchTableWidth}\raggedright
\hyperref[v8-section-common-native-family]{Common family};
\PaperSixEvidence{evidence/research_v8/reviews/common_native_crosscheck.md}{independent review}.
\end{minipage}\hfill%
\begin{minipage}[t]{0.29295\OfflineResearchTableWidth}\raggedright
\texttt{TruncatedProbeMenu} proves the equality matching row bound;
physical/source steps written.
\end{minipage}%
\OfflineResearchRow
\noindent%
\begin{minipage}[t]{0.16065\OfflineResearchTableWidth}\raggedright
C4
\end{minipage}\hfill%
\begin{minipage}[t]{0.31185\OfflineResearchTableWidth}\raggedright
All completed low-precision root slots, over every coordinate and
probe, are jointly determined by the stated sublinear number of base
bits.
\end{minipage}\hfill%
\begin{minipage}[t]{0.17955\OfflineResearchTableWidth}\raggedright
\hyperref[v8-section-pooled-menu-barrier]{Pooled menu};
\PaperSixEvidence{evidence/research_v8/reviews/pooled_menu_native_review.md}{native review}.
\end{minipage}\hfill%
\begin{minipage}[t]{0.29295\OfflineResearchTableWidth}\raggedright
Weighted \nolinkurl{contraction/evaluation} encoding cores in
\texttt{TruncatedProbeMenu}; pooling and source modes written.
\end{minipage}%
\OfflineResearchRow
\noindent%
\begin{minipage}[t]{0.16065\OfflineResearchTableWidth}\raggedright
C5--C6
\end{minipage}\hfill%
\begin{minipage}[t]{0.31185\OfflineResearchTableWidth}\raggedright
The caterpillar recurrence (C5) and a fixed cubic expander
completion give the exponential exact-count range (C6) on the same
physical CNF family.
\end{minipage}\hfill%
\begin{minipage}[t]{0.17955\OfflineResearchTableWidth}\raggedright
\hyperref[v8-section-pooled-menu-barrier]{Construction and recurrence};
\PaperSixEvidence{evidence/research_v8/reviews/caterpillar_range_independent.md}{range audit}.
\end{minipage}\hfill%
\begin{minipage}[t]{0.29295\OfflineResearchTableWidth}\raggedright
\texttt{selections\_card}, \nolinkurl{output_eq_selections},
\texttt{ratio\_injective}, \nolinkurl{output_range_square}; expander and
CNF completion written.
\end{minipage}%
\OfflineResearchRow
\noindent%
\begin{minipage}[t]{0.16065\OfflineResearchTableWidth}\raggedright
Pooled consequence of C4--C6
\end{minipage}\hfill%
\begin{minipage}[t]{0.31185\OfflineResearchTableWidth}\raggedright
Arbitrary decoding of the entire pooled
menu plus bounded extra information requires a linear total bit budget.
This is an exact-count interface impossibility, not a SAT or
high-precision time lower bound.
\end{minipage}\hfill%
\begin{minipage}[t]{0.17955\OfflineResearchTableWidth}\raggedright
\hyperref[v8-section-pooled-menu-barrier]{Pooled conclusion}.
\end{minipage}\hfill%
\begin{minipage}[t]{0.29295\OfflineResearchTableWidth}\raggedright
\nolinkurl{range_le_transcripts}, \nolinkurl{exact_decoder_bit_budget};
insert C4--C6 written identifications.
\end{minipage}%
\OfflineResearchRow
\noindent%
\begin{minipage}[t]{0.16065\OfflineResearchTableWidth}\raggedright
Raw-entry closure
\end{minipage}\hfill%
\begin{minipage}[t]{0.31185\OfflineResearchTableWidth}\raggedright
Legal complementary raw table entries and two-rail
arithmetic simulate Boolean circuits; an output-only definition cannot
inherit the positive-origin barrier.
\end{minipage}\hfill%
\begin{minipage}[t]{0.17955\OfflineResearchTableWidth}\raggedright
\hyperref[v8-section-native-closure]{Escape model}.
\end{minipage}\hfill%
\begin{minipage}[t]{0.29295\OfflineResearchTableWidth}\raggedright
AND rail identities
in \texttt{NonlinearProgram}; circuit simulation written.
\end{minipage}%
\OfflineResearchRow
\OfflineResearchHeaderRule
\end{OfflineResearchTable}

All newly checked declarations, as well as the preserved v6/v7 reports,
are listed with module, source hash and module-level standard axiom
unions in
\PaperSixEvidence{evidence/research_v8/formal/DECLARATION_INDEX.json}{\nolinkurl{DECLARATION_INDEX.json}}.
The \hyperref[v8-section-formal-status]{formal status} records the clean
replay and distinguishes formal ingredients from the full written
applications. The \hyperref[v8-section-coverage]{coverage matrix} states
which original execution mechanisms remain outside these theorems.

Classical antecedents and the closest source-specific comparisons appear
in the \hyperref[v8-section-contribution-audit]{contribution audit} and the
\hyperref[v8-section-final-targeted-audit]{targeted literature audit}. In
particular, Newton/factorial ideals, Smith and Frobenius arguments,
monotone support covers, rooted caterpillar ratios and the
pair-to-scalar counting step are not presented as inventions of this
project.

\section{V8 archive: The next unresolved interface}\label{v8-the-next-unresolved-interface}

The most consequential missing theorem concerns compressed evaluation of
the actual high-precision completed registers on immutable full-input
layouts. The two simple factors in the earlier multiplication escape can
be obtained from raw local entries or freshly chosen small networks.
Equality with efficient evaluations of the prescribed full registers has
not been established. Their printed acquisition routines and several
restricted replacements have lower bounds; a completely reorganized
evaluation may have a different cost.

There are concrete routes for further work.

\begin{enumerate}
\def\labelenumi{\arabic{enumi}.}
\tightlist
\item
  \textbf{Compute actual attainable-image profiles.} For a small family
  of genuine completed probes, form their evaluation lattice and test
  the required target columns using the minimum-loss criterion. A loss
  bound based only on the number of registers misses correlations; an
  unrestricted raw-bit encoding discards their source identity.
\item
  \textbf{Extend the coordinate menu to higher precision and mixed
  moments.} The proven low-digit menu kills every nonzero mode. At
  greater precision, products of different tile moments can supply
  additional functions. Determine their actual module and the cost of
  generating it, rather than treating every formal mode as independently
  acquired.
\item
  \textbf{Charge a signed computational interface.} A useful next
  invariant must track the original precision, canonical lifts, allowed
  routing and source-preserving operations. An additive gate potential
  would need an independently proved bound under multiplication and
  reuse. The present Smith profile measures achievable functions, not
  all such computational costs.
\item
  \textbf{Use the stronger CNF family without losing quantifiers.}
  Uniform hardness embeddings and positive contraction lower bounds do
  not provide an unconditional signed arithmetic lower bound. One common
  input family must force the selected native interfaces simultaneously,
  including their input-dependent layout rules and request signs.
\end{enumerate}

A meaningful next success may be a better exact evaluator for a genuine
register, a stronger source-specific impossibility certificate, or a
counterexample to a proposed cost transfer. The current results specify
what such an advance would have to change or preserve. No conclusion
about the validity of the entire ETH claim follows from the barriers in
this paper alone.

\subsection{Contributions and research
record}\label{v8-contributions-and-research-record}

The user initiated the problem, chose the research direction, insisted
on native interfaces and rigorous counterexamples, and set the delivery
and review requirements. ChatGPT coordinated the work and generated the
mathematical arguments, exact test programs, formal proofs and
manuscript with explicitly configured GPT-6.2 Astra Max agents. Separate
agents reviewed algebra, physical CNFs, source correspondence and model
scope; the coordinator integrated their corrections and ran independent
checks. This disclosure describes the actual workflow, not a substitute
for human mathematical responsibility or an authorship decision. No
external submission, repository publication or pull request was made.

\section{V9 archive: Scope and remaining obligations}\label{v9-section-coverage}

All conclusions refer to the archived source version and to the
interface definitions in this manuscript. A successful Lean replay
validates the formalized finite mathematical cores; it does not by
itself prove that an entire source execution satisfies a contract.

\begin{OfflineResearchTable}
\OfflineResearchHeaderRule
\noindent%
\begin{minipage}[t]{0.17955\OfflineResearchTableWidth}\raggedright
Interface
\end{minipage}\hfill%
\begin{minipage}[t]{0.22680\OfflineResearchTableWidth}\raggedright
Established result
\end{minipage}\hfill%
\begin{minipage}[t]{0.23625\OfflineResearchTableWidth}\raggedright
Layout and precision
\end{minipage}\hfill%
\begin{minipage}[t]{0.30240\OfflineResearchTableWidth}\raggedright
Unproved transfer
\end{minipage}%
\OfflineResearchRow
\nobreak
\OfflineResearchHeaderRule
\nobreak
\noindent%
\begin{minipage}[t]{0.17955\OfflineResearchTableWidth}\raggedright
Actual MAIN completed jets
\end{minipage}\hfill%
\begin{minipage}[t]{0.22680\OfflineResearchTableWidth}\raggedright
Adaptive exact and randomized query lower
bounds; arbitrary receiver operations.
\end{minipage}\hfill%
\begin{minipage}[t]{0.23625\OfflineResearchTableWidth}\raggedright
Constructed certified
diagonal-pin layout; exact rational values.
\end{minipage}\hfill%
\begin{minipage}[t]{0.30240\OfflineResearchTableWidth}\raggedright
Every possible MAIN
layout, raw input preprocessing, and acquisition of other
observables.
\end{minipage}%
\OfflineResearchRow
\noindent%
\begin{minipage}[t]{0.17955\OfflineResearchTableWidth}\raggedright
Actual individual MAIN groups
\end{minipage}\hfill%
\begin{minipage}[t]{0.22680\OfflineResearchTableWidth}\raggedright
Same support argument with per-group
exponent.
\end{minipage}\hfill%
\begin{minipage}[t]{0.23625\OfflineResearchTableWidth}\raggedright
All nodes/groups on that fixed layout.
\end{minipage}\hfill%
\begin{minipage}[t]{0.30240\OfflineResearchTableWidth}\raggedright
A joint producer
returning a different aggregate cannot be priced by merely counting
requested table entries.
\end{minipage}%
\OfflineResearchRow
\noindent%
\begin{minipage}[t]{0.17955\OfflineResearchTableWidth}\raggedright
Completed D products, guarded side
\end{minipage}\hfill%
\begin{minipage}[t]{0.22680\OfflineResearchTableWidth}\raggedright
Entire-menu erasure and coefficient
alphabet bounds.
\end{minipage}\hfill%
\begin{minipage}[t]{0.23625\OfflineResearchTableWidth}\raggedright
Subdivision columns; constructed or stated random
order; precision below the proved threshold.
\end{minipage}\hfill%
\begin{minipage}[t]{0.30240\OfflineResearchTableWidth}\raggedright
The printed larger
precision at every nonzero probe; individual D children.
\end{minipage}%
\OfflineResearchRow
\noindent%
\begin{minipage}[t]{0.17955\OfflineResearchTableWidth}\raggedright
Individual D source moments
\end{minipage}\hfill%
\begin{minipage}[t]{0.22680\OfflineResearchTableWidth}\raggedright
Explicit polynomial bit acquisition and
sign decoding.
\end{minipage}\hfill%
\begin{minipage}[t]{0.23625\OfflineResearchTableWidth}\raggedright
Subdivision rows; no-matched-pin layout; legal zero
probe.
\end{minipage}\hfill%
\begin{minipage}[t]{0.30240\OfflineResearchTableWidth}\raggedright
General CNFs, prescribed nonzero probes evaluated individually,
the source's fixed tie-break.
\end{minipage}%
\OfflineResearchRow
\noindent%
\begin{minipage}[t]{0.17955\OfflineResearchTableWidth}\raggedright
Completed D probes at an extreme node
\end{minipage}\hfill%
\begin{minipage}[t]{0.22680\OfflineResearchTableWidth}\raggedright
Polynomial queries and receiver
computation recover all signs.
\end{minipage}\hfill%
\begin{minipage}[t]{0.23625\OfflineResearchTableWidth}\raggedright
Same no-matched-pin layout; \(10M\)
bits; any source-sized unit-separated probe set.
\end{minipage}\hfill%
\begin{minipage}[t]{0.30240\OfflineResearchTableWidth}\raggedright
Polynomial cost of
separately materializing every completed probe.
\end{minipage}%
\OfflineResearchRow
\noindent%
\begin{minipage}[t]{0.17955\OfflineResearchTableWidth}\raggedright
Compiled completed-mode coefficients
\end{minipage}\hfill%
\begin{minipage}[t]{0.22680\OfflineResearchTableWidth}\raggedright
Exact polynomial bit evaluator of
the required native linear readouts.
\end{minipage}\hfill%
\begin{minipage}[t]{0.23625\OfflineResearchTableWidth}\raggedright
Same layout and family, with full
input and shared intermediate arithmetic.
\end{minipage}\hfill%
\begin{minipage}[t]{0.30240\OfflineResearchTableWidth}\raggedright
A general compiler for all
source readouts or general SAT instances.
\end{minipage}%
\OfflineResearchRow
\noindent%
\begin{minipage}[t]{0.17955\OfflineResearchTableWidth}\raggedright
Supplied-layout register evaluation
\end{minipage}\hfill%
\begin{minipage}[t]{0.22680\OfflineResearchTableWidth}\raggedright
Uniform counting hardness on
binary tensors and on literal sparse CNFs.
\end{minipage}\hfill%
\begin{minipage}[t]{0.23625\OfflineResearchTableWidth}\raggedright
Register and certified
order are part of the request; stated precision.
\end{minipage}\hfill%
\begin{minipage}[t]{0.30240\OfflineResearchTableWidth}\raggedright
Necessity of that
register for the origin count; hardness under a specified alternative
source order.
\end{minipage}%
\OfflineResearchRow
\noindent%
\begin{minipage}[t]{0.17955\OfflineResearchTableWidth}\raggedright
Raw entries
\end{minipage}\hfill%
\begin{minipage}[t]{0.22680\OfflineResearchTableWidth}\raggedright
Polynomial scans compute the pinned family and its two
factors.
\end{minipage}\hfill%
\begin{minipage}[t]{0.23625\OfflineResearchTableWidth}\raggedright
Original entries are directly accessible.
\end{minipage}\hfill%
\begin{minipage}[t]{0.30240\OfflineResearchTableWidth}\raggedright
The pinned family
cannot witness a superpolynomial lower bound for this interface.
\end{minipage}%
\OfflineResearchRow
\OfflineResearchHeaderRule
\end{OfflineResearchTable}

\subsection{Quantifiers that cannot be
interchanged}\label{v9-quantifiers-that-cannot-be-interchanged}

The guarded theorem proves that for every permitted set of vertex labels
there exists a suitable packet ordering, after which every switch input
obeys the erasure. Its random counterpart fixes a probability law on
orderings and gives a simultaneous high-probability conclusion. The
moment construction chooses another legal labeling and order before any
signs vary. Its upper bound is existential over layouts and uniform over
the entire sign cube. The lower half of the extreme-node precision
theorem quantifies over every certified layout in the stated
orientation.

The single-register hardness theorem quantifies over a uniform
evaluation problem with a supplied order. The order itself may encode
the hard graph. A separate general-tensor construction fixes the
topology and layout as a function of length and varies only entry bits.
Neither assertion proves that a chosen source algorithm generates such a
request. Preprocessing used in a uniform polynomial-time conclusion must
itself take polynomial time; a polynomial number of advice values from
expensive preprocessing changes the model.

The MAIN query bound has an exact input-interface condition. The
decoder's choices and initial state must depend only on the public
layout and earlier replies. A separately read clause entry, an
input-dependent layout, timing, address or cache can convey additional
information. The known promise even has a one-bit raw-entry escape.
These channels cannot be declared free and simultaneously excluded from
the proof.

\subsection{What remains for a complete framework
barrier}\label{v9-what-remains-for-a-complete-framework-barrier}

A full barrier would need an explicitly delimited source-faithful
algorithm class for which every allowed input acquisition, nonlinear
transformation, joint readout and recursive reuse is covered. A hard
family must resist the allowed transformations; this pinned family does
not resist the native moment implementation. The source's actual
deterministic layout choices would need to be checked, or a lower bound
proved uniformly over all legal choices. A quantitative connection
between the resulting information obstruction and real bit cost would
then be necessary. Finally the theorem would need to cover the full
printed precision and every actual interface the program may use.

These are unresolved mathematical obligations, not missing formatting or
additional routine Lean compilation. The present release supplies a
genuine native counterexample, several narrower barriers and precise
value-complexity classifications. It does not assert that the complete
barrier has been established.

\section{V9 archive: Prior work and contribution assessment}\label{v9-section-contribution-audit}

The v9 result most closely tied to the source is the recovery of all
\(N\) pin signs of the four-clone CNF from specified Appendix D
registers. The source itself defines the power sums and the product of
group contractions. The work here constructs a certified layout on which
these objects expose individual signs, proves that their multiplicative
factors can be computed, and accounts for the precision of the completed
registers. The appropriate priority status for these constructions and
the restricted lower bounds is \textbf{independently derived, priority
unestablished}. The original formulas, inherited results, and classical
steps can be attributed more precisely.
\href{https://github.com/openai/math/blob/main/preprints/The-Exponential-Time-Hypothesis-is-False-October-25-2026/build/quantitative-contraction.tex\#L402}{Source
formulas, lines 402--442 and 487--523} ·
\PaperSixEvidence{evidence/research_v7/reviews/hybrid_independent.md}{v7 pinned family}
· \hyperref[v9-section-moment-escape]{J7--J9}

\subsection{Recovering the pins}\label{v9-recovering-the-pins}

\begin{OfflineResearchTable}
\OfflineResearchHeaderRule
\noindent%
\begin{minipage}[t]{0.31500\OfflineResearchTableWidth}\raggedright
Finding
\end{minipage}\hfill%
\begin{minipage}[t]{0.31500\OfflineResearchTableWidth}\raggedright
Prior ingredient
\end{minipage}\hfill%
\begin{minipage}[t]{0.31500\OfflineResearchTableWidth}\raggedright
What the v9 proof establishes
\end{minipage}%
\OfflineResearchRow
\nobreak
\OfflineResearchHeaderRule
\nobreak
\noindent%
\begin{minipage}[t]{0.31500\OfflineResearchTableWidth}\raggedright
J7, individual source children
\end{minipage}\hfill%
\begin{minipage}[t]{0.31500\OfflineResearchTableWidth}\raggedright
With distinct known indices \(j_h\),
the observations \(s_b=\sum_h j_h^b x_h\), \(0\le b<d\), form a square
Vandermonde system. Björck and Pereyra explicitly treat both the matrix
in this orientation and its transpose.
\end{minipage}\hfill%
\begin{minipage}[t]{0.31500\OfflineResearchTableWidth}\raggedright
With subdivision equalities as
rows and old clause and clone vertices as columns, a sign-independent
certified placement puts every pin in a crossing tile and keeps all pins
off the two matched endpoint seats. At \(t=L-1,\zeta=0\), each actual
source child equals a positive graph-known normalizer times the tile's
power sum of local means. Nonpin contributions are computable. The
integer is recovered before any division. At most \(N\) singleton and 64
baseline child values suffice, with polynomial bit time for each
acquisition and for the decoder, when zero is an allowed probe.
\href{https://www.ams.org/mcom/1970-24-112/S0025-5718-1970-0290541-1/S0025-5718-1970-0290541-1.pdf}{Björck--Pereyra,
equations (1)--(3)} ·
\hyperref[v9-j7-polynomial-acquisition-from-individual-source-children]{J7}
\end{minipage}%
\OfflineResearchRow
\noindent%
\begin{minipage}[t]{0.31500\OfflineResearchTableWidth}\raggedright
J8, first coefficient of a completed singleton
\end{minipage}\hfill%
\begin{minipage}[t]{0.31500\OfflineResearchTableWidth}\raggedright
Differentiating a
factored sum can recover marginal information. For example, Darwiche's
Theorem 2 identifies the first partial derivative of a Bayesian network
polynomial with respect to an evidence indicator as a joint probability.
The product rule and interpolation at points with unit differences also
apply in the present setting.
\end{minipage}\hfill%
\begin{minipage}[t]{0.31500\OfflineResearchTableWidth}\raggedright
Here the derivative is with respect to
the \textbf{one native probe coordinate}. At zero the target child
vanishes. Its one-port derivatives have a common sign and at least one
responding row exists. Consequently the completed first coefficient is a
strictly positive, graph-known rational multiple of the source moment.
Odd-denominator clearing, a size bound, and interpolation in the
prescribed extension recover that rational number even if the prescribed
probe set omits zero.
\href{https://people.cs.pitt.edu/~milos/courses/cs3710/readings/differential-darwiche.pdf}{Darwiche,
Theorem 2} ·
\href{https://github.com/openai/math/blob/main/preprints/The-Exponential-Time-Hypothesis-is-False-October-25-2026/build/quantitative-contraction.tex\#L338}{Source
interpolation schedule} ·
\hyperref[v9-j8-polynomially-many-completed-full-input-probes]{J8}
\end{minipage}%
\OfflineResearchRow
\noindent%
\begin{minipage}[t]{0.31500\OfflineResearchTableWidth}\raggedright
J9, single-node precision
\end{minipage}\hfill%
\begin{minipage}[t]{0.31500\OfflineResearchTableWidth}\raggedright
Exact recovery of a bounded integer from a
residue and invertibility of the rational Vandermonde matrix are
elementary. The proof lifts the integers before inverting a matrix whose
determinant may be even.
\end{minipage}\hfill%
\begin{minipage}[t]{0.31500\OfflineResearchTableWidth}\raggedright
A bound on the cleared integer gives the
concrete sufficiency of \(10M\) bits and \(N(10M+1)\) completed
addresses at one fixed node, \(t=L-1\), on the constructed layout. In
the same orientation \textbf{every} certified packing has at least
\(M/6-3M/L\) unused equality columns at that node. Their complete menu
is zero through that precision. This yields the order \(\Theta(M)\) for
full-sign recovery there when the packing can be chosen. The leading
constant and the corresponding threshold with multiple nodes remain
open. \hyperref[v9-j9-the-single-node-precision-order-is-linear]{J9} ·
\hyperref[v9-what-unused-equality-columns-force]{Unused equalities,
P5--P7}
\end{minipage}%
\OfflineResearchRow
\OfflineResearchHeaderRule
\end{OfflineResearchTable}

The two outputs of the v7 one-multiplication identity are the
internal-edge product \(U_M\) and the exterior-edge product \(U_D\).
After J7 or J8 reveals the signs \textbf{entirely through D}, two public
edge scans compute these factors. A route-specific implementation must
additionally obtain \(U_M\) through genuine MAIN queries. J3 gives a
sufficient graph-only rank test and a polynomial evaluator on layouts
passing it. Existence of an asymptotic certified v7 MAIN layout passing
that test remains open. J2 supplies an obstruction on its own diagonal
MAIN layouts. On the completed-D side J8 gives a polynomial decoder from
the supplied slots and a polynomial evaluator of their required
first-coefficient linear combinations from the physical input.
Individual acquisition of every requested nonzero-probe completed slot
on this family remains open.
\PaperSixEvidence{evidence/research_v7/reviews/hybrid_independent.md}{Original factor
ownership} ·
\hyperref[v9-j2-adaptive-query-obstruction-and-a-mixed-interface]{J2 and
J3} · \hyperref[v9-j8-polynomially-many-completed-full-input-probes]{J8
acquisition statement}

\subsection{Contraction, support and
precision}\label{v9-contraction-support-and-precision}

\begin{itemize}
\tightlist
\item
  \textbf{Contraction width.} Multiplying the factors containing a
  variable, summing over it, and bounding the resulting table by the
  induced width is classical bucket elimination. Markov and Shi's
  Proposition 4.2 identifies graph contraction complexity with the
  treewidth of the line graph in their formulation. The supplied order
  in J4 is an order on the \emph{reduced matched-bit primal graph}. J4
  derives those factors from an actual outer zero-probe root. On the J7
  layout each subdivision row leaves at most one unary weight, so the
  remaining column sums separate completely.
  \href{https://ics.uci.edu/~dechter/publications/r48b.pdf}{Dechter,
  \emph{Bucket Elimination}} ·
  \href{https://www.mit.edu/~6.454/www_fall_2002/lizhong/factorgraph.pdf}{Kschischang--Frey--Loeliger,
  factorization and marginalization} ·
  \href{https://arxiv.org/pdf/quant-ph/0511069}{Markov--Shi, Proposition
  4.2} ·
  \hyperref[v9-j4-exact-contraction-of-an-outer-zero-probe-register]{J4}
\item
  \textbf{S1--S3.} Following one distinguished reply until its
  supporting sets cover the witness patterns is a decision-tree
  adversary argument. Fixing random coins and averaging on an explicit
  finite distribution is the elementary direction of the distributional
  method. The v9 work identifies the literal singleton measurement rows,
  certifies diagonal pins in every MAIN group, supplies two differently
  valued legal inputs at each pattern, and combines these with a
  separately certified constant D menu on the same CNF. S2 charges a
  completed jet at one address or one individual MAIN group. The free D
  replies are completed values at the stated precision, with subdivision
  equalities as columns and the guarded order. S3 counts the diagonal
  incidences permitted by the cell-load certificate and bounds that
  particular construction.
  \href{https://doi.org/10.1109/SFCS.1977.24}{Yao, FOCS 1977} ·
  \hyperref[v9-a-nonlinear-adaptive-obstruction-on-a-specified-native-interface]{S1--S3}
\item
  \textbf{S4--S5.} A Boolean monomial on \(h\) distinct free bits has
  sum \(2^{K-h}\) on the \(K\)-cube, and an integral-coefficient
  polynomial respects congruence. The source-specific step in S4 is
  proving that a residual coefficient of degree \(j\) can mention at
  most \(|\beta|+j\) of the unused original equality bits, with
  \emph{all} source moments included. The resulting divisibility holds
  for a completed register and is compatible with J7's informative
  individual child on the same kind of input. P9 translates this into
  all-node erasure for a chosen ordering. The first-new-digit witness
  P10 instead uses the general binary tensor entry point. S5 charges a
  common power of two clearing the coefficients of a rational polynomial
  decoder of \textbf{integral} observations. The charge gives no
  operation-count lower bound.
  \hyperref[v9-erasure-denominators-and-computation]{S4 and S5} ·
  \hyperref[v9-what-unused-equality-columns-force]{P5--P10, layouts and
  input classes}
\end{itemize}

P3--P6 count the shared base coefficients needed to encode all
extension-valued probes of a mode. Counting a new word for each probe
would miss that dependence. For the v8 physical family, P8 gives a
linear erasure range at outer nodes under every certified packing, with
an orientation-dependent node interval. P9 supplies the stronger
\emph{all-node} erasure of \((1-o(1))M/(2L)\) bits on a constructed
subdivision-column packing. Its leading coefficient is optimal for the
proof by counting unused subdivision vertices at \(t=\pm1\). P11
analyzes a specified random distribution of packet orders. P10 exhibits
linearly many recoverable bits in the first additional digit across one
coordinate on a general binary input whose origin count is constant. The
proofs establish distinct facts about the information in the completed
values, the number of their addresses, and which layouts realize them.
\hyperref[v9-section-precision-frontier]{P3--P11}

\subsection{Which counting problem is
hard}\label{v9-which-counting-problem-is-hard}

\begin{OfflineResearchTable}
\OfflineResearchHeaderRule
\noindent%
\begin{minipage}[t]{0.31500\OfflineResearchTableWidth}\raggedright
Primary counting result
\end{minipage}\hfill%
\begin{minipage}[t]{0.31500\OfflineResearchTableWidth}\raggedright
Verified statement
\end{minipage}\hfill%
\begin{minipage}[t]{0.31500\OfflineResearchTableWidth}\raggedright
Transfer to the native value
\end{minipage}%
\OfflineResearchRow
\nobreak
\OfflineResearchHeaderRule
\nobreak
\noindent%
\begin{minipage}[t]{0.31500\OfflineResearchTableWidth}\raggedright
Catherine Greenhill (2000), Theorem 3.1
\end{minipage}\hfill%
\begin{minipage}[t]{0.31500\OfflineResearchTableWidth}\raggedright
Counting all independent sets
of a 3-regular graph is \(\#\mathrm P\)-complete. This is stronger than
bounded maximum degree three. The proof uses counting reductions with
interpolation.
\end{minipage}\hfill%
\begin{minipage}[t]{0.31500\OfflineResearchTableWidth}\raggedright
A1 maps a cubic graph to one supplied certified
star-network layout with binary entries. Its actual completed zero-mode
value at \((-1,0,0)\) equals the independent-set count. This last map
preserves the count with one evaluation when zero is eligible. A2
separately gives general binary constraints, and a family with the
physical topology and packing fixed by input length.
\href{https://web.maths.unsw.edu.au/~csg/papers/complexity.pdf}{Greenhill,
original paper} ·
\hyperref[v9-a1-complexity-of-a-prescribed-root-register]{A1 and A2}
\end{minipage}%
\OfflineResearchRow
\noindent%
\begin{minipage}[t]{0.31500\OfflineResearchTableWidth}\raggedright
Antoine Amarilli, Mikaël Monet and Pierre Senellart (PODS 2017), full
version, Theorem D.1
\end{minipage}\hfill%
\begin{minipage}[t]{0.31500\OfflineResearchTableWidth}\raggedright
Counting \textbf{edge covers} is
\(\#\mathrm P\)-complete on undirected planar bipartite graphs with
every vertex on one side of degree two and every vertex on the other of
degree three. Appendix D explicitly identifies the local signatures
\([0,1,1]\) and \([0,1,1,1]\), and relates their complements by flipping
all edge bits. Its first proof sentence uses ``matching'' in error.
Checking that the selected edges meet every vertex establishes
membership, and the displayed signatures prove hardness for edge covers.
\end{minipage}\hfill%
\begin{minipage}[t]{0.31500\OfflineResearchTableWidth}\raggedright
A4 realizes the complemented local constraints in the \emph{frozen
sparse clause compiler}. The original CNF has disjoint variable sets
across clauses. The supplied certified order carries the desired
incidences. One prescribed zero-mode value equals the number of covers
times the computable positive integer \(\prod_c(1+2^{d_c})\). It fits
the native precision before division, including when factors are even.
\href{https://pierre.senellart.com/publications/amarilli2017conjunctive.pdf\#page=16}{Amarilli--Monet--Senellart,
Appendix D, p.~232} ·
\hyperref[v9-a4-a-single-fixed-mesh-point-on-the-literal-sparse-clause-branch]{A4}
\end{minipage}%
\OfflineResearchRow
\OfflineResearchHeaderRule
\end{OfflineResearchTable}

The additional content of A1, A2 and A4 lies in the native shift, graph
and order certification, binary physical tables and decoding cost. The
graph counting problems and their local Boolean predicates already occur
in the cited literature. The general register inclusion is in
\(\mathrm{FP}^{\#\mathrm P}\) under the stated polynomial-time
original-entry hypothesis. With a probe set omitting zero, the
single-register reductions use polynomially many actual probes for that
mode. Uniform hardness of the sum over \textbf{all} modes on a
source-chosen sparse-clause layout has a separate interpolation
reduction. A lower bound for the origin would also need to show that
every permitted execution must acquire costly information. Such a
running-time bound for an optimized full-precision hybrid remains open.
\hyperref[v9-section-acquisition-cost]{A1--A4} ·
\hyperref[v9-an-exact-value-can-be-hard-without-being-necessary]{Synthesis}

A3 uses the standard count of leaves of an adaptive answer tree. A bound
of at most \(2^{\rho c}\) \emph{new conditional replies} per operation
of cost \(c\) must be justified for the chosen interface. The diverse v8
terminal counts then give an order-\(N\) bit-read lower bound when the
free menu and all other incoming input-dependent state carry \(o(N)\)
bits in total, and each operation supplies a bounded number of new bits.
On the very witness used for those counts, a public sequence of \(O(N)\)
original Boolean reads suffices by a two-integer recurrence. Charging
one operation per high-precision reply requires a separate analysis of
its conditional range.
\hyperref[v9-a3-what-state-counting-prices-in-an-adaptive-execution]{A3 and
its explicit recurrence}

\subsection{Nearest source-specific antecedents and outstanding
provenance}\label{v9-nearest-source-specific-antecedents-and-outstanding-provenance}

\begin{enumerate}
\def\labelenumi{\arabic{enumi}.}
\tightlist
\item
  The frozen Appendix D already gives the Cauchy kernel, the partial
  shift at a mesh point, geometric mode coefficients, source power sums,
  interpolation at distinct binary-coordinate lifts, and the sum of
  products of child contractions. The source also invokes the classical
  conditional-expectation method for its packing. The v9 layout changes
  the distribution and the seats while proving the printed certificate.
  The archived manifest identifies each inspected source file by hash.
  \PaperSixEvidence{evidence/research_work/sources/source_manifest.json}{Frozen
  manifest} ·
  \href{https://github.com/openai/math/blob/main/preprints/The-Exponential-Time-Hypothesis-is-False-October-25-2026/build/quantitative-contraction.tex}{Original
  contraction source} ·
  \href{https://github.com/openai/math/blob/main/preprints/The-Exponential-Time-Hypothesis-is-False-October-25-2026/build/quantitative-reduction.tex}{Original
  sparse compiler}
\item
  Dias and Lima explicitly estimate the cost at Appendix D's prescribed
  precision in Proposition 4.1 and Appendix B.1. Definition 5.1 and
  Theorems 8.1/8.3 concern a rational clause-product reconstruction and
  specified evaluation costs. Their §5.1 explains why tapered two-adic
  Appendix D falls outside that definition. Tulsyan distinguishes native
  engines, including a static full-label MAIN template (Proposition
  5.2), a linear-functional substitution-algebra bound (Theorem 3.7),
  and a residue-query result on a different set family (Theorem 6.4).
  These results provide nearby models for assessing the meaning of a
  lower bound. Their displayed conclusions concern different queries
  from J7--J9's particular singleton children and completed derivatives.
  \href{https://arxiv.org/html/2611.00517v1}{Dias--Lima, §§4--5 and 8,
  Appendix B} · \href{https://hexagonmath.org/pdf/2611.00149v1}{Tulsyan,
  §§3.5, 5 and 6}
\item
  Nazarov's coefficient-problem argument maximizes a convex functional
  over sign choices and compares the maximizer with each one-sign flip.
  This is a precedent for the inherited v7 simultaneous correlation
  argument. The J7--J9 derivations use the v7 pinned graph and edge
  factors, together with the v9 matching construction, bounded local
  sums and native interpolation. The earlier audit also identifies the
  independent-set caterpillar construction underlying the v8 many-valued
  count witness used in A3.
  \href{https://www.mathnet.ru/php/archive.phtml?jrnid=aa&option_lang=eng&paperid=774&wshow=paper}{Nazarov,
  journal record} ·
  \PaperSixEvidence{evidence/research_v7/literature/novelty_matrix.md}{v7 attribution}
  · \PaperSixEvidence{evidence/research_v8/literature/final_targeted_audit.md}{v8
  targeted comparison}
\end{enumerate}

The targeted search checked the source paper and its cited formulas,
work explicitly discussing Appendix D and its precision, and primary
papers on moment inversion, derivatives of factored sums, contraction
width, cubic independent sets and planar biregular edge covers. The
closest identified exact moment formula for the \textbf{modes} is the
source's own \(L_Z(X)\). None of the checked primary results states the
certified full-cube J7/J8 recovery or the two bounds of J9 together.
Independent discovery of the same source-specific constructions
elsewhere and priority relative to unindexed or later work remain
unresolved. The same qualification applies to the source-specific
support construction and hardness embeddings.

\section{V9 archive: Independent review corrections}\label{v9-section-review-resolutions}

All changes in this record affect v9 only. Historical source and proof
files are retained byte-for-byte.

\begin{enumerate}
\def\labelenumi{\arabic{enumi}.}
\tightlist
\item
  \textbf{A2 clause width.} The direct maximum-degree-six simulation now
  starts from \#3SAT. Wider bounded clauses require an additional
  arity-reduction step; it is not silently included.
\item
  \textbf{A2 preprocessing.} The uniform counting-class implication
  explicitly permits polynomial-time graph-only preprocessing. Expensive
  offline preprocessing yielding a polynomial advice string is a
  different model.
\item
  \textbf{A2 inactive roots.} Empty clauses and empty formulas are
  handled directly; tautologies and unused variables are removed before
  computing source size. The known free-variable factor is restored. The
  active-root claim is now stated for the normalized nontrivial
  instance.
\item
  \textbf{A4 layout input.} The literal-CNF reduction explicitly
  supplies a certified layout that can encode the edge-cover incidence
  graph. It is not claimed to be a prescribed function of a canonically
  indexed read-once CNF topology alone.
\item
  \textbf{P8 padding.} The all-layout padding estimate explicitly uses
  the printed maximum-real-load, rounded common cell size. The abstract
  three packing guarantees alone do not imply negligible padding.
\item
  \textbf{J7/J8 biased labels.} The proof centers every wire at its own
  cell-indicator mean. Independent biased wires have zero covariance
  after this centering, but their agreement probability need not be the
  uniform-cell probability. The paired mean identity and bounded
  dependency graph are both used.
\item
  \textbf{Operand provenance.} The native moment decoder obtains every
  pin from D and then computes both older mathematical factors. It does
  not identify the MAIN factor with one actual MAIN register. J3
  provides a conditional MAIN reconstruction test, not an asymptotic
  existence result.
\item
  \textbf{Execution accounting.} Polynomially many completed probes and
  polynomial receiver work are separated from the cost of acquiring each
  completed value. The polynomial compiled first-coefficient evaluator
  is proved as an exact alternative execution of the required aggregate,
  not counted as having produced absent sample values.
\item
  \textbf{Sharpness.} The diagonal count scales are optimal for that
  mechanism; the unsupported-equality guard constant is optimal for that
  certificate; the one-node precision order is linear. None of these
  statements gives an optimal full-algorithm running time.
\end{enumerate}

The fresh mathematical reviewer independently checked a once-subdivided
K4 edge-cover example with 621 covers. The coverage reviewer verified
the original source semantics of J7--J9, including the exact integer
bounds and both row orientations. They also identified a stronger
full-precision zero-probe erasure on the guarded orientation, recorded
in their review. The public evidence catalogue identifies available
records of those reviews. The stated limitations remain in force.

\section{V9 archive: Formal verification and reproduction}\label{v9-section-formal-status}

The historical v9 build compiled the complete retained v6/v7/v8/v9
project import graph using Lean 4.24.0 and Mathlib commit
\nolinkurl{f897ebcf72cd16f89ab4577d0c826cd14afaafc7}. All \textbf{58
modules} compiled with warnings treated as errors. Their \textbf{450
requested axiom-dependency reports} were present and contained only the
ordinary foundations \texttt{propext}, \texttt{Classical.choice}, and
\texttt{Quot.sound}, when used. The five v9 modules contribute 46 of
those reports. No preexisting project object file was reused: local
modules were compiled in dependency order into a fresh temporary
directory. The Mathlib checkout and direct imported source provenance
were checked against the pinned commit.

\begin{OfflineResearchTable}
\OfflineResearchHeaderRule
\noindent%
\begin{minipage}[t]{0.23625\OfflineResearchTableWidth}\raggedright
Module
\end{minipage}\hfill%
\begin{minipage}[t]{0.08505\OfflineResearchTableWidth}\raggedright
Reports
\end{minipage}\hfill%
\begin{minipage}[t]{0.31185\OfflineResearchTableWidth}\raggedright
What the kernel verifies
\end{minipage}\hfill%
\begin{minipage}[t]{0.31185\OfflineResearchTableWidth}\raggedright
Written source correspondence
\end{minipage}%
\OfflineResearchRow
\nobreak
\OfflineResearchHeaderRule
\nobreak
\noindent%
\begin{minipage}[t]{0.23625\OfflineResearchTableWidth}\raggedright
BooleanSummation
\end{minipage}\hfill%
\begin{minipage}[t]{0.08505\OfflineResearchTableWidth}\raggedright
9
\end{minipage}\hfill%
\begin{minipage}[t]{0.31185\OfflineResearchTableWidth}\raggedright
Actual Boolean cube sums, ordinary monomial
powers, support/degree divisibility, graded coefficient sums and
arbitrary-probe zero in finite precision.
\end{minipage}\hfill%
\begin{minipage}[t]{0.31185\OfflineResearchTableWidth}\raggedright
The native expansion and
unused equality bits yield the required support bound; layouts and
asymptotics remain written.
\end{minipage}%
\OfflineResearchRow
\noindent%
\begin{minipage}[t]{0.23625\OfflineResearchTableWidth}\raggedright
DiagonalQueryTree
\end{minipage}\hfill%
\begin{minipage}[t]{0.08505\OfflineResearchTableWidth}\raggedright
7
\end{minipage}\hfill%
\begin{minipage}[t]{0.31185\OfflineResearchTableWidth}\raggedright
An adaptive tree with arbitrary replies, a
capacity lower bound on a real paid execution, exact cylinder sizes,
truncation and finite-seed error counting.
\end{minipage}\hfill%
\begin{minipage}[t]{0.31185\OfflineResearchTableWidth}\raggedright
Actual MAIN singleton rows,
certified pin counts, two CNF witnesses and erased D replies are
established separately.
\end{minipage}%
\OfflineResearchRow
\noindent%
\begin{minipage}[t]{0.23625\OfflineResearchTableWidth}\raggedright
ShiftedAcquisition
\end{minipage}\hfill%
\begin{minipage}[t]{0.08505\OfflineResearchTableWidth}\raggedright
7
\end{minipage}\hfill%
\begin{minipage}[t]{0.31185\OfflineResearchTableWidth}\raggedright
Signed one-port transfer, paired Boolean
identity, arbitrary row functions and full finite network partition
equality.
\end{minipage}\hfill%
\begin{minipage}[t]{0.31185\OfflineResearchTableWidth}\raggedright
Mesh specialization and certified physical placement remain
written.
\end{minipage}%
\OfflineResearchRow
\noindent%
\begin{minipage}[t]{0.23625\OfflineResearchTableWidth}\raggedright
PhysicalAcquisition
\end{minipage}\hfill%
\begin{minipage}[t]{0.08505\OfflineResearchTableWidth}\raggedright
10
\end{minipage}\hfill%
\begin{minipage}[t]{0.31185\OfflineResearchTableWidth}\raggedright
Even negative clauses, controls, residual
equality sums and the joint finite edge-cover cardinality identity over
a commutative ring.
\end{minipage}\hfill%
\begin{minipage}[t]{0.31185\OfflineResearchTableWidth}\raggedright
Real clause compilation, padding, parity of
odd-degree vertices, precision extraction and counting reductions remain
written.
\end{minipage}%
\OfflineResearchRow
\noindent%
\begin{minipage}[t]{0.23625\OfflineResearchTableWidth}\raggedright
\texttt{NativeMoment\allowbreak{}Reconstruction}
\end{minipage}\hfill%
\begin{minipage}[t]{0.08505\OfflineResearchTableWidth}\raggedright
13
\end{minipage}\hfill%
\begin{minipage}[t]{0.31185\OfflineResearchTableWidth}\raggedright
An explicit
inverse Vandermonde decoder from distinct coordinates; actual weighted
assignment moments, known \nolinkurl{backgrounds/scales,} pinned
marginals and heterogeneous product-table factorization.
\end{minipage}\hfill%
\begin{minipage}[t]{0.31185\OfflineResearchTableWidth}\raggedright
J7/J8 show
that genuine native values realize these expressions and prove
acquisition/bit bounds.
\end{minipage}%
\OfflineResearchRow
\OfflineResearchHeaderRule
\end{OfflineResearchTable}

The historical v9 report records partial formal coverage of the
source-specific CNF compiler, packet construction, covariance estimates,
native reconstruction program, asymptotic thresholds and Turing-machine
cost model. The finite
cores do not assume the lower bound or decoder correctness being
established; they do assume the structural hypotheses documented in
their statements and scope files. In particular, the actual tree proof
receives the off-support response law, and the native moment proof
receives distinct coordinates, nonzero scales and the stated
weighted-sum realization. Their applicability is a separate mathematical
proof.

\PaperSixEvidence{evidence/research_v9/formal/replay_results_v9.json}{Full clean replay report} ·
\PaperSixEvidence{evidence/research_v9/formal/replay_all.py}{Replay program} ·
\PaperSixEvidence{evidence/research_v9/formal/DECLARATION_INDEX.json}{Declaration index} ·
\PaperSixEvidence{evidence/research_v9/formal/BooleanSummation_scope.md}{Boolean scope} ·
\PaperSixEvidence{evidence/research_v9/formal/DiagonalQueryTree_scope.md}{Decision-tree scope} ·
\PaperSixEvidence{evidence/research_v9/formal/ShiftedAcquisition_scope.md}{Acquisition scope} ·
\PaperSixEvidence{evidence/research_v9/formal/NativeMomentReconstruction_scope.md}{Moment scope}

\subsection{Exact finite checks}\label{v9-exact-finite-checks}

Five independent executable check programs were rerun from a fresh
source-only directory. Each exited successfully and reproduced its saved
JSON exactly as data. The suite covers rational Cauchy coefficients,
extension probes, original Boolean assignment sums, guarded and matched
physical packings, source moment vectors, signed sparse-CNF reductions
and the independent edge-cover example. Its reference outputs include
recovery of all 4,096 sign words on a 288-wire four-clone instance and
both 41-point extension probe sets, one omitting zero, on a separate
eight-wire literal network. The independent once-subdivided K4 check
recovers 621 covers from a 1,101-bit actual shifted value.

The toy arrays have intentionally small parameters. The source's fixed
parameter schedule has very large asymptotic thresholds, so the finite
tests check identities and semantics rather than the entire asymptotic
source regime. Neither finite enumeration nor a successful Lean core
replay substitutes for the separately audited layout and cost proofs.

\PaperSixEvidence{evidence/research_v9/checks/replay_summary.json}{Finite replay report} ·
\PaperSixEvidence{evidence/research_v9/checks/replay_checks.py}{Fresh replay driver} ·
\PaperSixEvidence{evidence/research_v9/checks/joint_algorithms_moment_escape.json}{Moment/probe
reference} · \PaperSixEvidence{evidence/research_v9/reviews/independent_math_review.json}{Independent
physical reference}

\subsection{Preservation and replay}\label{v9-preservation-and-replay}

The parent v8 full archive was restored from its retained copy and
verified before research resumed. The v9 replay independently checked
every retained active v6, v7 and v8 Lean source hash against its
historical replay record. That distribution preserved those sources,
logs, new proof files and v9 reports. The
\PaperSixCatalogue{public evidence catalogue} records the original
paths and hashes and identifies the files available with this edition.
The historical reproduction guide describes the original build and replay.

\PaperSixEvidence{evidence/research_v9/REPRODUCE.md}{Reproduction guide} ·
\PaperSixEvidence{evidence/research_v9/SOURCE_MANIFEST.json}{Frozen source hashes} ·
\PaperSixEvidence{evidence/research_v9/BASELINE.json}{Parent provenance}

\section{V9 archive: Theorem index}\label{v9-section-theorem-index}

The new main positive result is J7--J9. Classical finite algebra and the
source formulas are attributed in the literature chapter; priority of
the source-specific constructions is not established.

\begin{OfflineResearchTable}
\OfflineResearchHeaderRule
\noindent%
\begin{minipage}[t]{0.06615\OfflineResearchTableWidth}\raggedright
ID
\end{minipage}\hfill%
\begin{minipage}[t]{0.27405\OfflineResearchTableWidth}\raggedright
Result
\end{minipage}\hfill%
\begin{minipage}[t]{0.29295\OfflineResearchTableWidth}\raggedright
Proof location
\end{minipage}\hfill%
\begin{minipage}[t]{0.31185\OfflineResearchTableWidth}\raggedright
Formal status
\end{minipage}%
\OfflineResearchRow
\nobreak
\OfflineResearchHeaderRule
\nobreak
\noindent%
\begin{minipage}[t]{0.06615\OfflineResearchTableWidth}\raggedright
J7
\end{minipage}\hfill%
\begin{minipage}[t]{0.27405\OfflineResearchTableWidth}\raggedright
Actual source-child moment recovery of every pin, with polynomial
bit acquisition and decoding.
\end{minipage}\hfill%
\begin{minipage}[t]{0.29295\OfflineResearchTableWidth}\raggedright
Native moment chapter, label law,
matching placement and rational moment systems.
\end{minipage}\hfill%
\begin{minipage}[t]{0.31185\OfflineResearchTableWidth}\raggedright
Weighted finite
assignment, distinct-coordinate decoder and product marginal core
verified.
\end{minipage}%
\OfflineResearchRow
\noindent%
\begin{minipage}[t]{0.06615\OfflineResearchTableWidth}\raggedright
J8
\end{minipage}\hfill%
\begin{minipage}[t]{0.27405\OfflineResearchTableWidth}\raggedright
All pins recovered from actual completed probes, including a
prescribed probe set without zero; polynomial compiled first-coefficient
evaluator.
\end{minipage}\hfill%
\begin{minipage}[t]{0.29295\OfflineResearchTableWidth}\raggedright
Native moment chapter, strict target sign, finite-precision
interpolation and exact rational extraction.
\end{minipage}\hfill%
\begin{minipage}[t]{0.31185\OfflineResearchTableWidth}\raggedright
Moment decoder core
verified; native derivative realization and source acquisition are
written.
\end{minipage}%
\OfflineResearchRow
\noindent%
\begin{minipage}[t]{0.06615\OfflineResearchTableWidth}\raggedright
J9
\end{minipage}\hfill%
\begin{minipage}[t]{0.27405\OfflineResearchTableWidth}\raggedright
Linear precision order at one extreme node: uniform erasure below
\(M/6-3M/L\), constructive recovery at \(10M\).
\end{minipage}\hfill%
\begin{minipage}[t]{0.29295\OfflineResearchTableWidth}\raggedright
Native moment chapter
and P5--P7.
\end{minipage}\hfill%
\begin{minipage}[t]{0.31185\OfflineResearchTableWidth}\raggedright
Graded Boolean sum and decoder cores verified; composition
and bit bounds written.
\end{minipage}%
\OfflineResearchRow
\noindent%
\begin{minipage}[t]{0.06615\OfflineResearchTableWidth}\raggedright
J1--J2
\end{minipage}\hfill%
\begin{minipage}[t]{0.27405\OfflineResearchTableWidth}\raggedright
Certified diagonal MAIN pins, actual-path adaptive query bound
and mixed erased D access.
\end{minipage}\hfill%
\begin{minipage}[t]{0.29295\OfflineResearchTableWidth}\raggedright
MAIN/joint chapter; synthesis S1--S3.
\end{minipage}\hfill%
\begin{minipage}[t]{0.31185\OfflineResearchTableWidth}\raggedright
Adaptive finite-tree capacity and bounded-error counting verified.
\end{minipage}%
\OfflineResearchRow
\noindent%
\begin{minipage}[t]{0.06615\OfflineResearchTableWidth}\raggedright
J3
\end{minipage}\hfill%
\begin{minipage}[t]{0.27405\OfflineResearchTableWidth}\raggedright
A graph-testable first-jet MAIN decoder condition.
\end{minipage}\hfill%
\begin{minipage}[t]{0.29295\OfflineResearchTableWidth}\raggedright
MAIN/joint
chapter.
\end{minipage}\hfill%
\begin{minipage}[t]{0.31185\OfflineResearchTableWidth}\raggedright
Written sufficient condition; no asymptotic certified example
claimed.
\end{minipage}%
\OfflineResearchRow
\noindent%
\begin{minipage}[t]{0.06615\OfflineResearchTableWidth}\raggedright
J4--J5
\end{minipage}\hfill%
\begin{minipage}[t]{0.27405\OfflineResearchTableWidth}\raggedright
Exact extreme-node contraction, matching-boundary polynomial
register, and joint source-moment dynamic programs.
\end{minipage}\hfill%
\begin{minipage}[t]{0.29295\OfflineResearchTableWidth}\raggedright
MAIN/joint
chapter.
\end{minipage}\hfill%
\begin{minipage}[t]{0.31185\OfflineResearchTableWidth}\raggedright
Written identities and bit algorithms, with direct finite
comparisons.
\end{minipage}%
\OfflineResearchRow
\noindent%
\begin{minipage}[t]{0.06615\OfflineResearchTableWidth}\raggedright
J6
\end{minipage}\hfill%
\begin{minipage}[t]{0.27405\OfflineResearchTableWidth}\raggedright
Multi-node counting reduction through one physical CNF and fixed
supplied layout.
\end{minipage}\hfill%
\begin{minipage}[t]{0.29295\OfflineResearchTableWidth}\raggedright
MAIN/joint chapter.
\end{minipage}\hfill%
\begin{minipage}[t]{0.31185\OfflineResearchTableWidth}\raggedright
Written construction and
arithmetic; supplemented by the separate single-mesh edge-cover reduction A4.
\end{minipage}%
\OfflineResearchRow
\noindent%
\begin{minipage}[t]{0.06615\OfflineResearchTableWidth}\raggedright
P1--P4
\end{minipage}\hfill%
\begin{minipage}[t]{0.27405\OfflineResearchTableWidth}\raggedright
Exact mode valuations, full-slot series and pooled
coefficient-alphabet budget.
\end{minipage}\hfill%
\begin{minipage}[t]{0.29295\OfflineResearchTableWidth}\raggedright
Precision chapter.
\end{minipage}\hfill%
\begin{minipage}[t]{0.31185\OfflineResearchTableWidth}\raggedright
Native series proof
written; finite calculations and inherited capacity foundations.
\end{minipage}%
\OfflineResearchRow
\noindent%
\begin{minipage}[t]{0.06615\OfflineResearchTableWidth}\raggedright
P5--P7
\end{minipage}\hfill%
\begin{minipage}[t]{0.27405\OfflineResearchTableWidth}\raggedright
Unsupported equality columns imply coefficient and all-probe
divisibility.
\end{minipage}\hfill%
\begin{minipage}[t]{0.29295\OfflineResearchTableWidth}\raggedright
Precision chapter.
\end{minipage}\hfill%
\begin{minipage}[t]{0.31185\OfflineResearchTableWidth}\raggedright
Actual monomial/cube sum and
receiving-semiring core verified.
\end{minipage}%
\OfflineResearchRow
\noindent%
\begin{minipage}[t]{0.06615\OfflineResearchTableWidth}\raggedright
P8--P9
\end{minipage}\hfill%
\begin{minipage}[t]{0.27405\OfflineResearchTableWidth}\raggedright
Every printed outer padding gives an erasure range; a
constructive all-node guard has sharp unsupported-equality constant.
\end{minipage}\hfill%
\begin{minipage}[t]{0.29295\OfflineResearchTableWidth}\raggedright
Precision chapter.
\end{minipage}\hfill%
\begin{minipage}[t]{0.31185\OfflineResearchTableWidth}\raggedright
Written covariance, conditional expectation,
packing and asymptotics.
\end{minipage}%
\OfflineResearchRow
\noindent%
\begin{minipage}[t]{0.06615\OfflineResearchTableWidth}\raggedright
P10--P11
\end{minipage}\hfill%
\begin{minipage}[t]{0.27405\OfflineResearchTableWidth}\raggedright
First-extra-digit linear information witness and
high-probability random-order physical erasure.
\end{minipage}\hfill%
\begin{minipage}[t]{0.29295\OfflineResearchTableWidth}\raggedright
Precision chapter.
\end{minipage}\hfill%
\begin{minipage}[t]{0.31185\OfflineResearchTableWidth}\raggedright
Written constructions; exact kernel and physical packing checks.
\end{minipage}%
\OfflineResearchRow
\noindent%
\begin{minipage}[t]{0.06615\OfflineResearchTableWidth}\raggedright
A1--A2
\end{minipage}\hfill%
\begin{minipage}[t]{0.27405\OfflineResearchTableWidth}\raggedright
General native register evaluation is in FP with a \#P oracle
under the entry-cost hypothesis; counting-hard specific and fixed-host
examples.
\end{minipage}\hfill%
\begin{minipage}[t]{0.29295\OfflineResearchTableWidth}\raggedright
Acquisition chapter.
\end{minipage}\hfill%
\begin{minipage}[t]{0.31185\OfflineResearchTableWidth}\raggedright
Signed finite row/network identity
verified; complexity reductions written.
\end{minipage}%
\OfflineResearchRow
\noindent%
\begin{minipage}[t]{0.06615\OfflineResearchTableWidth}\raggedright
A3
\end{minipage}\hfill%
\begin{minipage}[t]{0.27405\OfflineResearchTableWidth}\raggedright
Conditional information per charged event bounds exact adaptive
execution; raw-entry escape attains linear read cost on the stated
witness.
\end{minipage}\hfill%
\begin{minipage}[t]{0.29295\OfflineResearchTableWidth}\raggedright
Acquisition chapter.
\end{minipage}\hfill%
\begin{minipage}[t]{0.31185\OfflineResearchTableWidth}\raggedright
Written finite-tree entropy/count
argument and explicit recurrence.
\end{minipage}%
\OfflineResearchRow
\noindent%
\begin{minipage}[t]{0.06615\OfflineResearchTableWidth}\raggedright
A4
\end{minipage}\hfill%
\begin{minipage}[t]{0.27405\OfflineResearchTableWidth}\raggedright
A specified literal sparse-CNF native value encodes edge covers
with a supplied certified layout.
\end{minipage}\hfill%
\begin{minipage}[t]{0.29295\OfflineResearchTableWidth}\raggedright
Acquisition chapter.
\end{minipage}\hfill%
\begin{minipage}[t]{0.31185\OfflineResearchTableWidth}\raggedright
Full finite
signed partition/count identity verified; physical placement and bit
reduction written.
\end{minipage}%
\OfflineResearchRow
\noindent%
\begin{minipage}[t]{0.06615\OfflineResearchTableWidth}\raggedright
S1--S5
\end{minipage}\hfill%
\begin{minipage}[t]{0.27405\OfflineResearchTableWidth}\raggedright
Support cover, optimal diagonal count scale and
congruence-to-denominator budget.
\end{minipage}\hfill%
\begin{minipage}[t]{0.29295\OfflineResearchTableWidth}\raggedright
Synthesis chapter.
\end{minipage}\hfill%
\begin{minipage}[t]{0.31185\OfflineResearchTableWidth}\raggedright
S1's finite
execution core and S4's assignment sums verified; S3 and S5 written.
\end{minipage}%
\OfflineResearchRow
\OfflineResearchHeaderRule
\end{OfflineResearchTable}

\PaperSixEvidence{evidence/research_v9/formal/DECLARATION_INDEX.json}{Detailed finite declaration
index} · \hyperref[v9-section-formal-status]{Formal coverage} ·
\hyperref[v9-section-coverage]{Source-model coverage}

\section{V9 archive: Consequences for the next research step}\label{v9-consequences-for-the-next-research-step}

The native reconstruction changes the next useful question. The old
pinned family already has a polynomial execution inside a sufficiently
rich, source-derived interface. Proving a larger rank for it cannot
restore a complete nonlinear lower bound. A new difficult family would
need varying clause structure or unpinned logical freedom that survives
the moment readout and resists local factorization. Its source layout,
available readouts, preprocessing and total bit cost must be part of the
theorem from the outset.

Several narrower questions remain concrete. Can the matching-structured
first-coefficient evaluator be extended to bounded-width shifted factor
graphs at nonzero probes? Which higher coefficients admit a similarly
compact joint state? Does a specified deterministic source layout permit
efficient native pin recovery without changing its orientation or order?
Can an arbitrary prescribed family of completed values be replaced by a
low-cost compiled readout while still performing the required Hermite
reconstruction? These are algorithm questions, not consequences of the
present hardness reductions.

On the lower-bound side, the single-node completed interface has a
proved linear precision scale on the common family. Determining its
optimal constant, or the effect of pooling nonextreme nodes, would
sharpen a genuine information statement. The guarded and random-order
estimates can be studied for larger graph classes and under fixed
packing rules. MAIN's diagonal mechanism is optimal to constant factors
in the number of diagonal pins allowed by its cell bound, so a matching
time exponent would require additional structure.

The separate original Boolean-cube extremal problem also remains
separate: these results do not determine the fully nonsymmetric
second-order degree constant. No statement about that constant is
inferred from the native moment construction.

The strongest immediate research discipline is to test the next proposed
hard family against its actual native first moments and low-order joint
derivatives before investing in further representation lower bounds.
Both a successful fast readout and a proved failure with an explicit
acquisition model would be substantive outcomes.



\clearpage
\section{Two consequences of the v9 source review}\label{r47-v9-review-consequences}
This supplement records two consequences contained in the original
\PaperSixEvidence{evidence/closeout/v9/independent_coverage_review.md}{independent
v9 source review}. R47 covers finite uniform seed sets. The R119
probability endpoints verify the argument below for arbitrary seed
probability laws, including the grouped and complete-only native
oracles on one certified layout. The zero-probe result has
separate statements for the cutoff at a given layout and for the
existence of one layout covering all active centers.

\subsection{Worst-case expected queries in J2}
Suppose the success probability is at least \(2/3\) on every promised
input, and the supremum, over real inputs, of the expected number of
paid queries is \(Q<\infty\). Count the queries charged in J2, with the
same public inputs and delivered observations. In the grouped case let
\(d=k_0\), and in the complete-only case let \(d=k\).
Each seed fixes a deterministic decision tree on this finite paired
promise. Draw the seed from any probability space on which the paid query
count and the success event on each input are measurable up to null
sets. Put \(B=\lfloor12Q\rfloor\). Truncate immediately before the
\((B+1)\)-st paid query. On each input, if \(C\) is its original
integer query count, then \(\Pr(C>B)\le Q/(B+1)\le 1/12\).
The original procedure is therefore both correct and within budget
with probability at least \(7/12\). For each fixed seed these successes
also occur in the truncated tree. On the uniform paired promise the
fixed-seed counting bound is \(1/2+B/2^{d+1}\). Integrate the number
of correct within-budget inputs over the seed law and exchange the
finite input sum with this integral. This gives
\(7/12\le 1/2+B/2^{d+1}\), including when \(Q=0\).
Consequently \(B\ge 2^d/6\) and \(Q\ge 2^d/72\). This is an expected
\emph{query} lower bound; it is not a bound on all arithmetic or on
the cost of evaluating one native reply.

\subsection{Full-precision erasure at the zero probe}
Use the J2 orientation with subdivisions on the column side.
At \(\zeta=0\) every nonzero mode has zero \emph{completed product}.
For the zero mode, a positive binary OR row vanishes identically
if neither of its two actual ports has a surviving shifted argument:
its extension then evaluates at \((0,0)\). In the original hard
network there are \(3N/2=M/16\) such rows. With \(P\) packets of length
\(L\), at most \(P(L-|t|)\) row ports can have surviving predecessors.
Thus for any certified order with \(P(L-|t|)<M/16\), its completed
zero probe at \(t\) is exactly zero, independent of precision and
private signs. The bound \(P=(1+o(1))M/L\) in the source schedule
implies this condition whenever \(|t|\ge(15/16+\delta)L\) for fixed
\(\delta>0\) and all sufficiently large admissible inputs.

There is also one certified order realizing this erasure at
\emph{every} active signed center. Choose two distinct positive OR
rows. Reserve the last position of their respective packets for
both actual ports of the first row, and the first position for both
ports of the second. The four adjacent subdivision columns are
distinct. Excluding them from the P9 candidate pool removes only
\(O(1)\) marked columns. Mask balancing gives \(H_{\rm pad}/L\to\infty\)
endpoint seats in each occupied cell. Subtracting the four reservations
from the associated capacities before deleting overfull marked
columns increases the total overflow by at most four. The same P9
estimate therefore still leaves \((1-o(1))M/(2L)\) fully guarded
columns on each side; complete the reservations, the guarded columns
and then the remaining seats in their original cells.

At positive \(t\), the two last-position row ports cannot see a
shifted column. At negative \(t\), the same is true of the two
first-position row ports. This certifies exact zero-probe erasure
on this one selected layout. J2 on that layout may give the whole
completed \(\zeta=0\) menu at its native precision for free, while
keeping the original receiver cap at nonzero probes. An installed
list that omits zero delivers no such probe, and this observation
does not erase nonzero-probe information at native precision.


\clearpage
\section{Source correspondence proved in the later Lean checks}\label{r64-source-correspondence}
This supplement updates three specific items from the earlier source audit.
It leaves the hypotheses and statements in the preceding manuscript unchanged.
The earlier coverage tables give the status at their own evidence cutoff.
A later source check does not by itself establish a new global resource bound.

\subsection{J3: a single literal first-jet query}
For each admissible graph-fixed menu in J3, the source evaluation of the
all-one first jet, normalization, and fraction-free recovery are checked
on the same signed system. The checks bound the actual arithmetic steps
and intermediate integer bit heights. Under standard integer arithmetic
they certify the polynomial bit work asserted for this computation,
measured in the printed parameters. They do not find an admissible
layout or furnish a bit implementation of the complete MAIN pipeline.

\subsection{J8: the actual full-precision roster}
For the original four-clone parameters $M=24N$ and $P_0=4N+1\le M$,
the printed root regime gives $p_t+1\le M(2L-1)+1$ at the odd extreme centre.
For any legal ordered binary-lift roster and any $p_t+1$ distinct
selected native addresses, interpolation on the actual completed slots
at precision $p_t$ gives the exact recovered first coefficients and all
pin signs. One list of $N(p_t+1)$ requests suffices for all sign words;
the cost of producing the individual replies must be accounted for separately. The separate
$10M$-bit receiver proof uses only $10M+1$ selected addresses. For the
full-precision argument the inlet precision restriction matters: Ready
and Stable alone permit raising the working precision beyond the
unchanged native menu size.

\subsection{F5: the emitted instance and the printed source entrances}
Given a suitable backbone as in F5, the checked ordered construction
embeds the input SAT decision into the emitted F1 CNF. The same signed
clause list is unchanged by the MAIN deletion rules. Each private
physical variable occurs exactly three times, and each detached variable
occurs once. In the admitted MAIN regime the printed search selects
a sparsification threshold of at least $16$; no first branching
condition holds, and the unique leaf is the original clause list.
The separate Appendix D clause-density entrance also holds on the
same list. Neither statement charges the later computation.
External graph-existence and hardness results remain referenced in
their original role; no new formalization of those external theorems is
asserted.

\paragraph{Evidence.}
See the \PaperSixEvidence{evidence/later/P6_R63_ORIGINAL_CLOSURE.md}{independent
source crosswalk}, the \PaperSixEvidence{evidence/later/J8R64_SCOPE.md}{J8
full-precision update}, and the \PaperSixEvidence{evidence/later/F5R63Sparsify_SCOPE.md}{F5
source-entrance certificate}. The crosswalk predates the full-precision
update; its J8 limitation is superseded exactly in the scope stated above.

\clearpage
\section{Checked attachment of the retained physical execution}
\label{sec:retained-physical-execution}

The completed-slot oracle in J6 is part of the statement of the query
reduction. The checks below concern a separate implementation of that
interface. They do not turn a bound on the number of calls into a
bound on the time of an arbitrary oracle.

\paragraph*{Exact count on the same execution.}
For each original address \((i,j,g)\), the executable attachment keeps the
site descriptor and the bank returned by the actual successful admission.
The continuation operates on that site and bank. Its first physical
child receives the installed resident data and the already emitted owner
and label words. The positive stopping rule then reads the retained
data on its own live bank. Its decoded result is the contraction of
the original group. Consequently the prescribed interpolation and
readout return the exact model count of the submitted CNF. This
attachment applies under the typed guards, ready-source operands, and
ordered-byte certificate of its Lean statement. These are not additional
axioms.

\paragraph*{Actual fees and named live data.}
The continuation charges the same child executions used for its value.
The acquired site is kept instead of being constructed by the
continuation a second time. Probe-chart construction and mounting are
charged on their own meter; the chart returned by that mount is the
one used by the retained driver. The shared chart is counted once along
one group execution. This does not assert one mount for all groups.
The source-indexed peak quantity includes the named retained owner
and label tapes, the installed records, and the observed suspended
buffers. At one posted group, the workspace theorem takes the maximum
of the chart-build, pre-child and reached-path entry peaks. It uses a
certificate for the chart-build region and the stated entry and
suspended-frame bounds along the active path, with the named ambient
data included. Arithmetic temporaries, the outer allocator and the
lifetime of an earlier preparer's old chart lie outside this bound.


\paragraph*{The assembled clock.}
At a nonterminal exterior, write \(A\) for the number of actual first
addresses and \(d\) for the exterior's absolute depth. The retained
contraction has the bound
\[
 F_{\mathrm{ext}}+A U_{\mathrm{ext}}
 +A\left(S_{\mathrm{first}}+
     (\operatorname{pathLimit}-d)(F+U)\,2^{\kappa H}\right).
\]
Here \(F_{\mathrm{ext}}\) and \(U_{\mathrm{ext}}\) bound the actual
fixed exterior preparation and arithmetic per address;
\(S_{\mathrm{first}}\) is the bound on the actual first-slot setup, installation,
publication and restoration; \(F+U\) is the sum of the fixed and per-address
ceilings for the charged local operations on the visited suffix. The parameters \(\kappa,H\) are those of the printed
absolute-depth budget. At a stopped exterior the executable uses its
original terminal meter instead. The bound at the last prescribed
point uses the absolute-depth certificate directly: that point is not
assumed to be a native-prefix occurrence. The outer theorem includes
compiled admission, the measured source preparer and chart mount,
the paid interpolation and the startup and readout. Its hypotheses
are the explicit measured preparation bound and the local
source-width, layout and occurrence certificates; there is no
hypothesis replacing the measured whole run by a supplied cost.


\paragraph*{Which size governs the retained execution.}
Let \(G\) denote the number of negative-unary padding gadgets.
Each such gadget uses one fresh variable and one clause. The \(L\)
copies of the guarded formula \(H\) are disjoint. Counting variables
and clauses separately therefore gives
\[
 V_{\rm pad}=L\,|\mathrm{Var}(H)|+G,\qquad
 C_{\rm pad}=L\,|\mathrm{Clauses}(H)|+G,\qquad 2G=M-Lm .
\]
Thus the source size at startup is \(s=V_{\rm pad}+C_{\rm pad}
=M-Lm+L s_H\). The sparse-CNF instance takes initial precision
\(P_0=V_{\rm pad}+1\) and truncation order \(\ell=1\).
The stopping test at an exterior or an interior call with \(e\) edges is
\(e\le\lceil s/\lfloor\sqrt{\lambda}\rfloor\rceil\), where
\(\lambda=\lceil\log_2 s\rceil\) in the admitted large-input regime.
A field-degree equality that determines \(\lambda\) alone does not
determine this threshold. In particular, an execution correspondence
must identify the size returned by the preparation of each visited group.

The subdivision count \(q_H=m/2\) has its own role. It supplies the number
of free columns in the interpolation factor \(2^{tq_H}\) and the
coefficient of \(2m+1\) in the factor-width bound. For example, guarding
the two positive unit clauses on two declared variables gives three
variables, four clauses and eight incidences. The guard has four
occurrences, hence adds one internal equality-tree edge. The original
variables each occur twice. Consequently \(s_H=7\), \(q_H=9\), and
\(m=18\). Using \(q_H\) in the source size would add \(2L\) to \(s\);
using \(s_H\) in the factor width would subtract \(74\) bits from its
stated bound. The example distinguishes the parameter formulas without
asserting that either threshold is attained by a visited call in an
admitted run.

For a group reached at its prescribed address, let \(b\) be the charge
returned by its original preparer and \(u\) the actual charge of
generating and mounting the chart. The preparation bound in the
assembled clock applies to \(b+u+4\), with four steps charged by the
wrapper. The chart bound \(u\le 3660s^5\) uses the size of the returned
record. For the concrete preparer the posted-row theorem identifies
that size with the padded source size and preserves it under the mount.
The generic preparer interface requires the corresponding size and selector law.

\paragraph*{Lean anchors.}
The accompanying theorem index identifies the source-specialized J6
startup, group receipt, joint count and clock, and named workspace
bound. In the main route, the joint count and clock use the original
ordered-byte certificate. In the binary route, the selected distinct
points are the prefixes of the deterministic binary enumeration. The
paid table matches those actual ordered rows, and the binary count and
clock are taken from that same binary-order execution. Strict prefixes
of the original order are a separate choice. Both routes retain the
typed startup, the actual mounted preparation cap and the prescribed
root-contraction certificates.

The final independent R119 incremental replay passed for both routes.
It verified 3,136 main-route modules and 36 endpoint axiom reports,
and 1,483 binary-route modules and 35 endpoint axiom reports. Previously
built dependencies were reused after verification. The main-route
endpoint reports include the two J2 native-query guarantees and their
joint statement. The reported axiom dependencies of the stated
endpoints are drawn from the standard Lean axioms
\texttt{propext}, \texttt{Classical.choice} and \texttt{Quot.sound}.


\end{document}
